% =====================================================================
%  Fondamenti di Unreal Engine, lezione 3 — Texture, materiali e mesh:
%  la parte grafica dentro Unreal e la torretta.
%  Terza lezione del corso breve di introduzione a Unreal Engine, online
%  su Zoom, condotta da Davide al posto di Igor (che "aveva un
%  impiccio"). Davide terrà anche le ultime lezioni del corso, dedicate
%  al level design.
%  Contenuti, nell'ordine della lezione: l'attesa e le presentazioni
%  (chi fa game design, 3D art, character animation, narrative design);
%  che cosa fa un grafico dentro Unreal; il materiale didattico; la
%  texture come "carta regalo" e le UV; i tipi di mappa (Diffuse o
%  Albedo, Normal, Specular, la mappa di altezza per il displacement),
%  il texturing procedurale, il bianco come 1 e il nero come 0; i
%  canali RGB e le mappe impacchettate (AORM); il materiale come
%  vernice e l'editor dei materiali; Material Domain, Blend Mode e
%  Shading Model; i Material Expression Node; il problema della mappa
%  vuota; dentro una texture; Texture Sample e i canali; le costanti e
%  le scorciatoie (T, 1, 2, 3, A, M, D, S, U, L, P + clic); il colore e
%  l'emissive, Multiply e il parametro Glow Intensity; il PBR, Metallic,
%  Specular, Roughness, Normal; il Texture Coordinate e il tiling; il
%  pavimento; le Material Instance (i tavoli, il Tiling Multiplier); le
%  Material Function; Lerp e Panner; la pausa (il progetto che non si
%  trova, Git e il lavoro a distanza); l'import delle mesh (le opzioni
%  della finestra di import, la torretta in tre pezzi); vertici,
%  spigoli e facce; dimensioni, unità di misura e formati; il pivot del
%  cannone e il Reimport; le cartelle; lo Static Mesh Editor (gli slot
%  dei materiali, le modalità di visualizzazione, le collisioni simple
%  e complex); raggruppare e Merge Actors; i materiali della torretta;
%  il pivot con la Modeling Mode; il merge e gli slot; lo stretch, lo
%  z-fighting e i materiali World Aligned; l'esercizio della lava, il
%  nodo Noise e la chiusura.
%  Fonte: trascrizione automatica della terza lezione (UE3), circa tre
%  ore. Integrazioni dal libro "Blueprints Visual Scripting for Unreal
%  Engine 5" (Romero, Sewell), quinto capitolo, e nozioni tecniche
%  generali su Unreal Engine 5, sempre segnalate come approfondimenti.
%
%  Pacchetti richiesti nel preambolo del libro (già presenti nel
%  main.tex): \usepackage{booktabs}, \usepackage{amsmath},
%  \usepackage{tikz} con
%    \usetikzlibrary{shapes.geometric,shapes.symbols,shapes.misc,
%      arrows.meta,positioning,calc,fit}
%  (le librerie vanno caricate nel preambolo, non nel capitolo).
%  Consigliati, come per gli altri capitoli: \usepackage[italian]{babel}
%  e \usepackage[T1]{fontenc} (il capitolo compila anche senza).
%
%  Collocazione nel libro: parte "Fondamenti di Unreal Engine", terzo
%  capitolo (numerazione propria 1.3). Nel file principale del libro,
%  subito dopo la riga della seconda lezione:
%
%    % =====================================================================
%  Fondamenti di Unreal Engine, lezione 3 — Texture, materiali e mesh:
%  la parte grafica dentro Unreal e la torretta.
%  Terza lezione del corso breve di introduzione a Unreal Engine, online
%  su Zoom, condotta da Davide al posto di Igor (che "aveva un
%  impiccio"). Davide terrà anche le ultime lezioni del corso, dedicate
%  al level design.
%  Contenuti, nell'ordine della lezione: l'attesa e le presentazioni
%  (chi fa game design, 3D art, character animation, narrative design);
%  che cosa fa un grafico dentro Unreal; il materiale didattico; la
%  texture come "carta regalo" e le UV; i tipi di mappa (Diffuse o
%  Albedo, Normal, Specular, la mappa di altezza per il displacement),
%  il texturing procedurale, il bianco come 1 e il nero come 0; i
%  canali RGB e le mappe impacchettate (AORM); il materiale come
%  vernice e l'editor dei materiali; Material Domain, Blend Mode e
%  Shading Model; i Material Expression Node; il problema della mappa
%  vuota; dentro una texture; Texture Sample e i canali; le costanti e
%  le scorciatoie (T, 1, 2, 3, A, M, D, S, U, L, P + clic); il colore e
%  l'emissive, Multiply e il parametro Glow Intensity; il PBR, Metallic,
%  Specular, Roughness, Normal; il Texture Coordinate e il tiling; il
%  pavimento; le Material Instance (i tavoli, il Tiling Multiplier); le
%  Material Function; Lerp e Panner; la pausa (il progetto che non si
%  trova, Git e il lavoro a distanza); l'import delle mesh (le opzioni
%  della finestra di import, la torretta in tre pezzi); vertici,
%  spigoli e facce; dimensioni, unità di misura e formati; il pivot del
%  cannone e il Reimport; le cartelle; lo Static Mesh Editor (gli slot
%  dei materiali, le modalità di visualizzazione, le collisioni simple
%  e complex); raggruppare e Merge Actors; i materiali della torretta;
%  il pivot con la Modeling Mode; il merge e gli slot; lo stretch, lo
%  z-fighting e i materiali World Aligned; l'esercizio della lava, il
%  nodo Noise e la chiusura.
%  Fonte: trascrizione automatica della terza lezione (UE3), circa tre
%  ore. Integrazioni dal libro "Blueprints Visual Scripting for Unreal
%  Engine 5" (Romero, Sewell), quinto capitolo, e nozioni tecniche
%  generali su Unreal Engine 5, sempre segnalate come approfondimenti.
%
%  Pacchetti richiesti nel preambolo del libro (già presenti nel
%  main.tex): \usepackage{booktabs}, \usepackage{amsmath},
%  \usepackage{tikz} con
%    \usetikzlibrary{shapes.geometric,shapes.symbols,shapes.misc,
%      arrows.meta,positioning,calc,fit}
%  (le librerie vanno caricate nel preambolo, non nel capitolo).
%  Consigliati, come per gli altri capitoli: \usepackage[italian]{babel}
%  e \usepackage[T1]{fontenc} (il capitolo compila anche senza).
%
%  Collocazione nel libro: parte "Fondamenti di Unreal Engine", terzo
%  capitolo (numerazione propria 1.3). Nel file principale del libro,
%  subito dopo la riga della seconda lezione:
%
%    % =====================================================================
%  Fondamenti di Unreal Engine, lezione 3 — Texture, materiali e mesh:
%  la parte grafica dentro Unreal e la torretta.
%  Terza lezione del corso breve di introduzione a Unreal Engine, online
%  su Zoom, condotta da Davide al posto di Igor (che "aveva un
%  impiccio"). Davide terrà anche le ultime lezioni del corso, dedicate
%  al level design.
%  Contenuti, nell'ordine della lezione: l'attesa e le presentazioni
%  (chi fa game design, 3D art, character animation, narrative design);
%  che cosa fa un grafico dentro Unreal; il materiale didattico; la
%  texture come "carta regalo" e le UV; i tipi di mappa (Diffuse o
%  Albedo, Normal, Specular, la mappa di altezza per il displacement),
%  il texturing procedurale, il bianco come 1 e il nero come 0; i
%  canali RGB e le mappe impacchettate (AORM); il materiale come
%  vernice e l'editor dei materiali; Material Domain, Blend Mode e
%  Shading Model; i Material Expression Node; il problema della mappa
%  vuota; dentro una texture; Texture Sample e i canali; le costanti e
%  le scorciatoie (T, 1, 2, 3, A, M, D, S, U, L, P + clic); il colore e
%  l'emissive, Multiply e il parametro Glow Intensity; il PBR, Metallic,
%  Specular, Roughness, Normal; il Texture Coordinate e il tiling; il
%  pavimento; le Material Instance (i tavoli, il Tiling Multiplier); le
%  Material Function; Lerp e Panner; la pausa (il progetto che non si
%  trova, Git e il lavoro a distanza); l'import delle mesh (le opzioni
%  della finestra di import, la torretta in tre pezzi); vertici,
%  spigoli e facce; dimensioni, unità di misura e formati; il pivot del
%  cannone e il Reimport; le cartelle; lo Static Mesh Editor (gli slot
%  dei materiali, le modalità di visualizzazione, le collisioni simple
%  e complex); raggruppare e Merge Actors; i materiali della torretta;
%  il pivot con la Modeling Mode; il merge e gli slot; lo stretch, lo
%  z-fighting e i materiali World Aligned; l'esercizio della lava, il
%  nodo Noise e la chiusura.
%  Fonte: trascrizione automatica della terza lezione (UE3), circa tre
%  ore. Integrazioni dal libro "Blueprints Visual Scripting for Unreal
%  Engine 5" (Romero, Sewell), quinto capitolo, e nozioni tecniche
%  generali su Unreal Engine 5, sempre segnalate come approfondimenti.
%
%  Pacchetti richiesti nel preambolo del libro (già presenti nel
%  main.tex): \usepackage{booktabs}, \usepackage{amsmath},
%  \usepackage{tikz} con
%    \usetikzlibrary{shapes.geometric,shapes.symbols,shapes.misc,
%      arrows.meta,positioning,calc,fit}
%  (le librerie vanno caricate nel preambolo, non nel capitolo).
%  Consigliati, come per gli altri capitoli: \usepackage[italian]{babel}
%  e \usepackage[T1]{fontenc} (il capitolo compila anche senza).
%
%  Collocazione nel libro: parte "Fondamenti di Unreal Engine", terzo
%  capitolo (numerazione propria 1.3). Nel file principale del libro,
%  subito dopo la riga della seconda lezione:
%
%    % =====================================================================
%  Fondamenti di Unreal Engine, lezione 3 — Texture, materiali e mesh:
%  la parte grafica dentro Unreal e la torretta.
%  Terza lezione del corso breve di introduzione a Unreal Engine, online
%  su Zoom, condotta da Davide al posto di Igor (che "aveva un
%  impiccio"). Davide terrà anche le ultime lezioni del corso, dedicate
%  al level design.
%  Contenuti, nell'ordine della lezione: l'attesa e le presentazioni
%  (chi fa game design, 3D art, character animation, narrative design);
%  che cosa fa un grafico dentro Unreal; il materiale didattico; la
%  texture come "carta regalo" e le UV; i tipi di mappa (Diffuse o
%  Albedo, Normal, Specular, la mappa di altezza per il displacement),
%  il texturing procedurale, il bianco come 1 e il nero come 0; i
%  canali RGB e le mappe impacchettate (AORM); il materiale come
%  vernice e l'editor dei materiali; Material Domain, Blend Mode e
%  Shading Model; i Material Expression Node; il problema della mappa
%  vuota; dentro una texture; Texture Sample e i canali; le costanti e
%  le scorciatoie (T, 1, 2, 3, A, M, D, S, U, L, P + clic); il colore e
%  l'emissive, Multiply e il parametro Glow Intensity; il PBR, Metallic,
%  Specular, Roughness, Normal; il Texture Coordinate e il tiling; il
%  pavimento; le Material Instance (i tavoli, il Tiling Multiplier); le
%  Material Function; Lerp e Panner; la pausa (il progetto che non si
%  trova, Git e il lavoro a distanza); l'import delle mesh (le opzioni
%  della finestra di import, la torretta in tre pezzi); vertici,
%  spigoli e facce; dimensioni, unità di misura e formati; il pivot del
%  cannone e il Reimport; le cartelle; lo Static Mesh Editor (gli slot
%  dei materiali, le modalità di visualizzazione, le collisioni simple
%  e complex); raggruppare e Merge Actors; i materiali della torretta;
%  il pivot con la Modeling Mode; il merge e gli slot; lo stretch, lo
%  z-fighting e i materiali World Aligned; l'esercizio della lava, il
%  nodo Noise e la chiusura.
%  Fonte: trascrizione automatica della terza lezione (UE3), circa tre
%  ore. Integrazioni dal libro "Blueprints Visual Scripting for Unreal
%  Engine 5" (Romero, Sewell), quinto capitolo, e nozioni tecniche
%  generali su Unreal Engine 5, sempre segnalate come approfondimenti.
%
%  Pacchetti richiesti nel preambolo del libro (già presenti nel
%  main.tex): \usepackage{booktabs}, \usepackage{amsmath},
%  \usepackage{tikz} con
%    \usetikzlibrary{shapes.geometric,shapes.symbols,shapes.misc,
%      arrows.meta,positioning,calc,fit}
%  (le librerie vanno caricate nel preambolo, non nel capitolo).
%  Consigliati, come per gli altri capitoli: \usepackage[italian]{babel}
%  e \usepackage[T1]{fontenc} (il capitolo compila anche senza).
%
%  Collocazione nel libro: parte "Fondamenti di Unreal Engine", terzo
%  capitolo (numerazione propria 1.3). Nel file principale del libro,
%  subito dopo la riga della seconda lezione:
%
%    \input{UE/UE3_texture_materiali_e_mesh}
%
%  Tutte le etichette di questo capitolo hanno il prefisso "ue3:"
%  (ue3:cap:..., ue3:sec:..., ue3:fig:..., ue3:tab:...), e tutti gli
%  stili TikZ hanno il prefisso "ue3" (ue3node, ue3arr, ...) per non
%  collidere con chiavi native di TikZ né con quelli delle prime due
%  lezioni. Il capitolo è autonomo: non contiene \ref verso altri
%  capitoli (i rimandi alle lezioni precedenti sono a parole).
%  Le figure sono ricostruzioni schematiche di quanto mostrato e detto a
%  lezione, non copie delle slide o delle schermate del docente. Dove la
%  trascrizione è troppo rovinata per ricostruire un dettaglio, il
%  capitolo lo dice invece di inventarlo.
% =====================================================================

\chapter[Texture, materiali e mesh]{Texture, materiali e mesh: la grafica dentro
Unreal e la torretta}
\label{ue3:cap:materiali}

Questa è la \textbf{terza lezione} del corso breve di \textbf{introduzione a Unreal
Engine}, e cambia sia docente sia argomento. Al posto di Igor, che ``aveva un
impiccio'', c'è \textbf{Davide}, che si presenta così: ``io sono Davide e in questa
lezione, nello specifico, sostituisco Igor''. Davide tornerà più avanti: ``le ultime
lezioni del corso le fate con me'', e saranno lezioni ``un po' più basate su
\textbf{level design} e la creazione di mappe'', con qualche esercitazione insieme. Dopo
due lezioni sull'editor e sulla logica (piazzare oggetti, i \emph{transform}, i
Blueprint, la pallina sulla piattaforma), questa lezione guarda a ciò che ``fa il
grafico all'interno di Unreal'': le \textbf{mesh}, le \textbf{texture}, i
\textbf{materiali}, e ``qualche nodo, qualche cosa carina che potete fare'' combinando
mesh e materiali.

La lezione ha due metà separate da una pausa. La prima è sui \textbf{materiali}: che
cos'è una texture e come si avvolge su un oggetto attraverso le \textbf{UV}; quali tipi
di mappa esistono (\textbf{Diffuse}, \textbf{Normal}, \textbf{Specular}, la mappa di
altezza per il \emph{displacement}) e perché in queste mappe il bianco vale 1 e il nero
0; che cos'è un \textbf{materiale} e come si costruisce nell'\textbf{editor dei
materiali}, ``nello stesso modo dell'editor dei Blueprint''; le impostazioni principali
(\textbf{Material Domain}, \textbf{Blend Mode}, \textbf{Shading Model}); i nodi di base
(\textbf{Texture Sample}, le costanti, le operazioni, il \textbf{Texture Coordinate},
il \textbf{Lerp}, il \textbf{Panner}); il rendering \textbf{PBR} con Metallic,
Specular e Roughness; e le \textbf{Material Instance}, il modo ``bon ton'' di applicare i
materiali. La seconda metà è sulle \textbf{mesh}: si importa una \textbf{torretta} in
tre pezzi dal materiale didattico, si guarda la finestra di \textbf{import}, si
corregge il \textbf{pivot} del cannone con un \textbf{Reimport}, si esplorano lo Static
Mesh Editor e le \textbf{collisioni}, si fondono i pezzi con \textbf{Merge Actors}, si
creano a mano i materiali della torretta, e si chiude con i materiali \textbf{World
Aligned}, che non si deformano quando si stira l'oggetto. In fondo, un esercizio libero:
un materiale di \textbf{lava}.

La lezione è \textbf{online}, su Zoom, come le precedenti, e finisce poco prima delle
19 (``si sono fatte le 7''). La pausa, una decina di minuti, arriva dopo ``un'oretta e
mezza''. Davide, alla sua prima volta con questa classe, combatte per tutta la lezione
con la condivisione dello schermo (``io sto facendo a cazzotti da tutta la lezione con
la finestra di Zoom'', ``quanto non sono abituato a streamare'') e con un PC di
emergenza: ``il fisso mi è esploso letteralmente tre giorni fa'', e a metà della seconda
parte il computer ``inizia ad accusare la stream'', tanto che Davide gira il
ventilatore verso di lui. Chiede spesso conferma che gli studenti vedano quello che
mostra: ``mi impanicherò ogni tanto a caso e vi chiederò conferma delle robe''.

Tre avvertenze di lettura, le stesse delle prime due lezioni. La prima: la trascrizione
è automatica e storpia moltissimi nomi. Unreal diventa ``ArrayLink'', ``ArrayL4'',
``Arial Engine'' e ``a Real''; la \textbf{Material Instance} è quasi sempre
``materialistance'', ``Material Distance'' e una volta ``mie dire distanze''; il
\textbf{Panner} è ``panel'', ``planner'' e ``banner''; il \textbf{Lerp} è ``l'ERP'';
\emph{World Aligned} è ``word aligned'' e ``word.allocation''; la texture
\textbf{T\_Bark} (una corteccia) è ``Arc'', ``barca'', ``Barc'' e ``park''; ``strecciare''
e ``si straccia'' stanno per lo \emph{stretch}, lo stiramento; ``traccare'' sta per
\emph{tracciare} un bersaglio (\emph{to track}); la \textbf{mappa di altezza}
(\emph{height map}) è quasi sempre ``light map'' o ``lightmap''. Dove la ricostruzione
è sicura il capitolo usa la forma corretta e, la prima volta, lo segnala in nota; dove
non lo è riporta la forma della trascrizione e lo dice. La seconda: durante la pausa e
verso la fine il riconoscimento vocale produce lunghi blocchi di testo senza senso e
ripetuto (``qual è la risposta che ci può essere?'', discorsi sulla ``vita data per
scontata'' e su ``vittime''): sono artefatti, non parole di nessuno, e il capitolo li
ignora. La terza: gran parte della lezione è una dimostrazione dal vivo nell'editor dei
materiali e nello Static Mesh Editor. Il capitolo ricostruisce i grafi di nodi da ciò
che Davide dice mentre li costruisce, e li completa, dove serve, con i nomi esatti dei
nodi e dei comandi di Unreal Engine 5.

\paragraph{Come leggere questo capitolo.} Il capitolo segue l'ordine della lezione. Le
figure ricostruiscono gli schemi spiegati a voce, sulle slide o nell'editor: le mappe
della torretta, i canali di una texture, il rapporto tra texture, materiale e mesh, i
grafi dei materiali costruiti a lezione, le Material Instance, la struttura di una
mesh, il pivot del cannone. Le tabelle riassumono i tipi di mappa, le impostazioni del
materiale, le opzioni di import, le collisioni e le scorciatoie dell'editor dei
materiali. Qua e là, dei paragrafi di \textbf{approfondimento} collegano la lezione al
libro \emph{Blueprints Visual Scripting for Unreal Engine 5} di Marcos Romero e Brenden
Sewell, il testo di riferimento per questa parte del libro, e a qualche nozione tecnica
generale su Unreal Engine 5; sono segnalati come tali e non fanno parte della lezione.

Il percorso del capitolo:
\begin{enumerate}
  \item l'attesa, le presentazioni e il programma; che cosa fa un grafico dentro
        Unreal; il materiale didattico (sezioni \ref{ue3:sec:attesa}
        e~\ref{ue3:sec:didattico});
  \item le \textbf{texture}: la carta regalo, le \textbf{UV}, i tipi di mappa, il
        bianco e il nero, i \textbf{canali RGB} (sezioni da~\ref{ue3:sec:texture}
        a~\ref{ue3:sec:canali});
  \item i \textbf{materiali}: la vernice, l'editor, Material Domain, Blend Mode e
        Shading Model, i nodi (sezioni da~\ref{ue3:sec:materiale}
        a~\ref{ue3:sec:nodi});
  \item in Engine: la mappa vuota, dentro una texture, il \textbf{Texture Sample}, le
        costanti e le scorciatoie, il colore e l'\textbf{emissive} (sezioni
        da~\ref{ue3:sec:mappavuota} a~\ref{ue3:sec:emissive});
  \item il \textbf{PBR} e la Normal (sezione \ref{ue3:sec:pbr});
  \item il \textbf{Texture Coordinate}, il tiling e il pavimento (sezione
        \ref{ue3:sec:tiling});
  \item le \textbf{Material Instance} e le Material Function (sezioni
        \ref{ue3:sec:instance} e~\ref{ue3:sec:function});
  \item il \textbf{Lerp} e il \textbf{Panner}, poi la pausa (sezioni
        da~\ref{ue3:sec:lerp} a~\ref{ue3:sec:pausa});
  \item l'\textbf{import} e la torretta; che cos'è una \textbf{mesh}; dimensioni e
        formati (sezioni da~\ref{ue3:sec:import} a~\ref{ue3:sec:regole});
  \item il \textbf{pivot} del cannone e il \textbf{Reimport} (sezione
        \ref{ue3:sec:pivot});
  \item lo \textbf{Static Mesh Editor} e le \textbf{collisioni}; raggruppare e
        \textbf{Merge Actors} (sezioni \ref{ue3:sec:sme} e~\ref{ue3:sec:merge});
  \item i \textbf{materiali della torretta} e le domande sul Reimport e sul merge
        (sezioni \ref{ue3:sec:matturret} e~\ref{ue3:sec:domande});
  \item lo stretch, lo \textbf{z-fighting} e i materiali \textbf{World Aligned}
        (sezione \ref{ue3:sec:world});
  \item l'esercizio della \textbf{lava}, il nodo \textbf{Noise} e la chiusura
        (sezione \ref{ue3:sec:lava}).
\end{enumerate}

% ---------------------------------------------------------------------
\section{L'attesa, le presentazioni e il programma}
\label{ue3:sec:attesa}

``Ok, buonasera.'' La registrazione parte subito (``sta registrando? Sì, sta
registrando, perfetto''), e Davide aspetta ``qualche minutino'' che entrino tutti. Uno
studente avvisa che dovrà staccarsi venti minuti prima per andare a lavorare: ``nessun
problema'', tanto la lezione resta registrata, e anzi a Davide piacerebbe finire ``con
voi che fate qualcosina'', quindi andare via un po' in anticipo ``non è neanche un
problema''. Davide chiede quante persone c'erano alle lezioni precedenti (``una dozzina,
credo''), ringrazia chi accende la telecamera (``almeno so che non parlo da solo'') e
controlla che non si senta il ventilatore in sottofondo (``inizia a far caldo, inizia a
essere complicato'').

\paragraph{Che cosa avete fatto con Igor?} Prima di cominciare, Davide vuole sapere a
che punto è la classe, per non ripetere cose già viste. Dal Discord ha capito che si è
``ancora all'inizio'': le schede, le finestre, i \emph{transform} sugli oggetti, i
tavolini, ``qualcosa con gli input, con la pallina che si muove''. ``Se vediamo
qualcosa che avete già visto, ditemelo.''

\paragraph{Il programma.} La lezione di oggi guarda a quello che fa ``il grafico
all'interno di Unreal'': ``un po' di nozioni base per quanto riguarda le mesh, un po' di
nozioni base per quanto riguarda i materiali, le texture'', e ``già qualche nodo,
qualche cosa carina che potete fare'' usando mesh e materiali.

\paragraph{Chi c'è.} Davide chiede a che cosa sono iscritti gli studenti, ``per capire
chi ho davanti''. Le risposte: \textbf{game design} (più d'uno), \textbf{3D art}
(tre persone), \textbf{character animation} (due), \textbf{narrative design}. ``Beh, in
realtà siamo abbastanza assortiti.''

\paragraph{Che cosa fa un grafico dentro Unreal.} Ai futuri grafici 3D Davide mette
``le mani avanti'': la maggior parte del lavoro di un grafico \textbf{non} si svolge
dentro Unreal, ma in altri strumenti, che verranno spiegati nei rispettivi corsi.
Quello che si vede oggi è ciò che un grafico fa ``nel momento in cui va a lavorare con
altre persone'' dentro Unreal Engine: \textbf{importare} gli oggetti e creare gli
\textbf{shader} dentro l'engine.

\paragraph{Come si struttura la lezione.} Davide chiede come lavorava Igor: tutta
teoria e poi tutta pratica, o un po' e un po'? ``Andava prima con la presentazione e poi
con Unreal.'' Davide prova a replicare: ``ci facciamo prima una chiacchiera sulla parte
un po' più teorica e poi passiamo al progetto''. Nella pratica, come si vedrà, teoria
ed engine si alterneranno di continuo.

% ---------------------------------------------------------------------
\section{Il materiale didattico}
\label{ue3:sec:didattico}

Sul Discord c'è stato ``un po' di casino'' per il \textbf{materiale didattico}. ``Non è
molta roba'': è un piccolo progetto di Unreal, ``un progetto che Unreal usa per
spiegare alcune cose'', e probabilmente quello che Igor farà usare nelle lezioni
successive ogni volta che si implementa qualcosa. Gli altri file della cartella sono
\textbf{asset} da importare nel progetto poco alla volta: oggi la \textbf{torretta};
quando si farà l'audio, l'audio; ``poi vediamo se aggiungere altre cose''.

Dove si trova, Davide lo spiega più avanti a una studentessa: sul Discord, ``il 19
maggio, se non ricordo male'', c'è un link a una pagina di Drive con un file che si
chiama, più o meno, ``materiale didattico Unreal Engine Intro''\footnote{La trascrizione
dice ``IVE, materiale didattico, a Real Engine Intro'' e, poco prima, ``Hive Unreal
Engine 5 intro'': la parola iniziale non è ricostruibile con certezza.}. Lo si
scarica, lo si \textbf{decomprime}, e dentro c'è ``un progettino con un
\emph{character}'': sono soprattutto cose per le prossime lezioni con Igor, e per oggi
serve poco.

% ---------------------------------------------------------------------
\section{Che cos'è una texture}
\label{ue3:sec:texture}

\paragraph{La carta regalo.} ``Che cos'è una texture? La texture fondamentalmente è il
\emph{wrap}, è la \textbf{carta regalo} in cui è avvolto il nostro oggetto.'' Sulla
slide c'è la torretta del materiale didattico e, accanto, le sue texture: un'immagine
``tutta scomposta e segmentata e divisa in varie parti''. Perché divisa in parti?
Perché quelle parti poi ``si avvolgeranno sopra la mesh''. Le texture sono ``una serie
di mappe, una serie di informazioni che vanno a dare delle qualità specifiche
all'oggetto''. Con un'altra immagine: la texture ``non è altro che una
\textbf{decalcomania}'' dell'oggetto, della \emph{mesh} che si mette in scena.

\paragraph{Le UV.} Perché una texture funzioni bene su un oggetto, deve essere pensata
``con in mente le \textbf{UV} dei modelli su cui si lavora''. Le UV ``sono le
coordinate'', tutto quello che riguarda ``la struttura dell'oggetto'': la \textbf{mappa
delle UV} è la \textbf{prima parte} del processo di creazione di una texture. Dal
modello 3D si estraggono informazioni che poi servono a creare ``tutti i vari tipi di
texture che esistono'': sulla slide le texture della torretta sono diverse, e ognuna
porta informazioni diverse, ``ma sono tutte texture pensate per la stessa torretta''.

\paragraph{L'UV mapping.} L'\textbf{UV mapping} è la tecnica che serve a
\textbf{mappare le texture 2D su modelli 3D}: si assegnano le coordinate dette UV ``ai
vertici e agli \emph{edge} del modello'', e tutte le texture applicate dopo usano quelle
informazioni ``per precisione''. È un lavoro che si fa \textbf{prima} di entrare in
Unreal, in un altro software ``come Blender, come Maya''. Se si usa una texture che non
è stata pensata per le UV di quell'oggetto, ``la texture verrà strecciata''
(\emph{stretched}, stirata) o comunque non darà il risultato voluto.

\begin{figure}[htbp]
\centering
\begin{tikzpicture}[
  font=\small,
  ue3box/.style={draw, thick, align=center, minimum height=0.9cm},
  ue3arr/.style={-{Stealth[length=2.5mm]}, thick}
]
  % il cubo 3D
  \draw[thick, fill=orange!25] (0,0) -- (1.6,0) -- (1.6,1.6) -- (0,1.6) -- cycle;
  \draw[thick, fill=orange!40] (0,1.6) -- (0.7,2.2) -- (2.3,2.2) -- (1.6,1.6) -- cycle;
  \draw[thick, fill=orange!15] (1.6,0) -- (2.3,0.6) -- (2.3,2.2) -- (1.6,1.6) -- cycle;
  \node[font=\footnotesize, align=center] at (1.15,-0.5) {modello 3D\\(la mesh)};
  % l'apertura
  \draw[ue3arr] (2.8,1.1) -- node[above, font=\footnotesize]{si ``apre''}
    node[below, font=\footnotesize]{(UV mapping)} (5.0,1.1);
  % la mappa UV a croce
  \begin{scope}[shift={(5.4,-0.6)}]
    \draw[step=0.75, gray!40, thin] (0,0) grid (3.75,3.0);
    \foreach \x/\y in {0.75/0, 0/0.75, 0.75/0.75, 1.5/0.75, 2.25/0.75, 0.75/1.5}
      \draw[thick, fill=orange!25] (\x,\y) rectangle ++(0.75,0.75);
    \draw[->, thick] (0,0) -- (4.0,0) node[right, font=\footnotesize]{U};
    \draw[->, thick] (0,0) -- (0,3.25) node[above, font=\footnotesize]{V};
    \node[font=\footnotesize, align=center] at (1.9,-0.5) {mappa UV: le facce\\
      stese sul piano $[0,1]^2$};
  \end{scope}
  % la texture
  \draw[ue3arr] (9.4,1.1) -- node[above, font=\footnotesize]{si dipinge} (10.6,1.1);
  \node[ue3box, minimum width=2.3cm, minimum height=2.3cm, fill=orange!10,
    text width=2.1cm] (tx) at (12.0,1.1) {\footnotesize texture 2D\\
    \footnotesize (Diffuse, Normal,\\\footnotesize Roughness\ldots)\\[3pt]
    \footnotesize si ``riavvolge''\\\footnotesize sulla mesh};
\end{tikzpicture}
\caption{La texture come ``carta regalo'': le facce del modello vengono stese su un
piano (la mappa UV, fatta in Blender o in Maya), le texture si dipingono su quel piano,
e poi si ``riavvolgono'' sull'oggetto. È per questo che la texture della torretta, sulla
slide, appare ``tutta scomposta e segmentata''. Schema ricostruito; le lettere U e V
indicano i due assi della texture, per distinguerli da X, Y e Z del mondo.}
\label{ue3:fig:uv}
\end{figure}

\paragraph{Approfondimento: perché U e V.} Le lettere U e V sono scelte solo per non
confondere gli assi della texture con gli assi X, Y e Z dello spazio 3D. Una coordinata
UV è una coppia di numeri, di solito tra 0 e 1: $(0,0)$ è un angolo della texture e
$(1,1)$ l'angolo opposto. Ogni vertice della mesh riceve una coppia $(u,v)$, e per i
punti interni a una faccia il motore interpola le coordinate dei vertici. Se una mesh
non ha UV (o ne ha di sbagliate), il motore non sa ``quale pezzo di carta regalo'' va
su quale faccia: di qui le texture stirate di cui parla Davide.

% ---------------------------------------------------------------------
\section{I tipi di mappa}
\label{ue3:sec:mappe}

Davide passa in rassegna, sulla slide della torretta, le texture una per una, ``piano
piano'', con le informazioni che portano.

\paragraph{La Diffuse, o Albedo.} La prima è la \textbf{Diffuse}, ``in alcuni casi
detta anche \textbf{Albedo}''. Nei file la si riconosce dalla \textbf{nomenclatura}:
``T maiuscolo, \emph{underscore}, nome dell'oggetto, \emph{underscore}, A'', o D: A per
\emph{albedo}, D per \emph{diffuse} (per esempio \texttt{T\_Turret\_D}). La Diffuse è
il \textbf{colore}: ``tutti i dettagli estetici colorati''. Sulla slide è quella in alto
a destra, con l'arancione: è la texture che dà alla torretta l'arancione sulla base, il
nero negli altri dettagli e il bianco in alcune zone. Più avanti Davide ne completa la
descrizione: la Diffuse applica ``i colori e i dettagli di base'' sulle superfici 3D, e
dà l'informazione del colore, del tipo di materiale, dei dettagli (``queste linee che
vedete qua'') e di quanto un oggetto sia ruvido o levigato, se per esempio si applica la
texture di una roccia o di un metallo.

\paragraph{La Normal.} Poi c'è la \textbf{Normal}, la \textbf{Normal Map}, che ``serve
a simulare una tridimensionalità'' dentro un'informazione che è 2D. Si riconosce
subito: è \textbf{blu}, ``sempre mappe blu con delle sfumature di rosso e di verde''.
L'esempio classico è un \textbf{muro di mattoni}: in gioco il muro è \textbf{piatto}, ma
quando gli si applica la Normal il calcolo delle luci e delle ombre ``prende in
considerazione la normal map'' e proietta le ombre in base a lei, dando l'illusione che
i mattoni siano davvero tridimensionali. La Normal dà informazioni sulla
tridimensionalità ``senza appesantire particolarmente il calcolo'', perché non deforma
davvero l'oggetto come fa una mappa di altezza.

Più avanti Davide ci torna con un esempio sulla torretta. Su una delle gambe della base
servono dei \textbf{graffi}. Si potrebbe cambiare il modello aggiungendo i graffi, ma
questo aggiungerebbe \textbf{poligoni}, e quindi ``potenza di calcolo non
necessaria''. Si può invece aggiungere il dettaglio ``direttamente in 2D'', con la
Normal: protuberanze, graffi, rughe. ``Per un singolo oggetto non è un appesantimento
così importante. Ma considerate quanti oggetti ci sono all'interno di una mappa quando
andate a creare un livello'': se ogni oggetto non è creato con la giusta accortezza,
``si fa molto velocemente ad appesantire il progetto più di quanto dovrebbe essere
pesante''.

\paragraph{La Specular.} Lo \textbf{Specular mapping} riguarda ``la brillantezza e la
riflettività dell'oggetto'': come l'oggetto gestisce la luce, come la luce ``rimbalza''
su di lui. Un oggetto può avere parti metalliche, parti in pietra e altre ancora, e ``a
seconda del tipo di materiale di cui l'oggetto è fatto, rifrangerà la luce in una
maniera diversa''. Queste mappe sono quelle ``in bianco e in nero''. La texture
determina come si calcolano i riflessi ``sui singoli pixel proprio dell'oggetto'',
quanta luce viene riflessa; tornerà utile quando si vedranno le luci in Unreal, perché
``cambiando un materiale la luce ha un effetto completamente diverso''.

\paragraph{La mappa di altezza.} Poi c'è la mappa che serve a creare
\textbf{displacement}\footnote{Nella trascrizione è sempre ``light map'' o
``lightmap'', ma la descrizione (rialza e abbassa zone dell'oggetto, crea
\emph{tessellation}, è ``una normal sotto steroidi'') è quella della \textbf{height
map}, la mappa di altezza o di \emph{displacement}. Le \emph{lightmap} vere sono
un'altra cosa: texture in cui il motore salva l'illuminazione precalcolata.}: crea
``una \emph{tessellation} nell'oggetto'', rialza alcune zone, ne abbassa altre,
``modifica leggermente l'oggetto''. Davide la chiama ``una normal sotto steroidi'',
perché fa davvero \emph{displacement}: aggiunge profondità e dettagli tridimensionali
alle superfici piane, e usa le scale di grigio per creare scanalature e protuberanze.
``Nelle parti in cui vedete il bianco è sempre il sì'', la protuberanza; nel nero,
l'effetto opposto.

\paragraph{Il texturing procedurale.} Infine il \textbf{procedural texturing}: il
processo di creare texture, materiali o \textbf{shader} ``attraverso degli
algoritmi''. ``Andando ad utilizzare dei calcoli matematici, noi andiamo a modificare
delle informazioni per avere un risultato.'' È esattamente quello che fa l'editor dei
materiali di Unreal.

\begin{table}[htbp]
\centering
\small
\begin{tabular}{@{}p{2.5cm}p{2.4cm}p{8.1cm}@{}}
\toprule
\textbf{Mappa} & \textbf{Aspetto} & \textbf{Che cosa porta} \\
\midrule
Diffuse (Albedo) & a colori & il colore e i dettagli estetici; suffisso \_D o \_A \\
Normal & blu, con sfumature rosse e verdi & l'illusione del rilievo: la luce e le
  ombre si calcolano come se la superficie avesse graffi, mattoni, rughe, senza
  aggiungere poligoni \\
Specular & bianco e nero & quanto e come la superficie riflette la luce \\
Roughness & bianco e nero & quanto la superficie è ruvida (1) o liscia (0) \\
Metallic & bianco e nero & dove la superficie è metallo (1) e dove no (0) \\
Ambient Occlusion & bianco e nero & le zone in ombra perché ``chiuse'' (pieghe,
  angoli) \\
Height (displacement) & bianco e nero & deforma davvero la superficie: il bianco
  rialza, il nero abbassa \\
\bottomrule
\end{tabular}
\caption{I tipi di mappa nominati a lezione. Diffuse, Normal, Specular e la mappa di
altezza sono quelle della slide; Roughness, Metallic e Ambient Occlusion compaiono
nella parte sulle mappe impacchettate e nei materiali della torretta.}
\label{ue3:tab:mappe}
\end{table}

\paragraph{Il bianco è 1, il nero è 0.} La cosa importante da capire, insiste Davide, è
che nelle mappe in bianco e nero ``dovete sempre considerare il \textbf{bianco come un
1}, il \textbf{nero come uno 0} e il \textbf{grigio come dei valori nel mezzo} fra lo 0 e
l'1''. Che cosa vuol dire? Che in Unreal, per esempio, il valore dello Specular va
soltanto tra 0 e 1: a 1 riflette la luce, a 0 no. Una mappa di Specular, quindi, va
letta come una griglia di numeri: ``ci sono degli 1 in queste parti e degli 0 nel resto
delle parti''. È la chiave per capire quasi tutto il resto della lezione.

\paragraph{Approfondimento: la mappa di altezza, la Normal e Unreal Engine 5.} La
differenza tra Normal e mappa di altezza è la differenza tra \emph{sembrare} e
\emph{essere}. La Normal Map memorizza, per ogni pixel, la direzione in cui ``guarda''
la superficie (la sua \emph{normale}): i tre numeri X, Y e Z della direzione sono
salvati nei canali rosso, verde e blu, e siccome la maggior parte dei pixel guarda
``dritto fuori'' dalla superficie (Z alto) la mappa appare blu. La geometria non
cambia: di profilo, un muro con Normal Map resta piatto. La mappa di altezza invece
sposta davvero i vertici, e per farlo servono vertici da spostare: di qui la
\emph{tessellation}, che suddivide la superficie in molti triangoli. In Unreal Engine 5
il \emph{displacement} vero si fa soprattutto con le mesh \textbf{Nanite}, ed è più
costoso della Normal: per questo, per graffi e piccoli dettagli, la Normal resta la
scelta di default, come dice Davide.

% ---------------------------------------------------------------------
\section{I canali RGB e le mappe impacchettate}
\label{ue3:sec:canali}

``Chiaramente, queste texture utilizzano l'\textbf{RGB}'': non hanno un solo canale, ma
\textbf{tre canali} su cui passano le informazioni. Sulla slide c'è una Diffuse
scomposta in rosso, verde e blu: ognuno porta informazioni separate, che combinate
danno il colore che serve.

Ma c'è un altro modo di vedere la cosa. \textbf{Singolarmente}, ogni canale appare come
un'immagine in \textbf{bianco, nero e grigio}. E bianchi, neri e grigi ``sono
combinazioni di 0 e di 1'': quindi ogni canale si può usare per trasmettere
\textbf{un'informazione diversa}. È il \textbf{multichannel}. Una pratica comune,
nell'esportare e importare texture, è avere una singola mappa ``che spesso e volentieri
viene siglata come \textbf{AORM}'', per \textbf{Ambient Occlusion, Roughness e
Metallic}: tre mappe in bianco e nero, ognuna in un canale. ``Magari noi vogliamo
utilizzare questa mappa qui sopra per l'ambient occlusion, questa mappa qui sotto per
le metallic'', e così via. Una sola mappa con più informazioni invece di tante mappe
diverse: ``va ad alleggerire il progetto, va a rendere più veloce, più modulare''
l'export delle texture.

\begin{figure}[htbp]
\centering
\begin{tikzpicture}[
  font=\small,
  ue3ch/.style={draw, thick, minimum width=2.4cm, minimum height=1.0cm,
                align=center},
  ue3arr/.style={-{Stealth[length=2.5mm]}, thick}
]
  \node[draw, very thick, minimum width=2.6cm, minimum height=3.6cm, align=center,
    fill=gray!10] (tex) at (0,0) {\textbf{una texture}\\\textbf{RGB}\\[4pt]
    \footnotesize es. T\_Turret\_ORM};
  \node[ue3ch, fill=red!15] (r) at (4.2,1.3) {canale \textbf{R}};
  \node[ue3ch, fill=green!15] (g) at (4.2,0) {canale \textbf{G}};
  \node[ue3ch, fill=blue!15] (b) at (4.2,-1.3) {canale \textbf{B}};
  \draw[ue3arr] (tex.east) ++(0,1.0) -- (r.west);
  \draw[ue3arr] (tex.east) -- (g.west);
  \draw[ue3arr] (tex.east) ++(0,-1.0) -- (b.west);
  \node[ue3ch, text width=3.4cm] (ao) at (8.6,1.3) {\textbf{Ambient Occlusion}\\
    \footnotesize grigi tra 0 e 1};
  \node[ue3ch, text width=3.4cm] (ro) at (8.6,0) {\textbf{Roughness}\\
    \footnotesize grigi tra 0 e 1};
  \node[ue3ch, text width=3.4cm] (me) at (8.6,-1.3) {\textbf{Metallic}\\
    \footnotesize grigi tra 0 e 1};
  \draw[ue3arr] (r) -- (ao);
  \draw[ue3arr] (g) -- (ro);
  \draw[ue3arr] (b) -- (me);
  \node[font=\footnotesize, align=center, text width=11cm] at (4.3,-2.5)
    {ogni canale, da solo, è un'immagine in bianco e nero: tre ``mappe'' diverse in
    un solo file};
\end{tikzpicture}
\caption{Una mappa impacchettata (\emph{channel packing}). La lezione la chiama AORM;
nella pratica di Unreal la convenzione più diffusa è ORM, con l'Ambient Occlusion nel
rosso, la Roughness nel verde e il Metallic nel blu, ma l'ordine dipende da chi
esporta le texture.}
\label{ue3:fig:canali}
\end{figure}

\paragraph{Approfondimento: sRGB e compressione.} Quando si importa in Unreal una
texture impacchettata (o qualunque mappa ``di dati'', come Roughness o Metallic), conviene
aprirla e \textbf{togliere la spunta a sRGB}: l'opzione sRGB dice al motore che
l'immagine è un \emph{colore} da correggere per la percezione dell'occhio, mentre una
mappa di dati contiene numeri da leggere così come sono. Nelle stesse impostazioni, il
campo \emph{Compression Settings} si mette su \textbf{Masks} per le mappe impacchettate
e su \textbf{Normalmap} per le Normal (Unreal di solito lo riconosce da solo per le
texture con suffisso \_N). Se ci si dimentica, il materiale funziona lo stesso, ma i
grigi vengono letti un po' sbagliati e la superficie risulta più lucida o più opaca del
previsto.

% ---------------------------------------------------------------------
\section{Il materiale: la vernice}
\label{ue3:sec:materiale}

\paragraph{Le texture non si modificano.} ``Quello che possiamo fare con le texture in
Unreal è relativamente poco, perché non possiamo andare a modificare effettivamente le
texture all'interno dell'engine.'' E allora come si danno qualità visive alle mesh?
``Lo facciamo con i \textbf{materiali}.''

\paragraph{Il materiale è uno shader.} I materiali ``sono degli \textbf{shader}''. E
l'\textbf{editor dei materiali}, quello che si vedrà ``principalmente'' oggi, ``funziona
praticamente nello stesso modo dell'editor dei Blueprint'': una serie di \textbf{nodi}
collegati, che trasportano informazioni fino a un \textbf{nodo risultante}.

\paragraph{Vernice, non carta regalo.} La distinzione è netta: ``il materiale è
effettivamente la \textbf{vernice} che viene applicata sulla mesh, non la texture''. Sugli
oggetti si applicano \textbf{materiali}. Le texture servono a \textbf{costruire} i
materiali: si combinano le texture in un unico materiale, che poi si applica
sull'oggetto finale. Il materiale dice all'oggetto ``come dovrebbe interagire con la
luce, che tipo di informazioni visive dovrebbe portare'', dove ci sono
tridimensionalità e così via: ``tutto quello che abbiamo visto dalle varie texture
precedentemente viene tutto convogliato all'interno dei materiali''. Questo, precisa
Davide, se si parla di \textbf{oggetti fisici}. I materiali hanno anche altre funzioni:
servono per esempio a fare \textbf{effetti speciali}, e un materiale per un effetto
speciale ``dovrà essere costruito in maniera diversa'' da quello che riveste una
torretta.

\begin{figure}[htbp]
\centering
\begin{tikzpicture}[
  font=\small,
  ue3tex/.style={draw, thick, minimum width=2.3cm, minimum height=0.7cm,
                 align=center, fill=gray!10},
  ue3big/.style={draw, very thick, rounded corners=4pt, align=center},
  ue3arr/.style={-{Stealth[length=2.5mm]}, thick}
]
  \node[ue3tex] (t1) at (0,1.5) {T\_Turret\_D};
  \node[ue3tex] (t2) at (0,0.5) {T\_Turret\_N};
  \node[ue3tex] (t3) at (0,-0.5) {T\_Turret\_R};
  \node[ue3tex] (t4) at (0,-1.5) {T\_Turret\_M};
  \node[font=\footnotesize, align=center] at (0,2.4) {\textbf{texture}\\(immagini)};
  \node[ue3big, minimum width=3.4cm, minimum height=3.6cm, fill=blue!6,
    text width=3.2cm] (mat) at (4.4,0) {\textbf{materiale}\\M\_Turret\\[6pt]
    \footnotesize un grafo di nodi\\\footnotesize (uno \emph{shader})\\[6pt]
    \footnotesize la ``vernice''};
  \foreach \tn in {t1,t2,t3,t4} \draw[ue3arr] (\tn.east) -- (mat.west |- \tn);
  \node[ue3big, minimum width=2.8cm, minimum height=1.6cm, fill=orange!12,
    text width=2.6cm] (mi) at (8.4,0) {\textbf{Material}\\\textbf{Instance}\\
    \footnotesize MI\_Turret\\\footnotesize (parametri)};
  \draw[ue3arr] (mat) -- (mi);
  \node[ue3big, minimum width=2.4cm, minimum height=1.6cm, fill=green!8,
    text width=2.2cm] (mesh) at (11.8,0) {\textbf{mesh}\\\footnotesize SM\_Turret\\
    \footnotesize slot del\\\footnotesize materiale};
  \draw[ue3arr] (mi) -- node[above, font=\footnotesize]{si applica} (mesh);
\end{tikzpicture}
\caption{Texture, materiale, istanza e mesh. Le texture non si applicano direttamente
agli oggetti: entrano in un materiale, che è la ``vernice''; dal materiale si crea
un'istanza, ed è l'istanza che si applica allo slot del materiale della mesh
(sezione~\ref{ue3:sec:instance}). I nomi seguono la nomenclatura suggerita a lezione
(T\_, M\_, MI\_, SM\_).}
\label{ue3:fig:pipeline}
\end{figure}

\paragraph{Approfondimento: il materiale nel libro di Romero e Sewell.} Il quinto
capitolo del libro (``Object Interaction with Blueprints'') usa la stessa immagine di
Davide: un materiale è un \emph{asset} che si aggiunge alla mesh di un actor per
crearne l'aspetto, ``come uno strato di vernice''. Il libro crea un materiale rosso,
\texttt{M\_TargetRed}, in una cartella Materials (``non è necessario, ma è buona
pratica tenere ordinata la gerarchia del progetto''), e nota che l'editor dei
materiali ``condivide molte caratteristiche e convenzioni con i Blueprint'': al centro
c'è il grafo; il nodo grande con il nome del materiale è il \textbf{result node}; con il
tasto destro si apre un menu di nodi con una ricerca sensibile al contesto; i fili si
tirano trascinando dai pin, gli ingressi a sinistra e le uscite a destra. Il colore lo
dà un \textbf{VectorParameter} collegato a Base Color; due \textbf{ScalarParameter} a
0,1 danno un po' di Metallic e di Roughness. Poi, in un Blueprint, il nodo \textbf{Set
Material} cambia il materiale di un bersaglio quando viene colpito: un primo esempio di
come la logica (le lezioni di Igor) e la grafica (questa lezione) si incontrino.

% ---------------------------------------------------------------------
\section{Material Domain, Blend Mode e Shading Model}
\label{ue3:sec:dettagli}

Come gli oggetti nella \emph{viewport}, anche il materiale ha i suoi
\textbf{Details}. Dentro l'editor si possono selezionare i singoli nodi, per i dettagli
del nodo, oppure vedere i \textbf{dettagli generali del materiale}. Tra questi ce ne
sono alcuni, evidenziati sulla slide, che ``cambiano proprio effettivamente come
funziona il materiale'' e ne definiscono la struttura e le possibilità.

\paragraph{Material Domain.} Il \textbf{Material Domain} dice \textbf{per che cosa}
verrà usato il materiale: una \textbf{superficie}, l'\textbf{interfaccia} (la UI), il
\textbf{post process}, un effetto.

\paragraph{Blend Mode.} La \textbf{Blend Mode} è il modo in cui il materiale
``\emph{blenda}'' (si mescola) con i pixel calcolati nel resto della scena. La modalità
\textbf{Opaque}, quella di default, \textbf{occlude completamente} ciò che c'è dietro.
Per un effetto speciale, per esempio un \textbf{fuoco}, serve invece la Blend Mode
\textbf{Translucent}, che permette di aggiungere parametri per l'\textbf{opacità} del
materiale.

\paragraph{Shading Model.} Lo \textbf{Shading Model}\footnote{Nella trascrizione una
volta ``scaling model''.} definisce come il materiale \textbf{reagisce alla luce}.
\textbf{Default Lit} vuol dire che segue le regole di base dell'illuminazione: il
materiale viene illuminato. Per un effetto speciale, in cui non si vuole che la luce
influenzi la resa, lo si mette invece su \textbf{Unlit}\footnote{Nella trascrizione
``lo metterò un lit''.}. Ci sono ``diverse tipologie'' di Shading Model, a seconda di
cosa si vuole realizzare.

\paragraph{I pin cambiano.} Queste impostazioni si applicano sul nodo risultante, e
cambiarle cambia i \textbf{pin} del nodo. L'\textbf{Opacity} non si può modificare in
Blend Mode Opaque, perché un materiale opaco ``dà per scontato che l'oggetto sia sempre
opaco'': non c'è bisogno di dire dove e quando sarà trasparente. Viceversa, la Blend
Mode Translucent toglie la possibilità di modificare Metallic, Specular, Roughness e
così via: se il materiale deve essere potenzialmente trasparente, ``non ha senso andare
a calcolare come rifletterà la luce''.

\begin{table}[htbp]
\centering
\small
\begin{tabular}{@{}p{2.6cm}p{3.3cm}p{7.1cm}@{}}
\toprule
\textbf{Impostazione} & \textbf{Valori citati} & \textbf{Che cosa decide} \\
\midrule
Material Domain & Surface, User Interface, Post Process\ldots & per che cosa si usa
  il materiale: una superficie, la UI, il post process, un effetto \\
Blend Mode & \textbf{Opaque} (default), Masked, Translucent & come il materiale si
  mescola con ciò che ha dietro: Opaque nasconde tutto; Masked ritaglia con una
  maschera; Translucent sfuma con l'Opacity (fuoco, fili d'erba) \\
Shading Model & \textbf{Default Lit} (default), Unlit\ldots & come reagisce alla luce:
  Default Lit viene illuminato, Unlit no (effetti speciali) \\
\bottomrule
\end{tabular}
\caption{Le tre impostazioni ``più importanti'' dei dettagli di un materiale. Cambiarle
accende e spegne i pin del nodo risultante.}
\label{ue3:tab:dettagli}
\end{table}

\paragraph{La prova in Engine.} Per far vedere la cosa, Davide apre il progetto del
materiale didattico e invita chi vuole a seguirlo: nel \textbf{Content Browser} si crea
un nuovo materiale (``va bene anche chiamarlo New Material''), lo si apre con un
\textbf{doppio clic} e si entra nell'editor. Le impostazioni di cui si parlava sono nei
dettagli, in basso a sinistra. Se la Blend Mode passa da Opaque a \textbf{Translucent},
il nodo risultante cambia: spariscono Metallic, Specular, Roughness, Anisotropy, e
compare \textbf{Opacity}. È utile, per esempio, per una \textbf{fiamma} che verso l'alto
diventa trasparente (le informazioni grafiche alla base, e in alto ``piano piano va a
dissolversi''), o per un \textbf{filo d'erba} che parte con il suo colore e poi va in
trasparenza. Con la Blend Mode \textbf{Masked}, ``mascherata'', si crea invece una
\textbf{maschera}, un \emph{layer} che funziona in modo diverso dal resto della
superficie. Per ora si resta su \textbf{Opaque}, perché i materiali di oggi andranno su
oggetti 3D. A metà spiegazione arriva uno studente in ritardo: ``Chat, buon pomeriggio,
Lorenzo. Noi qui abbiamo già cominciato, ma in realtà abbiamo parlato di cose
principalmente teoriche.''

\paragraph{Approfondimento: Masked o Translucent?} Le due modalità sembrano simili ma
lavorano in modo opposto. \textbf{Masked} decide pixel per pixel se il materiale
\emph{c'è o non c'è}: l'ingresso Opacity Mask viene confrontato con una soglia (di
default 0,333), e sotto la soglia il pixel non viene disegnato affatto. È la modalità
delle foglie, delle reti e dell'erba ``ritagliata'': economica, compatibile con la
luce e con le ombre, ma senza sfumature. \textbf{Translucent} invece \emph{mescola} il
colore del materiale con ciò che sta dietro, secondo un valore di Opacity tra 0 e 1:
permette le sfumature (fumo, vetro, fiamme), ma costa di più, perché ogni strato
traslucido va sommato a quelli dietro, e gestisce peggio luce e ordinamento. Regola
pratica: Masked quando basta un ``sì o no'', Translucent solo quando serve davvero la
sfumatura.

% ---------------------------------------------------------------------
\section{I nodi: le Material Expression}
\label{ue3:sec:nodi}

``Come per i Blueprint, anche l'editor dei materiali ha dei nodi, delle
\textbf{espressioni}.'' I \textbf{Material Expression Node} sono ``dei piccoli set di
istruzioni preprogrammate che eseguono delle cose nello specifico, esattamente come i
nodi dei Blueprint'', o, con un'altra formula, ``pezzi di codice statico''.

\paragraph{Il primo grafo.} Sulla slide c'è un piccolo grafo
(figura~\ref{ue3:fig:primografo}), e Davide lo legge nodo per nodo:
\begin{itemize}
  \item il \textbf{Texture Sample} comunica l'informazione della texture: ``sta dicendo
        che la texture che c'è in questo nodo è la texture che determina il
        \textbf{base color}''; permette di usare un'immagine e farla diventare la
        texture del materiale;
  \item il \textbf{Texture Coordinate}, il nodo \textbf{rosso}, restituisce le
        coordinate U e V del materiale, e serve a decidere ``quanto è grande il
        \textbf{tiling} della mia texture, quindi quante volte si ripete'';
  \item il \textbf{Multiply} è programmato per moltiplicare due valori: ``moltiplica A
        per B e dà una risultante'';
  \item il quadratino è una \textbf{costante}, ``un singolo \emph{float}'', una variabile
        singola.
\end{itemize}
Il grafo dice: ``moltiplica le coordinate per uno e comunica queste cose alle UV della
texture''. Così com'è, quindi, ``è come se non ci fosse'': moltiplicare per 1 lascia le
coordinate di default. Ma è così che si costruiscono i grafi: combinando nodi si
ottengono risultati molto diversi. Si può moltiplicare, si può fare
un'\textbf{interpolazione}, si possono far \textbf{muovere} le texture per suggerire un
movimento nel materiale, ``eccetera, eccetera, eccetera''.

\begin{figure}[htbp]
\centering
\begin{tikzpicture}[
  font=\small,
  ue3node/.style={draw, thick, rounded corners=4pt, align=center,
                  minimum height=0.9cm},
  ue3uv/.style={-{Stealth[length=2mm]}, thick, red!60!black},
  ue3flt/.style={-{Stealth[length=2mm]}, thick, green!50!black},
  ue3col/.style={-{Stealth[length=2mm]}, thick, blue!60!black}
]
  \node[ue3node, fill=red!12, text width=2.2cm] (tc) at (0,0.6)
    {\textbf{TexCoord}\\\footnotesize UTiling 1, VTiling 1};
  \node[ue3node, fill=green!8, minimum width=1.0cm] (k) at (0,-1.0) {\textbf{1}};
  \node[font=\footnotesize] at (0,-1.75) {costante};
  \node[ue3node, fill=gray!10, text width=1.6cm] (mul) at (3.0,0)
    {\textbf{Multiply}\\\footnotesize A $\times$ B};
  \node[ue3node, fill=gray!10, text width=2.4cm, minimum height=1.8cm] (ts)
    at (6.4,0) {\textbf{Texture}\\\textbf{Sample}\\\footnotesize (T\_Bark)};
  \node[ue3node, fill=blue!6, text width=2.4cm, minimum height=2.6cm] (res)
    at (10.4,0) {\textbf{M\_NewMaterial}\\[3pt]\footnotesize\textbf{Base Color}\\
    \footnotesize Metallic\\\footnotesize Specular\\\footnotesize Roughness\\
    \footnotesize Emissive Color\\\footnotesize Normal\ldots};
  \draw[ue3uv] (tc.east) -- ++(0.4,0) |- node[pos=0.75, above, font=\scriptsize]{A}
    ([yshift=2mm]mul.west);
  \draw[ue3flt] (k.east) -- ++(0.6,0) |- node[pos=0.75, below, font=\scriptsize]{B}
    ([yshift=-2mm]mul.west);
  \draw[ue3uv] (mul.east) -- node[above, font=\scriptsize]{UVs}
    (ts.west |- mul.east);
  \draw[ue3col] (ts.east) -- node[above, font=\scriptsize]{RGB}
    (res.west |- ts.east);
\end{tikzpicture}
\caption{Il primo grafo della slide: le coordinate UV, moltiplicate per una costante,
entrano nel Texture Sample, il cui colore va nel Base Color del nodo risultante. Con la
costante a 1 il Multiply ``è come se non ci fosse''; cambiando la costante cambia il
tiling (sezione~\ref{ue3:sec:tiling}). Schema ricostruito.}
\label{ue3:fig:primografo}
\end{figure}

\paragraph{I parametri.} Non solo le costanti, ma i nodi in generale ``possono essere
utilizzati come \textbf{parametri}''. Per rendere un nodo un parametro si fa
\textbf{tasto destro} sul nodo e \textbf{Convert to Parameter}, gli si dà un nome, e
questo permette di modificare in seguito il valore ``senza dover ricalcolare da capo il
materiale''. I parametri torneranno utili con le \textbf{Material Instance}
(sezione~\ref{ue3:sec:instance}).

\paragraph{Come si chiama un nodo.} ``Tasto destro, scriviamo il nome'', per esempio
\emph{texture sample}, ed ecco il nodo. C'è poi ``tutta una serie di \emph{shortcut}''
per i nodi più usati, che si vedranno tra poco.

% ---------------------------------------------------------------------
\section{Il problema della mappa vuota}
\label{ue3:sec:mappavuota}

Una studentessa interrompe: non riesce più ad aprire la mappa, ``me la apre totalmente
vuota, senza titolo''. Davide le chiede di condividere lo schermo, ma lei segue da due
computer diversi. Prima ipotesi: ha creato un \textbf{nuovo livello}, e un \emph{New
Level} appena caricato è tutto nero, perché non c'è niente in scena (``potremmo andare
a popolarla mettendoci cose''). Ma il problema è un altro: ha aperto il progetto del
materiale didattico, cliccato sulla mappa del personaggio in terza persona
(\emph{Third Person Map}), e non succede niente. Davide sospetta che il progetto sia
ancora in una \textbf{cartella zippata} (``di solito questa cosa succede quando il
progetto è una cartella che è ancora zippata''), e le consiglia di
\textbf{riscaricarlo}. Poco dopo il problema è risolto: ``sono riuscita a risolvere
adesso''. ``Meno male, meno male.''

\paragraph{Approfondimento: aprire un progetto da una cartella.} Un progetto di Unreal
è una cartella che contiene un file \texttt{.uproject} e le sottocartelle
\texttt{Content}, \texttt{Config} e altre. Unreal non può aprire un progetto che sta
ancora dentro un archivio \texttt{.zip}: Windows mostra il contenuto dello zip come se
fosse una cartella, ma l'editor non ci può scrivere, e il risultato tipico è proprio
una mappa vuota o mancante. La soluzione è estrarre tutto (``Estrai tutto'') e aprire il
\texttt{.uproject} estratto con un doppio clic. Se il progetto non compare nell'elenco
dell'Epic Games Launcher o nel Project Browser, basta aprirlo dal \texttt{.uproject}
la prima volta: da lì in poi compare tra i progetti recenti.

% ---------------------------------------------------------------------
\section{Dentro una texture e il Texture Sample}
\label{ue3:sec:sample}

``Facciamo le cose in ordine, se no io mi confondo e mi impiccio.'' Davide torna al
punto saltato: com'è fatta \textbf{dentro} una texture.

\paragraph{La finestra della texture.} Con un doppio clic su una texture nel Content
Browser (Davide ne prende una a caso, una corteccia, \textbf{T\_Bark}) si apre una
finestra in cui ``non possiamo fare grandi modifiche'', ma si possono vedere i
\textbf{canali}: il blu, il verde, il rosso, ciascuno fatto a modo suo. ``Questa texture
non è particolarmente interessante'', ma si rivedrà meglio con la torretta.

\paragraph{Il Texture Sample.} Nel grafo del materiale si mette un \textbf{Texture
Sample}, si seleziona il nodo, e nei dettagli, in basso, si sceglie quale texture
metterci. Cercando T\_Bark e mettendola nel nodo, la \emph{preview} mostra la corteccia:
``il nostro materiale sarà così, perché è effettivamente la texture che ci abbiamo
messo''.

\paragraph{Le uscite separate.} Il Texture Sample non ha una sola uscita: oltre
all'\textbf{RGB} ci sono \textbf{R}, \textbf{G} e \textbf{B} separati (e l'alfa, A). Si
può quindi usare, per esempio, solo il canale \textbf{rosso} (``non quella completa, ma
questa qui'') per trasportare un'informazione: collegandolo allo \textbf{Specular} si
ottiene un materiale che in alcune parti riflette la luce più che in altre. ``Con
questo materiale in realtà è un po' complicato farlo vedere'', ma il principio è quello
dei canali visto sulle slide.

\paragraph{Come si mette una texture nel nodo.} Una studentessa non ha capito come si
aggiunge la texture ``al quadratino nero'' del Texture Sample. Davide rifà il
percorso: si clicca il nodo, e nei suoi dettagli c'è il campo \textbf{Texture}, con il
quadratino e una lista. Due modi:
\begin{itemize}
  \item se una texture è già \textbf{selezionata nel Content Browser}, la freccetta
        \textbf{Use Selected Asset from Content Browser} la inserisce automaticamente
        nel nodo;
  \item più spesso, si clicca sul \textbf{nome} della texture nel campo e si fa un
        \emph{browse}: si cerca il nome (``T underscore Bark'') e lo si sceglie dalla
        lista di tutte le texture disponibili.
\end{itemize}
``Voi lo vedete questo pop-up che esce?'' Sì: perfetto.

% ---------------------------------------------------------------------
\section{Le costanti, le operazioni e le scorciatoie}
\label{ue3:sec:costanti}

``Chiaramente ne vediamo un po', non li vediamo tutti, perché sono tanti, sono
veramente tanti.'' Ci sono ``tantissime possibilità, tante cose che vanno studiate'',
ma per ora servono le basi.

\paragraph{La regola delle scorciatoie.} Il Texture Sample si può chiamare cercandolo
con il tasto destro, oppure \textbf{tenendo premuto il tasto T e facendo clic
sinistro} nel grafo. ``Questa struttura dello shortcut funziona per ogni singolo nodo
che ha uno shortcut'': un tasto, tenuto premuto, più clic sinistro. Il tasto è
``molto spesso l'iniziale del nome del nodo''.

\paragraph{Le costanti.} Cercando \emph{vector} compare, per esempio, il Vector
Noise; cercando \emph{constant} compaiono le costanti: \textbf{Constant},
\textbf{Constant2Vector}, \textbf{Constant3Vector} (e Constant4Vector). Le scorciatoie
sono i numeri: \textbf{1 + clic} per la costante singola, \textbf{2 + clic} per quella
a due valori, \textbf{3 + clic} per quella a tre. La costante a due ``la utilizzerete
relativamente poco''. Le altre due moltissimo:
\begin{itemize}
  \item la \textbf{costante singola} è un \textbf{valore numerico}, che si usa ``per
        andare a fare qualsiasi tipo di calcolo, qualsiasi tipo di costruzione'';
  \item la \textbf{costante a tre} è un \textbf{colore}: l'engine usa i tre valori come
        \textbf{R}, \textbf{G} e \textbf{B}. Si sceglie un colore (Davide prende un
        arancione) e l'engine modifica X, Y e Z per ottenerlo. ``Qualcuno di voi l'aveva
        anche già fatto nel test con la pallina.''
\end{itemize}
Collegando l'arancione al \textbf{Base Color} si ottiene un materiale arancione,
``perché effettivamente è l'unica cosa che gli sto passando''. Ma ``molto raramente''
si usa un colore piatto come Base Color: lo si usa piuttosto per fare
\textbf{operazioni}. Più avanti una studentessa chiederà di nuovo come si crea il colore:
``è come la costante singola, che è uno clic sinistro, ma è a tre, quindi è tre clic
sinistro''.

\paragraph{Staccare i fili.} Per rompere un collegamento: \textbf{Alt + clic sinistro}
sul pin, ``e si rompono''.

\paragraph{Le operazioni.} Le operazioni utili sono ``le classiche'': addizione,
moltiplicazione, divisione, sottrazione. Le scorciatoie: \textbf{A + clic} per
\textbf{Add}, \textbf{M + clic} per \textbf{Multiply}, \textbf{D + clic} per
\textbf{Divide}. La sottrazione (\textbf{Subtract}) non ha scorciatoia: si cerca. Il
motivo lo spiega Davide più avanti: \textbf{S + clic} crea automaticamente un
\textbf{parametro} (uno Scalar Parameter), ``quindi questo potrebbe tornarvi utile''.

\paragraph{Quali si usano davvero.} In chiusura di lezione Davide aggiunge un consiglio
d'esperienza: ``l'operazione che utilizzerete di più in assoluto è la
\textbf{moltiplicazione}''. Moltiplicazione e addizione sono le più usate; divisione e
sottrazione servono soprattutto per cose più complesse, ``tipo sugli effetti
speciali''.

% ---------------------------------------------------------------------
\section{Il colore che brilla: Emissive, Multiply e Glow Intensity}
\label{ue3:sec:emissive}

``A cosa ci servono'' le operazioni? Per esempio ad \textbf{aumentare il valore} di
qualcosa. Il nodo risultante ha diversi pin, ``diverse cose su cui andare a collegare
dei valori'': il \textbf{Base Color}, e più in basso l'\textbf{Emissive Color}, che
controlla ``quali parti del materiale emetteranno una piccola luce'', un
\textbf{\emph{glow}}.

\paragraph{L'arancione nell'Emissive.} Davide stacca l'arancione dal Base Color e lo
collega all'Emissive: il materiale ``sarà leggermente luminoso''. ``Adesso è molto poco
luminoso, ma è comunque un emissive.''

\paragraph{Moltiplicare.} Come si ottiene un'emissione maggiore? ``Moltiplichiamo.''
L'arancione entra nella A di un \textbf{Multiply}, il Multiply va nell'Emissive. Il
secondo valore si può scrivere direttamente nel nodo Multiply (``invece di 1 mettiamo
100''): ``verrà decisamente più luminoso''. Ma così ``è meno controllabile''. Meglio
attaccare alla B una \textbf{costante} separata, darle il valore 100 e convertirla in
\textbf{parametro}, con un nome come \textbf{Glow Intensity}: quando più avanti ci sarà
una Material Instance, il valore si potrà cambiare ``direttamente da lì''. ``Ci andiamo
a estrapolare un parametro.''

\begin{figure}[htbp]
\centering
\begin{tikzpicture}[
  font=\small,
  ue3node/.style={draw, thick, rounded corners=4pt, align=center,
                  minimum height=0.9cm},
  ue3col/.style={-{Stealth[length=2mm]}, thick, blue!60!black},
  ue3flt/.style={-{Stealth[length=2mm]}, thick, green!50!black}
]
  \node[ue3node, fill=orange!30, text width=2.5cm] (col) at (0,0.7)
    {\textbf{Constant3Vector}\\\footnotesize arancione (3 + clic)};
  \node[ue3node, fill=green!8, text width=2.5cm] (gi) at (0,-0.9)
    {\textbf{Glow Intensity}\\\footnotesize parametro = 100};
  \node[ue3node, fill=gray!10, text width=1.6cm] (mul) at (3.6,0)
    {\textbf{Multiply}\\\footnotesize (M + clic)};
  \node[ue3node, fill=blue!6, text width=2.6cm, minimum height=2.2cm] (res)
    at (7.4,0) {\textbf{nodo risultante}\\[3pt]\footnotesize Base Color\\
    \footnotesize Metallic, Specular\\\footnotesize Roughness\\
    \footnotesize\textbf{Emissive Color}};
  \draw[ue3col] (col.east) -- ++(0.4,0) |- node[pos=0.75, above, font=\scriptsize]{A}
    ([yshift=2mm]mul.west);
  \draw[ue3flt] (gi.east) -- ++(0.4,0) |- node[pos=0.75, below, font=\scriptsize]{B}
    ([yshift=-2mm]mul.west);
  \draw[ue3col] (mul.east) -- ++(0.6,0) |- ([yshift=-7mm]res.west);
\end{tikzpicture}
\caption{Un materiale che brilla. Il colore, collegato all'Emissive Color, emette ``una
piccola luce''; moltiplicato per 100 diventa molto più luminoso. Il valore 100 sta in
una costante convertita in parametro (Glow Intensity), così che si possa cambiare da
una Material Instance senza toccare il grafo. Ricostruzione dalla dimostrazione.}
\label{ue3:fig:emissive}
\end{figure}

\paragraph{Approfondimento: perché 100 funziona.} I colori di un materiale non sono
limitati tra 0 e 1. Nel Base Color valori oltre 1 non hanno senso fisico (nessuna
superficie riflette più luce di quella che riceve), ma nell'Emissive sì: un'emissione di
$(1; 0{,}5; 0)$ moltiplicata per 100 diventa $(100; 50; 0)$, una sorgente molto intensa
che il \emph{post process} di Unreal trasforma in un alone luminoso (il \emph{bloom}).
Il Multiply tra un colore e un numero moltiplica ciascun canale per quel numero, quindi
la tinta resta la stessa e cresce solo l'intensità. Con Lumen, il sistema di
illuminazione globale di Unreal Engine 5, un materiale emissivo può anche illuminare un
po' le superfici vicine.

% ---------------------------------------------------------------------
\section{Il PBR: Metallic, Specular, Roughness e Normal}
\label{ue3:sec:pbr}

``Mi rendo conto che sto andando un po' a destra e a sinistra'', ammette Davide, e
torna ai pin del nodo risultante. Base Color ed Emissive si sono visti. Restano
\textbf{Metallic}, \textbf{Specular} e \textbf{Roughness}, ``degli input abbastanza
interessanti'', perché sono quelli che \textbf{determinano la superficie} del
materiale.

\paragraph{Il PBR.} Il rendering dei materiali in Unreal è \textbf{PBR},
\emph{Physically Based Rendering}: per decidere come si comporta una superficie, il
motore mette insieme una serie di informazioni fisiche. Tutti questi parametri vanno
\textbf{da 0 a 1}.

\paragraph{Roughness.} La \textbf{Roughness} dice quanto una superficie è
\textbf{liscia} o \textbf{ruvida}. Una \textbf{roccia} vuole una Roughness molto alta,
cioè 1: il materiale sarà ruvido ``e non rifletterà la luce in nessun modo''. Uno
\textbf{specchio} vuole una Roughness a 0: il più liscio possibile, il meno ruvido, il
più riflettente.

\paragraph{Specular e Metallic.} ``Per aiutare in questa situazione'' ci sono Specular e
Metallic, che ``si alternano a seconda del tipo di materiale''. Lo
\textbf{Specular}, come si è visto con le mappe, determina quanta luce viene riflessa e
l'angolazione a cui viene riflessa. Per creare uno \textbf{specchio} serve un materiale
molto, molto liscio e molto, molto riflettente: Davide abbassa la Roughness e alza la
riflessione, e la \emph{preview} comincia ad assomigliare ``a una parete riflettente'';
con altre cose che si vedranno in seguito ``iniziamo ad avere un riflesso''.

\paragraph{Emissive e Normal.} L'\textbf{Emissive Color} restituisce un colore che è
``effettivamente una luce''. La \textbf{Normal} è l'ultimo parametro utile, e richiede
la texture blu vista prima: si usa sempre un Texture Sample, ma con
\textbf{T\_Bark\_N}, la Normal della corteccia, invece della Diffuse. Davide toglie
tutto il resto e mette a confronto il materiale con e senza Normal: ``vedete la
differenza?'' La Normal aggiunge un dettaglio tridimensionale ``molto semplice da
vedere'', eppure ``la mappa è 2D, non c'è tridimensionalità'': crea solo l'illusione.

\begin{figure}[htbp]
\centering
\begin{tikzpicture}[font=\small]
  \foreach \i/\rr in {0/0, 1/0.25, 2/0.5, 3/0.75, 4/1} {
    \pgfmathsetmacro{\hs}{0.12+0.35*\rr}
    \pgfmathsetmacro{\op}{1-0.85*\rr}
    \shade[ball color=gray!70] (\i*2.4,0) circle (0.8);
    \fill[white, opacity=\op] (\i*2.4-0.28,0.3) circle (\hs);
    \node[font=\footnotesize] at (\i*2.4,-1.2) {Roughness $=\rr$};
  }
  \draw[-{Stealth[length=2.5mm]}, thick] (-0.6,-1.8) -- (10.2,-1.8);
  \node[font=\footnotesize, anchor=north west] at (-0.6,-1.85) {liscio, riflessi netti
    (lo specchio)};
  \node[font=\footnotesize, anchor=north east] at (10.2,-1.85) {ruvido, nessun
    riflesso netto (la roccia)};
\end{tikzpicture}
\caption{La Roughness, da 0 a 1, in uno schema qualitativo: più la superficie è liscia,
più il riflesso della luce è piccolo e netto; più è ruvida, più la luce si sparpaglia
fino a sparire come riflesso. Disegno indicativo, non un rendering di Unreal.}
\label{ue3:fig:roughness}
\end{figure}

\paragraph{Approfondimento: Metallic e Specular nella pratica.} Nel modello PBR di
Unreal i due parametri non hanno lo stesso peso. \textbf{Metallic} è quasi sempre
\emph{binario}: 1 per i metalli, 0 per tutto il resto (legno, pietra, plastica, pelle);
i valori intermedi servono soprattutto per i bordi tra metallo e vernice o per lo
sporco. Lo \textbf{Specular} ha un default di 0,5, che va bene per la grande maggioranza
dei materiali non metallici, e la documentazione di Unreal consiglia di toccarlo poco:
per i metalli viene ignorato, perché il loro riflesso dipende dal Base Color. Lo
specchio della lezione, in termini PBR, è quindi Metallic $=1$, Roughness $=0$, Base
Color chiaro. È il senso dell'``alternarsi'' di cui parla Davide: nei materiali metallici
si lavora con il Metallic, negli altri, se proprio serve, con lo Specular, e quasi
sempre con la Roughness, che è il parametro che cambia di più l'aspetto di una
superficie.

% ---------------------------------------------------------------------
\section{Il Texture Coordinate e il tiling}
\label{ue3:sec:tiling}

``Chiaramente ogni nodo funziona a sé.'' Il \textbf{Texture Coordinate}, il nodo rosso
della prima slide, ha i suoi parametri. La scorciatoia è ``un po' \emph{tricky} da
ricordare'', perché T + clic è già preso dal Texture Sample: è \textbf{U + clic},
``perché il Texture Coordinate lavora sulle UV''.

\paragraph{UTiling e VTiling.} Selezionando il nodo, nei dettagli compaiono
\textbf{UTiling} e \textbf{VTiling}: i parametri ``molto comodi'' con cui si aumenta o
si diminuisce il tiling di un oggetto. Davide collega il Texture Coordinate alle UV
della Normal e della Diffuse, e chiede: se metto UTiling a 2, che cosa succede? ``Vediamo
due volte la texture applicata all'oggetto'', risponde uno studente. ``Sì, però c'è un
\emph{catch}, perché è soltanto sulle U'': la texture si ripete due volte
sull'\textbf{asse orizzontale} e una sola volta su quello verticale. Con anche VTiling a
2, la texture si ripete \textbf{quattro} volte: ``ce n'è una qua, una qua, una qui e una
qua''.

\paragraph{Lo zoom.} Si può fare anche il contrario, per ``zoomare'' sulla texture: con
0,5 la si vede ``la metà delle volte. Anzi, non la metà delle volte, scusate: un
quarto'' (0,5 su entrambi gli assi). E questo non vuol dire che una parte della
texture ``viene mangiata'': viene semplicemente ``applicata in uno spazio differente''.

\begin{figure}[htbp]
\centering
\begin{tikzpicture}[font=\small]
  \newcommand{\uetreTile}[5]{% x0, y0, nu, nv, etichetta
    \begin{scope}[shift={(#1,#2)}]
      \pgfmathsetmacro{\du}{2.2/#3}
      \pgfmathsetmacro{\dv}{2.2/#4}
      \foreach \i in {1,...,#3} \foreach \j in {1,...,#4} {
        \draw[fill=brown!25] ({(\i-1)*\du},{(\j-1)*\dv}) rectangle
          ({\i*\du},{\j*\dv});
        \draw[brown!70!black, thick] ({(\i-1)*\du+0.15*\du},{(\j-1)*\dv+0.2*\dv})
          -- ({(\i-1)*\du+0.85*\du},{(\j-1)*\dv+0.8*\dv});
      }
      \draw[very thick] (0,0) rectangle (2.2,2.2);
      \node[font=\footnotesize, align=center] at (1.1,-0.5) {#5};
    \end{scope}}
  \uetreTile{0}{0}{1}{1}{U = 1, V = 1}
  \uetreTile{3.0}{0}{2}{1}{U = 2, V = 1}
  \uetreTile{6.0}{0}{2}{2}{U = 2, V = 2}
  % zoom: un quarto della texture, ingrandito
  \begin{scope}[shift={(9.0,0)}]
    \draw[fill=brown!25] (0,0) rectangle (2.2,2.2);
    \draw[brown!70!black, thick] (0.66,0.88) -- (2.2,2.2);
    \draw[very thick] (0,0) rectangle (2.2,2.2);
    \node[font=\footnotesize, align=center] at (1.1,-0.5) {U = 0,5, V = 0,5};
  \end{scope}
\end{tikzpicture}
\caption{Il tiling sulla stessa faccia. Con UTiling a 2 la texture si ripete solo in
orizzontale; con entrambi a 2 si ripete quattro volte; con entrambi a 0,5 si vede un
quarto della texture, ingrandito: non è ``mangiata'', è applicata in uno spazio più
grande. Schema ricostruito dalla dimostrazione.}
\label{ue3:fig:tiling}
\end{figure}

\paragraph{Il pavimento.} ``Il tiling, però, va visto.'' Nella mappa c'è un
\textbf{pavimento}, e Davide ne cambia il materiale: lo seleziona, e nei
\textbf{Details}, sotto \textbf{Materials}, al posto della \emph{material instance}
della griglia grigia di prototipazione\footnote{Nella trascrizione ``Materialistance
Prototype Grid Gray'': è l'istanza della griglia grigia dei progetti modello di
Unreal.} trascina il \textbf{master material} appena creato. Il risultato: la corteccia
è ``molto, molto grande'' (``io non so neanche qual è l'utilizzo di questa texture che
sto usando, ci ho preso Bark a caso, ma va bene così''). Davide porta il tiling a 3,
poi a 10: ``iniziamo a vederla più spesso, questa corteccia. Sembrano anche delle dune,
vagamente.'' Il pavimento è probabilmente un cubo allargato, con UV ``molto, molto
piccole'' rispetto alle sue dimensioni: per questo nella \emph{preview} sferica del
materiale la texture si ripete ``diecimila volte'', mentre sul pavimento ``abbiamo un
pavimento che è più dimensionato''.

\paragraph{Una domanda che apre il tema successivo.} Uno studente chiede: si possono
modificare U e V in base all'oggetto su cui è applicata la texture, per esempio 10 su un
pavimento e 0,5 su un altro oggetto? ``Sì, è possibile, non facendolo dallo stesso
materiale. Questo è il momento in cui entrano in gioco le \textbf{Material
Instance}.''

\paragraph{Approfondimento: il tiling in una formula.} Il Texture Coordinate restituisce
per ogni punto la sua coppia $(u, v)$; il Multiply la scala: $(u', v') = (k_u\,u,\;
k_v\,v)$. Una texture è definita sull'intervallo $[0,1]$ e si ripete fuori da lì (è il
comportamento di default, \emph{wrap}): quando $u'$ va da 0 a 2, la texture viene letta
due volte. Con $k < 1$ la faccia ``vede'' solo un pezzo dell'intervallo, e quel pezzo
viene ingrandito. Per questo il tiling dipende dalle UV della mesh: se le UV di un
pavimento enorme occupano comunque tutto il quadrato $[0,1]^2$, una texture con tiling
1 si stende una volta sola su tutta la superficie, ed è ``molto, molto grande''.
Moltiplicare per una costante e usare UTiling e VTiling nel nodo sono due modi di fare
la stessa cosa; il primo ha il vantaggio che la costante può diventare un parametro.

% ---------------------------------------------------------------------
\section{Le Material Instance}
\label{ue3:sec:instance}

Davide torna alle slide (``passiamo un attimo alle Material Instance, visto che ne
stavamo parlando''). Quello costruito finora è un \textbf{master material}: tutto ciò
che sta nel grafo. Se si modifica il grafo, si modifica il materiale. Per avere lo
\textbf{stesso materiale con una piccola differenza} si possono creare due
\textbf{istanze} e cambiare in ciascuna il parametro collegato al Texture Coordinate.

\paragraph{Che cos'è un'istanza.} Una \textbf{Material Instance} è ``un'istanza di un
materiale che noi possiamo utilizzare per modificare quel materiale \textbf{senza
modificare la sorgente}''. Ha diversi vantaggi, ``oltre a essere \emph{bon ton}'', cioè
il modo in cui si dovrebbero applicare i materiali in generale: ``non dovreste prendere
il master material e applicarlo agli oggetti in scena. È sempre consigliabile applicare
un'istanza.''

\paragraph{Perché è più leggera: i tavoli.} L'istanza è ``molto più leggera''. Davide
lo spiega con i \textbf{tavolini} della prima lezione. Tre tavoli in scena con lo
stesso materiale del legno: per ogni tavolo l'engine carica la mesh e carica il
materiale, quindi ``tavolo, tavolo, tavolo'' e ``materiale, materiale, materiale''. Se
invece sui tavoli ci sono tre Material Instance, l'engine carica le tre mesh, ma trova
tre istanze che puntano \textbf{tutte allo stesso master material}, e ``va a
renderizzare un solo materiale''. ``Chiaramente, su tre tavoli non cambia niente. Ora
immaginate che invece di tre tavoli ci siano diecimila oggetti in scena'' e che mille
condividano lo stesso materiale: si capisce perché ``torna comodo avere un'istanza
piuttosto che avere un materiale generale''.

\begin{figure}[htbp]
\centering
\begin{tikzpicture}[
  font=\small,
  ue3tab/.style={draw, thick, minimum width=1.5cm, minimum height=0.7cm,
                 fill=brown!20, align=center},
  ue3mat/.style={draw, thick, rounded corners=3pt, minimum width=1.5cm,
                 minimum height=0.7cm, align=center},
  ue3arr/.style={-{Stealth[length=2mm]}, thick}
]
  % senza istanze
  \node[font=\footnotesize\bfseries] at (2.0,2.2) {come la racconta Davide: master
    material su ogni tavolo};
  \foreach \i in {0,1,2} {
    \node[ue3tab] (t\i) at (\i*2.0,1.2) {tavolo};
    \node[ue3mat, fill=blue!15] (m\i) at (\i*2.0,0) {M\_Legno};
    \draw[ue3arr] (t\i) -- (m\i);
  }
  \node[font=\footnotesize, align=center] at (2.0,-0.8) {tre mesh e ``materiale,
    materiale, materiale''};
  % con istanze
  \begin{scope}[shift={(7.6,0)}]
    \node[font=\footnotesize\bfseries] at (2.0,2.2) {con le Material Instance};
    \foreach \i in {0,1,2} {
      \node[ue3tab] (u\i) at (\i*2.0,1.2) {tavolo};
      \node[ue3mat, fill=orange!15, font=\footnotesize] (i\i) at (\i*2.0,0)
        {MI\_Legno\_\i};
      \draw[ue3arr] (u\i) -- (i\i);
    }
    \node[ue3mat, fill=blue!15] (mm) at (2.0,-1.3) {\textbf{M\_Legno}};
    \foreach \i in {0,1,2} \draw[ue3arr] (i\i) -- (mm);
    \node[font=\footnotesize, align=center] at (2.0,-2.1) {tre mesh, tre istanze, un
      solo materiale da compilare};
  \end{scope}
\end{tikzpicture}
\caption{L'esempio dei tavoli. A sinistra, il master material applicato direttamente;
a destra, tre istanze (che possono avere parametri diversi: un legno più chiaro, un
tiling più fitto) che rimandano tutte allo stesso master material. Schema ricostruito
dalla spiegazione; sul costo reale si veda l'approfondimento.}
\label{ue3:fig:tavoli}
\end{figure}

\paragraph{Come si crea.} ``Come lo facciamo su Unreal?'' Si fa \textbf{tasto destro}
sul materiale nel Content Browser e \textbf{Create Material Instance}, ``qui in
alto''. L'istanza appare \textbf{vuota}: ``non c'è un grafo per modificarlo, perché
chiaramente non posso modificare nessun grafo quando sono qui dentro''. Per avere
qualcosa da cambiare, bisogna estrarre dei \textbf{parametri} dal master material.
Davide torna nel grafo, fa tasto destro sul Texture Sample del colore, \textbf{Convert
to Parameter}, e lo chiama più o meno \emph{Diffuse Color}\footnote{Nella trascrizione
``discolor'': il nome esatto non si ricostruisce, ma è il parametro della texture del
colore.}; lo stesso per l'altro Texture Sample, chiamato \textbf{Normal Map}. Salvato il
materiale, nell'istanza compaiono \textbf{due parametri} selezionabili.

\paragraph{La nomenclatura.} ``Normalmente dovreste chiamarlo'' non \emph{New Material}
ma \textbf{M\_NewMaterial}, e l'istanza \textbf{MI\_NewMaterial}: ``per nomenclatura,
insomma, per ritrovare facilmente tutto quello che ci serve''.

\paragraph{Cambiare la texture dall'istanza.} Davide applica MI\_NewMaterial al
pavimento: è identico a prima. Ma se nell'istanza si cambia il parametro della texture,
da Bark a un'altra, il materiale cambia, senza toccare il grafo. Si potrebbe avere un
master material senza texture, ``delle texture assenti o base'', e tre o quattro
istanze in cui si cambiano le texture: un muro arancione, uno blu, uno rosso e uno verde
sono quattro istanze. ``Quando io renderizzo la scena, non andrò a chiamare cinque
materiali diversi, chiamerò sempre M\_NewMaterial, perché sono tutte istanze dello
stesso materiale.'' ``Non vorrei stare complicando troppo la cosa, però insomma, ditemi
voi se è chiaro il discorso.''

\paragraph{Il Tiling Multiplier.} Ed ecco la risposta alla domanda sul tiling diverso
per ogni oggetto. Nel master material Davide prende un \textbf{Multiply}, ci collega il
Texture Coordinate, collega il Multiply alle UV dei due Texture Sample e attacca alla B
una \textbf{costante}, che converte in parametro con il nome \textbf{Tiling
Multiplier}. Nell'istanza, il Tiling Multiplier portato a 10 dà ``lo stesso risultato di
prima''; ma ora è un parametro. Una copia dell'istanza con il Tiling Multiplier a 5
darebbe un tiling diverso con lo stesso materiale. ``È chiaro il discorso?''
``Perfetto.''

\begin{figure}[htbp]
\centering
\begin{tikzpicture}[
  font=\small,
  ue3node/.style={draw, thick, rounded corners=4pt, align=center,
                  minimum height=0.9cm},
  ue3uv/.style={-{Stealth[length=2mm]}, thick, red!60!black},
  ue3flt/.style={-{Stealth[length=2mm]}, thick, green!50!black},
  ue3col/.style={-{Stealth[length=2mm]}, thick, blue!60!black}
]
  \node[ue3node, fill=red!12, text width=1.9cm] (tc) at (0,0.7)
    {\textbf{TexCoord}\\\footnotesize (U + clic)};
  \node[ue3node, fill=green!8, text width=1.9cm] (tm) at (0,-0.9)
    {\textbf{Tiling}\\\textbf{Multiplier}\\\footnotesize parametro};
  \node[ue3node, fill=gray!10, text width=1.5cm] (mul) at (3.0,0)
    {\textbf{Multiply}};
  \node[ue3node, fill=orange!12, text width=2.4cm] (ts1) at (6.2,1.0)
    {\textbf{Diffuse Color}\\\footnotesize Texture Sample,\\
    \footnotesize parametro};
  \node[ue3node, fill=orange!12, text width=2.4cm] (ts2) at (6.2,-1.0)
    {\textbf{Normal Map}\\\footnotesize Texture Sample,\\\footnotesize parametro};
  \node[ue3node, fill=blue!6, text width=2.2cm, minimum height=2.6cm] (res)
    at (10.0,0) {\textbf{M\_NewMaterial}\\[4pt]\footnotesize Base Color\\[14pt]
    \footnotesize Normal};
  \draw[ue3uv] (tc.east) -- ++(0.4,0) |- ([yshift=2mm]mul.west);
  \draw[ue3flt] (tm.east) -- ++(0.4,0) |- ([yshift=-2mm]mul.west);
  \draw[ue3uv] (mul.east) -- ++(0.5,0) |- node[pos=0.8, above, font=\scriptsize]{UVs}
    (ts1.west);
  \draw[ue3uv] (mul.east) -- ++(0.5,0) |- node[pos=0.8, below, font=\scriptsize]{UVs}
    (ts2.west);
  \draw[ue3col] (ts1.east) -- ++(0.6,0) |- ([yshift=6mm]res.west);
  \draw[ue3col] (ts2.east) -- ++(0.6,0) |- ([yshift=-6mm]res.west);
  % istanze
  \node[ue3node, fill=yellow!15, text width=3.6cm, font=\footnotesize] (mi1)
    at (3.0,-2.9) {\textbf{MI\_NewMaterial}\\Tiling Multiplier = 10\\Diffuse = T\_Bark};
  \node[ue3node, fill=yellow!15, text width=3.6cm, font=\footnotesize] (mi2)
    at (8.0,-2.9) {\textbf{MI\_NewMaterial\_2}\\Tiling Multiplier = 5\\
    Diffuse = un'altra texture};
  \draw[dashed, thick] (mi1.north) -- ++(0,0.4) -| (res.south);
  \draw[dashed, thick] (mi2.north) -- (mi2.north |- res.south);
\end{tikzpicture}
\caption{Il master material con tre parametri (le due texture e il Tiling Multiplier)
e due istanze che li impostano in modo diverso. Le istanze non hanno un grafo: hanno
solo le caselle dei parametri. Ricostruzione dalla dimostrazione.}
\label{ue3:fig:instance}
\end{figure}

\paragraph{Il vantaggio principale.} Ricapitolando dalle slide: le Material Instance
permettono di \textbf{personalizzare velocemente} le superfici ``senza andare a
modificare il grafico che avete fatto'', e le modifiche si vedono
\textbf{immediatamente}, ``senza dover aspettare la ricompilazione del materiale''.

\paragraph{Approfondimento: che cosa si risparmia davvero.} La spiegazione dei tavoli
va presa come un'immagine. Il costo principale di un materiale non è ``caricarlo'' una
volta per oggetto, ma \textbf{compilarlo}: ogni master material diventa uno o più
\emph{shader}, programmi per la scheda video che Unreal deve compilare (e che fanno
comparire, nella prima apertura di un progetto, il famoso contatore ``Compiling
Shaders''). Un'istanza \textbf{non} genera un nuovo shader: riusa quello del genitore e
cambia solo i valori dei parametri. Cento istanze dello stesso master material sono
quindi un solo shader con cento ``schede'' di valori; cento materiali diversi sono cento
shader. Per questo cambiare un parametro in un'istanza è istantaneo, mentre cambiare un
nodo nel master material richiede la ricompilazione, come dice la slide. Le istanze del
Content Browser si chiamano \textbf{Material Instance Constant}; esistono anche le
\textbf{Material Instance Dynamic}, create durante il gioco da un Blueprint (nodo
\emph{Create Dynamic Material Instance}) per cambiare i parametri a \emph{runtime}, per
esempio con \emph{Set Scalar Parameter Value}: un materiale che si accende quando la
torretta di Igor ha un bersaglio nel mirino.

% ---------------------------------------------------------------------
\section{Le Material Function}
\label{ue3:sec:function}

``Per chiudere il discorso materiali'', l'ultima voce delle slide: le \textbf{Material
Function}, che oggi non si useranno. Sono funzioni che raccolgono ``una serie di calcoli
un po' più complessi'': invece di rifare ogni volta dei calcoli che si fanno spesso, li
si mette dentro una funzione e si usa la funzione ``a vostra discrezione''. ``È come se
fossero delle \textbf{parentesi}: voi prendete queste parentesi, le spostate in giro. È
un pacchetto che potete portarvi in giro.'' Servono a velocizzare lo sviluppo quando ci
sono cose da ripetere, e possono essere usate e condivise ``fra membri del team, fra
materiali diversi''.

\paragraph{Approfondimento: Material Function e Material Instance.} Le due cose
risolvono problemi diversi. L'istanza riusa \emph{un intero materiale} cambiandone i
valori; la funzione riusa \emph{un pezzo di grafo} dentro materiali diversi. Unreal
stesso è pieno di Material Function già pronte (si trovano nel menu dei nodi con
l'icona della funzione), e i nodi World Aligned Texture e World Aligned Normal che
compaiono alla fine della lezione (sezione~\ref{ue3:sec:world}) sono, appunto, Material
Function del motore.

% ---------------------------------------------------------------------
\section{Il Lerp: interpolare}
\label{ue3:sec:lerp}

Restano ``un paio di nodi che possono essere carini''. Davide ricapitola quelli visti
(la costante singola, il colore, il Texture Sample, le operazioni, il Texture
Coordinate) e ne aggiunge due: il \textbf{Lerp} e il \textbf{Panner}. ``Se abbiamo
capito la regola degli shortcut, questi saranno quale lettera più clic sinistro?''
``L.'' ``Esatto'': \textbf{L + clic} per il Lerp, \textbf{P + clic} per il Panner. ``Volevo
essere sicuro che era arrivata la cosa.''

\paragraph{Due colori.} Il \textbf{Lerp} si occupa dell'\textbf{interpolazione
lineare} (\emph{linear interpolation}). Per farlo vedere Davide usa due colori, ``una
roba che non farete mai, non con due colori perlomeno'': un \textbf{blu} nella
\textbf{A} e un \textbf{rosso} nella \textbf{B}. Il Lerp, collegato al Base Color, dà il
\textbf{viola}. Perché è inutile farlo con i colori? ``Perché potevo cercarmi il viola
direttamente dal colore.'' Ma mostra come funziona l'interpolazione. Il Lerp si può fare
anche con i numeri, ``per fare operazioni matematiche'', e anche con le texture.

\paragraph{L'alfa.} L'operazione è ``gestita interamente dall'\textbf{alfa}'', il terzo
ingresso. L'alfa va \textbf{da 0 a 1} e decide quale colore prevale, con A che vale 0 e
B che vale 1: con l'alfa a 0 si ha il blu, a 1 il rosso, e in mezzo ``una percentuale di
quel colore'', per esempio 0,7. Modificando l'alfa si va sempre più verso il rosso, o
sempre più verso il blu.

\paragraph{Una texture come alfa.} Ma l'alfa non deve essere un numero fisso. Davide
prende ``il nostro fidato Bark'' e collega \textbf{un suo canale} all'alfa: ``vi ricordate
il discorso dei canali sulle texture?'' Ora quella sequenza di 0 e di 1, di bianco e di
nero, ``è l'\textbf{arbitro} nel decidere dove ci sarà più A e dove ci sarà più B'':
nelle zone più nere il Lerp mette più blu, in quelle più bianche più rosso. Questa
texture è ``estremamente bianca'', quindi il risultato tende al rosso; con altre
texture, o lavorando su questa, si otterrebbero risultati diversi. E collegando il
Multiply delle coordinate anche al Texture Sample dell'alfa, si ottiene un tiling
anche dell'alfa, cioè di dove va il blu e dove il rosso. ``Chiaro questo processo?''
``Sì.''

\begin{figure}[htbp]
\centering
\begin{tikzpicture}[
  font=\small,
  ue3node/.style={draw, thick, rounded corners=4pt, align=center,
                  minimum height=0.9cm},
  ue3col/.style={-{Stealth[length=2mm]}, thick, blue!60!black},
  ue3flt/.style={-{Stealth[length=2mm]}, thick, green!50!black}
]
  \node[ue3node, fill=blue!30, text width=1.9cm] (a) at (0,1.4) {\textbf{blu}\\
    \footnotesize Constant3Vector};
  \node[ue3node, fill=red!30, text width=1.9cm] (b) at (0,0) {\textbf{rosso}\\
    \footnotesize Constant3Vector};
  \node[ue3node, fill=gray!10, text width=1.9cm] (al) at (0,-1.5)
    {\textbf{Alpha}\\\footnotesize 0,7 oppure\\\footnotesize un canale di T\_Bark};
  \node[ue3node, fill=gray!10, text width=1.6cm, minimum height=2.2cm] (lerp)
    at (3.6,0) {\textbf{Lerp}\\\footnotesize (L + clic)\\[4pt]\footnotesize A\\
    \footnotesize B\\\footnotesize Alpha};
  \node[ue3node, fill=violet!30, text width=2.0cm] (res) at (7.0,0)
    {\textbf{Base Color}\\\footnotesize viola, più o\\\footnotesize meno rosso};
  \draw[ue3col] (a.east) -- ++(0.6,0) |- ([yshift=5mm]lerp.west);
  \draw[ue3col] (b.east) -- ([yshift=0mm]lerp.west);
  \draw[ue3flt] (al.east) -- ++(0.6,0) |- ([yshift=-5mm]lerp.west);
  \draw[ue3col] (lerp) -- (res);
  % la scala dell'alfa
  \begin{scope}[shift={(9.4,-0.6)}]
    \shade[left color=blue!60, right color=red!60] (0,0) rectangle (3.4,0.6);
    \draw[thick] (0,0) rectangle (3.4,0.6);
    \node[font=\footnotesize] at (0,-0.3) {0};
    \node[font=\footnotesize] at (3.4,-0.3) {1};
    \node[font=\footnotesize] at (1.7,-0.3) {Alpha};
    \node[font=\footnotesize] at (0,0.95) {A};
    \node[font=\footnotesize] at (3.4,0.95) {B};
    \draw[thick] (2.38,-0.05) -- (2.38,0.65);
    \node[font=\scriptsize] at (2.38,0.95) {0,7};
  \end{scope}
\end{tikzpicture}
\caption{Il Lerp della dimostrazione: con l'alfa a 0 esce A (il blu), a 1 esce B (il
rosso), in mezzo una miscela. Se l'alfa è un canale di una texture, ogni pixel ha la sua
miscela: blu dove la texture è nera, rosso dove è bianca.}
\label{ue3:fig:lerp}
\end{figure}

\paragraph{L'esercizio della pausa.} ``Provateci voi.'' Durante la pausa, chi vuole può
prendere \textbf{due texture}, collegarle a un Lerp e, giocando con l'alfa, vedere quale
parte delle due texture viene presa di più e quale di meno.

\paragraph{Approfondimento: la formula del Lerp.} L'interpolazione lineare tra A e B
con peso $\alpha$ è
\[
  \operatorname{lerp}(A, B, \alpha) = A\,(1-\alpha) + B\,\alpha .
\]
Con colori, il calcolo si fa canale per canale. Blu $=(0,0,1)$ e rosso $=(1,0,0)$ con
$\alpha = 0{,}5$ danno $(0{,}5;\,0;\,0{,}5)$, il viola della lezione; con $\alpha =
0{,}7$ danno $(0{,}7;\,0;\,0{,}3)$, un viola più rosso. Quando l'alfa è una texture, la
formula si applica pixel per pixel, ed è così che si mescolano due materiali su una
stessa superficie: erba e terra su un terreno, vernice e ruggine su un metallo, con una
maschera in bianco e nero a fare da ``arbitro''.

% ---------------------------------------------------------------------
\section{Il Panner: far muovere le texture}
\label{ue3:sec:panner}

Il \textbf{Panner}\footnote{Nella trascrizione ``panel'', ``planner'' e ``banner''.}
``si occupa di far muovere le texture''. Davide riprende il Base Color, stacca il resto,
e collega il Panner alle \textbf{UV} del Texture Sample, ``perché lui lavora sulle UV
della texture''. Selezionando il nodo (o direttamente sul nodo) si vedono
\textbf{Speed X} e \textbf{Speed Y}. Con Speed X a 1 la \emph{preview} comincia a
muoversi sull'asse X, ``anche un po' troppo veloce per i miei gusti''; con 0,3 in X e
0,1 in Y si ha un Panner che si muove su entrambi gli assi.

\paragraph{A che cosa serve.} ``Che materiale potremmo volerci fare con questo nodo?''
``Dell'acqua.'' ``Dell'acqua. Qualsiasi fluido, qualsiasi cosa che si muove.'' Il Panner
ha il parametro di velocità e l'ingresso delle \textbf{coordinate}: per combinarlo con il
tiling si collega il Texture Coordinate (o il Multiply) all'ingresso Coordinate. C'è
poi l'ingresso \textbf{Time}, con cui si può decidere ``quando e come deve muoversi'';
altrimenti, di default, usa il \emph{game time}, ``cioè è sempre in movimento''.

\paragraph{Il master material cambia, l'istanza segue.} Mentre il Panner lavora, il
pavimento nella scena \textbf{si muove}: sul pavimento c'è la Material Instance, e
``nonostante stiamo facendo delle modifiche al materiale originale e non alla Material
Instance, la Material Instance continua a modificarsi''. È il rovescio del rapporto tra
master e istanze: quello che cambia nel genitore cambia in tutte le istanze.

\paragraph{Approfondimento: il Panner in una formula.} Il Panner non fa altro che
sommare alle coordinate un tempo moltiplicato per una velocità:
\[
  (u', v') = (u + s_x\,t,\; v + s_y\,t),
\]
dove $t$ è il tempo in secondi e $s_x, s_y$ le velocità in ``texture al secondo'': con
Speed X a 1 la texture scorre di una larghezza intera ogni secondo, per questo era
troppo veloce. Siccome la texture si ripete, lo scorrimento è infinito e senza salti,
purché la texture sia \emph{tileable}, cioè con i bordi che combaciano. Due Panner con
velocità e tiling diversi sommati tra loro sono il trucco classico per l'acqua e per la
lava: il movimento non si ripete mai uguale.

% ---------------------------------------------------------------------
\section{La pausa: il progetto che non si trova, e Git}
\label{ue3:sec:pausa}

``Visto che comunque ci siamo fatti un'oretta e mezza'', Davide chiude per ora il
discorso materiali e propone una pausa di una decina di minuti. Al ritorno: l'import,
le mesh, e ``una cosina con la torretta che c'è nei file che vi abbiamo passato'', così
che Igor, quando torna, potrà farci ``un piccolo sistema di puntamento''.

\paragraph{Il progetto che Unreal non vede.} Al rientro, in attesa degli altri, una
studentessa racconta il suo problema: ha scaricato e decompresso la cartella del Drive,
ma aprendo Unreal il progetto non compare, e deve aprirlo ``a mano'' da un'altra parte.
Davide spiega che dipende probabilmente dal percorso in cui è salvata la cartella: se la
si mette dove sono salvati gli altri progetti, Unreal dovrebbe trovarla; in ogni caso,
aprire il progetto dal file eseguibile nella cartella funziona. ``Se non è un problema
che ti impedisce di accedere al progetto, fa niente.''

\paragraph{Come si lavora insieme a distanza?} Uno studente, programmatore, fa una
domanda ``fuori dall'argomento'': i programmatori usano \textbf{Git}, ma i file di un
progetto di Unreal sembrano troppo pesanti per lavorarci in contemporanea da lontano.
Davide conferma che \textbf{Git}, ``o diciamo repository di sorta'', è ``la cosa
principale che utilizzerai''. Con \textbf{Unity} ha sempre avuto meno problemi, perché
Unreal ``va gestito in maniera un po' più delicata'', ma entrambi supportano senza
problemi repository e \emph{source control}. Esistono anche altri metodi, che Davide non
ha mai provato: per Unreal, ``una cosa che io sento da una vita'', un modo per lavorare
\textbf{contemporaneamente} sulla stessa cosa, tramite \textbf{VPN} o in \textbf{LAN} e
con dei \emph{plugin} specifici, in modo che due colleghi modifichino attivamente lo
stesso livello. Con Git invece ``ognuno lavora sul suo e poi si \emph{merge} tutto
insieme''. ``Git è quello che ho utilizzato sempre, è il più comodo, quindi dovresti
trovarti abbastanza a casa.''

\paragraph{Approfondimento: Git LFS, Perforce e Multi-User Editing.} I timori dello
studente sono fondati. Git nasce per file di testo, e gli \emph{asset} di Unreal
(\texttt{.uasset}, \texttt{.umap}) sono file \textbf{binari}, spesso grandi, che Git non
sa fondere: se due persone modificano lo stesso livello, una delle due modifiche va
persa. Per questo con Unreal si usa \textbf{Git LFS} (\emph{Large File Storage}), che
tiene i file pesanti fuori dalla cronologia normale e permette di \textbf{bloccare}
(\emph{lock}) un file mentre lo si modifica. Nell'industria, per i progetti Unreal, è
molto diffuso \textbf{Perforce} (Helix Core), il sistema che Epic stessa usa e che
l'editor supporta direttamente, proprio per la gestione dei file binari e dei blocchi.
Il sistema ``sentito da una vita'' di cui parla Davide esiste: è il plugin
\textbf{Multi-User Editing}, con cui più editor connessi in rete (in LAN o via VPN)
condividono una stessa sessione e vedono le modifiche degli altri in tempo reale; nato
per la produzione virtuale nel cinema, si può usare anche nello sviluppo di giochi, ma
non sostituisce il \emph{source control}, a cui le modifiche vanno comunque consegnate
alla fine.

% ---------------------------------------------------------------------
\section{L'import e la torretta}
\label{ue3:sec:import}

Passato ``il nostro quarto d'ora'', si riparte: si importano le prime mesh, la
\textbf{torretta}, ``le tre parti della torretta, perché l'abbiamo divisa in tre
parti''.

\paragraph{La finestra di import, sulla slide.} Quando si importa una mesh compare una
finestra (``che può essere più larga, può essere più stretta'') piena di parametri. Le
cose principali che si possono fare da lì:
\begin{itemize}
  \item \textbf{generare automaticamente le collisioni}, se non sono presenti nel file
        FBX, perché magari gli oggetti non sono stati creati con le collisioni;
  \item \textbf{modificare il transform} dell'oggetto importato ``proprio di base'':
        ruotarlo sul suo \textbf{pivot}, traslarlo. ``Avete già parlato di cos'è il
        pivot con Igor, giusto? Il fulcro che vi permette di traslare e ruotare gli
        oggetti all'interno della scena.'' Con la funzione di \textbf{reimport}, si
        vedrà, questa possibilità diventa molto utile;
  \item adattare le \textbf{unità} e gli \textbf{assi} tra il software di origine e
        Unreal: \textbf{convertire la scena} dalle coordinate FBX a quelle di Unreal,
        \textbf{convertire le unità} (Unreal lavora in centimetri), oppure usare
        \textbf{Force Front X Axis}, che cambia l'asse che Unreal considera come
        ``fronte'' dell'oggetto. Davide l'ha usato poche volte: una volta il modello di
        un \emph{character}, importato normalmente, ``veniva rivoltato, come se il
        dentro fosse fuori'', e cambiare l'asse frontale gli ha restituito la mesh
        com'era stata esportata (``mi pare fosse Blender'');
  \item decidere se importare anche i \textbf{materiali} e le \textbf{texture}
        collegati alla mesh, e \textbf{dove} cercarli: in locale, in un'altra cartella,
        e così via.
\end{itemize}

\paragraph{Il pulsante Import.} Prima di importare, Davide spegne il Panner (``facciamo
che smetto di far muovere il pavimento'') e il pavimento torna normale. Nel Content
Browser c'è ``questo pulsantone con \textbf{Import}'', che apre l'esplora risorse del
computer per cercare i file. Nella cartella del materiale didattico c'è una cartella
\textbf{torretta}, con tre \textbf{static mesh}: \textbf{SM\_Turret\_Base},
\textbf{SM\_Turret\_Cannon} e \textbf{SM\_Turret\_Pivot}\footnote{Nella trascrizione
``Turret Base, Turret Cannon e Turret Pivo''; il prefisso SM\_ è detto a voce.}.

\paragraph{La finestra di import, dal vivo.} Importandole si apre la finestra (``nel
mio caso è così enorme, ma perché io tendo a usarla così''), e si ritrovano le opzioni
della slide: la possibilità di modificare traslazione e rotazione, ``anche
singolarmente sui vari assi'' (``questo è qualcosa che faremo più tardi''); i
\textbf{LOD} (``decidere il gruppo dei LOD''); le collisioni, se e come importarle. Se si
importa una \textbf{Skeletal Mesh}\footnote{Nella trascrizione ``scheda dal Mesh'' e
``Scheletal Mesh'', come nella seconda lezione.}, come un personaggio, la finestra
permette di collegarla a uno \textbf{scheletro} esistente, di generarne uno, oppure di
importare solo delle \textbf{animazioni}: ``faremo una lezione specifica sulle
animazioni, quindi per il momento glissiamo''. In alto ci sono anche le sezioni per
andare più veloci, tra cui \textbf{Import Materials} e \textbf{Import Textures}, con la
\textbf{Search Location} in cui cercarli.

\paragraph{Niente materiali e texture, per ora.} La consegna di Davide: \textbf{togliere
la spunta} a Import Materials e Import Textures, e fare \textbf{Import}. Unreal importa i
tre oggetti (Davide elimina quelli della prova precedente e salva). Perché togliere le
spunte? Se si fossero importati gli oggetti con materiali e texture, Unreal avrebbe
importato automaticamente tutte le texture, creato i materiali adatti e li avrebbe
anche \textbf{applicati}, ``in modo da averli già precisi sulla torretta''. Ed è
proprio quello che si farà, ma \textbf{a mano}, ``così vediamo effettivamente anche il
processo''. ``Ci sono stati problemi? Tutto a posto con la torretta?'' ``Sì, nessun
problema.''

\begin{table}[htbp]
\centering
\small
\begin{tabular}{@{}p{4.2cm}p{8.8cm}@{}}
\toprule
\textbf{Opzione} & \textbf{A che cosa serve} \\
\midrule
Generate Missing Collision & genera le collisioni se il file FBX non le ha \\
Import Translation / Rotation / Uniform Scale & sposta, ruota, scala l'oggetto ``alla
  radice'': cambia dove sta il pivot rispetto alla geometria \\
LOD Group & sceglie un gruppo di livelli di dettaglio (\emph{Level Of Detail}) \\
Convert Scene, Convert Scene Unit & adattano assi e unità di misura del file FBX a
  quelli di Unreal (centimetri) \\
Force Front X Axis & cambia l'asse considerato come ``fronte'' \\
Skeleton / Import Animations & per le Skeletal Mesh: scheletro esistente o nuovo, solo
  animazioni \\
Import Materials, Import Textures & importano e applicano materiali e texture
  collegati \\
Search Location & dove cercare materiali e texture da importare \\
\bottomrule
\end{tabular}
\caption{Le opzioni della finestra di import citate a lezione, con i nomi della
finestra di Unreal Engine 5 (FBX Import Options; il raggruppamento può variare tra le
versioni). Per la torretta: Import Materials e Import Textures senza spunta.}
\label{ue3:tab:import}
\end{table}

% ---------------------------------------------------------------------
\section{Che cos'è una mesh}
\label{ue3:sec:mesh}

Mentre la torretta è importata, Davide torna sulle slide per ``le ultime cose sulle
mesh''. Sulla slide c'è un \textbf{vaso} in \emph{wireframe} e, sotto, una
semplificazione. Una mesh è ``un insieme di \textbf{vertici}, \textbf{spigoli} e
\textbf{facce} che vanno a creare un oggetto in tre dimensioni''.
\begin{itemize}
  \item I \textbf{vertici} sono i \textbf{punti} che costituiscono gli angoli
        dell'oggetto, punti nello spazio.
  \item Gli \textbf{spigoli} (\emph{edge}) sono i \textbf{segmenti} che collegano due
        vertici.
  \item Le \textbf{facce} sono le \textbf{superfici piatte} definite da più segmenti:
        sulla slide, ``tutti i triangoli''.
\end{itemize}
Il \textbf{triangolo} è ``la forma più comoda che esiste a livello di calcolo'', e
tendenzialmente tutti gli oggetti saranno composti da ``un insieme di triangolini''. Più
triangoli insieme, più facce, formano dei \textbf{poligoni}; i poligoni combinati formano
delle \textbf{superfici}; e la superficie determina la mesh. Tutto questo si può vedere
anche nell'engine, con la visualizzazione \textbf{wireframe}: l'oggetto non ``come
appare'', ma ``come griglia, come insieme di tutti i poligoni, vertici e spigoli''.

\begin{figure}[htbp]
\centering
\begin{tikzpicture}[font=\small]
  \coordinate (A) at (0,0);
  \coordinate (B) at (2.6,0.3);
  \coordinate (C) at (3.2,2.2);
  \coordinate (D) at (0.6,2.6);
  \coordinate (E) at (1.7,1.3);
  \fill[orange!15] (A) -- (B) -- (E) -- cycle;
  \fill[orange!30] (B) -- (C) -- (E) -- cycle;
  \fill[orange!20] (C) -- (D) -- (E) -- cycle;
  \fill[orange!10] (D) -- (A) -- (E) -- cycle;
  \draw[very thick] (A) -- (B) -- (C) -- (D) -- cycle;
  \draw[very thick] (A) -- (E) (B) -- (E) (C) -- (E) (D) -- (E);
  \foreach \p in {A,B,C,D,E} \fill (\p) circle (2.5pt);
  \draw[-{Stealth[length=2mm]}] (4.3,2.8) node[right, font=\footnotesize]
    {\textbf{vertice}: un punto} -- (C);
  \draw[-{Stealth[length=2mm]}] (4.3,1.6) node[right, font=\footnotesize]
    {\textbf{spigolo} (\emph{edge}): un segmento tra due vertici} -- ($(B)!0.5!(C)$);
  \draw[-{Stealth[length=2mm]}] (4.3,0.4) node[right, font=\footnotesize]
    {\textbf{faccia}: un triangolo} -- (1.45,0.55);
  \node[font=\footnotesize, align=center, text width=4cm] at (1.6,-0.8)
    {4 facce triangolari\\formano un poligono};
\end{tikzpicture}
\caption{Vertici, spigoli e facce, come nella semplificazione della slide. Il triangolo
è l'unità di calcolo: anche le facce a quattro lati dei programmi di modellazione,
quando arrivano in Unreal, vengono divise in triangoli.}
\label{ue3:fig:mesh}
\end{figure}

% ---------------------------------------------------------------------
\section{Dimensioni, unità di misura e formati}
\label{ue3:sec:regole}

Quando si porta una mesh in Unreal, mentre la si produce o mentre la si cerca, bisogna
rispettare alcune cose.

\paragraph{Le dimensioni.} La cosa più comune: si decide di fare ``una mappa molto
carina'' su Unreal, si cercano mesh su internet o su \textbf{Fab} (``su Fab questo
problema c'è meno, però utilizzando siti esterni può capitare''), e la mesh importata è
\textbf{gigantesca} o \textbf{incredibilmente piccola}, perché chi ha fatto l'export, o
chi ha fatto l'import, non ha rispettato ``le misure di adattamento e transizione da un
software all'altro''. Unreal usa come unità di misura i \textbf{centimetri}; Blender,
per esempio, ``usa altra roba''. Nel passaggio da Blender a Unreal bisogna controllare
che la conversione delle dimensioni sia corretta.

\paragraph{FBX.} ``Altra cosa molto importante'': Unreal ``tendenzialmente legge
soltanto \textbf{FBX}''. Sono supportati anche altri formati, come gli \textbf{OBJ}, ``ma
non assicurano un corretto funzionamento all'interno dell'engine''. Il consiglio:
``importate sempre unicamente FBX''.

\paragraph{I formati delle texture.} Per le texture i formati supportati sono molti
(sulla slide c'è un elenco), ma nella pratica si usano soprattutto \textbf{JPEG} e
\textbf{PNG}, ``i più comuni, a meno che non state facendo qualcosa di molto
specifico''.

\paragraph{Approfondimento: Blender, i metri e il fattore 100.} Blender lavora di default
in metri, con una ``unità'' uguale a un metro; Unreal lavora in centimetri. Un cubo di
lato 2 in Blender è quindi un cubo di 2~m, cioè 200 unità di Unreal: se l'export non
converte, l'oggetto arriva cento volte troppo piccolo (o, con altre impostazioni, cento
volte troppo grande). La soluzione standard è nell'esportatore FBX di Blender
(l'opzione \emph{Apply Scalings} su \emph{FBX Units Scale}) o, lato Unreal, nella
scala uniforme e nella conversione delle unità della finestra di import. Un controllo
rapido: un personaggio umano in Unreal è alto circa 180 unità, come il manichino dei
progetti modello; se la mesh appena importata è molto diversa, c'è un problema di scala.
Oltre a FBX, le versioni recenti di Unreal Engine 5 importano anche \textbf{glTF} e
\textbf{USD}, ma FBX resta il formato più collaudato, come dice Davide.

% ---------------------------------------------------------------------
\section{Il pivot del cannone e il Reimport}
\label{ue3:sec:pivot}

Davide apre le tre mesh (la base, la torretta e il pivot) e le piazza in scena. Sono
state esportate con il \textbf{pivot nello stesso punto}: per sovrapporle senza
spostarle a mano basta \textbf{copiare la X e la Y} di un pezzo negli altri. Ed ecco la
torretta, ``a parte il fatto che non c'ha il materiale''.

\paragraph{I problemi della torretta.} Ma ``sta torretta ha dei problemi: quali sono?''
Il primo si vede: per qualche motivo \textbf{non è allineata}, ed è ``scomodamente''
ruotata: trovare la rotazione giusta sarà scomodo. Il secondo è più sottile. Si
supponga di mettere i tre pezzi insieme in un \textbf{Blueprint}. Che cosa deve fare una
torretta? Deve \textbf{tracciare dei bersagli}: ruotare su sé stessa, ma anche andare su
e giù, se il bersaglio è in alto o in basso. E qui c'è il problema: il \textbf{pivot del
cannone} non è in una posizione comoda per la rotazione verticale, perché sta
\textbf{in basso}, in fondo. Ruotare sul piano orizzontale non dà problemi, ma per
alzare e abbassare la canna il pivot dovrebbe essere ``più o meno all'altezza di questo
\textbf{bullone}'', il fulcro, e da lì il cannone dovrebbe essere manipolato e ruotato
(figura~\ref{ue3:fig:pivot}).

\begin{figure}[htbp]
\centering
\begin{tikzpicture}[font=\small,
  ue3arr/.style={-{Stealth[length=2.5mm]}, thick}]
  \foreach \xs/\titolo/\px/\py/\ang in {0/{pivot in basso (com'era)}/1.25/0.0/0,
                                         6.6/{pivot nel fulcro (dopo il Reimport)}/1.25/1.75/0} {
    \begin{scope}[shift={(\xs,0)}]
      % base
      \draw[thick, fill=orange!35] (0.3,0) -- (2.2,0) -- (1.9,0.5) -- (0.6,0.5) -- cycle;
      % pivot (supporto)
      \draw[thick, fill=gray!30] (0.8,0.5) rectangle (1.7,1.5);
      % cannone
      \draw[thick, fill=orange!20] (0.7,1.5) rectangle (1.8,2.1);
      \draw[thick, fill=orange!20] (1.8,1.6) rectangle (4.0,1.95);
      % bullone
      \fill[black] (1.25,1.75) circle (2pt);
      % il pivot del cannone
      \draw[red, very thick] (\px-0.18,\py) -- (\px+0.18,\py) (\px,\py-0.18) -- (\px,\py+0.18);
      \draw[red, thick] (\px,\py) circle (0.18);
      \node[font=\footnotesize, align=center] at (2.0,-0.6) {\titolo};
    \end{scope}
  }
  % archi di rotazione
  \draw[ue3arr, red] (3.4,0.5) arc[start angle=-20, end angle=25, radius=2.6];
  \node[font=\scriptsize, red, align=center] at (4.6,0.55) {la canna ruota\\attorno al
    piede};
  \draw[ue3arr, red] (6.6+3.6,1.0) arc[start angle=-30, end angle=30, radius=1.6];
  \node[font=\scriptsize, red, align=center] at (6.6+4.75,1.75) {la canna ruota\\
    attorno al bullone};
  % assi
  \draw[ue3arr, blue] (-0.6,0) -- (-0.6,1.2) node[above, font=\scriptsize]{Z};
\end{tikzpicture}
\caption{Il problema del pivot del cannone, visto di lato (schema non in scala). Con il
pivot (il mirino rosso) in fondo, alzare la canna la farebbe girare attorno al piede,
staccandola dal supporto; con il pivot nel fulcro, all'altezza del bullone, la canna
si alza e si abbassa come una vera torretta. La rotazione orizzontale, attorno all'asse
Z, funziona invece in entrambi i casi.}
\label{ue3:fig:pivot}
\end{figure}

\paragraph{Come si sposta un pivot.} ``In realtà possiamo farlo in tutta una miriade di
modi.'' Con la \textbf{Modeling Mode} si potrebbe rimodellare l'oggetto. Ma lo si può
fare anche \textbf{nell'import}: si fa tasto destro sulla mesh e \textbf{Reimport Mesh
With Dialog}, che riapre la finestra di import vista prima. Unreal importa lo stesso
oggetto e lo \textbf{sovrascrive}, ma permette di cambiare come verrà posizionato:
l'\textbf{offset della rotazione} e l'\textbf{offset della traslazione}.

\paragraph{I numeri.} Davide se li era scritti, ma non li ricorda, e li ricava. Per
sapere quale asse toccare, guarda il \emph{gizmo}: il cannone deve essere allineato con
il resto della torretta (e lo stesso il pivot), e il pivot del cannone deve stare non
``qui giù in fondo'', ma allineato con il fulcro, sull'\textbf{asse Z}, ``perché è
questo blu''. Per questa torretta, ``quindi anche per voi'', i parametri sono:
\begin{itemize}
  \item per il \textbf{cannone}: traslazione Z $= -135$, che reimporta il cannone ``più
        attaccato al pivot'', e rotazione $= 113$, che lo riallinea\footnote{La
        trascrizione dice prima ``113 per la traslazione'' e poi si corregge: ``e 113
        sulla rotazione''. L'asse della rotazione non è detto; trattandosi di
        riallineare la torretta nel piano orizzontale, è con ogni probabilità la
        rotazione attorno a Z (il \emph{yaw}).};
  \item per il \textbf{pivot}: ``questa volta non serve modificare la traslazione'', ma
        solo la rotazione, sempre 113.
\end{itemize}
Dopo il reimport il cannone ``sta sotto'' (la sua origine ora è nel fulcro, quindi
nella posizione di prima risulta più basso) e va rimesso a posto; ma ora, ruotandolo, si
alza e si abbassa nel modo giusto. ``E chiaramente poi questo farete più avanti, con dei
Blueprint.'' Si salva tutto.

\paragraph{Approfondimento: perché il pivot conta così tanto.} In Unreal ogni rotazione
e ogni scala di un componente avvengono \textbf{attorno alla sua origine}, cioè al
pivot della mesh. Un Blueprint di torretta, di solito, ha la base come componente
radice, il supporto come figlio che ruota attorno a Z (il \emph{yaw}, per seguire il
bersaglio a destra e a sinistra) e il cannone come figlio del supporto che ruota attorno
a Y (il \emph{pitch}, su e giù). Per questo ogni pezzo deve avere l'origine esattamente
nel punto in cui gira nella realtà: se l'origine è sbagliata, nessun nodo di rotazione
potrà correggerla senza calcoli aggiuntivi. Il ``piccolo sistema di puntamento'' che
Igor farà con questa torretta (con nodi come \emph{Find Look at Rotation}, che calcola
la rotazione necessaria per guardare un punto) si appoggia interamente a questo lavoro
preparatorio.

% ---------------------------------------------------------------------
\section{Mettere ordine: le cartelle}
\label{ue3:sec:cartelle}

Prima di importare le \textbf{texture}, ``per dare un po' di ordine'', Davide crea nel
Content Browser una nuova cartella, \textbf{Torretta}, ci sposta dentro tutto
(``muoviamo, così è un po' più pulito''), e dentro crea un'ulteriore cartella,
\textbf{Texture}, in cui importare le texture della torretta. ``Questo non è
obbligatorio, è solo una questione di \emph{best practice}, cioè è come il discorso
delle Material Instance. Non vi vengo a trovare a casa se non lo fate, però vi risulta
tutto più comodo.''

% ---------------------------------------------------------------------
\section{Lo Static Mesh Editor e le collisioni}
\label{ue3:sec:sme}

Con un doppio clic su una mesh si apre lo \textbf{Static Mesh Editor}: ``andiamoci a
vedere effettivamente un po' qui dentro cosa possiamo vedere delle mesh''.

\paragraph{Gli slot dei materiali.} A destra, nei \textbf{Material Slots}, queste mesh
hanno \textbf{due slot}: ``non è la normalità''. (Una delle tre ne ha uno solo.) Il modo
in cui la torretta è stata suddivisa fa sì che, per rendere bene i materiali, servano
due slot: la base e il cannone condividono in parte il materiale del \textbf{pivot}. Nei
due slot si inseriranno i materiali più avanti. Altre impostazioni (Davide accenna a
quelle sui LOD e sulla \emph{Nanite}\footnote{La trascrizione dice ``per quanto riguarda
la night'': la ricostruzione come Nanite, la tecnologia di geometria virtualizzata di
Unreal Engine 5 che ha una sua sezione nei dettagli della mesh, è probabile ma non
certa.}) non si vedono oggi.

\paragraph{Le modalità di visualizzazione.} Dal menu della \emph{viewport} (la
prospettiva ``non ci interessa'') si cambia la modalità: \textbf{Lit}, quella normale;
\textbf{Unlit}, ``com'è il mio oggetto se non è colpito dalla luce?''; e
\textbf{Wireframe}: ``vedete quello che dicevamo prima?'' Vertici, facce ed \emph{edge}
della mesh, che si potrebbero modificare con la \textbf{Modeling Mode} (``se avanzerà
il tempo con le lezioni, sarebbe carino vederla più avanti''). Ci sono anche
\textbf{Lighting Only}, le \textbf{Reflections} e altre ancora. Dal menu \textbf{Show}
si possono mostrare i \textbf{vertici}: ``dove sono i vertici di questo oggetto?''

\paragraph{Simple e complex.} Un'altra cosa utile ``da un punto di vista di gameplay''
sono le \textbf{collisioni}. Dal menu Show si attiva la \textbf{Simple Collision}: ``queste
sono le collisioni del mio oggetto''. In scena, il \emph{character} interagirà con
queste: ci salterà sopra, o lo dovrà scavalcare. Ma oltre alle \textbf{simple
collision} ci sono le \textbf{complex collision}. La differenza:
\begin{itemize}
  \item le \textbf{simple collision} sono forme che si possono posizionare a mano,
        anche da qui, e rifare da capo;
  \item le \textbf{complex collision} usano la struttura in vertici della mesh stessa
        come collisione. ``È qualcosa che è molto più pesante rispetto a una collisione
        semplice, e torna utile solo in alcuni casi specifici.''
\end{itemize}

\paragraph{Modificare le collisioni.} La collisione semplice si può selezionare,
spostare, anche eliminare; e se ne possono \textbf{aggiungere} altre: una
\textbf{sfera}, di cui si decide la grandezza; una \textbf{capsula}, che ``è
tendenzialmente la collisione che usano i vostri \emph{character}'': il personaggio di
solito non ha una collisione super precisa, ``è un capsulone'', a meno che non si voglia
davvero il contrario; una \textbf{box}. E una box, per questa mesh, ``è più ottimizzata
rispetto alla collisione che questa aveva di base'': all'atto pratico fa la stessa cosa
``con molti meno vertici''. Su una singola mesh cambia poco; su molte mesh aiuta le
prestazioni. Ci sono anche collisioni più complesse, ma ``idealmente, il modo migliore
per fare delle simple collision è letteralmente prendere delle \textbf{forme base}'',
magari più d'una, ``piuttosto che avere tanti vertici sulla stessa collisione''.

\begin{table}[htbp]
\centering
\small
\begin{tabular}{@{}p{3.3cm}p{4.4cm}p{5.3cm}@{}}
\toprule
\textbf{Collisione} & \textbf{Com'è fatta} & \textbf{Quando} \\
\midrule
Box & un parallelepipedo & oggetti squadrati; spesso basta anche per oggetti più
  complessi, ed è la più leggera \\
Sphere & una sfera & oggetti tondi, palle \\
Capsule & un cilindro con due semisfere & i \emph{character}: ``un capsulone'' \\
Convex (più forme) & poliedri convessi, magari più d'uno & quando una forma sola non
  basta; meglio poche forme base che tanti vertici \\
Complex & i triangoli della mesh stessa & casi specifici: precisione massima, costo
  molto più alto \\
\bottomrule
\end{tabular}
\caption{Le collisioni di una static mesh viste nello Static Mesh Editor. Le prime
quattro sono \emph{simple collision}, che si aggiungono e si modificano dal menu
Collision dell'editor; l'ultima è la \emph{complex collision}.}
\label{ue3:tab:collisioni}
\end{table}

\paragraph{Approfondimento: chi usa quale collisione.} Unreal tiene le due collisioni
separate perché le usa per cose diverse. La \textbf{simple} è quella della
\textbf{fisica}: la pallina della seconda lezione che rotola sulla piattaforma, il
\emph{character} che cammina e salta, gli oggetti che cadono; i calcoli fisici su forme
semplici sono molto più veloci. La \textbf{complex} si usa soprattutto per le
\textbf{interrogazioni} precise, come un \emph{line trace} (un ``raggio'' lanciato per
sapere che cosa colpisce un proiettile o che cosa sta guardando il giocatore), quando
la precisione al triangolo serve davvero. Nei dettagli della mesh, il campo
\emph{Collision Complexity} decide come combinarle; l'opzione \emph{Use Complex
Collision As Simple} usa i triangoli anche per la fisica, ed è comoda per oggetti
statici di forma strana, ma va usata con giudizio, per il costo di cui parla Davide.

% ---------------------------------------------------------------------
\section{Raggruppare e fondere: Merge Actors}
\label{ue3:sec:merge}

``Se ad esempio vogliamo avere tutti e tre gli oggetti assieme'', o semplicemente
spostarli insieme, ci sono due strade.

\paragraph{Raggruppare.} Si selezionano tutti e tre e li si \textbf{raggruppa} con
\textbf{Ctrl+G}; si spostano insieme; con \textbf{Shift+G} si toglie il gruppo.

\paragraph{Fondere.} Oppure li si \textbf{fonde} (\emph{merge}). Si evidenziano i tre
pezzi in scena, si va nel menu \textbf{Actor} in alto e si sceglie \textbf{Merge
Actors}. Unreal chiede dove salvare il nuovo oggetto e con che nome (per esempio
\textbf{SM\_Turret}\footnote{Nella trascrizione ``smTourette''.}), e crea una
\textbf{nuova mesh}, ``che è semplicemente la fusione delle tre mesh precedenti'': una
somma, un \emph{merge}, appunto.

\paragraph{Il merge, di nuovo.} Più tardi uno studente dice di essersi perso come si fa
il merge, e Davide lo rifà con calma. I tre oggetti (il cannone, la base e il pivot)
sono in scena. Con il \textbf{tasto destro} su di essi ci sono delle opzioni, ``però
queste sono opzioni limitate''. Per averle tutte si usa il menu \textbf{Actor} in alto:
mostra gran parte delle stesse opzioni, ma in basso, nella sezione degli strumenti di
Unreal Engine, c'è \textbf{Merge Actors}. Al merge Unreal chiede dove salvare e con che
nome, e ``ti fonde i tre oggetti''. Nello stesso menu c'è anche \textbf{Chaos}, che
``serve per gli oggetti distruttibili'': cose più avanzate, che probabilmente non si
vedranno nel corso, ma da approfondire ``se vi capiterà di voler mettere distruttibilità
all'interno del vostro gioco''.

\paragraph{Approfondimento: gruppo, merge o Blueprint?} Le tre soluzioni non sono
equivalenti. Il \textbf{gruppo} è solo una comodità dell'editor: gli actor restano
separati, e in gioco non cambia nulla. Il \textbf{merge} crea una mesh nuova, una sola:
meno oggetti da disegnare (utile per le prestazioni, per esempio con tanti pezzi di
arredo statici), ma i pezzi \textbf{non si possono più muovere uno rispetto
all'altro}. Per la torretta di Igor, che deve ruotare il supporto e alzare la canna, il
merge è quindi la scelta sbagliata: serve un \textbf{Blueprint} con i tre pezzi come
\textbf{componenti} separati, imparentati tra loro, come si è visto nella seconda
lezione. Il merge della lezione è una dimostrazione dello strumento, e serve anche a
rispondere alla domanda sui materiali (sezione~\ref{ue3:sec:domande}).

% ---------------------------------------------------------------------
\section{I materiali della torretta}
\label{ue3:sec:matturret}

``Ora abbiamo importato le texture. Creiamoci i materiali di queste torrette.'' Tre
materiali, uno per pezzo: \textbf{M\_Turret\_Base}, \textbf{M\_Turret\_Pivot} e
\textbf{M\_Turret\_Cannon}. Si potrebbero fare molto elaborati, ``ma partiamo
semplicemente andando a utilizzare le texture che ci sono state date''.

\paragraph{Quattro Texture Sample.} Nel materiale della base Davide prende \textbf{quattro
Texture Sample} e ci mette le texture della base: \textbf{Base Color},
\textbf{Metallic}, \textbf{Roughness} e \textbf{Normal} (Turret Base Color, Turret Base
Metallic, Turret Base Roughness, Turret Base Normal). Poi le collega ai pin
corrispondenti del nodo risultante: il colore al Base Color, la Metallic al Metallic, la
Roughness alla Roughness, la Normal alla Normal (figura~\ref{ue3:fig:turret}).

\begin{figure}[htbp]
\centering
\begin{tikzpicture}[
  font=\small,
  ue3node/.style={draw, thick, rounded corners=4pt, align=center,
                  minimum height=0.85cm, text width=3.6cm},
  ue3col/.style={-{Stealth[length=2mm]}, thick, blue!60!black},
  ue3flt/.style={-{Stealth[length=2mm]}, thick, green!50!black}
]
  \node[ue3node, fill=orange!20] (c) at (0,2.1) {\textbf{Texture Sample}\\
    \footnotesize T\_Turret\_Base\_Color};
  \node[ue3node, fill=gray!15] (m) at (0,0.7) {\textbf{Texture Sample}\\
    \footnotesize T\_Turret\_Base\_Metallic};
  \node[ue3node, fill=gray!15] (r) at (0,-0.7) {\textbf{Texture Sample}\\
    \footnotesize T\_Turret\_Base\_Roughness};
  \node[ue3node, fill=blue!15] (n) at (0,-2.1) {\textbf{Texture Sample}\\
    \footnotesize T\_Turret\_Base\_Normal};
  \node[draw, very thick, rounded corners=4pt, fill=blue!6, minimum width=3.0cm,
    minimum height=5.2cm] (res) at (6.2,0) {};
  \node[font=\small\bfseries] at (6.2,2.3) {M\_Turret\_Base};
  \foreach \lab/\y in {Base Color/2.1, Metallic/0.7, Specular/0.0, Roughness/-0.7,
                       Emissive Color/-1.4, Normal/-2.1}
    \node[font=\footnotesize, anchor=west] at (4.85,\y) {\lab};
  \draw[ue3col] (c.east) -- (4.8,2.1);
  \draw[ue3flt] (m.east) -- node[above, font=\scriptsize]{R} (4.8,0.7);
  \draw[ue3flt] (r.east) -- node[above, font=\scriptsize]{R} (4.8,-0.7);
  \draw[ue3col] (n.east) -- (4.8,-2.1);
\end{tikzpicture}
\caption{Il materiale della base della torretta, fatto a mano: quattro Texture Sample,
quattro pin. Metallic e Roughness sono mappe in bianco e nero, e basta un canale (per
esempio R). M\_Turret\_Pivot e M\_Turret\_Cannon sono identici, con le loro texture.
Ricostruzione dalla dimostrazione; i nomi esatti dei file sono indicativi.}
\label{ue3:fig:turret}
\end{figure}

\paragraph{Le preview.} ``Mi ero dimenticato di farvi vedere una cosa.'' Nella
\emph{preview} del materiale si può scegliere la forma: il \textbf{cilindro}, la
\textbf{sfera}, il \textbf{piano}, il \textbf{cubo} e il \textbf{custom}. Con il custom,
se si seleziona una mesh nel Content Browser e si clicca il pulsante, la preview mostra
\textbf{quella mesh} con il materiale applicato. La prima volta non la mostra: ``sta
compilando''. Poi compare (Davide aveva collegato una texture ``dalla parte sbagliata'',
e la sistema).

\paragraph{Gli slot.} Il materiale è salvato e compilato. Sulla mesh della torretta ci
sono i \textbf{due slot}: in uno c'è M\_Turret\_Base, ma ``per come è stata strutturata
la torretta'' nell'altro va il materiale del \textbf{pivot}. ``Sta rallentando
veramente tanto il PC.''

\paragraph{Un grafico, un engine.} Davide si ferma su che cosa si è fatto: ``qui non
abbiamo fatto altro che utilizzare le mesh che ci sono state fornite da un grafico''.
Idealmente un grafico, in un software a parte, ha realizzato la mesh e le texture, e ce
le ha passate; noi le abbiamo inserite. Si potrebbe lavorare ancora sull'estetica della
torretta modificando i materiali, ``tuttavia, in questo caso non è necessario''.

\paragraph{Il pivot e il cannone, più furbi.} Per il materiale del pivot Davide è ``un
po' più furbo'': cerca la texture \emph{Turret Pivot Color}, e le altre le copia da lì,
perché nella lista ``stanno direttamente dopo''. Collega Metallic, Roughness e Normal,
salva, ``e abbiamo anche il materiale del nostro pivot''. Poi lo stesso per il cannone:
la texture \emph{Cannon Color}, \textbf{Ctrl+C}, \textbf{Ctrl+V}. Lì Davide ricorda le
scorciatoie classiche: ``tutti i classici shortcut di Control'', Ctrl+C, Ctrl+V, Ctrl+Z,
Ctrl+Y, Ctrl+A, ``sono applicabili anche a Unreal nella maggior parte dei casi''.

\paragraph{Salvate spesso.} Uno studente chiede perché il suo non si è applicato:
``Salvalo.'' E Davide ne fa una regola: ``puntualmente mi capita di vedere gente che sta
lavorando a qualcosa, non salva per qualche motivo, Unreal crasha e perde tipo tutto il
lavoro che stava facendo. E quindi salvate spesso.''

\paragraph{La torretta finita.} Nei due slot vanno il pivot e il cannone. ``Oh, eccoci
qua. Bello lucido.'' Salvato tutto, la torretta ha i suoi materiali.

\paragraph{Ritoccare il colore con un Lerp.} ``Se volete che la torretta sia meno
lucida, o che sia meno nitido l'arancione'', si può intervenire sui singoli materiali.
Davide improvvisa sul cannone: non è sicuro di volere il colore ``così netto'', quindi
prende un \textbf{nero} e un \textbf{Lerp}, collega il nero e la texture del colore alle
due entrate e mette nell'alfa un parametro con cui decidere ``quanto si deve vedere'',
per esempio 0,6. Si può convertire in parametro, creare una \textbf{Material
Instance}, e regolarlo da lì. Al posto del nero si può mettere un altro colore, un
verde: ``viene un po' più verdastra, è troppo verde'', e con l'alfa a 0,8 viene ``questo
verde mezzo, insomma, un po' scuro, un po' cromato''. ``Qui chiaramente sto improvvisando
cose, ma cose anche stupide. Ma nel momento in cui voi avete un'idea un po' più chiara,
un po' più specifica di quello che volete, potete andare a fare un po' più di cose.''

\paragraph{Approfondimento: lo stesso lavoro con un master material.} I tre materiali
della torretta hanno lo stesso identico grafo: cambiano solo le texture. È il caso da
manuale per le Material Instance della prima parte della lezione: un solo master
material (per esempio \texttt{M\_Turret}) con quattro Texture Parameter (colore,
metallic, roughness, normal) e magari un Lerp per la tinta, e tre istanze
(\texttt{MI\_Turret\_Base}, \texttt{MI\_Turret\_Pivot}, \texttt{MI\_Turret\_Cannon}) che
impostano le texture di ciascun pezzo. Un solo shader da compilare, e il ritocco
verdastro del cannone diventa una casella da regolare nell'istanza. Se le texture
fossero state importate insieme alle mesh (Import Materials e Import Textures), Unreal
avrebbe creato da solo i materiali, uno per slot, con un grafo simile a quello della
figura~\ref{ue3:fig:turret}.

% ---------------------------------------------------------------------
\section{Le domande: il Reimport e il merge}
\label{ue3:sec:domande}

``Ora tutto chiaro il discorso dell'import e dei materiali?'' Due domande.

\paragraph{Perché il Reimport con la Z?} Una studentessa non ha capito il reimport
modificando la Z. Davide riapre il cannone. Il reimport è servito a \textbf{modificare la
mesh} alla radice. L'oggetto importato, modellato da qualcun altro, era fatto in un
certo modo, e non è un problema in sé: lo si può abbassare e girare a mano, in scena.
``Però poi, tutte le volte che io prendo e rimetto in scena l'oggetto, l'oggetto sarà
sempre posizionato in quel modo e ruotato in quel modo.'' Il reimport invece ``modifica
proprio alla radice il modo in cui funziona l'oggetto''. E il problema principale non
era la rotazione (``so che l'angolo è 113, allora lo rimodifico''), ma il
\textbf{pivot}, ``in una posizione scomoda'': il reimport è un modo per spostarlo.

\paragraph{La Modeling Mode.} ``Non è l'unico modo.'' Davide apre la \textbf{Modeling
Mode}, cerca lo strumento del \textbf{pivot}, e mostra che lo si può mettere al centro,
a sinistra, a destra, oppure ``nel punto specifico che dico io: lo prendo e me lo
sposto''. Ma questa è ``una soluzione che hanno aggiunto con Unreal Engine 5'': prima la
Modeling Mode non c'era, e in Unreal Engine 4 reimportare era il modo comodo per farlo
``senza avere bisogno di rimodellare l'oggetto a posteriori''. ``Fai la stessa cosa,
semplicemente non hai fatto un'operazione extra.''

\paragraph{Una prova.} Davide stesso è curioso: se si reimporta di nuovo con $-135$ e
113? Il reimport \textbf{sovrascrive} quello che c'era prima, ma il cannone in scena
conserva la rotazione che gli era stata data a mano. Davide lo rimette dove stava (``io
sto facendo a cazzotti da tutta la lezione con la finestra di Zoom, se non si fosse
capito'').

\paragraph{Che cosa succede ai materiali con il merge?} Era la domanda arrivata in chat
durante la costruzione dei materiali: se faccio il merge e i pezzi della torretta hanno
texture diverse? Davide prova: seleziona i tre oggetti, \textbf{Actor}, \textbf{Merge
Actors}, nome \emph{Turret}, salva. Il risultato è un oggetto con \textbf{tre slot} per i
materiali: lo slot 0 è tutta la parte del cannone, lo slot 1 il pivot ``e tutta questa
zona qui'', e così via. Se in uno slot si mette un materiale vuoto, ``diventa tutto
nero'': ``vedete, tutta questa zona adesso viene influenzata da quel materiale lì''.
Risposta: con il merge l'actor viene fuso, e i materiali vengono \textbf{riposizionati
negli slot} della nuova mesh.

% ---------------------------------------------------------------------
\section{Lo stretch, lo z-fighting e i materiali World Aligned}
\label{ue3:sec:world}

``Mi è venuto in mente un'ultima cosa che vi potrei far vedere per quanto riguarda i
materiali, e poi vi vorrei far fare una cosina.'' Manca una mezz'oretta, ``visto che
Igor vi ha fatto fare delle cose, ho visto che siete molto partecipi. Quindi vorrei
fomentare questa cosa.''

\paragraph{M\_Test.} Davide ricrea il materiale della corteccia, cancellato per sbaglio:
\textbf{M\_Test}, con due Texture Sample (Bark e la sua Normal) e un Texture
Coordinate, ``proprio la roba più facile del mondo'', con il tiling a 10.

\paragraph{Ingrandire e stirare.} ``Ora vediamo se abbiamo capito come funziona tutto
quanto.'' I materiali sono collegati alle \textbf{UV} di un oggetto. Quindi se
l'oggetto si \textbf{ingrandisce}, al materiale ``non dovrebbe succedere granché'': si
ingrandisce anche lui (``la risoluzione non ne giova particolarmente''). E se l'oggetto
si \textbf{stira} (\emph{stretch}), tirandolo verso destra? ``Credo si stira anche
lui.'' Esatto: il materiale si allunga. ``In alcuni casi non è neanche particolarmente
brutto'', ma la corteccia ``non sembra più corteccia''; e tirando di più, ``così è
brutto''.

\paragraph{Il pavimento grande.} Se si vuole fare un pavimento per una zona molto
ampia, dover pensare a questo ``può essere una noia'': per non stirare il materiale si
dovrebbero fare tante copie dello stesso pezzo, una attaccata all'altra, ``che
tendenzialmente non è quello che voglio'', perché è scomodo.

\paragraph{Una digressione: lo z-fighting.} ``Visto che è una domanda che arriverà,
anticipo'': se due piani si trovano \textbf{nella stessa posizione}, ``clippano'': la
superficie sfarfalla a chiazze, perché l'engine ``non sa se utilizzare il materiale
dell'oggetto a destra o dell'oggetto a sinistra''. Le soluzioni: separarli davvero,
oppure spostarne uno di pochissimo in \textbf{Z}, per esempio di 0,1: ``è
impercettibile, non si nota, però uno è leggermente più sopra rispetto all'altro'', e il
problema sparisce.

\paragraph{Allineato al mondo.} ``C'è un modo per far sì che il mio materiale non si
stiri quando stiro l'oggetto? Chiaramente sì, perché se no non avrei fatto questa
domanda.'' Bisogna fare un materiale che non è allineato con le coordinate
dell'oggetto, ma con quelle del \textbf{mondo}. Davide lo spiega con il \emph{gizmo}
di spostamento: un oggetto ruotato si può muovere lungo gli assi del mondo o lungo i
suoi assi locali, e il pulsante in alto passa dall'uno all'altro. ``Più o meno il
discorso è lo stesso'': nel caso del materiale non si parla di assi, ma di UV.

\paragraph{I nodi World Aligned.} Davide copia il materiale (\textbf{M\_Test1}) e usa
due nodi, \textbf{World Aligned Texture} e \textbf{World Aligned Normal}. Non hanno
scorciatoie: si cercano per nome. Hanno ``un sacco di parametri'', ma per ora ne servono
due: \textbf{Texture Object} e \textbf{Texture Size}. World Aligned Texture va nel
\textbf{Base Color}, World Aligned Normal nella \textbf{Normal}, usando l'uscita
\textbf{XYZ}, ``perché almeno prendo tutti gli assi'': il materiale è allineato al
mondo su tutti e tre.

\paragraph{Texture Size e Texture Object.} Il \textbf{Texture Size}, un vettore a tre
valori, ``fa la stessa cosa del Texture Coordinate'': regola la grandezza della
ripetizione. Perché allora non tenere il Texture Coordinate sulle UV del Texture
Sample? Perché i nodi World Aligned \textbf{non prendono} un Texture Sample: vogliono un
\textbf{Texture Object}. Si converte: tasto destro sul Texture Sample, \textbf{Convert to
Texture Object}. Il Texture Object ``è un oggetto molto più semplice'': non ha
ingressi, ha una sola uscita, niente RGB; ``è molto più limitato in quello che può
fare'', ma è il ``messaggio'' che serve ai nodi World Aligned.

\paragraph{La prova.} Davide prova diversi valori di Texture Size (1, 2, 10) per avvicinarsi
al risultato dell'altro materiale, con qualche esitazione su quale sia ``troppo'' e quale
``troppo poco'' (``è proprio piccola la texture, allora''), e applica il nuovo
materiale a un secondo oggetto. Poi \textbf{stira} entrambi: il primo si stira; il
secondo ``semplicemente genera nuove texture'', rimane orientato con il mondo,
``rimane sempre della stessa proporzione''.

\paragraph{Dov'è la fregatura.} ``Se questo rimane sempre proporzionato anche quando
stiriamo le cose, perché non dovremmo avere tutti i materiali World Aligned? Dov'è la
fregatura?'' Eccola: essendo allineato al mondo, se l'oggetto si \textbf{muove} o
\textbf{ruota}, il materiale \textbf{resta fermo}. Una \textbf{porta} con un materiale
normale si apre e il materiale resta attaccato alla porta; una porta con un materiale
World Aligned si apre e ``il materiale sembra scorrere sulla porta'', ``come se il
materiale fosse animato''. ``Ci siamo? Ho rotto le palle? Datemi un feedback.''

\begin{figure}[htbp]
\centering
\begin{tikzpicture}[font=\small]
  % materiale UV: stiramento
  \node[font=\footnotesize\bfseries] at (2.6,2.6) {materiale normale (UV)};
  \draw[thick, fill=brown!20] (0,1.0) rectangle (1.2,2.2);
  \foreach \i in {0.2,0.6,1.0} \draw[brown!70!black, thick] (\i,1.15) -- (\i,2.05);
  \draw[-{Stealth[length=2mm]}, thick] (1.5,1.6) -- (2.1,1.6);
  \draw[thick, fill=brown!20] (2.4,1.0) rectangle (5.4,2.2);
  \foreach \i in {0.5,1.5,2.5} \draw[brown!70!black, thick] (2.4+\i,1.15) --
    (2.4+\i,2.05);
  \node[font=\footnotesize, align=center] at (2.7,0.55) {stirato: le venature si
    allargano};
  % materiale world aligned
  \begin{scope}[shift={(6.6,0)}]
    \node[font=\footnotesize\bfseries] at (2.6,2.6) {materiale World Aligned};
    \draw[thick, fill=brown!20] (0,1.0) rectangle (1.2,2.2);
    \foreach \i in {0.2,0.6,1.0} \draw[brown!70!black, thick] (\i,1.15) -- (\i,2.05);
    \draw[-{Stealth[length=2mm]}, thick] (1.5,1.6) -- (2.1,1.6);
    \draw[thick, fill=brown!20] (2.4,1.0) rectangle (5.4,2.2);
    \foreach \i in {0.2,0.6,...,3.0} \draw[brown!70!black, thick] (2.4+\i,1.15) --
      (2.4+\i,2.05);
    \node[font=\footnotesize, align=center] at (2.7,0.55) {stirato: ``genera nuove
      texture'', stessa scala};
  \end{scope}
  % la porta
  \node[font=\footnotesize, align=center, text width=12.5cm] at (5.9,-0.3)
    {la fregatura: se l'oggetto si muove o ruota (una porta che si apre), il materiale
    World Aligned resta fermo nel mondo e ``sembra scorrere'' sull'oggetto};
\end{tikzpicture}
\caption{Materiale legato alle UV e materiale World Aligned a confronto. Il primo
segue l'oggetto e si deforma con lui; il secondo non si deforma, ma non segue
l'oggetto quando si muove. Va bene per pavimenti, muri e rocce statici; non per porte,
veicoli o la torretta.}
\label{ue3:fig:world}
\end{figure}

\paragraph{Le linee nere.} Una studentessa vede delle \textbf{linee nere} al confine tra
una ripetizione della texture e l'altra, sul pavimento. Davide guarda il suo grafo:
stessa texture, Bark, ma ha collegato l'uscita \textbf{XY} invece di \textbf{XYZ}. ``In
quel modo lui aveva due assi che erano allineati col mondo, e l'altro no'', e il
risultato era diverso. ``Ok, grazie.''

\paragraph{Approfondimento: come funziona un World Aligned.} Il nodo World Aligned
Texture, una Material Function del motore, non usa le UV della mesh: per ogni pixel
prende la sua \textbf{posizione nel mondo} e la usa come coordinata della texture. Lo fa
tre volte, proiettando la texture lungo i tre assi (dall'alto, di lato, di fronte), e
mescola le tre proiezioni secondo la direzione della superficie: è la tecnica detta
\textbf{triplanar mapping}. L'uscita XY usa una sola proiezione (dall'alto, buona per i
pavimenti), l'uscita XYZ le mescola tutte; per questo, su un oggetto con facce in tutte
le direzioni, con XY le pareti laterali ricevono una texture ``stirata'' a strisce, che
può dare proprio le linee scure viste dalla studentessa. Il \textbf{Texture Size} è
espresso in \textbf{unità del mondo}, cioè in centimetri: è la grandezza di una
ripetizione della texture. Funziona quindi al contrario del tiling: un tiling più alto
fa ripetere di più la texture, un Texture Size più alto la fa ripetere di meno (una
texture di 100 si ripete ogni metro, una di 10 ogni dieci centimetri). È la ragione
delle esitazioni di Davide su quale valore fosse ``troppo''. Il costo: tre letture della
texture invece di una, quindi un materiale un po' più pesante.

% ---------------------------------------------------------------------
\section{L'esercizio della lava, il nodo Noise e la chiusura}
\label{ue3:sec:lava}

``Allora, ci siamo visti un po' di cose, fatto un miscuglione di cose per quanto riguarda
materiali, texture e import delle mesh.'' Nessuna domanda: ``sembra tutto a posto''.

\paragraph{La consegna.} Per esercitarsi con i materiali, e magari scoprire qualche nodo
non usato a lezione (il nodo \textbf{Noise}, qualche \textbf{maschera}), Davide propone
una prova di una decina di minuti, da continuare poi a casa come le cose di Igor:
\textbf{un materiale di lava}. Di che componenti avrà bisogno? Deve \textbf{muoversi}
un po', quindi bisognerà gestire la velocità del \textbf{Panner}, ``a meno che non vi
cercate proprio la texture della lava''. Ma la richiesta è di improvvisare: ``provate a
vedere cosa riuscite a combinare andando a fare delle operazioni fra colori e fra altre
texture'', con un po' di movimento, un \textbf{Base Color}, ``ma anche un
\textbf{Emissive}, perché la lava comunque deve emettere luce''.

\paragraph{Il ripasso dei nodi.} Una studentessa non ricorda il nome della funzione
``che abbiamo usato per collegare rosso e blu'': il Lerp. Davide ricondivide lo schermo
e fa il riepilogo dei nodi della giornata: la \textbf{costante}, il \textbf{colore}, la
\textbf{texture}, il \textbf{nodo risultante}, il \textbf{Texture Coordinate}, i
\textbf{World Aligned}, le \textbf{operazioni} (con il consiglio sulla moltiplicazione
visto nella sezione~\ref{ue3:sec:costanti}), il \textbf{Panner}, il \textbf{Lerp}.
Divisione e sottrazione servono soprattutto per gli effetti speciali, e Davide vorrebbe
fare una lezione su \textbf{Niagara}, ``il tool che usa Unreal per fare gli effetti
speciali'', se ci sarà tempo.

\paragraph{Le prossime lezioni con Davide.} Davide annuncia come saranno le sue
lezioni finali: ci sarà ``una lezione che sarà un'esercitazione intera'' in cui si
proverà a fare dei \textbf{livelli}, e chiederà agli studenti di condividere lo schermo
e di \emph{streamare} a turno (``se per voi è un problema, magari lo facciamo solo con le
persone per cui non è un problema'', e lo dirà in anticipo). Per esempio, ``un livello
\textbf{platform}'', con le interazioni che si faranno nelle lezioni di Igor e quello che
si vedrà insieme per il \emph{level design}.

\paragraph{Il nodo Noise.} Un nodo ``un po' più complicato'', per chi vuole provare: il
\textbf{Noise}, che ``serve per simulare un \emph{noise}'': il \emph{Voronoi}, il
\emph{White Noise}, un'alternanza di bianco e nero ``che si usa molto spesso per fare
delle distorsioni''. Ha molti parametri: la scala del \emph{noise}, e la possibilità di
filtrarlo con dei vettori o con delle mappe. ``Potete smanettarci un pochino.''

\paragraph{Lo stream della studentessa.} ``C'è qualcuno che, nel mentre sta provando a
fare qualcosa, vuole streamare in questi ultimi cinque minuti?'' Una studentessa ha
provato, ma non riesce a trasformare la texture in World Aligned: è perché usa un
Texture Sample, e serviva un Texture Object. ``Prova a non usare il World Aligned, e
farlo così, diciamo, semplice.'' Senza World Aligned le viene ``questa cosa'', sempre
con il Bark. Davide la guida:
\begin{itemize}
  \item nella lava servono ``delle zone più scure e delle zone più luminose'': quindi
        un \textbf{Lerp} tra un \textbf{colore} e un \textbf{nero}, gestito da un'alfa
        che sia una texture con dei bianchi e dei neri ``posizionati come vuoi'', da
        mandare nell'\textbf{Emissive Color};
  \item l'alfa va cercata tra le texture del progetto, ``qualcosa per dare un effetto
        tipo di distorsione''; le \textbf{Normal}, invece, ``lasciale perdere'': per
        un'alfa si usano le Diffuse o le mappe come Specular e Roughness;
  \item se una texture fa da alfa, nelle entrate del Lerp ci vogliono \textbf{due
        colori}, ``uno che sia molto scuro, uno che sia più chiaro''; ricordando che
        l'alfa ``prende i bianchi e i neri e dice: nei bianchi vai più verso B, nei neri
        vai più verso A'';
  \item il risultato va ``forse più nell'Emissive che non nel Base Color'', perché la
        lava la si vuole ``infuocata, accesa, luccicante''; il Base Color si può anche
        non usare, o usarlo con una texture adatta (``il Bark, per come è fatto, forse
        non è proprio la texture più adatta per fare una lava'');
  \item nell'Emissive si può \textbf{moltiplicare} il risultato, in modo che le zone
        accese siano molto luminose.
\end{itemize}
La studentessa nota che il nero, nell'Emissive, quasi non si vede: ``esatto'', il nero
non emette niente. E usandolo come colore ``non si vede molto'', perché il Bark come alfa
``è molto bianco'', e finisce per scurire un po' il tutto. ``È comunque una cosa
complicata, quindi vedete effettivamente se riuscite a tirar fuori qualcosa.''

\begin{figure}[htbp]
\centering
\begin{tikzpicture}[
  font=\small,
  ue3node/.style={draw, thick, rounded corners=4pt, align=center,
                  minimum height=0.85cm},
  ue3uv/.style={-{Stealth[length=2mm]}, thick, red!60!black},
  ue3flt/.style={-{Stealth[length=2mm]}, thick, green!50!black},
  ue3col/.style={-{Stealth[length=2mm]}, thick, blue!60!black}
]
  \node[ue3node, fill=red!12, text width=1.7cm] (tc) at (0,0) {\textbf{TexCoord}};
  \node[ue3node, fill=gray!10, text width=1.9cm] (pan) at (2.6,0) {\textbf{Panner}\\
    \footnotesize Speed 0,05; 0,02};
  \node[ue3node, fill=gray!15, text width=2.3cm] (msk) at (5.6,0)
    {\textbf{Texture Sample}\\\footnotesize maschera in b/n\\\footnotesize (o un
    Noise)};
  \node[ue3node, fill=black!70, text=white, text width=1.9cm] (dark) at (5.6,2.0)
    {\textbf{scuro}\\\footnotesize quasi nero};
  \node[ue3node, fill=orange!60, text width=1.9cm] (hot) at (5.6,-2.0)
    {\textbf{chiaro}\\\footnotesize arancione};
  \node[ue3node, fill=gray!10, text width=1.3cm, minimum height=2.0cm] (lerp)
    at (8.6,0) {\textbf{Lerp}\\[3pt]\footnotesize A\\\footnotesize B\\
    \footnotesize Alpha};
  \node[ue3node, fill=green!8, text width=1.6cm] (k) at (8.6,-2.4)
    {\footnotesize intensità\\\footnotesize (parametro)};
  \node[ue3node, fill=gray!10, text width=1.4cm] (mul) at (10.9,-0.6)
    {\textbf{Multiply}};
  \node[ue3node, fill=blue!6, text width=2.0cm, minimum height=1.6cm] (res)
    at (13.4,0) {\footnotesize Base Color\\[6pt]\footnotesize\textbf{Emissive}};
  \draw[ue3uv] (tc) -- (pan);
  \draw[ue3uv] (pan) -- node[above, font=\scriptsize]{UVs} (msk);
  \draw[ue3col] (dark.east) -- ++(0.7,0) |- ([yshift=5mm]lerp.west);
  \draw[ue3col] (hot.east) -- ++(0.9,0) |- ([yshift=0mm]lerp.west);
  \draw[ue3flt] (msk.east) -- ++(0.35,0) |- ([yshift=-5mm]lerp.west);
  \draw[ue3col] (lerp.east) -- ++(0.3,0) |- ([yshift=2mm]mul.west);
  \draw[ue3flt] (k.east) -- ++(0.5,0) |- ([yshift=-2mm]mul.west);
  \draw[ue3col] (mul.east) -- ++(0.4,0) |- ([yshift=-3mm]res.west);
  \draw[ue3col, dashed] (lerp.north) |- ([yshift=4mm]res.west);
\end{tikzpicture}
\caption{Una possibile soluzione dell'esercizio della lava, costruita con i soli nodi
della lezione e con i consigli dati alla studentessa: una maschera in bianco e nero che
scorre con il Panner fa da alfa a un Lerp tra uno scuro e un arancione; il risultato,
moltiplicato, va nell'Emissive (e, tratteggiato, eventualmente anche nel Base Color).
Non è un grafo mostrato a lezione: i valori sono indicativi.}
\label{ue3:fig:lava}
\end{figure}

\paragraph{La chiusura.} ``Vabbè, io in realtà adesso ho gettato questa bomba, e penso
che vi lascerò fare. Prendetevi tutto il tempo che volete, nessuna pressione.'' Davide
rivedrà la classe ``a luglio, probabilmente'', per le ultime quattro o cinque lezioni del
corso, sul \textbf{level design} e sulla \textbf{creazione di livelli}. ``Si sono fatte
le 7.'' ``Direi che è stato un piacere, mi sembrate tutti abbastanza interessati,
sicuramente belli, belli, belli attivi.'' E forse ci si rivedrà prima, ``se Igor ha
bisogno di un'altra sostituzione''. ``Ciao a tutti, arrivederci.''

\paragraph{Approfondimento: una lava più convincente.} Partendo dallo schema della
figura~\ref{ue3:fig:lava}, tre accorgimenti, tutti con nodi già visti, rendono la lava
molto più viva. Primo: \textbf{due Panner} con velocità e tiling diversi sulla stessa
maschera, sommati o moltiplicati, così che il movimento non si ripeta mai uguale (lo
stesso trucco dell'acqua). Secondo: un Multiply sull'Emissive con valori ben oltre 1
(10, 20), perché solo così il \emph{bloom} fa ``bruciare'' le zone calde, come nel
materiale con il Glow Intensity. Terzo: usare la \textbf{stessa} maschera per la
\textbf{Roughness}, con un Lerp tra 0,3 (la lava liquida, lucida) e 1 (la crosta
raffreddata, opaca): il materiale cambia così non solo colore ma anche superficie, che è
il senso del PBR. Il nodo Noise di Davide può sostituire la texture della maschera, ma è
calcolato a ogni pixel e costa più di una texture: per un oggetto grande in scena, una
texture di \emph{noise} è di solito la scelta più leggera.

\begin{table}[htbp]
\centering
\small
\begin{tabular}{@{}p{4.4cm}p{8.6cm}@{}}
\toprule
\textbf{Comando} & \textbf{Che cosa fa} \\
\midrule
Tasto destro nel grafo & menu dei nodi, con ricerca \\
T + clic & Texture Sample \\
1 + clic / 2 + clic / 3 + clic & Constant, Constant2Vector, Constant3Vector (il
  colore) \\
A / M / D + clic & Add, Multiply, Divide \\
S + clic & Scalar Parameter (per questo Subtract non ha scorciatoia) \\
U + clic & Texture Coordinate (``lavora sulle UV'') \\
L + clic & Lerp (Linear Interpolate) \\
P + clic & Panner \\
Alt + clic su un pin & stacca i collegamenti \\
Tasto destro su un nodo, Convert to Parameter & trasforma il nodo in un parametro,
  regolabile dalle Material Instance \\
Tasto destro su un Texture Sample, Convert to Texture Object & per i nodi World Aligned \\
Use Selected Asset from Content Browser & mette nel nodo la texture selezionata nel
  Content Browser \\
Tasto destro su un materiale, Create Material Instance & crea un'istanza \\
Reimport Mesh With Dialog & reimporta una mesh riaprendo la finestra di import \\
Ctrl+G / Shift+G & raggruppa e separa gli actor selezionati \\
Actor, Merge Actors & fonde gli actor selezionati in una nuova mesh \\
Ctrl+C, Ctrl+V, Ctrl+Z, Ctrl+Y, Ctrl+A & le scorciatoie classiche, valide quasi
  ovunque \\
Ctrl+S & salvare: ``salvate spesso'' \\
\bottomrule
\end{tabular}
\caption{I comandi della terza lezione nell'editor dei materiali, nel Content Browser e
nel Level Editor di Unreal Engine 5. Le scorciatoie con il clic si usano tenendo premuto
il tasto e cliccando nel grafo; World Aligned Texture e World Aligned Normal non ne
hanno.}
\label{ue3:tab:scorciatoie}
\end{table}

% ---------------------------------------------------------------------
\section*{Punti chiave del capitolo}
\addcontentsline{toc}{section}{Punti chiave del capitolo}

\begin{itemize}
  \item Il grafico lavora soprattutto fuori da Unreal (Blender, Maya); dentro Unreal
        \textbf{importa} gli oggetti e costruisce gli \textbf{shader}.
  \item Una \textbf{texture} è la ``carta regalo'' dell'oggetto: si avvolge sulla mesh
        attraverso le \textbf{UV}, che si creano prima, nel software di modellazione.
        Una texture non pensata per quelle UV si stira.
  \item Le mappe principali: \textbf{Diffuse/Albedo} (il colore, suffisso \_D o \_A),
        \textbf{Normal} (blu: l'illusione del rilievo senza poligoni in più),
        \textbf{Specular}, \textbf{Roughness}, \textbf{Metallic}, la mappa di
        \textbf{altezza} (\emph{displacement} vero, ``una normal sotto steroidi'').
  \item Nelle mappe in bianco e nero il \textbf{bianco è 1}, il \textbf{nero è 0}, il
        grigio sta in mezzo. Ogni canale RGB è una mappa a sé: una texture
        \textbf{AORM}/ORM impacchetta tre mappe in un file.
  \item Il \textbf{materiale} è la ``vernice'', uno \textbf{shader}: si costruisce con
        nodi, come un Blueprint, e combina le texture. Sugli oggetti si applicano
        materiali, non texture.
  \item \textbf{Material Domain} (per che cosa), \textbf{Blend Mode} (Opaque, Masked,
        Translucent) e \textbf{Shading Model} (Default Lit, Unlit) cambiano i pin del
        nodo risultante.
  \item Nodi di base: \textbf{Texture Sample} (T), costanti (1, 2, 3), \textbf{Add},
        \textbf{Multiply}, \textbf{Divide} (A, M, D), Scalar Parameter (S),
        \textbf{Texture Coordinate} (U), \textbf{Lerp} (L), \textbf{Panner} (P). La
        moltiplicazione è l'operazione più usata.
  \item Il rendering è \textbf{PBR}: Roughness 1 è la roccia, 0 lo specchio; Metallic
        e Specular decidono come la superficie riflette; l'\textbf{Emissive}
        moltiplicato fa brillare il materiale.
  \item \textbf{UTiling} e \textbf{VTiling} (o un Multiply sulle coordinate) decidono
        quante volte si ripete la texture: 2 e 2 sono quattro ripetizioni, 0,5 e 0,5 un
        quarto della texture.
  \item Un nodo diventa un \textbf{parametro} con Convert to Parameter. Le
        \textbf{Material Instance} (MI\_) cambiano i parametri senza toccare il grafo
        del \textbf{master material} (M\_), si aggiornano subito e sono più leggere:
        sugli oggetti si applicano sempre le istanze.
  \item Le \textbf{Material Function} sono ``parentesi'' di calcoli riusabili tra
        materiali e tra membri del team.
  \item Il \textbf{Lerp} interpola tra A e B secondo un'\textbf{alfa} da 0 a 1, che può
        essere una texture (``l'arbitro''); il \textbf{Panner} fa scorrere la texture
        (acqua, fluidi, lava).
  \item Importare \textbf{FBX}, controllare le \textbf{unità} (Unreal usa i
        centimetri), texture in PNG o JPEG. La finestra di import genera collisioni,
        corregge il transform, importa materiali e texture.
  \item Una mesh è fatta di \textbf{vertici}, \textbf{spigoli} e \textbf{facce}
        (triangoli). Il \textbf{pivot} deve stare dove il pezzo ruota: con
        \textbf{Reimport Mesh With Dialog} (cannone: Z $-135$, rotazione 113) o con la
        \textbf{Modeling Mode} di Unreal Engine 5.
  \item Le \textbf{simple collision} (box, sfera, capsula: il ``capsulone'' dei
        \emph{character}) sono leggere; le \textbf{complex} usano i triangoli della mesh
        e costano molto di più.
  \item \textbf{Ctrl+G} raggruppa; \textbf{Merge Actors} fonde in una nuova mesh, con
        tanti slot quanti i materiali dei pezzi.
  \item Due piani nella stessa posizione sfarfallano (\textbf{z-fighting}): basta
        alzarne uno di 0,1. I materiali \textbf{World Aligned} (Texture Object, Texture
        Size, uscita XYZ) non si stirano, ma restano fermi se l'oggetto si muove.
  \item \textbf{Salvate spesso.} E per esercizio: un materiale di \textbf{lava}, con
        Panner, Lerp ed Emissive.
\end{itemize}

%
%  Tutte le etichette di questo capitolo hanno il prefisso "ue3:"
%  (ue3:cap:..., ue3:sec:..., ue3:fig:..., ue3:tab:...), e tutti gli
%  stili TikZ hanno il prefisso "ue3" (ue3node, ue3arr, ...) per non
%  collidere con chiavi native di TikZ né con quelli delle prime due
%  lezioni. Il capitolo è autonomo: non contiene \ref verso altri
%  capitoli (i rimandi alle lezioni precedenti sono a parole).
%  Le figure sono ricostruzioni schematiche di quanto mostrato e detto a
%  lezione, non copie delle slide o delle schermate del docente. Dove la
%  trascrizione è troppo rovinata per ricostruire un dettaglio, il
%  capitolo lo dice invece di inventarlo.
% =====================================================================

\chapter[Texture, materiali e mesh]{Texture, materiali e mesh: la grafica dentro
Unreal e la torretta}
\label{ue3:cap:materiali}

Questa è la \textbf{terza lezione} del corso breve di \textbf{introduzione a Unreal
Engine}, e cambia sia docente sia argomento. Al posto di Igor, che ``aveva un
impiccio'', c'è \textbf{Davide}, che si presenta così: ``io sono Davide e in questa
lezione, nello specifico, sostituisco Igor''. Davide tornerà più avanti: ``le ultime
lezioni del corso le fate con me'', e saranno lezioni ``un po' più basate su
\textbf{level design} e la creazione di mappe'', con qualche esercitazione insieme. Dopo
due lezioni sull'editor e sulla logica (piazzare oggetti, i \emph{transform}, i
Blueprint, la pallina sulla piattaforma), questa lezione guarda a ciò che ``fa il
grafico all'interno di Unreal'': le \textbf{mesh}, le \textbf{texture}, i
\textbf{materiali}, e ``qualche nodo, qualche cosa carina che potete fare'' combinando
mesh e materiali.

La lezione ha due metà separate da una pausa. La prima è sui \textbf{materiali}: che
cos'è una texture e come si avvolge su un oggetto attraverso le \textbf{UV}; quali tipi
di mappa esistono (\textbf{Diffuse}, \textbf{Normal}, \textbf{Specular}, la mappa di
altezza per il \emph{displacement}) e perché in queste mappe il bianco vale 1 e il nero
0; che cos'è un \textbf{materiale} e come si costruisce nell'\textbf{editor dei
materiali}, ``nello stesso modo dell'editor dei Blueprint''; le impostazioni principali
(\textbf{Material Domain}, \textbf{Blend Mode}, \textbf{Shading Model}); i nodi di base
(\textbf{Texture Sample}, le costanti, le operazioni, il \textbf{Texture Coordinate},
il \textbf{Lerp}, il \textbf{Panner}); il rendering \textbf{PBR} con Metallic,
Specular e Roughness; e le \textbf{Material Instance}, il modo ``bon ton'' di applicare i
materiali. La seconda metà è sulle \textbf{mesh}: si importa una \textbf{torretta} in
tre pezzi dal materiale didattico, si guarda la finestra di \textbf{import}, si
corregge il \textbf{pivot} del cannone con un \textbf{Reimport}, si esplorano lo Static
Mesh Editor e le \textbf{collisioni}, si fondono i pezzi con \textbf{Merge Actors}, si
creano a mano i materiali della torretta, e si chiude con i materiali \textbf{World
Aligned}, che non si deformano quando si stira l'oggetto. In fondo, un esercizio libero:
un materiale di \textbf{lava}.

La lezione è \textbf{online}, su Zoom, come le precedenti, e finisce poco prima delle
19 (``si sono fatte le 7''). La pausa, una decina di minuti, arriva dopo ``un'oretta e
mezza''. Davide, alla sua prima volta con questa classe, combatte per tutta la lezione
con la condivisione dello schermo (``io sto facendo a cazzotti da tutta la lezione con
la finestra di Zoom'', ``quanto non sono abituato a streamare'') e con un PC di
emergenza: ``il fisso mi è esploso letteralmente tre giorni fa'', e a metà della seconda
parte il computer ``inizia ad accusare la stream'', tanto che Davide gira il
ventilatore verso di lui. Chiede spesso conferma che gli studenti vedano quello che
mostra: ``mi impanicherò ogni tanto a caso e vi chiederò conferma delle robe''.

Tre avvertenze di lettura, le stesse delle prime due lezioni. La prima: la trascrizione
è automatica e storpia moltissimi nomi. Unreal diventa ``ArrayLink'', ``ArrayL4'',
``Arial Engine'' e ``a Real''; la \textbf{Material Instance} è quasi sempre
``materialistance'', ``Material Distance'' e una volta ``mie dire distanze''; il
\textbf{Panner} è ``panel'', ``planner'' e ``banner''; il \textbf{Lerp} è ``l'ERP'';
\emph{World Aligned} è ``word aligned'' e ``word.allocation''; la texture
\textbf{T\_Bark} (una corteccia) è ``Arc'', ``barca'', ``Barc'' e ``park''; ``strecciare''
e ``si straccia'' stanno per lo \emph{stretch}, lo stiramento; ``traccare'' sta per
\emph{tracciare} un bersaglio (\emph{to track}); la \textbf{mappa di altezza}
(\emph{height map}) è quasi sempre ``light map'' o ``lightmap''. Dove la ricostruzione
è sicura il capitolo usa la forma corretta e, la prima volta, lo segnala in nota; dove
non lo è riporta la forma della trascrizione e lo dice. La seconda: durante la pausa e
verso la fine il riconoscimento vocale produce lunghi blocchi di testo senza senso e
ripetuto (``qual è la risposta che ci può essere?'', discorsi sulla ``vita data per
scontata'' e su ``vittime''): sono artefatti, non parole di nessuno, e il capitolo li
ignora. La terza: gran parte della lezione è una dimostrazione dal vivo nell'editor dei
materiali e nello Static Mesh Editor. Il capitolo ricostruisce i grafi di nodi da ciò
che Davide dice mentre li costruisce, e li completa, dove serve, con i nomi esatti dei
nodi e dei comandi di Unreal Engine 5.

\paragraph{Come leggere questo capitolo.} Il capitolo segue l'ordine della lezione. Le
figure ricostruiscono gli schemi spiegati a voce, sulle slide o nell'editor: le mappe
della torretta, i canali di una texture, il rapporto tra texture, materiale e mesh, i
grafi dei materiali costruiti a lezione, le Material Instance, la struttura di una
mesh, il pivot del cannone. Le tabelle riassumono i tipi di mappa, le impostazioni del
materiale, le opzioni di import, le collisioni e le scorciatoie dell'editor dei
materiali. Qua e là, dei paragrafi di \textbf{approfondimento} collegano la lezione al
libro \emph{Blueprints Visual Scripting for Unreal Engine 5} di Marcos Romero e Brenden
Sewell, il testo di riferimento per questa parte del libro, e a qualche nozione tecnica
generale su Unreal Engine 5; sono segnalati come tali e non fanno parte della lezione.

Il percorso del capitolo:
\begin{enumerate}
  \item l'attesa, le presentazioni e il programma; che cosa fa un grafico dentro
        Unreal; il materiale didattico (sezioni \ref{ue3:sec:attesa}
        e~\ref{ue3:sec:didattico});
  \item le \textbf{texture}: la carta regalo, le \textbf{UV}, i tipi di mappa, il
        bianco e il nero, i \textbf{canali RGB} (sezioni da~\ref{ue3:sec:texture}
        a~\ref{ue3:sec:canali});
  \item i \textbf{materiali}: la vernice, l'editor, Material Domain, Blend Mode e
        Shading Model, i nodi (sezioni da~\ref{ue3:sec:materiale}
        a~\ref{ue3:sec:nodi});
  \item in Engine: la mappa vuota, dentro una texture, il \textbf{Texture Sample}, le
        costanti e le scorciatoie, il colore e l'\textbf{emissive} (sezioni
        da~\ref{ue3:sec:mappavuota} a~\ref{ue3:sec:emissive});
  \item il \textbf{PBR} e la Normal (sezione \ref{ue3:sec:pbr});
  \item il \textbf{Texture Coordinate}, il tiling e il pavimento (sezione
        \ref{ue3:sec:tiling});
  \item le \textbf{Material Instance} e le Material Function (sezioni
        \ref{ue3:sec:instance} e~\ref{ue3:sec:function});
  \item il \textbf{Lerp} e il \textbf{Panner}, poi la pausa (sezioni
        da~\ref{ue3:sec:lerp} a~\ref{ue3:sec:pausa});
  \item l'\textbf{import} e la torretta; che cos'è una \textbf{mesh}; dimensioni e
        formati (sezioni da~\ref{ue3:sec:import} a~\ref{ue3:sec:regole});
  \item il \textbf{pivot} del cannone e il \textbf{Reimport} (sezione
        \ref{ue3:sec:pivot});
  \item lo \textbf{Static Mesh Editor} e le \textbf{collisioni}; raggruppare e
        \textbf{Merge Actors} (sezioni \ref{ue3:sec:sme} e~\ref{ue3:sec:merge});
  \item i \textbf{materiali della torretta} e le domande sul Reimport e sul merge
        (sezioni \ref{ue3:sec:matturret} e~\ref{ue3:sec:domande});
  \item lo stretch, lo \textbf{z-fighting} e i materiali \textbf{World Aligned}
        (sezione \ref{ue3:sec:world});
  \item l'esercizio della \textbf{lava}, il nodo \textbf{Noise} e la chiusura
        (sezione \ref{ue3:sec:lava}).
\end{enumerate}

% ---------------------------------------------------------------------
\section{L'attesa, le presentazioni e il programma}
\label{ue3:sec:attesa}

``Ok, buonasera.'' La registrazione parte subito (``sta registrando? Sì, sta
registrando, perfetto''), e Davide aspetta ``qualche minutino'' che entrino tutti. Uno
studente avvisa che dovrà staccarsi venti minuti prima per andare a lavorare: ``nessun
problema'', tanto la lezione resta registrata, e anzi a Davide piacerebbe finire ``con
voi che fate qualcosina'', quindi andare via un po' in anticipo ``non è neanche un
problema''. Davide chiede quante persone c'erano alle lezioni precedenti (``una dozzina,
credo''), ringrazia chi accende la telecamera (``almeno so che non parlo da solo'') e
controlla che non si senta il ventilatore in sottofondo (``inizia a far caldo, inizia a
essere complicato'').

\paragraph{Che cosa avete fatto con Igor?} Prima di cominciare, Davide vuole sapere a
che punto è la classe, per non ripetere cose già viste. Dal Discord ha capito che si è
``ancora all'inizio'': le schede, le finestre, i \emph{transform} sugli oggetti, i
tavolini, ``qualcosa con gli input, con la pallina che si muove''. ``Se vediamo
qualcosa che avete già visto, ditemelo.''

\paragraph{Il programma.} La lezione di oggi guarda a quello che fa ``il grafico
all'interno di Unreal'': ``un po' di nozioni base per quanto riguarda le mesh, un po' di
nozioni base per quanto riguarda i materiali, le texture'', e ``già qualche nodo,
qualche cosa carina che potete fare'' usando mesh e materiali.

\paragraph{Chi c'è.} Davide chiede a che cosa sono iscritti gli studenti, ``per capire
chi ho davanti''. Le risposte: \textbf{game design} (più d'uno), \textbf{3D art}
(tre persone), \textbf{character animation} (due), \textbf{narrative design}. ``Beh, in
realtà siamo abbastanza assortiti.''

\paragraph{Che cosa fa un grafico dentro Unreal.} Ai futuri grafici 3D Davide mette
``le mani avanti'': la maggior parte del lavoro di un grafico \textbf{non} si svolge
dentro Unreal, ma in altri strumenti, che verranno spiegati nei rispettivi corsi.
Quello che si vede oggi è ciò che un grafico fa ``nel momento in cui va a lavorare con
altre persone'' dentro Unreal Engine: \textbf{importare} gli oggetti e creare gli
\textbf{shader} dentro l'engine.

\paragraph{Come si struttura la lezione.} Davide chiede come lavorava Igor: tutta
teoria e poi tutta pratica, o un po' e un po'? ``Andava prima con la presentazione e poi
con Unreal.'' Davide prova a replicare: ``ci facciamo prima una chiacchiera sulla parte
un po' più teorica e poi passiamo al progetto''. Nella pratica, come si vedrà, teoria
ed engine si alterneranno di continuo.

% ---------------------------------------------------------------------
\section{Il materiale didattico}
\label{ue3:sec:didattico}

Sul Discord c'è stato ``un po' di casino'' per il \textbf{materiale didattico}. ``Non è
molta roba'': è un piccolo progetto di Unreal, ``un progetto che Unreal usa per
spiegare alcune cose'', e probabilmente quello che Igor farà usare nelle lezioni
successive ogni volta che si implementa qualcosa. Gli altri file della cartella sono
\textbf{asset} da importare nel progetto poco alla volta: oggi la \textbf{torretta};
quando si farà l'audio, l'audio; ``poi vediamo se aggiungere altre cose''.

Dove si trova, Davide lo spiega più avanti a una studentessa: sul Discord, ``il 19
maggio, se non ricordo male'', c'è un link a una pagina di Drive con un file che si
chiama, più o meno, ``materiale didattico Unreal Engine Intro''\footnote{La trascrizione
dice ``IVE, materiale didattico, a Real Engine Intro'' e, poco prima, ``Hive Unreal
Engine 5 intro'': la parola iniziale non è ricostruibile con certezza.}. Lo si
scarica, lo si \textbf{decomprime}, e dentro c'è ``un progettino con un
\emph{character}'': sono soprattutto cose per le prossime lezioni con Igor, e per oggi
serve poco.

% ---------------------------------------------------------------------
\section{Che cos'è una texture}
\label{ue3:sec:texture}

\paragraph{La carta regalo.} ``Che cos'è una texture? La texture fondamentalmente è il
\emph{wrap}, è la \textbf{carta regalo} in cui è avvolto il nostro oggetto.'' Sulla
slide c'è la torretta del materiale didattico e, accanto, le sue texture: un'immagine
``tutta scomposta e segmentata e divisa in varie parti''. Perché divisa in parti?
Perché quelle parti poi ``si avvolgeranno sopra la mesh''. Le texture sono ``una serie
di mappe, una serie di informazioni che vanno a dare delle qualità specifiche
all'oggetto''. Con un'altra immagine: la texture ``non è altro che una
\textbf{decalcomania}'' dell'oggetto, della \emph{mesh} che si mette in scena.

\paragraph{Le UV.} Perché una texture funzioni bene su un oggetto, deve essere pensata
``con in mente le \textbf{UV} dei modelli su cui si lavora''. Le UV ``sono le
coordinate'', tutto quello che riguarda ``la struttura dell'oggetto'': la \textbf{mappa
delle UV} è la \textbf{prima parte} del processo di creazione di una texture. Dal
modello 3D si estraggono informazioni che poi servono a creare ``tutti i vari tipi di
texture che esistono'': sulla slide le texture della torretta sono diverse, e ognuna
porta informazioni diverse, ``ma sono tutte texture pensate per la stessa torretta''.

\paragraph{L'UV mapping.} L'\textbf{UV mapping} è la tecnica che serve a
\textbf{mappare le texture 2D su modelli 3D}: si assegnano le coordinate dette UV ``ai
vertici e agli \emph{edge} del modello'', e tutte le texture applicate dopo usano quelle
informazioni ``per precisione''. È un lavoro che si fa \textbf{prima} di entrare in
Unreal, in un altro software ``come Blender, come Maya''. Se si usa una texture che non
è stata pensata per le UV di quell'oggetto, ``la texture verrà strecciata''
(\emph{stretched}, stirata) o comunque non darà il risultato voluto.

\begin{figure}[htbp]
\centering
\begin{tikzpicture}[
  font=\small,
  ue3box/.style={draw, thick, align=center, minimum height=0.9cm},
  ue3arr/.style={-{Stealth[length=2.5mm]}, thick}
]
  % il cubo 3D
  \draw[thick, fill=orange!25] (0,0) -- (1.6,0) -- (1.6,1.6) -- (0,1.6) -- cycle;
  \draw[thick, fill=orange!40] (0,1.6) -- (0.7,2.2) -- (2.3,2.2) -- (1.6,1.6) -- cycle;
  \draw[thick, fill=orange!15] (1.6,0) -- (2.3,0.6) -- (2.3,2.2) -- (1.6,1.6) -- cycle;
  \node[font=\footnotesize, align=center] at (1.15,-0.5) {modello 3D\\(la mesh)};
  % l'apertura
  \draw[ue3arr] (2.8,1.1) -- node[above, font=\footnotesize]{si ``apre''}
    node[below, font=\footnotesize]{(UV mapping)} (5.0,1.1);
  % la mappa UV a croce
  \begin{scope}[shift={(5.4,-0.6)}]
    \draw[step=0.75, gray!40, thin] (0,0) grid (3.75,3.0);
    \foreach \x/\y in {0.75/0, 0/0.75, 0.75/0.75, 1.5/0.75, 2.25/0.75, 0.75/1.5}
      \draw[thick, fill=orange!25] (\x,\y) rectangle ++(0.75,0.75);
    \draw[->, thick] (0,0) -- (4.0,0) node[right, font=\footnotesize]{U};
    \draw[->, thick] (0,0) -- (0,3.25) node[above, font=\footnotesize]{V};
    \node[font=\footnotesize, align=center] at (1.9,-0.5) {mappa UV: le facce\\
      stese sul piano $[0,1]^2$};
  \end{scope}
  % la texture
  \draw[ue3arr] (9.4,1.1) -- node[above, font=\footnotesize]{si dipinge} (10.6,1.1);
  \node[ue3box, minimum width=2.3cm, minimum height=2.3cm, fill=orange!10,
    text width=2.1cm] (tx) at (12.0,1.1) {\footnotesize texture 2D\\
    \footnotesize (Diffuse, Normal,\\\footnotesize Roughness\ldots)\\[3pt]
    \footnotesize si ``riavvolge''\\\footnotesize sulla mesh};
\end{tikzpicture}
\caption{La texture come ``carta regalo'': le facce del modello vengono stese su un
piano (la mappa UV, fatta in Blender o in Maya), le texture si dipingono su quel piano,
e poi si ``riavvolgono'' sull'oggetto. È per questo che la texture della torretta, sulla
slide, appare ``tutta scomposta e segmentata''. Schema ricostruito; le lettere U e V
indicano i due assi della texture, per distinguerli da X, Y e Z del mondo.}
\label{ue3:fig:uv}
\end{figure}

\paragraph{Approfondimento: perché U e V.} Le lettere U e V sono scelte solo per non
confondere gli assi della texture con gli assi X, Y e Z dello spazio 3D. Una coordinata
UV è una coppia di numeri, di solito tra 0 e 1: $(0,0)$ è un angolo della texture e
$(1,1)$ l'angolo opposto. Ogni vertice della mesh riceve una coppia $(u,v)$, e per i
punti interni a una faccia il motore interpola le coordinate dei vertici. Se una mesh
non ha UV (o ne ha di sbagliate), il motore non sa ``quale pezzo di carta regalo'' va
su quale faccia: di qui le texture stirate di cui parla Davide.

% ---------------------------------------------------------------------
\section{I tipi di mappa}
\label{ue3:sec:mappe}

Davide passa in rassegna, sulla slide della torretta, le texture una per una, ``piano
piano'', con le informazioni che portano.

\paragraph{La Diffuse, o Albedo.} La prima è la \textbf{Diffuse}, ``in alcuni casi
detta anche \textbf{Albedo}''. Nei file la si riconosce dalla \textbf{nomenclatura}:
``T maiuscolo, \emph{underscore}, nome dell'oggetto, \emph{underscore}, A'', o D: A per
\emph{albedo}, D per \emph{diffuse} (per esempio \texttt{T\_Turret\_D}). La Diffuse è
il \textbf{colore}: ``tutti i dettagli estetici colorati''. Sulla slide è quella in alto
a destra, con l'arancione: è la texture che dà alla torretta l'arancione sulla base, il
nero negli altri dettagli e il bianco in alcune zone. Più avanti Davide ne completa la
descrizione: la Diffuse applica ``i colori e i dettagli di base'' sulle superfici 3D, e
dà l'informazione del colore, del tipo di materiale, dei dettagli (``queste linee che
vedete qua'') e di quanto un oggetto sia ruvido o levigato, se per esempio si applica la
texture di una roccia o di un metallo.

\paragraph{La Normal.} Poi c'è la \textbf{Normal}, la \textbf{Normal Map}, che ``serve
a simulare una tridimensionalità'' dentro un'informazione che è 2D. Si riconosce
subito: è \textbf{blu}, ``sempre mappe blu con delle sfumature di rosso e di verde''.
L'esempio classico è un \textbf{muro di mattoni}: in gioco il muro è \textbf{piatto}, ma
quando gli si applica la Normal il calcolo delle luci e delle ombre ``prende in
considerazione la normal map'' e proietta le ombre in base a lei, dando l'illusione che
i mattoni siano davvero tridimensionali. La Normal dà informazioni sulla
tridimensionalità ``senza appesantire particolarmente il calcolo'', perché non deforma
davvero l'oggetto come fa una mappa di altezza.

Più avanti Davide ci torna con un esempio sulla torretta. Su una delle gambe della base
servono dei \textbf{graffi}. Si potrebbe cambiare il modello aggiungendo i graffi, ma
questo aggiungerebbe \textbf{poligoni}, e quindi ``potenza di calcolo non
necessaria''. Si può invece aggiungere il dettaglio ``direttamente in 2D'', con la
Normal: protuberanze, graffi, rughe. ``Per un singolo oggetto non è un appesantimento
così importante. Ma considerate quanti oggetti ci sono all'interno di una mappa quando
andate a creare un livello'': se ogni oggetto non è creato con la giusta accortezza,
``si fa molto velocemente ad appesantire il progetto più di quanto dovrebbe essere
pesante''.

\paragraph{La Specular.} Lo \textbf{Specular mapping} riguarda ``la brillantezza e la
riflettività dell'oggetto'': come l'oggetto gestisce la luce, come la luce ``rimbalza''
su di lui. Un oggetto può avere parti metalliche, parti in pietra e altre ancora, e ``a
seconda del tipo di materiale di cui l'oggetto è fatto, rifrangerà la luce in una
maniera diversa''. Queste mappe sono quelle ``in bianco e in nero''. La texture
determina come si calcolano i riflessi ``sui singoli pixel proprio dell'oggetto'',
quanta luce viene riflessa; tornerà utile quando si vedranno le luci in Unreal, perché
``cambiando un materiale la luce ha un effetto completamente diverso''.

\paragraph{La mappa di altezza.} Poi c'è la mappa che serve a creare
\textbf{displacement}\footnote{Nella trascrizione è sempre ``light map'' o
``lightmap'', ma la descrizione (rialza e abbassa zone dell'oggetto, crea
\emph{tessellation}, è ``una normal sotto steroidi'') è quella della \textbf{height
map}, la mappa di altezza o di \emph{displacement}. Le \emph{lightmap} vere sono
un'altra cosa: texture in cui il motore salva l'illuminazione precalcolata.}: crea
``una \emph{tessellation} nell'oggetto'', rialza alcune zone, ne abbassa altre,
``modifica leggermente l'oggetto''. Davide la chiama ``una normal sotto steroidi'',
perché fa davvero \emph{displacement}: aggiunge profondità e dettagli tridimensionali
alle superfici piane, e usa le scale di grigio per creare scanalature e protuberanze.
``Nelle parti in cui vedete il bianco è sempre il sì'', la protuberanza; nel nero,
l'effetto opposto.

\paragraph{Il texturing procedurale.} Infine il \textbf{procedural texturing}: il
processo di creare texture, materiali o \textbf{shader} ``attraverso degli
algoritmi''. ``Andando ad utilizzare dei calcoli matematici, noi andiamo a modificare
delle informazioni per avere un risultato.'' È esattamente quello che fa l'editor dei
materiali di Unreal.

\begin{table}[htbp]
\centering
\small
\begin{tabular}{@{}p{2.5cm}p{2.4cm}p{8.1cm}@{}}
\toprule
\textbf{Mappa} & \textbf{Aspetto} & \textbf{Che cosa porta} \\
\midrule
Diffuse (Albedo) & a colori & il colore e i dettagli estetici; suffisso \_D o \_A \\
Normal & blu, con sfumature rosse e verdi & l'illusione del rilievo: la luce e le
  ombre si calcolano come se la superficie avesse graffi, mattoni, rughe, senza
  aggiungere poligoni \\
Specular & bianco e nero & quanto e come la superficie riflette la luce \\
Roughness & bianco e nero & quanto la superficie è ruvida (1) o liscia (0) \\
Metallic & bianco e nero & dove la superficie è metallo (1) e dove no (0) \\
Ambient Occlusion & bianco e nero & le zone in ombra perché ``chiuse'' (pieghe,
  angoli) \\
Height (displacement) & bianco e nero & deforma davvero la superficie: il bianco
  rialza, il nero abbassa \\
\bottomrule
\end{tabular}
\caption{I tipi di mappa nominati a lezione. Diffuse, Normal, Specular e la mappa di
altezza sono quelle della slide; Roughness, Metallic e Ambient Occlusion compaiono
nella parte sulle mappe impacchettate e nei materiali della torretta.}
\label{ue3:tab:mappe}
\end{table}

\paragraph{Il bianco è 1, il nero è 0.} La cosa importante da capire, insiste Davide, è
che nelle mappe in bianco e nero ``dovete sempre considerare il \textbf{bianco come un
1}, il \textbf{nero come uno 0} e il \textbf{grigio come dei valori nel mezzo} fra lo 0 e
l'1''. Che cosa vuol dire? Che in Unreal, per esempio, il valore dello Specular va
soltanto tra 0 e 1: a 1 riflette la luce, a 0 no. Una mappa di Specular, quindi, va
letta come una griglia di numeri: ``ci sono degli 1 in queste parti e degli 0 nel resto
delle parti''. È la chiave per capire quasi tutto il resto della lezione.

\paragraph{Approfondimento: la mappa di altezza, la Normal e Unreal Engine 5.} La
differenza tra Normal e mappa di altezza è la differenza tra \emph{sembrare} e
\emph{essere}. La Normal Map memorizza, per ogni pixel, la direzione in cui ``guarda''
la superficie (la sua \emph{normale}): i tre numeri X, Y e Z della direzione sono
salvati nei canali rosso, verde e blu, e siccome la maggior parte dei pixel guarda
``dritto fuori'' dalla superficie (Z alto) la mappa appare blu. La geometria non
cambia: di profilo, un muro con Normal Map resta piatto. La mappa di altezza invece
sposta davvero i vertici, e per farlo servono vertici da spostare: di qui la
\emph{tessellation}, che suddivide la superficie in molti triangoli. In Unreal Engine 5
il \emph{displacement} vero si fa soprattutto con le mesh \textbf{Nanite}, ed è più
costoso della Normal: per questo, per graffi e piccoli dettagli, la Normal resta la
scelta di default, come dice Davide.

% ---------------------------------------------------------------------
\section{I canali RGB e le mappe impacchettate}
\label{ue3:sec:canali}

``Chiaramente, queste texture utilizzano l'\textbf{RGB}'': non hanno un solo canale, ma
\textbf{tre canali} su cui passano le informazioni. Sulla slide c'è una Diffuse
scomposta in rosso, verde e blu: ognuno porta informazioni separate, che combinate
danno il colore che serve.

Ma c'è un altro modo di vedere la cosa. \textbf{Singolarmente}, ogni canale appare come
un'immagine in \textbf{bianco, nero e grigio}. E bianchi, neri e grigi ``sono
combinazioni di 0 e di 1'': quindi ogni canale si può usare per trasmettere
\textbf{un'informazione diversa}. È il \textbf{multichannel}. Una pratica comune,
nell'esportare e importare texture, è avere una singola mappa ``che spesso e volentieri
viene siglata come \textbf{AORM}'', per \textbf{Ambient Occlusion, Roughness e
Metallic}: tre mappe in bianco e nero, ognuna in un canale. ``Magari noi vogliamo
utilizzare questa mappa qui sopra per l'ambient occlusion, questa mappa qui sotto per
le metallic'', e così via. Una sola mappa con più informazioni invece di tante mappe
diverse: ``va ad alleggerire il progetto, va a rendere più veloce, più modulare''
l'export delle texture.

\begin{figure}[htbp]
\centering
\begin{tikzpicture}[
  font=\small,
  ue3ch/.style={draw, thick, minimum width=2.4cm, minimum height=1.0cm,
                align=center},
  ue3arr/.style={-{Stealth[length=2.5mm]}, thick}
]
  \node[draw, very thick, minimum width=2.6cm, minimum height=3.6cm, align=center,
    fill=gray!10] (tex) at (0,0) {\textbf{una texture}\\\textbf{RGB}\\[4pt]
    \footnotesize es. T\_Turret\_ORM};
  \node[ue3ch, fill=red!15] (r) at (4.2,1.3) {canale \textbf{R}};
  \node[ue3ch, fill=green!15] (g) at (4.2,0) {canale \textbf{G}};
  \node[ue3ch, fill=blue!15] (b) at (4.2,-1.3) {canale \textbf{B}};
  \draw[ue3arr] (tex.east) ++(0,1.0) -- (r.west);
  \draw[ue3arr] (tex.east) -- (g.west);
  \draw[ue3arr] (tex.east) ++(0,-1.0) -- (b.west);
  \node[ue3ch, text width=3.4cm] (ao) at (8.6,1.3) {\textbf{Ambient Occlusion}\\
    \footnotesize grigi tra 0 e 1};
  \node[ue3ch, text width=3.4cm] (ro) at (8.6,0) {\textbf{Roughness}\\
    \footnotesize grigi tra 0 e 1};
  \node[ue3ch, text width=3.4cm] (me) at (8.6,-1.3) {\textbf{Metallic}\\
    \footnotesize grigi tra 0 e 1};
  \draw[ue3arr] (r) -- (ao);
  \draw[ue3arr] (g) -- (ro);
  \draw[ue3arr] (b) -- (me);
  \node[font=\footnotesize, align=center, text width=11cm] at (4.3,-2.5)
    {ogni canale, da solo, è un'immagine in bianco e nero: tre ``mappe'' diverse in
    un solo file};
\end{tikzpicture}
\caption{Una mappa impacchettata (\emph{channel packing}). La lezione la chiama AORM;
nella pratica di Unreal la convenzione più diffusa è ORM, con l'Ambient Occlusion nel
rosso, la Roughness nel verde e il Metallic nel blu, ma l'ordine dipende da chi
esporta le texture.}
\label{ue3:fig:canali}
\end{figure}

\paragraph{Approfondimento: sRGB e compressione.} Quando si importa in Unreal una
texture impacchettata (o qualunque mappa ``di dati'', come Roughness o Metallic), conviene
aprirla e \textbf{togliere la spunta a sRGB}: l'opzione sRGB dice al motore che
l'immagine è un \emph{colore} da correggere per la percezione dell'occhio, mentre una
mappa di dati contiene numeri da leggere così come sono. Nelle stesse impostazioni, il
campo \emph{Compression Settings} si mette su \textbf{Masks} per le mappe impacchettate
e su \textbf{Normalmap} per le Normal (Unreal di solito lo riconosce da solo per le
texture con suffisso \_N). Se ci si dimentica, il materiale funziona lo stesso, ma i
grigi vengono letti un po' sbagliati e la superficie risulta più lucida o più opaca del
previsto.

% ---------------------------------------------------------------------
\section{Il materiale: la vernice}
\label{ue3:sec:materiale}

\paragraph{Le texture non si modificano.} ``Quello che possiamo fare con le texture in
Unreal è relativamente poco, perché non possiamo andare a modificare effettivamente le
texture all'interno dell'engine.'' E allora come si danno qualità visive alle mesh?
``Lo facciamo con i \textbf{materiali}.''

\paragraph{Il materiale è uno shader.} I materiali ``sono degli \textbf{shader}''. E
l'\textbf{editor dei materiali}, quello che si vedrà ``principalmente'' oggi, ``funziona
praticamente nello stesso modo dell'editor dei Blueprint'': una serie di \textbf{nodi}
collegati, che trasportano informazioni fino a un \textbf{nodo risultante}.

\paragraph{Vernice, non carta regalo.} La distinzione è netta: ``il materiale è
effettivamente la \textbf{vernice} che viene applicata sulla mesh, non la texture''. Sugli
oggetti si applicano \textbf{materiali}. Le texture servono a \textbf{costruire} i
materiali: si combinano le texture in un unico materiale, che poi si applica
sull'oggetto finale. Il materiale dice all'oggetto ``come dovrebbe interagire con la
luce, che tipo di informazioni visive dovrebbe portare'', dove ci sono
tridimensionalità e così via: ``tutto quello che abbiamo visto dalle varie texture
precedentemente viene tutto convogliato all'interno dei materiali''. Questo, precisa
Davide, se si parla di \textbf{oggetti fisici}. I materiali hanno anche altre funzioni:
servono per esempio a fare \textbf{effetti speciali}, e un materiale per un effetto
speciale ``dovrà essere costruito in maniera diversa'' da quello che riveste una
torretta.

\begin{figure}[htbp]
\centering
\begin{tikzpicture}[
  font=\small,
  ue3tex/.style={draw, thick, minimum width=2.3cm, minimum height=0.7cm,
                 align=center, fill=gray!10},
  ue3big/.style={draw, very thick, rounded corners=4pt, align=center},
  ue3arr/.style={-{Stealth[length=2.5mm]}, thick}
]
  \node[ue3tex] (t1) at (0,1.5) {T\_Turret\_D};
  \node[ue3tex] (t2) at (0,0.5) {T\_Turret\_N};
  \node[ue3tex] (t3) at (0,-0.5) {T\_Turret\_R};
  \node[ue3tex] (t4) at (0,-1.5) {T\_Turret\_M};
  \node[font=\footnotesize, align=center] at (0,2.4) {\textbf{texture}\\(immagini)};
  \node[ue3big, minimum width=3.4cm, minimum height=3.6cm, fill=blue!6,
    text width=3.2cm] (mat) at (4.4,0) {\textbf{materiale}\\M\_Turret\\[6pt]
    \footnotesize un grafo di nodi\\\footnotesize (uno \emph{shader})\\[6pt]
    \footnotesize la ``vernice''};
  \foreach \tn in {t1,t2,t3,t4} \draw[ue3arr] (\tn.east) -- (mat.west |- \tn);
  \node[ue3big, minimum width=2.8cm, minimum height=1.6cm, fill=orange!12,
    text width=2.6cm] (mi) at (8.4,0) {\textbf{Material}\\\textbf{Instance}\\
    \footnotesize MI\_Turret\\\footnotesize (parametri)};
  \draw[ue3arr] (mat) -- (mi);
  \node[ue3big, minimum width=2.4cm, minimum height=1.6cm, fill=green!8,
    text width=2.2cm] (mesh) at (11.8,0) {\textbf{mesh}\\\footnotesize SM\_Turret\\
    \footnotesize slot del\\\footnotesize materiale};
  \draw[ue3arr] (mi) -- node[above, font=\footnotesize]{si applica} (mesh);
\end{tikzpicture}
\caption{Texture, materiale, istanza e mesh. Le texture non si applicano direttamente
agli oggetti: entrano in un materiale, che è la ``vernice''; dal materiale si crea
un'istanza, ed è l'istanza che si applica allo slot del materiale della mesh
(sezione~\ref{ue3:sec:instance}). I nomi seguono la nomenclatura suggerita a lezione
(T\_, M\_, MI\_, SM\_).}
\label{ue3:fig:pipeline}
\end{figure}

\paragraph{Approfondimento: il materiale nel libro di Romero e Sewell.} Il quinto
capitolo del libro (``Object Interaction with Blueprints'') usa la stessa immagine di
Davide: un materiale è un \emph{asset} che si aggiunge alla mesh di un actor per
crearne l'aspetto, ``come uno strato di vernice''. Il libro crea un materiale rosso,
\texttt{M\_TargetRed}, in una cartella Materials (``non è necessario, ma è buona
pratica tenere ordinata la gerarchia del progetto''), e nota che l'editor dei
materiali ``condivide molte caratteristiche e convenzioni con i Blueprint'': al centro
c'è il grafo; il nodo grande con il nome del materiale è il \textbf{result node}; con il
tasto destro si apre un menu di nodi con una ricerca sensibile al contesto; i fili si
tirano trascinando dai pin, gli ingressi a sinistra e le uscite a destra. Il colore lo
dà un \textbf{VectorParameter} collegato a Base Color; due \textbf{ScalarParameter} a
0,1 danno un po' di Metallic e di Roughness. Poi, in un Blueprint, il nodo \textbf{Set
Material} cambia il materiale di un bersaglio quando viene colpito: un primo esempio di
come la logica (le lezioni di Igor) e la grafica (questa lezione) si incontrino.

% ---------------------------------------------------------------------
\section{Material Domain, Blend Mode e Shading Model}
\label{ue3:sec:dettagli}

Come gli oggetti nella \emph{viewport}, anche il materiale ha i suoi
\textbf{Details}. Dentro l'editor si possono selezionare i singoli nodi, per i dettagli
del nodo, oppure vedere i \textbf{dettagli generali del materiale}. Tra questi ce ne
sono alcuni, evidenziati sulla slide, che ``cambiano proprio effettivamente come
funziona il materiale'' e ne definiscono la struttura e le possibilità.

\paragraph{Material Domain.} Il \textbf{Material Domain} dice \textbf{per che cosa}
verrà usato il materiale: una \textbf{superficie}, l'\textbf{interfaccia} (la UI), il
\textbf{post process}, un effetto.

\paragraph{Blend Mode.} La \textbf{Blend Mode} è il modo in cui il materiale
``\emph{blenda}'' (si mescola) con i pixel calcolati nel resto della scena. La modalità
\textbf{Opaque}, quella di default, \textbf{occlude completamente} ciò che c'è dietro.
Per un effetto speciale, per esempio un \textbf{fuoco}, serve invece la Blend Mode
\textbf{Translucent}, che permette di aggiungere parametri per l'\textbf{opacità} del
materiale.

\paragraph{Shading Model.} Lo \textbf{Shading Model}\footnote{Nella trascrizione una
volta ``scaling model''.} definisce come il materiale \textbf{reagisce alla luce}.
\textbf{Default Lit} vuol dire che segue le regole di base dell'illuminazione: il
materiale viene illuminato. Per un effetto speciale, in cui non si vuole che la luce
influenzi la resa, lo si mette invece su \textbf{Unlit}\footnote{Nella trascrizione
``lo metterò un lit''.}. Ci sono ``diverse tipologie'' di Shading Model, a seconda di
cosa si vuole realizzare.

\paragraph{I pin cambiano.} Queste impostazioni si applicano sul nodo risultante, e
cambiarle cambia i \textbf{pin} del nodo. L'\textbf{Opacity} non si può modificare in
Blend Mode Opaque, perché un materiale opaco ``dà per scontato che l'oggetto sia sempre
opaco'': non c'è bisogno di dire dove e quando sarà trasparente. Viceversa, la Blend
Mode Translucent toglie la possibilità di modificare Metallic, Specular, Roughness e
così via: se il materiale deve essere potenzialmente trasparente, ``non ha senso andare
a calcolare come rifletterà la luce''.

\begin{table}[htbp]
\centering
\small
\begin{tabular}{@{}p{2.6cm}p{3.3cm}p{7.1cm}@{}}
\toprule
\textbf{Impostazione} & \textbf{Valori citati} & \textbf{Che cosa decide} \\
\midrule
Material Domain & Surface, User Interface, Post Process\ldots & per che cosa si usa
  il materiale: una superficie, la UI, il post process, un effetto \\
Blend Mode & \textbf{Opaque} (default), Masked, Translucent & come il materiale si
  mescola con ciò che ha dietro: Opaque nasconde tutto; Masked ritaglia con una
  maschera; Translucent sfuma con l'Opacity (fuoco, fili d'erba) \\
Shading Model & \textbf{Default Lit} (default), Unlit\ldots & come reagisce alla luce:
  Default Lit viene illuminato, Unlit no (effetti speciali) \\
\bottomrule
\end{tabular}
\caption{Le tre impostazioni ``più importanti'' dei dettagli di un materiale. Cambiarle
accende e spegne i pin del nodo risultante.}
\label{ue3:tab:dettagli}
\end{table}

\paragraph{La prova in Engine.} Per far vedere la cosa, Davide apre il progetto del
materiale didattico e invita chi vuole a seguirlo: nel \textbf{Content Browser} si crea
un nuovo materiale (``va bene anche chiamarlo New Material''), lo si apre con un
\textbf{doppio clic} e si entra nell'editor. Le impostazioni di cui si parlava sono nei
dettagli, in basso a sinistra. Se la Blend Mode passa da Opaque a \textbf{Translucent},
il nodo risultante cambia: spariscono Metallic, Specular, Roughness, Anisotropy, e
compare \textbf{Opacity}. È utile, per esempio, per una \textbf{fiamma} che verso l'alto
diventa trasparente (le informazioni grafiche alla base, e in alto ``piano piano va a
dissolversi''), o per un \textbf{filo d'erba} che parte con il suo colore e poi va in
trasparenza. Con la Blend Mode \textbf{Masked}, ``mascherata'', si crea invece una
\textbf{maschera}, un \emph{layer} che funziona in modo diverso dal resto della
superficie. Per ora si resta su \textbf{Opaque}, perché i materiali di oggi andranno su
oggetti 3D. A metà spiegazione arriva uno studente in ritardo: ``Chat, buon pomeriggio,
Lorenzo. Noi qui abbiamo già cominciato, ma in realtà abbiamo parlato di cose
principalmente teoriche.''

\paragraph{Approfondimento: Masked o Translucent?} Le due modalità sembrano simili ma
lavorano in modo opposto. \textbf{Masked} decide pixel per pixel se il materiale
\emph{c'è o non c'è}: l'ingresso Opacity Mask viene confrontato con una soglia (di
default 0,333), e sotto la soglia il pixel non viene disegnato affatto. È la modalità
delle foglie, delle reti e dell'erba ``ritagliata'': economica, compatibile con la
luce e con le ombre, ma senza sfumature. \textbf{Translucent} invece \emph{mescola} il
colore del materiale con ciò che sta dietro, secondo un valore di Opacity tra 0 e 1:
permette le sfumature (fumo, vetro, fiamme), ma costa di più, perché ogni strato
traslucido va sommato a quelli dietro, e gestisce peggio luce e ordinamento. Regola
pratica: Masked quando basta un ``sì o no'', Translucent solo quando serve davvero la
sfumatura.

% ---------------------------------------------------------------------
\section{I nodi: le Material Expression}
\label{ue3:sec:nodi}

``Come per i Blueprint, anche l'editor dei materiali ha dei nodi, delle
\textbf{espressioni}.'' I \textbf{Material Expression Node} sono ``dei piccoli set di
istruzioni preprogrammate che eseguono delle cose nello specifico, esattamente come i
nodi dei Blueprint'', o, con un'altra formula, ``pezzi di codice statico''.

\paragraph{Il primo grafo.} Sulla slide c'è un piccolo grafo
(figura~\ref{ue3:fig:primografo}), e Davide lo legge nodo per nodo:
\begin{itemize}
  \item il \textbf{Texture Sample} comunica l'informazione della texture: ``sta dicendo
        che la texture che c'è in questo nodo è la texture che determina il
        \textbf{base color}''; permette di usare un'immagine e farla diventare la
        texture del materiale;
  \item il \textbf{Texture Coordinate}, il nodo \textbf{rosso}, restituisce le
        coordinate U e V del materiale, e serve a decidere ``quanto è grande il
        \textbf{tiling} della mia texture, quindi quante volte si ripete'';
  \item il \textbf{Multiply} è programmato per moltiplicare due valori: ``moltiplica A
        per B e dà una risultante'';
  \item il quadratino è una \textbf{costante}, ``un singolo \emph{float}'', una variabile
        singola.
\end{itemize}
Il grafo dice: ``moltiplica le coordinate per uno e comunica queste cose alle UV della
texture''. Così com'è, quindi, ``è come se non ci fosse'': moltiplicare per 1 lascia le
coordinate di default. Ma è così che si costruiscono i grafi: combinando nodi si
ottengono risultati molto diversi. Si può moltiplicare, si può fare
un'\textbf{interpolazione}, si possono far \textbf{muovere} le texture per suggerire un
movimento nel materiale, ``eccetera, eccetera, eccetera''.

\begin{figure}[htbp]
\centering
\begin{tikzpicture}[
  font=\small,
  ue3node/.style={draw, thick, rounded corners=4pt, align=center,
                  minimum height=0.9cm},
  ue3uv/.style={-{Stealth[length=2mm]}, thick, red!60!black},
  ue3flt/.style={-{Stealth[length=2mm]}, thick, green!50!black},
  ue3col/.style={-{Stealth[length=2mm]}, thick, blue!60!black}
]
  \node[ue3node, fill=red!12, text width=2.2cm] (tc) at (0,0.6)
    {\textbf{TexCoord}\\\footnotesize UTiling 1, VTiling 1};
  \node[ue3node, fill=green!8, minimum width=1.0cm] (k) at (0,-1.0) {\textbf{1}};
  \node[font=\footnotesize] at (0,-1.75) {costante};
  \node[ue3node, fill=gray!10, text width=1.6cm] (mul) at (3.0,0)
    {\textbf{Multiply}\\\footnotesize A $\times$ B};
  \node[ue3node, fill=gray!10, text width=2.4cm, minimum height=1.8cm] (ts)
    at (6.4,0) {\textbf{Texture}\\\textbf{Sample}\\\footnotesize (T\_Bark)};
  \node[ue3node, fill=blue!6, text width=2.4cm, minimum height=2.6cm] (res)
    at (10.4,0) {\textbf{M\_NewMaterial}\\[3pt]\footnotesize\textbf{Base Color}\\
    \footnotesize Metallic\\\footnotesize Specular\\\footnotesize Roughness\\
    \footnotesize Emissive Color\\\footnotesize Normal\ldots};
  \draw[ue3uv] (tc.east) -- ++(0.4,0) |- node[pos=0.75, above, font=\scriptsize]{A}
    ([yshift=2mm]mul.west);
  \draw[ue3flt] (k.east) -- ++(0.6,0) |- node[pos=0.75, below, font=\scriptsize]{B}
    ([yshift=-2mm]mul.west);
  \draw[ue3uv] (mul.east) -- node[above, font=\scriptsize]{UVs}
    (ts.west |- mul.east);
  \draw[ue3col] (ts.east) -- node[above, font=\scriptsize]{RGB}
    (res.west |- ts.east);
\end{tikzpicture}
\caption{Il primo grafo della slide: le coordinate UV, moltiplicate per una costante,
entrano nel Texture Sample, il cui colore va nel Base Color del nodo risultante. Con la
costante a 1 il Multiply ``è come se non ci fosse''; cambiando la costante cambia il
tiling (sezione~\ref{ue3:sec:tiling}). Schema ricostruito.}
\label{ue3:fig:primografo}
\end{figure}

\paragraph{I parametri.} Non solo le costanti, ma i nodi in generale ``possono essere
utilizzati come \textbf{parametri}''. Per rendere un nodo un parametro si fa
\textbf{tasto destro} sul nodo e \textbf{Convert to Parameter}, gli si dà un nome, e
questo permette di modificare in seguito il valore ``senza dover ricalcolare da capo il
materiale''. I parametri torneranno utili con le \textbf{Material Instance}
(sezione~\ref{ue3:sec:instance}).

\paragraph{Come si chiama un nodo.} ``Tasto destro, scriviamo il nome'', per esempio
\emph{texture sample}, ed ecco il nodo. C'è poi ``tutta una serie di \emph{shortcut}''
per i nodi più usati, che si vedranno tra poco.

% ---------------------------------------------------------------------
\section{Il problema della mappa vuota}
\label{ue3:sec:mappavuota}

Una studentessa interrompe: non riesce più ad aprire la mappa, ``me la apre totalmente
vuota, senza titolo''. Davide le chiede di condividere lo schermo, ma lei segue da due
computer diversi. Prima ipotesi: ha creato un \textbf{nuovo livello}, e un \emph{New
Level} appena caricato è tutto nero, perché non c'è niente in scena (``potremmo andare
a popolarla mettendoci cose''). Ma il problema è un altro: ha aperto il progetto del
materiale didattico, cliccato sulla mappa del personaggio in terza persona
(\emph{Third Person Map}), e non succede niente. Davide sospetta che il progetto sia
ancora in una \textbf{cartella zippata} (``di solito questa cosa succede quando il
progetto è una cartella che è ancora zippata''), e le consiglia di
\textbf{riscaricarlo}. Poco dopo il problema è risolto: ``sono riuscita a risolvere
adesso''. ``Meno male, meno male.''

\paragraph{Approfondimento: aprire un progetto da una cartella.} Un progetto di Unreal
è una cartella che contiene un file \texttt{.uproject} e le sottocartelle
\texttt{Content}, \texttt{Config} e altre. Unreal non può aprire un progetto che sta
ancora dentro un archivio \texttt{.zip}: Windows mostra il contenuto dello zip come se
fosse una cartella, ma l'editor non ci può scrivere, e il risultato tipico è proprio
una mappa vuota o mancante. La soluzione è estrarre tutto (``Estrai tutto'') e aprire il
\texttt{.uproject} estratto con un doppio clic. Se il progetto non compare nell'elenco
dell'Epic Games Launcher o nel Project Browser, basta aprirlo dal \texttt{.uproject}
la prima volta: da lì in poi compare tra i progetti recenti.

% ---------------------------------------------------------------------
\section{Dentro una texture e il Texture Sample}
\label{ue3:sec:sample}

``Facciamo le cose in ordine, se no io mi confondo e mi impiccio.'' Davide torna al
punto saltato: com'è fatta \textbf{dentro} una texture.

\paragraph{La finestra della texture.} Con un doppio clic su una texture nel Content
Browser (Davide ne prende una a caso, una corteccia, \textbf{T\_Bark}) si apre una
finestra in cui ``non possiamo fare grandi modifiche'', ma si possono vedere i
\textbf{canali}: il blu, il verde, il rosso, ciascuno fatto a modo suo. ``Questa texture
non è particolarmente interessante'', ma si rivedrà meglio con la torretta.

\paragraph{Il Texture Sample.} Nel grafo del materiale si mette un \textbf{Texture
Sample}, si seleziona il nodo, e nei dettagli, in basso, si sceglie quale texture
metterci. Cercando T\_Bark e mettendola nel nodo, la \emph{preview} mostra la corteccia:
``il nostro materiale sarà così, perché è effettivamente la texture che ci abbiamo
messo''.

\paragraph{Le uscite separate.} Il Texture Sample non ha una sola uscita: oltre
all'\textbf{RGB} ci sono \textbf{R}, \textbf{G} e \textbf{B} separati (e l'alfa, A). Si
può quindi usare, per esempio, solo il canale \textbf{rosso} (``non quella completa, ma
questa qui'') per trasportare un'informazione: collegandolo allo \textbf{Specular} si
ottiene un materiale che in alcune parti riflette la luce più che in altre. ``Con
questo materiale in realtà è un po' complicato farlo vedere'', ma il principio è quello
dei canali visto sulle slide.

\paragraph{Come si mette una texture nel nodo.} Una studentessa non ha capito come si
aggiunge la texture ``al quadratino nero'' del Texture Sample. Davide rifà il
percorso: si clicca il nodo, e nei suoi dettagli c'è il campo \textbf{Texture}, con il
quadratino e una lista. Due modi:
\begin{itemize}
  \item se una texture è già \textbf{selezionata nel Content Browser}, la freccetta
        \textbf{Use Selected Asset from Content Browser} la inserisce automaticamente
        nel nodo;
  \item più spesso, si clicca sul \textbf{nome} della texture nel campo e si fa un
        \emph{browse}: si cerca il nome (``T underscore Bark'') e lo si sceglie dalla
        lista di tutte le texture disponibili.
\end{itemize}
``Voi lo vedete questo pop-up che esce?'' Sì: perfetto.

% ---------------------------------------------------------------------
\section{Le costanti, le operazioni e le scorciatoie}
\label{ue3:sec:costanti}

``Chiaramente ne vediamo un po', non li vediamo tutti, perché sono tanti, sono
veramente tanti.'' Ci sono ``tantissime possibilità, tante cose che vanno studiate'',
ma per ora servono le basi.

\paragraph{La regola delle scorciatoie.} Il Texture Sample si può chiamare cercandolo
con il tasto destro, oppure \textbf{tenendo premuto il tasto T e facendo clic
sinistro} nel grafo. ``Questa struttura dello shortcut funziona per ogni singolo nodo
che ha uno shortcut'': un tasto, tenuto premuto, più clic sinistro. Il tasto è
``molto spesso l'iniziale del nome del nodo''.

\paragraph{Le costanti.} Cercando \emph{vector} compare, per esempio, il Vector
Noise; cercando \emph{constant} compaiono le costanti: \textbf{Constant},
\textbf{Constant2Vector}, \textbf{Constant3Vector} (e Constant4Vector). Le scorciatoie
sono i numeri: \textbf{1 + clic} per la costante singola, \textbf{2 + clic} per quella
a due valori, \textbf{3 + clic} per quella a tre. La costante a due ``la utilizzerete
relativamente poco''. Le altre due moltissimo:
\begin{itemize}
  \item la \textbf{costante singola} è un \textbf{valore numerico}, che si usa ``per
        andare a fare qualsiasi tipo di calcolo, qualsiasi tipo di costruzione'';
  \item la \textbf{costante a tre} è un \textbf{colore}: l'engine usa i tre valori come
        \textbf{R}, \textbf{G} e \textbf{B}. Si sceglie un colore (Davide prende un
        arancione) e l'engine modifica X, Y e Z per ottenerlo. ``Qualcuno di voi l'aveva
        anche già fatto nel test con la pallina.''
\end{itemize}
Collegando l'arancione al \textbf{Base Color} si ottiene un materiale arancione,
``perché effettivamente è l'unica cosa che gli sto passando''. Ma ``molto raramente''
si usa un colore piatto come Base Color: lo si usa piuttosto per fare
\textbf{operazioni}. Più avanti una studentessa chiederà di nuovo come si crea il colore:
``è come la costante singola, che è uno clic sinistro, ma è a tre, quindi è tre clic
sinistro''.

\paragraph{Staccare i fili.} Per rompere un collegamento: \textbf{Alt + clic sinistro}
sul pin, ``e si rompono''.

\paragraph{Le operazioni.} Le operazioni utili sono ``le classiche'': addizione,
moltiplicazione, divisione, sottrazione. Le scorciatoie: \textbf{A + clic} per
\textbf{Add}, \textbf{M + clic} per \textbf{Multiply}, \textbf{D + clic} per
\textbf{Divide}. La sottrazione (\textbf{Subtract}) non ha scorciatoia: si cerca. Il
motivo lo spiega Davide più avanti: \textbf{S + clic} crea automaticamente un
\textbf{parametro} (uno Scalar Parameter), ``quindi questo potrebbe tornarvi utile''.

\paragraph{Quali si usano davvero.} In chiusura di lezione Davide aggiunge un consiglio
d'esperienza: ``l'operazione che utilizzerete di più in assoluto è la
\textbf{moltiplicazione}''. Moltiplicazione e addizione sono le più usate; divisione e
sottrazione servono soprattutto per cose più complesse, ``tipo sugli effetti
speciali''.

% ---------------------------------------------------------------------
\section{Il colore che brilla: Emissive, Multiply e Glow Intensity}
\label{ue3:sec:emissive}

``A cosa ci servono'' le operazioni? Per esempio ad \textbf{aumentare il valore} di
qualcosa. Il nodo risultante ha diversi pin, ``diverse cose su cui andare a collegare
dei valori'': il \textbf{Base Color}, e più in basso l'\textbf{Emissive Color}, che
controlla ``quali parti del materiale emetteranno una piccola luce'', un
\textbf{\emph{glow}}.

\paragraph{L'arancione nell'Emissive.} Davide stacca l'arancione dal Base Color e lo
collega all'Emissive: il materiale ``sarà leggermente luminoso''. ``Adesso è molto poco
luminoso, ma è comunque un emissive.''

\paragraph{Moltiplicare.} Come si ottiene un'emissione maggiore? ``Moltiplichiamo.''
L'arancione entra nella A di un \textbf{Multiply}, il Multiply va nell'Emissive. Il
secondo valore si può scrivere direttamente nel nodo Multiply (``invece di 1 mettiamo
100''): ``verrà decisamente più luminoso''. Ma così ``è meno controllabile''. Meglio
attaccare alla B una \textbf{costante} separata, darle il valore 100 e convertirla in
\textbf{parametro}, con un nome come \textbf{Glow Intensity}: quando più avanti ci sarà
una Material Instance, il valore si potrà cambiare ``direttamente da lì''. ``Ci andiamo
a estrapolare un parametro.''

\begin{figure}[htbp]
\centering
\begin{tikzpicture}[
  font=\small,
  ue3node/.style={draw, thick, rounded corners=4pt, align=center,
                  minimum height=0.9cm},
  ue3col/.style={-{Stealth[length=2mm]}, thick, blue!60!black},
  ue3flt/.style={-{Stealth[length=2mm]}, thick, green!50!black}
]
  \node[ue3node, fill=orange!30, text width=2.5cm] (col) at (0,0.7)
    {\textbf{Constant3Vector}\\\footnotesize arancione (3 + clic)};
  \node[ue3node, fill=green!8, text width=2.5cm] (gi) at (0,-0.9)
    {\textbf{Glow Intensity}\\\footnotesize parametro = 100};
  \node[ue3node, fill=gray!10, text width=1.6cm] (mul) at (3.6,0)
    {\textbf{Multiply}\\\footnotesize (M + clic)};
  \node[ue3node, fill=blue!6, text width=2.6cm, minimum height=2.2cm] (res)
    at (7.4,0) {\textbf{nodo risultante}\\[3pt]\footnotesize Base Color\\
    \footnotesize Metallic, Specular\\\footnotesize Roughness\\
    \footnotesize\textbf{Emissive Color}};
  \draw[ue3col] (col.east) -- ++(0.4,0) |- node[pos=0.75, above, font=\scriptsize]{A}
    ([yshift=2mm]mul.west);
  \draw[ue3flt] (gi.east) -- ++(0.4,0) |- node[pos=0.75, below, font=\scriptsize]{B}
    ([yshift=-2mm]mul.west);
  \draw[ue3col] (mul.east) -- ++(0.6,0) |- ([yshift=-7mm]res.west);
\end{tikzpicture}
\caption{Un materiale che brilla. Il colore, collegato all'Emissive Color, emette ``una
piccola luce''; moltiplicato per 100 diventa molto più luminoso. Il valore 100 sta in
una costante convertita in parametro (Glow Intensity), così che si possa cambiare da
una Material Instance senza toccare il grafo. Ricostruzione dalla dimostrazione.}
\label{ue3:fig:emissive}
\end{figure}

\paragraph{Approfondimento: perché 100 funziona.} I colori di un materiale non sono
limitati tra 0 e 1. Nel Base Color valori oltre 1 non hanno senso fisico (nessuna
superficie riflette più luce di quella che riceve), ma nell'Emissive sì: un'emissione di
$(1; 0{,}5; 0)$ moltiplicata per 100 diventa $(100; 50; 0)$, una sorgente molto intensa
che il \emph{post process} di Unreal trasforma in un alone luminoso (il \emph{bloom}).
Il Multiply tra un colore e un numero moltiplica ciascun canale per quel numero, quindi
la tinta resta la stessa e cresce solo l'intensità. Con Lumen, il sistema di
illuminazione globale di Unreal Engine 5, un materiale emissivo può anche illuminare un
po' le superfici vicine.

% ---------------------------------------------------------------------
\section{Il PBR: Metallic, Specular, Roughness e Normal}
\label{ue3:sec:pbr}

``Mi rendo conto che sto andando un po' a destra e a sinistra'', ammette Davide, e
torna ai pin del nodo risultante. Base Color ed Emissive si sono visti. Restano
\textbf{Metallic}, \textbf{Specular} e \textbf{Roughness}, ``degli input abbastanza
interessanti'', perché sono quelli che \textbf{determinano la superficie} del
materiale.

\paragraph{Il PBR.} Il rendering dei materiali in Unreal è \textbf{PBR},
\emph{Physically Based Rendering}: per decidere come si comporta una superficie, il
motore mette insieme una serie di informazioni fisiche. Tutti questi parametri vanno
\textbf{da 0 a 1}.

\paragraph{Roughness.} La \textbf{Roughness} dice quanto una superficie è
\textbf{liscia} o \textbf{ruvida}. Una \textbf{roccia} vuole una Roughness molto alta,
cioè 1: il materiale sarà ruvido ``e non rifletterà la luce in nessun modo''. Uno
\textbf{specchio} vuole una Roughness a 0: il più liscio possibile, il meno ruvido, il
più riflettente.

\paragraph{Specular e Metallic.} ``Per aiutare in questa situazione'' ci sono Specular e
Metallic, che ``si alternano a seconda del tipo di materiale''. Lo
\textbf{Specular}, come si è visto con le mappe, determina quanta luce viene riflessa e
l'angolazione a cui viene riflessa. Per creare uno \textbf{specchio} serve un materiale
molto, molto liscio e molto, molto riflettente: Davide abbassa la Roughness e alza la
riflessione, e la \emph{preview} comincia ad assomigliare ``a una parete riflettente'';
con altre cose che si vedranno in seguito ``iniziamo ad avere un riflesso''.

\paragraph{Emissive e Normal.} L'\textbf{Emissive Color} restituisce un colore che è
``effettivamente una luce''. La \textbf{Normal} è l'ultimo parametro utile, e richiede
la texture blu vista prima: si usa sempre un Texture Sample, ma con
\textbf{T\_Bark\_N}, la Normal della corteccia, invece della Diffuse. Davide toglie
tutto il resto e mette a confronto il materiale con e senza Normal: ``vedete la
differenza?'' La Normal aggiunge un dettaglio tridimensionale ``molto semplice da
vedere'', eppure ``la mappa è 2D, non c'è tridimensionalità'': crea solo l'illusione.

\begin{figure}[htbp]
\centering
\begin{tikzpicture}[font=\small]
  \foreach \i/\rr in {0/0, 1/0.25, 2/0.5, 3/0.75, 4/1} {
    \pgfmathsetmacro{\hs}{0.12+0.35*\rr}
    \pgfmathsetmacro{\op}{1-0.85*\rr}
    \shade[ball color=gray!70] (\i*2.4,0) circle (0.8);
    \fill[white, opacity=\op] (\i*2.4-0.28,0.3) circle (\hs);
    \node[font=\footnotesize] at (\i*2.4,-1.2) {Roughness $=\rr$};
  }
  \draw[-{Stealth[length=2.5mm]}, thick] (-0.6,-1.8) -- (10.2,-1.8);
  \node[font=\footnotesize, anchor=north west] at (-0.6,-1.85) {liscio, riflessi netti
    (lo specchio)};
  \node[font=\footnotesize, anchor=north east] at (10.2,-1.85) {ruvido, nessun
    riflesso netto (la roccia)};
\end{tikzpicture}
\caption{La Roughness, da 0 a 1, in uno schema qualitativo: più la superficie è liscia,
più il riflesso della luce è piccolo e netto; più è ruvida, più la luce si sparpaglia
fino a sparire come riflesso. Disegno indicativo, non un rendering di Unreal.}
\label{ue3:fig:roughness}
\end{figure}

\paragraph{Approfondimento: Metallic e Specular nella pratica.} Nel modello PBR di
Unreal i due parametri non hanno lo stesso peso. \textbf{Metallic} è quasi sempre
\emph{binario}: 1 per i metalli, 0 per tutto il resto (legno, pietra, plastica, pelle);
i valori intermedi servono soprattutto per i bordi tra metallo e vernice o per lo
sporco. Lo \textbf{Specular} ha un default di 0,5, che va bene per la grande maggioranza
dei materiali non metallici, e la documentazione di Unreal consiglia di toccarlo poco:
per i metalli viene ignorato, perché il loro riflesso dipende dal Base Color. Lo
specchio della lezione, in termini PBR, è quindi Metallic $=1$, Roughness $=0$, Base
Color chiaro. È il senso dell'``alternarsi'' di cui parla Davide: nei materiali metallici
si lavora con il Metallic, negli altri, se proprio serve, con lo Specular, e quasi
sempre con la Roughness, che è il parametro che cambia di più l'aspetto di una
superficie.

% ---------------------------------------------------------------------
\section{Il Texture Coordinate e il tiling}
\label{ue3:sec:tiling}

``Chiaramente ogni nodo funziona a sé.'' Il \textbf{Texture Coordinate}, il nodo rosso
della prima slide, ha i suoi parametri. La scorciatoia è ``un po' \emph{tricky} da
ricordare'', perché T + clic è già preso dal Texture Sample: è \textbf{U + clic},
``perché il Texture Coordinate lavora sulle UV''.

\paragraph{UTiling e VTiling.} Selezionando il nodo, nei dettagli compaiono
\textbf{UTiling} e \textbf{VTiling}: i parametri ``molto comodi'' con cui si aumenta o
si diminuisce il tiling di un oggetto. Davide collega il Texture Coordinate alle UV
della Normal e della Diffuse, e chiede: se metto UTiling a 2, che cosa succede? ``Vediamo
due volte la texture applicata all'oggetto'', risponde uno studente. ``Sì, però c'è un
\emph{catch}, perché è soltanto sulle U'': la texture si ripete due volte
sull'\textbf{asse orizzontale} e una sola volta su quello verticale. Con anche VTiling a
2, la texture si ripete \textbf{quattro} volte: ``ce n'è una qua, una qua, una qui e una
qua''.

\paragraph{Lo zoom.} Si può fare anche il contrario, per ``zoomare'' sulla texture: con
0,5 la si vede ``la metà delle volte. Anzi, non la metà delle volte, scusate: un
quarto'' (0,5 su entrambi gli assi). E questo non vuol dire che una parte della
texture ``viene mangiata'': viene semplicemente ``applicata in uno spazio differente''.

\begin{figure}[htbp]
\centering
\begin{tikzpicture}[font=\small]
  \newcommand{\uetreTile}[5]{% x0, y0, nu, nv, etichetta
    \begin{scope}[shift={(#1,#2)}]
      \pgfmathsetmacro{\du}{2.2/#3}
      \pgfmathsetmacro{\dv}{2.2/#4}
      \foreach \i in {1,...,#3} \foreach \j in {1,...,#4} {
        \draw[fill=brown!25] ({(\i-1)*\du},{(\j-1)*\dv}) rectangle
          ({\i*\du},{\j*\dv});
        \draw[brown!70!black, thick] ({(\i-1)*\du+0.15*\du},{(\j-1)*\dv+0.2*\dv})
          -- ({(\i-1)*\du+0.85*\du},{(\j-1)*\dv+0.8*\dv});
      }
      \draw[very thick] (0,0) rectangle (2.2,2.2);
      \node[font=\footnotesize, align=center] at (1.1,-0.5) {#5};
    \end{scope}}
  \uetreTile{0}{0}{1}{1}{U = 1, V = 1}
  \uetreTile{3.0}{0}{2}{1}{U = 2, V = 1}
  \uetreTile{6.0}{0}{2}{2}{U = 2, V = 2}
  % zoom: un quarto della texture, ingrandito
  \begin{scope}[shift={(9.0,0)}]
    \draw[fill=brown!25] (0,0) rectangle (2.2,2.2);
    \draw[brown!70!black, thick] (0.66,0.88) -- (2.2,2.2);
    \draw[very thick] (0,0) rectangle (2.2,2.2);
    \node[font=\footnotesize, align=center] at (1.1,-0.5) {U = 0,5, V = 0,5};
  \end{scope}
\end{tikzpicture}
\caption{Il tiling sulla stessa faccia. Con UTiling a 2 la texture si ripete solo in
orizzontale; con entrambi a 2 si ripete quattro volte; con entrambi a 0,5 si vede un
quarto della texture, ingrandito: non è ``mangiata'', è applicata in uno spazio più
grande. Schema ricostruito dalla dimostrazione.}
\label{ue3:fig:tiling}
\end{figure}

\paragraph{Il pavimento.} ``Il tiling, però, va visto.'' Nella mappa c'è un
\textbf{pavimento}, e Davide ne cambia il materiale: lo seleziona, e nei
\textbf{Details}, sotto \textbf{Materials}, al posto della \emph{material instance}
della griglia grigia di prototipazione\footnote{Nella trascrizione ``Materialistance
Prototype Grid Gray'': è l'istanza della griglia grigia dei progetti modello di
Unreal.} trascina il \textbf{master material} appena creato. Il risultato: la corteccia
è ``molto, molto grande'' (``io non so neanche qual è l'utilizzo di questa texture che
sto usando, ci ho preso Bark a caso, ma va bene così''). Davide porta il tiling a 3,
poi a 10: ``iniziamo a vederla più spesso, questa corteccia. Sembrano anche delle dune,
vagamente.'' Il pavimento è probabilmente un cubo allargato, con UV ``molto, molto
piccole'' rispetto alle sue dimensioni: per questo nella \emph{preview} sferica del
materiale la texture si ripete ``diecimila volte'', mentre sul pavimento ``abbiamo un
pavimento che è più dimensionato''.

\paragraph{Una domanda che apre il tema successivo.} Uno studente chiede: si possono
modificare U e V in base all'oggetto su cui è applicata la texture, per esempio 10 su un
pavimento e 0,5 su un altro oggetto? ``Sì, è possibile, non facendolo dallo stesso
materiale. Questo è il momento in cui entrano in gioco le \textbf{Material
Instance}.''

\paragraph{Approfondimento: il tiling in una formula.} Il Texture Coordinate restituisce
per ogni punto la sua coppia $(u, v)$; il Multiply la scala: $(u', v') = (k_u\,u,\;
k_v\,v)$. Una texture è definita sull'intervallo $[0,1]$ e si ripete fuori da lì (è il
comportamento di default, \emph{wrap}): quando $u'$ va da 0 a 2, la texture viene letta
due volte. Con $k < 1$ la faccia ``vede'' solo un pezzo dell'intervallo, e quel pezzo
viene ingrandito. Per questo il tiling dipende dalle UV della mesh: se le UV di un
pavimento enorme occupano comunque tutto il quadrato $[0,1]^2$, una texture con tiling
1 si stende una volta sola su tutta la superficie, ed è ``molto, molto grande''.
Moltiplicare per una costante e usare UTiling e VTiling nel nodo sono due modi di fare
la stessa cosa; il primo ha il vantaggio che la costante può diventare un parametro.

% ---------------------------------------------------------------------
\section{Le Material Instance}
\label{ue3:sec:instance}

Davide torna alle slide (``passiamo un attimo alle Material Instance, visto che ne
stavamo parlando''). Quello costruito finora è un \textbf{master material}: tutto ciò
che sta nel grafo. Se si modifica il grafo, si modifica il materiale. Per avere lo
\textbf{stesso materiale con una piccola differenza} si possono creare due
\textbf{istanze} e cambiare in ciascuna il parametro collegato al Texture Coordinate.

\paragraph{Che cos'è un'istanza.} Una \textbf{Material Instance} è ``un'istanza di un
materiale che noi possiamo utilizzare per modificare quel materiale \textbf{senza
modificare la sorgente}''. Ha diversi vantaggi, ``oltre a essere \emph{bon ton}'', cioè
il modo in cui si dovrebbero applicare i materiali in generale: ``non dovreste prendere
il master material e applicarlo agli oggetti in scena. È sempre consigliabile applicare
un'istanza.''

\paragraph{Perché è più leggera: i tavoli.} L'istanza è ``molto più leggera''. Davide
lo spiega con i \textbf{tavolini} della prima lezione. Tre tavoli in scena con lo
stesso materiale del legno: per ogni tavolo l'engine carica la mesh e carica il
materiale, quindi ``tavolo, tavolo, tavolo'' e ``materiale, materiale, materiale''. Se
invece sui tavoli ci sono tre Material Instance, l'engine carica le tre mesh, ma trova
tre istanze che puntano \textbf{tutte allo stesso master material}, e ``va a
renderizzare un solo materiale''. ``Chiaramente, su tre tavoli non cambia niente. Ora
immaginate che invece di tre tavoli ci siano diecimila oggetti in scena'' e che mille
condividano lo stesso materiale: si capisce perché ``torna comodo avere un'istanza
piuttosto che avere un materiale generale''.

\begin{figure}[htbp]
\centering
\begin{tikzpicture}[
  font=\small,
  ue3tab/.style={draw, thick, minimum width=1.5cm, minimum height=0.7cm,
                 fill=brown!20, align=center},
  ue3mat/.style={draw, thick, rounded corners=3pt, minimum width=1.5cm,
                 minimum height=0.7cm, align=center},
  ue3arr/.style={-{Stealth[length=2mm]}, thick}
]
  % senza istanze
  \node[font=\footnotesize\bfseries] at (2.0,2.2) {come la racconta Davide: master
    material su ogni tavolo};
  \foreach \i in {0,1,2} {
    \node[ue3tab] (t\i) at (\i*2.0,1.2) {tavolo};
    \node[ue3mat, fill=blue!15] (m\i) at (\i*2.0,0) {M\_Legno};
    \draw[ue3arr] (t\i) -- (m\i);
  }
  \node[font=\footnotesize, align=center] at (2.0,-0.8) {tre mesh e ``materiale,
    materiale, materiale''};
  % con istanze
  \begin{scope}[shift={(7.6,0)}]
    \node[font=\footnotesize\bfseries] at (2.0,2.2) {con le Material Instance};
    \foreach \i in {0,1,2} {
      \node[ue3tab] (u\i) at (\i*2.0,1.2) {tavolo};
      \node[ue3mat, fill=orange!15, font=\footnotesize] (i\i) at (\i*2.0,0)
        {MI\_Legno\_\i};
      \draw[ue3arr] (u\i) -- (i\i);
    }
    \node[ue3mat, fill=blue!15] (mm) at (2.0,-1.3) {\textbf{M\_Legno}};
    \foreach \i in {0,1,2} \draw[ue3arr] (i\i) -- (mm);
    \node[font=\footnotesize, align=center] at (2.0,-2.1) {tre mesh, tre istanze, un
      solo materiale da compilare};
  \end{scope}
\end{tikzpicture}
\caption{L'esempio dei tavoli. A sinistra, il master material applicato direttamente;
a destra, tre istanze (che possono avere parametri diversi: un legno più chiaro, un
tiling più fitto) che rimandano tutte allo stesso master material. Schema ricostruito
dalla spiegazione; sul costo reale si veda l'approfondimento.}
\label{ue3:fig:tavoli}
\end{figure}

\paragraph{Come si crea.} ``Come lo facciamo su Unreal?'' Si fa \textbf{tasto destro}
sul materiale nel Content Browser e \textbf{Create Material Instance}, ``qui in
alto''. L'istanza appare \textbf{vuota}: ``non c'è un grafo per modificarlo, perché
chiaramente non posso modificare nessun grafo quando sono qui dentro''. Per avere
qualcosa da cambiare, bisogna estrarre dei \textbf{parametri} dal master material.
Davide torna nel grafo, fa tasto destro sul Texture Sample del colore, \textbf{Convert
to Parameter}, e lo chiama più o meno \emph{Diffuse Color}\footnote{Nella trascrizione
``discolor'': il nome esatto non si ricostruisce, ma è il parametro della texture del
colore.}; lo stesso per l'altro Texture Sample, chiamato \textbf{Normal Map}. Salvato il
materiale, nell'istanza compaiono \textbf{due parametri} selezionabili.

\paragraph{La nomenclatura.} ``Normalmente dovreste chiamarlo'' non \emph{New Material}
ma \textbf{M\_NewMaterial}, e l'istanza \textbf{MI\_NewMaterial}: ``per nomenclatura,
insomma, per ritrovare facilmente tutto quello che ci serve''.

\paragraph{Cambiare la texture dall'istanza.} Davide applica MI\_NewMaterial al
pavimento: è identico a prima. Ma se nell'istanza si cambia il parametro della texture,
da Bark a un'altra, il materiale cambia, senza toccare il grafo. Si potrebbe avere un
master material senza texture, ``delle texture assenti o base'', e tre o quattro
istanze in cui si cambiano le texture: un muro arancione, uno blu, uno rosso e uno verde
sono quattro istanze. ``Quando io renderizzo la scena, non andrò a chiamare cinque
materiali diversi, chiamerò sempre M\_NewMaterial, perché sono tutte istanze dello
stesso materiale.'' ``Non vorrei stare complicando troppo la cosa, però insomma, ditemi
voi se è chiaro il discorso.''

\paragraph{Il Tiling Multiplier.} Ed ecco la risposta alla domanda sul tiling diverso
per ogni oggetto. Nel master material Davide prende un \textbf{Multiply}, ci collega il
Texture Coordinate, collega il Multiply alle UV dei due Texture Sample e attacca alla B
una \textbf{costante}, che converte in parametro con il nome \textbf{Tiling
Multiplier}. Nell'istanza, il Tiling Multiplier portato a 10 dà ``lo stesso risultato di
prima''; ma ora è un parametro. Una copia dell'istanza con il Tiling Multiplier a 5
darebbe un tiling diverso con lo stesso materiale. ``È chiaro il discorso?''
``Perfetto.''

\begin{figure}[htbp]
\centering
\begin{tikzpicture}[
  font=\small,
  ue3node/.style={draw, thick, rounded corners=4pt, align=center,
                  minimum height=0.9cm},
  ue3uv/.style={-{Stealth[length=2mm]}, thick, red!60!black},
  ue3flt/.style={-{Stealth[length=2mm]}, thick, green!50!black},
  ue3col/.style={-{Stealth[length=2mm]}, thick, blue!60!black}
]
  \node[ue3node, fill=red!12, text width=1.9cm] (tc) at (0,0.7)
    {\textbf{TexCoord}\\\footnotesize (U + clic)};
  \node[ue3node, fill=green!8, text width=1.9cm] (tm) at (0,-0.9)
    {\textbf{Tiling}\\\textbf{Multiplier}\\\footnotesize parametro};
  \node[ue3node, fill=gray!10, text width=1.5cm] (mul) at (3.0,0)
    {\textbf{Multiply}};
  \node[ue3node, fill=orange!12, text width=2.4cm] (ts1) at (6.2,1.0)
    {\textbf{Diffuse Color}\\\footnotesize Texture Sample,\\
    \footnotesize parametro};
  \node[ue3node, fill=orange!12, text width=2.4cm] (ts2) at (6.2,-1.0)
    {\textbf{Normal Map}\\\footnotesize Texture Sample,\\\footnotesize parametro};
  \node[ue3node, fill=blue!6, text width=2.2cm, minimum height=2.6cm] (res)
    at (10.0,0) {\textbf{M\_NewMaterial}\\[4pt]\footnotesize Base Color\\[14pt]
    \footnotesize Normal};
  \draw[ue3uv] (tc.east) -- ++(0.4,0) |- ([yshift=2mm]mul.west);
  \draw[ue3flt] (tm.east) -- ++(0.4,0) |- ([yshift=-2mm]mul.west);
  \draw[ue3uv] (mul.east) -- ++(0.5,0) |- node[pos=0.8, above, font=\scriptsize]{UVs}
    (ts1.west);
  \draw[ue3uv] (mul.east) -- ++(0.5,0) |- node[pos=0.8, below, font=\scriptsize]{UVs}
    (ts2.west);
  \draw[ue3col] (ts1.east) -- ++(0.6,0) |- ([yshift=6mm]res.west);
  \draw[ue3col] (ts2.east) -- ++(0.6,0) |- ([yshift=-6mm]res.west);
  % istanze
  \node[ue3node, fill=yellow!15, text width=3.6cm, font=\footnotesize] (mi1)
    at (3.0,-2.9) {\textbf{MI\_NewMaterial}\\Tiling Multiplier = 10\\Diffuse = T\_Bark};
  \node[ue3node, fill=yellow!15, text width=3.6cm, font=\footnotesize] (mi2)
    at (8.0,-2.9) {\textbf{MI\_NewMaterial\_2}\\Tiling Multiplier = 5\\
    Diffuse = un'altra texture};
  \draw[dashed, thick] (mi1.north) -- ++(0,0.4) -| (res.south);
  \draw[dashed, thick] (mi2.north) -- (mi2.north |- res.south);
\end{tikzpicture}
\caption{Il master material con tre parametri (le due texture e il Tiling Multiplier)
e due istanze che li impostano in modo diverso. Le istanze non hanno un grafo: hanno
solo le caselle dei parametri. Ricostruzione dalla dimostrazione.}
\label{ue3:fig:instance}
\end{figure}

\paragraph{Il vantaggio principale.} Ricapitolando dalle slide: le Material Instance
permettono di \textbf{personalizzare velocemente} le superfici ``senza andare a
modificare il grafico che avete fatto'', e le modifiche si vedono
\textbf{immediatamente}, ``senza dover aspettare la ricompilazione del materiale''.

\paragraph{Approfondimento: che cosa si risparmia davvero.} La spiegazione dei tavoli
va presa come un'immagine. Il costo principale di un materiale non è ``caricarlo'' una
volta per oggetto, ma \textbf{compilarlo}: ogni master material diventa uno o più
\emph{shader}, programmi per la scheda video che Unreal deve compilare (e che fanno
comparire, nella prima apertura di un progetto, il famoso contatore ``Compiling
Shaders''). Un'istanza \textbf{non} genera un nuovo shader: riusa quello del genitore e
cambia solo i valori dei parametri. Cento istanze dello stesso master material sono
quindi un solo shader con cento ``schede'' di valori; cento materiali diversi sono cento
shader. Per questo cambiare un parametro in un'istanza è istantaneo, mentre cambiare un
nodo nel master material richiede la ricompilazione, come dice la slide. Le istanze del
Content Browser si chiamano \textbf{Material Instance Constant}; esistono anche le
\textbf{Material Instance Dynamic}, create durante il gioco da un Blueprint (nodo
\emph{Create Dynamic Material Instance}) per cambiare i parametri a \emph{runtime}, per
esempio con \emph{Set Scalar Parameter Value}: un materiale che si accende quando la
torretta di Igor ha un bersaglio nel mirino.

% ---------------------------------------------------------------------
\section{Le Material Function}
\label{ue3:sec:function}

``Per chiudere il discorso materiali'', l'ultima voce delle slide: le \textbf{Material
Function}, che oggi non si useranno. Sono funzioni che raccolgono ``una serie di calcoli
un po' più complessi'': invece di rifare ogni volta dei calcoli che si fanno spesso, li
si mette dentro una funzione e si usa la funzione ``a vostra discrezione''. ``È come se
fossero delle \textbf{parentesi}: voi prendete queste parentesi, le spostate in giro. È
un pacchetto che potete portarvi in giro.'' Servono a velocizzare lo sviluppo quando ci
sono cose da ripetere, e possono essere usate e condivise ``fra membri del team, fra
materiali diversi''.

\paragraph{Approfondimento: Material Function e Material Instance.} Le due cose
risolvono problemi diversi. L'istanza riusa \emph{un intero materiale} cambiandone i
valori; la funzione riusa \emph{un pezzo di grafo} dentro materiali diversi. Unreal
stesso è pieno di Material Function già pronte (si trovano nel menu dei nodi con
l'icona della funzione), e i nodi World Aligned Texture e World Aligned Normal che
compaiono alla fine della lezione (sezione~\ref{ue3:sec:world}) sono, appunto, Material
Function del motore.

% ---------------------------------------------------------------------
\section{Il Lerp: interpolare}
\label{ue3:sec:lerp}

Restano ``un paio di nodi che possono essere carini''. Davide ricapitola quelli visti
(la costante singola, il colore, il Texture Sample, le operazioni, il Texture
Coordinate) e ne aggiunge due: il \textbf{Lerp} e il \textbf{Panner}. ``Se abbiamo
capito la regola degli shortcut, questi saranno quale lettera più clic sinistro?''
``L.'' ``Esatto'': \textbf{L + clic} per il Lerp, \textbf{P + clic} per il Panner. ``Volevo
essere sicuro che era arrivata la cosa.''

\paragraph{Due colori.} Il \textbf{Lerp} si occupa dell'\textbf{interpolazione
lineare} (\emph{linear interpolation}). Per farlo vedere Davide usa due colori, ``una
roba che non farete mai, non con due colori perlomeno'': un \textbf{blu} nella
\textbf{A} e un \textbf{rosso} nella \textbf{B}. Il Lerp, collegato al Base Color, dà il
\textbf{viola}. Perché è inutile farlo con i colori? ``Perché potevo cercarmi il viola
direttamente dal colore.'' Ma mostra come funziona l'interpolazione. Il Lerp si può fare
anche con i numeri, ``per fare operazioni matematiche'', e anche con le texture.

\paragraph{L'alfa.} L'operazione è ``gestita interamente dall'\textbf{alfa}'', il terzo
ingresso. L'alfa va \textbf{da 0 a 1} e decide quale colore prevale, con A che vale 0 e
B che vale 1: con l'alfa a 0 si ha il blu, a 1 il rosso, e in mezzo ``una percentuale di
quel colore'', per esempio 0,7. Modificando l'alfa si va sempre più verso il rosso, o
sempre più verso il blu.

\paragraph{Una texture come alfa.} Ma l'alfa non deve essere un numero fisso. Davide
prende ``il nostro fidato Bark'' e collega \textbf{un suo canale} all'alfa: ``vi ricordate
il discorso dei canali sulle texture?'' Ora quella sequenza di 0 e di 1, di bianco e di
nero, ``è l'\textbf{arbitro} nel decidere dove ci sarà più A e dove ci sarà più B'':
nelle zone più nere il Lerp mette più blu, in quelle più bianche più rosso. Questa
texture è ``estremamente bianca'', quindi il risultato tende al rosso; con altre
texture, o lavorando su questa, si otterrebbero risultati diversi. E collegando il
Multiply delle coordinate anche al Texture Sample dell'alfa, si ottiene un tiling
anche dell'alfa, cioè di dove va il blu e dove il rosso. ``Chiaro questo processo?''
``Sì.''

\begin{figure}[htbp]
\centering
\begin{tikzpicture}[
  font=\small,
  ue3node/.style={draw, thick, rounded corners=4pt, align=center,
                  minimum height=0.9cm},
  ue3col/.style={-{Stealth[length=2mm]}, thick, blue!60!black},
  ue3flt/.style={-{Stealth[length=2mm]}, thick, green!50!black}
]
  \node[ue3node, fill=blue!30, text width=1.9cm] (a) at (0,1.4) {\textbf{blu}\\
    \footnotesize Constant3Vector};
  \node[ue3node, fill=red!30, text width=1.9cm] (b) at (0,0) {\textbf{rosso}\\
    \footnotesize Constant3Vector};
  \node[ue3node, fill=gray!10, text width=1.9cm] (al) at (0,-1.5)
    {\textbf{Alpha}\\\footnotesize 0,7 oppure\\\footnotesize un canale di T\_Bark};
  \node[ue3node, fill=gray!10, text width=1.6cm, minimum height=2.2cm] (lerp)
    at (3.6,0) {\textbf{Lerp}\\\footnotesize (L + clic)\\[4pt]\footnotesize A\\
    \footnotesize B\\\footnotesize Alpha};
  \node[ue3node, fill=violet!30, text width=2.0cm] (res) at (7.0,0)
    {\textbf{Base Color}\\\footnotesize viola, più o\\\footnotesize meno rosso};
  \draw[ue3col] (a.east) -- ++(0.6,0) |- ([yshift=5mm]lerp.west);
  \draw[ue3col] (b.east) -- ([yshift=0mm]lerp.west);
  \draw[ue3flt] (al.east) -- ++(0.6,0) |- ([yshift=-5mm]lerp.west);
  \draw[ue3col] (lerp) -- (res);
  % la scala dell'alfa
  \begin{scope}[shift={(9.4,-0.6)}]
    \shade[left color=blue!60, right color=red!60] (0,0) rectangle (3.4,0.6);
    \draw[thick] (0,0) rectangle (3.4,0.6);
    \node[font=\footnotesize] at (0,-0.3) {0};
    \node[font=\footnotesize] at (3.4,-0.3) {1};
    \node[font=\footnotesize] at (1.7,-0.3) {Alpha};
    \node[font=\footnotesize] at (0,0.95) {A};
    \node[font=\footnotesize] at (3.4,0.95) {B};
    \draw[thick] (2.38,-0.05) -- (2.38,0.65);
    \node[font=\scriptsize] at (2.38,0.95) {0,7};
  \end{scope}
\end{tikzpicture}
\caption{Il Lerp della dimostrazione: con l'alfa a 0 esce A (il blu), a 1 esce B (il
rosso), in mezzo una miscela. Se l'alfa è un canale di una texture, ogni pixel ha la sua
miscela: blu dove la texture è nera, rosso dove è bianca.}
\label{ue3:fig:lerp}
\end{figure}

\paragraph{L'esercizio della pausa.} ``Provateci voi.'' Durante la pausa, chi vuole può
prendere \textbf{due texture}, collegarle a un Lerp e, giocando con l'alfa, vedere quale
parte delle due texture viene presa di più e quale di meno.

\paragraph{Approfondimento: la formula del Lerp.} L'interpolazione lineare tra A e B
con peso $\alpha$ è
\[
  \operatorname{lerp}(A, B, \alpha) = A\,(1-\alpha) + B\,\alpha .
\]
Con colori, il calcolo si fa canale per canale. Blu $=(0,0,1)$ e rosso $=(1,0,0)$ con
$\alpha = 0{,}5$ danno $(0{,}5;\,0;\,0{,}5)$, il viola della lezione; con $\alpha =
0{,}7$ danno $(0{,}7;\,0;\,0{,}3)$, un viola più rosso. Quando l'alfa è una texture, la
formula si applica pixel per pixel, ed è così che si mescolano due materiali su una
stessa superficie: erba e terra su un terreno, vernice e ruggine su un metallo, con una
maschera in bianco e nero a fare da ``arbitro''.

% ---------------------------------------------------------------------
\section{Il Panner: far muovere le texture}
\label{ue3:sec:panner}

Il \textbf{Panner}\footnote{Nella trascrizione ``panel'', ``planner'' e ``banner''.}
``si occupa di far muovere le texture''. Davide riprende il Base Color, stacca il resto,
e collega il Panner alle \textbf{UV} del Texture Sample, ``perché lui lavora sulle UV
della texture''. Selezionando il nodo (o direttamente sul nodo) si vedono
\textbf{Speed X} e \textbf{Speed Y}. Con Speed X a 1 la \emph{preview} comincia a
muoversi sull'asse X, ``anche un po' troppo veloce per i miei gusti''; con 0,3 in X e
0,1 in Y si ha un Panner che si muove su entrambi gli assi.

\paragraph{A che cosa serve.} ``Che materiale potremmo volerci fare con questo nodo?''
``Dell'acqua.'' ``Dell'acqua. Qualsiasi fluido, qualsiasi cosa che si muove.'' Il Panner
ha il parametro di velocità e l'ingresso delle \textbf{coordinate}: per combinarlo con il
tiling si collega il Texture Coordinate (o il Multiply) all'ingresso Coordinate. C'è
poi l'ingresso \textbf{Time}, con cui si può decidere ``quando e come deve muoversi'';
altrimenti, di default, usa il \emph{game time}, ``cioè è sempre in movimento''.

\paragraph{Il master material cambia, l'istanza segue.} Mentre il Panner lavora, il
pavimento nella scena \textbf{si muove}: sul pavimento c'è la Material Instance, e
``nonostante stiamo facendo delle modifiche al materiale originale e non alla Material
Instance, la Material Instance continua a modificarsi''. È il rovescio del rapporto tra
master e istanze: quello che cambia nel genitore cambia in tutte le istanze.

\paragraph{Approfondimento: il Panner in una formula.} Il Panner non fa altro che
sommare alle coordinate un tempo moltiplicato per una velocità:
\[
  (u', v') = (u + s_x\,t,\; v + s_y\,t),
\]
dove $t$ è il tempo in secondi e $s_x, s_y$ le velocità in ``texture al secondo'': con
Speed X a 1 la texture scorre di una larghezza intera ogni secondo, per questo era
troppo veloce. Siccome la texture si ripete, lo scorrimento è infinito e senza salti,
purché la texture sia \emph{tileable}, cioè con i bordi che combaciano. Due Panner con
velocità e tiling diversi sommati tra loro sono il trucco classico per l'acqua e per la
lava: il movimento non si ripete mai uguale.

% ---------------------------------------------------------------------
\section{La pausa: il progetto che non si trova, e Git}
\label{ue3:sec:pausa}

``Visto che comunque ci siamo fatti un'oretta e mezza'', Davide chiude per ora il
discorso materiali e propone una pausa di una decina di minuti. Al ritorno: l'import,
le mesh, e ``una cosina con la torretta che c'è nei file che vi abbiamo passato'', così
che Igor, quando torna, potrà farci ``un piccolo sistema di puntamento''.

\paragraph{Il progetto che Unreal non vede.} Al rientro, in attesa degli altri, una
studentessa racconta il suo problema: ha scaricato e decompresso la cartella del Drive,
ma aprendo Unreal il progetto non compare, e deve aprirlo ``a mano'' da un'altra parte.
Davide spiega che dipende probabilmente dal percorso in cui è salvata la cartella: se la
si mette dove sono salvati gli altri progetti, Unreal dovrebbe trovarla; in ogni caso,
aprire il progetto dal file eseguibile nella cartella funziona. ``Se non è un problema
che ti impedisce di accedere al progetto, fa niente.''

\paragraph{Come si lavora insieme a distanza?} Uno studente, programmatore, fa una
domanda ``fuori dall'argomento'': i programmatori usano \textbf{Git}, ma i file di un
progetto di Unreal sembrano troppo pesanti per lavorarci in contemporanea da lontano.
Davide conferma che \textbf{Git}, ``o diciamo repository di sorta'', è ``la cosa
principale che utilizzerai''. Con \textbf{Unity} ha sempre avuto meno problemi, perché
Unreal ``va gestito in maniera un po' più delicata'', ma entrambi supportano senza
problemi repository e \emph{source control}. Esistono anche altri metodi, che Davide non
ha mai provato: per Unreal, ``una cosa che io sento da una vita'', un modo per lavorare
\textbf{contemporaneamente} sulla stessa cosa, tramite \textbf{VPN} o in \textbf{LAN} e
con dei \emph{plugin} specifici, in modo che due colleghi modifichino attivamente lo
stesso livello. Con Git invece ``ognuno lavora sul suo e poi si \emph{merge} tutto
insieme''. ``Git è quello che ho utilizzato sempre, è il più comodo, quindi dovresti
trovarti abbastanza a casa.''

\paragraph{Approfondimento: Git LFS, Perforce e Multi-User Editing.} I timori dello
studente sono fondati. Git nasce per file di testo, e gli \emph{asset} di Unreal
(\texttt{.uasset}, \texttt{.umap}) sono file \textbf{binari}, spesso grandi, che Git non
sa fondere: se due persone modificano lo stesso livello, una delle due modifiche va
persa. Per questo con Unreal si usa \textbf{Git LFS} (\emph{Large File Storage}), che
tiene i file pesanti fuori dalla cronologia normale e permette di \textbf{bloccare}
(\emph{lock}) un file mentre lo si modifica. Nell'industria, per i progetti Unreal, è
molto diffuso \textbf{Perforce} (Helix Core), il sistema che Epic stessa usa e che
l'editor supporta direttamente, proprio per la gestione dei file binari e dei blocchi.
Il sistema ``sentito da una vita'' di cui parla Davide esiste: è il plugin
\textbf{Multi-User Editing}, con cui più editor connessi in rete (in LAN o via VPN)
condividono una stessa sessione e vedono le modifiche degli altri in tempo reale; nato
per la produzione virtuale nel cinema, si può usare anche nello sviluppo di giochi, ma
non sostituisce il \emph{source control}, a cui le modifiche vanno comunque consegnate
alla fine.

% ---------------------------------------------------------------------
\section{L'import e la torretta}
\label{ue3:sec:import}

Passato ``il nostro quarto d'ora'', si riparte: si importano le prime mesh, la
\textbf{torretta}, ``le tre parti della torretta, perché l'abbiamo divisa in tre
parti''.

\paragraph{La finestra di import, sulla slide.} Quando si importa una mesh compare una
finestra (``che può essere più larga, può essere più stretta'') piena di parametri. Le
cose principali che si possono fare da lì:
\begin{itemize}
  \item \textbf{generare automaticamente le collisioni}, se non sono presenti nel file
        FBX, perché magari gli oggetti non sono stati creati con le collisioni;
  \item \textbf{modificare il transform} dell'oggetto importato ``proprio di base'':
        ruotarlo sul suo \textbf{pivot}, traslarlo. ``Avete già parlato di cos'è il
        pivot con Igor, giusto? Il fulcro che vi permette di traslare e ruotare gli
        oggetti all'interno della scena.'' Con la funzione di \textbf{reimport}, si
        vedrà, questa possibilità diventa molto utile;
  \item adattare le \textbf{unità} e gli \textbf{assi} tra il software di origine e
        Unreal: \textbf{convertire la scena} dalle coordinate FBX a quelle di Unreal,
        \textbf{convertire le unità} (Unreal lavora in centimetri), oppure usare
        \textbf{Force Front X Axis}, che cambia l'asse che Unreal considera come
        ``fronte'' dell'oggetto. Davide l'ha usato poche volte: una volta il modello di
        un \emph{character}, importato normalmente, ``veniva rivoltato, come se il
        dentro fosse fuori'', e cambiare l'asse frontale gli ha restituito la mesh
        com'era stata esportata (``mi pare fosse Blender'');
  \item decidere se importare anche i \textbf{materiali} e le \textbf{texture}
        collegati alla mesh, e \textbf{dove} cercarli: in locale, in un'altra cartella,
        e così via.
\end{itemize}

\paragraph{Il pulsante Import.} Prima di importare, Davide spegne il Panner (``facciamo
che smetto di far muovere il pavimento'') e il pavimento torna normale. Nel Content
Browser c'è ``questo pulsantone con \textbf{Import}'', che apre l'esplora risorse del
computer per cercare i file. Nella cartella del materiale didattico c'è una cartella
\textbf{torretta}, con tre \textbf{static mesh}: \textbf{SM\_Turret\_Base},
\textbf{SM\_Turret\_Cannon} e \textbf{SM\_Turret\_Pivot}\footnote{Nella trascrizione
``Turret Base, Turret Cannon e Turret Pivo''; il prefisso SM\_ è detto a voce.}.

\paragraph{La finestra di import, dal vivo.} Importandole si apre la finestra (``nel
mio caso è così enorme, ma perché io tendo a usarla così''), e si ritrovano le opzioni
della slide: la possibilità di modificare traslazione e rotazione, ``anche
singolarmente sui vari assi'' (``questo è qualcosa che faremo più tardi''); i
\textbf{LOD} (``decidere il gruppo dei LOD''); le collisioni, se e come importarle. Se si
importa una \textbf{Skeletal Mesh}\footnote{Nella trascrizione ``scheda dal Mesh'' e
``Scheletal Mesh'', come nella seconda lezione.}, come un personaggio, la finestra
permette di collegarla a uno \textbf{scheletro} esistente, di generarne uno, oppure di
importare solo delle \textbf{animazioni}: ``faremo una lezione specifica sulle
animazioni, quindi per il momento glissiamo''. In alto ci sono anche le sezioni per
andare più veloci, tra cui \textbf{Import Materials} e \textbf{Import Textures}, con la
\textbf{Search Location} in cui cercarli.

\paragraph{Niente materiali e texture, per ora.} La consegna di Davide: \textbf{togliere
la spunta} a Import Materials e Import Textures, e fare \textbf{Import}. Unreal importa i
tre oggetti (Davide elimina quelli della prova precedente e salva). Perché togliere le
spunte? Se si fossero importati gli oggetti con materiali e texture, Unreal avrebbe
importato automaticamente tutte le texture, creato i materiali adatti e li avrebbe
anche \textbf{applicati}, ``in modo da averli già precisi sulla torretta''. Ed è
proprio quello che si farà, ma \textbf{a mano}, ``così vediamo effettivamente anche il
processo''. ``Ci sono stati problemi? Tutto a posto con la torretta?'' ``Sì, nessun
problema.''

\begin{table}[htbp]
\centering
\small
\begin{tabular}{@{}p{4.2cm}p{8.8cm}@{}}
\toprule
\textbf{Opzione} & \textbf{A che cosa serve} \\
\midrule
Generate Missing Collision & genera le collisioni se il file FBX non le ha \\
Import Translation / Rotation / Uniform Scale & sposta, ruota, scala l'oggetto ``alla
  radice'': cambia dove sta il pivot rispetto alla geometria \\
LOD Group & sceglie un gruppo di livelli di dettaglio (\emph{Level Of Detail}) \\
Convert Scene, Convert Scene Unit & adattano assi e unità di misura del file FBX a
  quelli di Unreal (centimetri) \\
Force Front X Axis & cambia l'asse considerato come ``fronte'' \\
Skeleton / Import Animations & per le Skeletal Mesh: scheletro esistente o nuovo, solo
  animazioni \\
Import Materials, Import Textures & importano e applicano materiali e texture
  collegati \\
Search Location & dove cercare materiali e texture da importare \\
\bottomrule
\end{tabular}
\caption{Le opzioni della finestra di import citate a lezione, con i nomi della
finestra di Unreal Engine 5 (FBX Import Options; il raggruppamento può variare tra le
versioni). Per la torretta: Import Materials e Import Textures senza spunta.}
\label{ue3:tab:import}
\end{table}

% ---------------------------------------------------------------------
\section{Che cos'è una mesh}
\label{ue3:sec:mesh}

Mentre la torretta è importata, Davide torna sulle slide per ``le ultime cose sulle
mesh''. Sulla slide c'è un \textbf{vaso} in \emph{wireframe} e, sotto, una
semplificazione. Una mesh è ``un insieme di \textbf{vertici}, \textbf{spigoli} e
\textbf{facce} che vanno a creare un oggetto in tre dimensioni''.
\begin{itemize}
  \item I \textbf{vertici} sono i \textbf{punti} che costituiscono gli angoli
        dell'oggetto, punti nello spazio.
  \item Gli \textbf{spigoli} (\emph{edge}) sono i \textbf{segmenti} che collegano due
        vertici.
  \item Le \textbf{facce} sono le \textbf{superfici piatte} definite da più segmenti:
        sulla slide, ``tutti i triangoli''.
\end{itemize}
Il \textbf{triangolo} è ``la forma più comoda che esiste a livello di calcolo'', e
tendenzialmente tutti gli oggetti saranno composti da ``un insieme di triangolini''. Più
triangoli insieme, più facce, formano dei \textbf{poligoni}; i poligoni combinati formano
delle \textbf{superfici}; e la superficie determina la mesh. Tutto questo si può vedere
anche nell'engine, con la visualizzazione \textbf{wireframe}: l'oggetto non ``come
appare'', ma ``come griglia, come insieme di tutti i poligoni, vertici e spigoli''.

\begin{figure}[htbp]
\centering
\begin{tikzpicture}[font=\small]
  \coordinate (A) at (0,0);
  \coordinate (B) at (2.6,0.3);
  \coordinate (C) at (3.2,2.2);
  \coordinate (D) at (0.6,2.6);
  \coordinate (E) at (1.7,1.3);
  \fill[orange!15] (A) -- (B) -- (E) -- cycle;
  \fill[orange!30] (B) -- (C) -- (E) -- cycle;
  \fill[orange!20] (C) -- (D) -- (E) -- cycle;
  \fill[orange!10] (D) -- (A) -- (E) -- cycle;
  \draw[very thick] (A) -- (B) -- (C) -- (D) -- cycle;
  \draw[very thick] (A) -- (E) (B) -- (E) (C) -- (E) (D) -- (E);
  \foreach \p in {A,B,C,D,E} \fill (\p) circle (2.5pt);
  \draw[-{Stealth[length=2mm]}] (4.3,2.8) node[right, font=\footnotesize]
    {\textbf{vertice}: un punto} -- (C);
  \draw[-{Stealth[length=2mm]}] (4.3,1.6) node[right, font=\footnotesize]
    {\textbf{spigolo} (\emph{edge}): un segmento tra due vertici} -- ($(B)!0.5!(C)$);
  \draw[-{Stealth[length=2mm]}] (4.3,0.4) node[right, font=\footnotesize]
    {\textbf{faccia}: un triangolo} -- (1.45,0.55);
  \node[font=\footnotesize, align=center, text width=4cm] at (1.6,-0.8)
    {4 facce triangolari\\formano un poligono};
\end{tikzpicture}
\caption{Vertici, spigoli e facce, come nella semplificazione della slide. Il triangolo
è l'unità di calcolo: anche le facce a quattro lati dei programmi di modellazione,
quando arrivano in Unreal, vengono divise in triangoli.}
\label{ue3:fig:mesh}
\end{figure}

% ---------------------------------------------------------------------
\section{Dimensioni, unità di misura e formati}
\label{ue3:sec:regole}

Quando si porta una mesh in Unreal, mentre la si produce o mentre la si cerca, bisogna
rispettare alcune cose.

\paragraph{Le dimensioni.} La cosa più comune: si decide di fare ``una mappa molto
carina'' su Unreal, si cercano mesh su internet o su \textbf{Fab} (``su Fab questo
problema c'è meno, però utilizzando siti esterni può capitare''), e la mesh importata è
\textbf{gigantesca} o \textbf{incredibilmente piccola}, perché chi ha fatto l'export, o
chi ha fatto l'import, non ha rispettato ``le misure di adattamento e transizione da un
software all'altro''. Unreal usa come unità di misura i \textbf{centimetri}; Blender,
per esempio, ``usa altra roba''. Nel passaggio da Blender a Unreal bisogna controllare
che la conversione delle dimensioni sia corretta.

\paragraph{FBX.} ``Altra cosa molto importante'': Unreal ``tendenzialmente legge
soltanto \textbf{FBX}''. Sono supportati anche altri formati, come gli \textbf{OBJ}, ``ma
non assicurano un corretto funzionamento all'interno dell'engine''. Il consiglio:
``importate sempre unicamente FBX''.

\paragraph{I formati delle texture.} Per le texture i formati supportati sono molti
(sulla slide c'è un elenco), ma nella pratica si usano soprattutto \textbf{JPEG} e
\textbf{PNG}, ``i più comuni, a meno che non state facendo qualcosa di molto
specifico''.

\paragraph{Approfondimento: Blender, i metri e il fattore 100.} Blender lavora di default
in metri, con una ``unità'' uguale a un metro; Unreal lavora in centimetri. Un cubo di
lato 2 in Blender è quindi un cubo di 2~m, cioè 200 unità di Unreal: se l'export non
converte, l'oggetto arriva cento volte troppo piccolo (o, con altre impostazioni, cento
volte troppo grande). La soluzione standard è nell'esportatore FBX di Blender
(l'opzione \emph{Apply Scalings} su \emph{FBX Units Scale}) o, lato Unreal, nella
scala uniforme e nella conversione delle unità della finestra di import. Un controllo
rapido: un personaggio umano in Unreal è alto circa 180 unità, come il manichino dei
progetti modello; se la mesh appena importata è molto diversa, c'è un problema di scala.
Oltre a FBX, le versioni recenti di Unreal Engine 5 importano anche \textbf{glTF} e
\textbf{USD}, ma FBX resta il formato più collaudato, come dice Davide.

% ---------------------------------------------------------------------
\section{Il pivot del cannone e il Reimport}
\label{ue3:sec:pivot}

Davide apre le tre mesh (la base, la torretta e il pivot) e le piazza in scena. Sono
state esportate con il \textbf{pivot nello stesso punto}: per sovrapporle senza
spostarle a mano basta \textbf{copiare la X e la Y} di un pezzo negli altri. Ed ecco la
torretta, ``a parte il fatto che non c'ha il materiale''.

\paragraph{I problemi della torretta.} Ma ``sta torretta ha dei problemi: quali sono?''
Il primo si vede: per qualche motivo \textbf{non è allineata}, ed è ``scomodamente''
ruotata: trovare la rotazione giusta sarà scomodo. Il secondo è più sottile. Si
supponga di mettere i tre pezzi insieme in un \textbf{Blueprint}. Che cosa deve fare una
torretta? Deve \textbf{tracciare dei bersagli}: ruotare su sé stessa, ma anche andare su
e giù, se il bersaglio è in alto o in basso. E qui c'è il problema: il \textbf{pivot del
cannone} non è in una posizione comoda per la rotazione verticale, perché sta
\textbf{in basso}, in fondo. Ruotare sul piano orizzontale non dà problemi, ma per
alzare e abbassare la canna il pivot dovrebbe essere ``più o meno all'altezza di questo
\textbf{bullone}'', il fulcro, e da lì il cannone dovrebbe essere manipolato e ruotato
(figura~\ref{ue3:fig:pivot}).

\begin{figure}[htbp]
\centering
\begin{tikzpicture}[font=\small,
  ue3arr/.style={-{Stealth[length=2.5mm]}, thick}]
  \foreach \xs/\titolo/\px/\py/\ang in {0/{pivot in basso (com'era)}/1.25/0.0/0,
                                         6.6/{pivot nel fulcro (dopo il Reimport)}/1.25/1.75/0} {
    \begin{scope}[shift={(\xs,0)}]
      % base
      \draw[thick, fill=orange!35] (0.3,0) -- (2.2,0) -- (1.9,0.5) -- (0.6,0.5) -- cycle;
      % pivot (supporto)
      \draw[thick, fill=gray!30] (0.8,0.5) rectangle (1.7,1.5);
      % cannone
      \draw[thick, fill=orange!20] (0.7,1.5) rectangle (1.8,2.1);
      \draw[thick, fill=orange!20] (1.8,1.6) rectangle (4.0,1.95);
      % bullone
      \fill[black] (1.25,1.75) circle (2pt);
      % il pivot del cannone
      \draw[red, very thick] (\px-0.18,\py) -- (\px+0.18,\py) (\px,\py-0.18) -- (\px,\py+0.18);
      \draw[red, thick] (\px,\py) circle (0.18);
      \node[font=\footnotesize, align=center] at (2.0,-0.6) {\titolo};
    \end{scope}
  }
  % archi di rotazione
  \draw[ue3arr, red] (3.4,0.5) arc[start angle=-20, end angle=25, radius=2.6];
  \node[font=\scriptsize, red, align=center] at (4.6,0.55) {la canna ruota\\attorno al
    piede};
  \draw[ue3arr, red] (6.6+3.6,1.0) arc[start angle=-30, end angle=30, radius=1.6];
  \node[font=\scriptsize, red, align=center] at (6.6+4.75,1.75) {la canna ruota\\
    attorno al bullone};
  % assi
  \draw[ue3arr, blue] (-0.6,0) -- (-0.6,1.2) node[above, font=\scriptsize]{Z};
\end{tikzpicture}
\caption{Il problema del pivot del cannone, visto di lato (schema non in scala). Con il
pivot (il mirino rosso) in fondo, alzare la canna la farebbe girare attorno al piede,
staccandola dal supporto; con il pivot nel fulcro, all'altezza del bullone, la canna
si alza e si abbassa come una vera torretta. La rotazione orizzontale, attorno all'asse
Z, funziona invece in entrambi i casi.}
\label{ue3:fig:pivot}
\end{figure}

\paragraph{Come si sposta un pivot.} ``In realtà possiamo farlo in tutta una miriade di
modi.'' Con la \textbf{Modeling Mode} si potrebbe rimodellare l'oggetto. Ma lo si può
fare anche \textbf{nell'import}: si fa tasto destro sulla mesh e \textbf{Reimport Mesh
With Dialog}, che riapre la finestra di import vista prima. Unreal importa lo stesso
oggetto e lo \textbf{sovrascrive}, ma permette di cambiare come verrà posizionato:
l'\textbf{offset della rotazione} e l'\textbf{offset della traslazione}.

\paragraph{I numeri.} Davide se li era scritti, ma non li ricorda, e li ricava. Per
sapere quale asse toccare, guarda il \emph{gizmo}: il cannone deve essere allineato con
il resto della torretta (e lo stesso il pivot), e il pivot del cannone deve stare non
``qui giù in fondo'', ma allineato con il fulcro, sull'\textbf{asse Z}, ``perché è
questo blu''. Per questa torretta, ``quindi anche per voi'', i parametri sono:
\begin{itemize}
  \item per il \textbf{cannone}: traslazione Z $= -135$, che reimporta il cannone ``più
        attaccato al pivot'', e rotazione $= 113$, che lo riallinea\footnote{La
        trascrizione dice prima ``113 per la traslazione'' e poi si corregge: ``e 113
        sulla rotazione''. L'asse della rotazione non è detto; trattandosi di
        riallineare la torretta nel piano orizzontale, è con ogni probabilità la
        rotazione attorno a Z (il \emph{yaw}).};
  \item per il \textbf{pivot}: ``questa volta non serve modificare la traslazione'', ma
        solo la rotazione, sempre 113.
\end{itemize}
Dopo il reimport il cannone ``sta sotto'' (la sua origine ora è nel fulcro, quindi
nella posizione di prima risulta più basso) e va rimesso a posto; ma ora, ruotandolo, si
alza e si abbassa nel modo giusto. ``E chiaramente poi questo farete più avanti, con dei
Blueprint.'' Si salva tutto.

\paragraph{Approfondimento: perché il pivot conta così tanto.} In Unreal ogni rotazione
e ogni scala di un componente avvengono \textbf{attorno alla sua origine}, cioè al
pivot della mesh. Un Blueprint di torretta, di solito, ha la base come componente
radice, il supporto come figlio che ruota attorno a Z (il \emph{yaw}, per seguire il
bersaglio a destra e a sinistra) e il cannone come figlio del supporto che ruota attorno
a Y (il \emph{pitch}, su e giù). Per questo ogni pezzo deve avere l'origine esattamente
nel punto in cui gira nella realtà: se l'origine è sbagliata, nessun nodo di rotazione
potrà correggerla senza calcoli aggiuntivi. Il ``piccolo sistema di puntamento'' che
Igor farà con questa torretta (con nodi come \emph{Find Look at Rotation}, che calcola
la rotazione necessaria per guardare un punto) si appoggia interamente a questo lavoro
preparatorio.

% ---------------------------------------------------------------------
\section{Mettere ordine: le cartelle}
\label{ue3:sec:cartelle}

Prima di importare le \textbf{texture}, ``per dare un po' di ordine'', Davide crea nel
Content Browser una nuova cartella, \textbf{Torretta}, ci sposta dentro tutto
(``muoviamo, così è un po' più pulito''), e dentro crea un'ulteriore cartella,
\textbf{Texture}, in cui importare le texture della torretta. ``Questo non è
obbligatorio, è solo una questione di \emph{best practice}, cioè è come il discorso
delle Material Instance. Non vi vengo a trovare a casa se non lo fate, però vi risulta
tutto più comodo.''

% ---------------------------------------------------------------------
\section{Lo Static Mesh Editor e le collisioni}
\label{ue3:sec:sme}

Con un doppio clic su una mesh si apre lo \textbf{Static Mesh Editor}: ``andiamoci a
vedere effettivamente un po' qui dentro cosa possiamo vedere delle mesh''.

\paragraph{Gli slot dei materiali.} A destra, nei \textbf{Material Slots}, queste mesh
hanno \textbf{due slot}: ``non è la normalità''. (Una delle tre ne ha uno solo.) Il modo
in cui la torretta è stata suddivisa fa sì che, per rendere bene i materiali, servano
due slot: la base e il cannone condividono in parte il materiale del \textbf{pivot}. Nei
due slot si inseriranno i materiali più avanti. Altre impostazioni (Davide accenna a
quelle sui LOD e sulla \emph{Nanite}\footnote{La trascrizione dice ``per quanto riguarda
la night'': la ricostruzione come Nanite, la tecnologia di geometria virtualizzata di
Unreal Engine 5 che ha una sua sezione nei dettagli della mesh, è probabile ma non
certa.}) non si vedono oggi.

\paragraph{Le modalità di visualizzazione.} Dal menu della \emph{viewport} (la
prospettiva ``non ci interessa'') si cambia la modalità: \textbf{Lit}, quella normale;
\textbf{Unlit}, ``com'è il mio oggetto se non è colpito dalla luce?''; e
\textbf{Wireframe}: ``vedete quello che dicevamo prima?'' Vertici, facce ed \emph{edge}
della mesh, che si potrebbero modificare con la \textbf{Modeling Mode} (``se avanzerà
il tempo con le lezioni, sarebbe carino vederla più avanti''). Ci sono anche
\textbf{Lighting Only}, le \textbf{Reflections} e altre ancora. Dal menu \textbf{Show}
si possono mostrare i \textbf{vertici}: ``dove sono i vertici di questo oggetto?''

\paragraph{Simple e complex.} Un'altra cosa utile ``da un punto di vista di gameplay''
sono le \textbf{collisioni}. Dal menu Show si attiva la \textbf{Simple Collision}: ``queste
sono le collisioni del mio oggetto''. In scena, il \emph{character} interagirà con
queste: ci salterà sopra, o lo dovrà scavalcare. Ma oltre alle \textbf{simple
collision} ci sono le \textbf{complex collision}. La differenza:
\begin{itemize}
  \item le \textbf{simple collision} sono forme che si possono posizionare a mano,
        anche da qui, e rifare da capo;
  \item le \textbf{complex collision} usano la struttura in vertici della mesh stessa
        come collisione. ``È qualcosa che è molto più pesante rispetto a una collisione
        semplice, e torna utile solo in alcuni casi specifici.''
\end{itemize}

\paragraph{Modificare le collisioni.} La collisione semplice si può selezionare,
spostare, anche eliminare; e se ne possono \textbf{aggiungere} altre: una
\textbf{sfera}, di cui si decide la grandezza; una \textbf{capsula}, che ``è
tendenzialmente la collisione che usano i vostri \emph{character}'': il personaggio di
solito non ha una collisione super precisa, ``è un capsulone'', a meno che non si voglia
davvero il contrario; una \textbf{box}. E una box, per questa mesh, ``è più ottimizzata
rispetto alla collisione che questa aveva di base'': all'atto pratico fa la stessa cosa
``con molti meno vertici''. Su una singola mesh cambia poco; su molte mesh aiuta le
prestazioni. Ci sono anche collisioni più complesse, ma ``idealmente, il modo migliore
per fare delle simple collision è letteralmente prendere delle \textbf{forme base}'',
magari più d'una, ``piuttosto che avere tanti vertici sulla stessa collisione''.

\begin{table}[htbp]
\centering
\small
\begin{tabular}{@{}p{3.3cm}p{4.4cm}p{5.3cm}@{}}
\toprule
\textbf{Collisione} & \textbf{Com'è fatta} & \textbf{Quando} \\
\midrule
Box & un parallelepipedo & oggetti squadrati; spesso basta anche per oggetti più
  complessi, ed è la più leggera \\
Sphere & una sfera & oggetti tondi, palle \\
Capsule & un cilindro con due semisfere & i \emph{character}: ``un capsulone'' \\
Convex (più forme) & poliedri convessi, magari più d'uno & quando una forma sola non
  basta; meglio poche forme base che tanti vertici \\
Complex & i triangoli della mesh stessa & casi specifici: precisione massima, costo
  molto più alto \\
\bottomrule
\end{tabular}
\caption{Le collisioni di una static mesh viste nello Static Mesh Editor. Le prime
quattro sono \emph{simple collision}, che si aggiungono e si modificano dal menu
Collision dell'editor; l'ultima è la \emph{complex collision}.}
\label{ue3:tab:collisioni}
\end{table}

\paragraph{Approfondimento: chi usa quale collisione.} Unreal tiene le due collisioni
separate perché le usa per cose diverse. La \textbf{simple} è quella della
\textbf{fisica}: la pallina della seconda lezione che rotola sulla piattaforma, il
\emph{character} che cammina e salta, gli oggetti che cadono; i calcoli fisici su forme
semplici sono molto più veloci. La \textbf{complex} si usa soprattutto per le
\textbf{interrogazioni} precise, come un \emph{line trace} (un ``raggio'' lanciato per
sapere che cosa colpisce un proiettile o che cosa sta guardando il giocatore), quando
la precisione al triangolo serve davvero. Nei dettagli della mesh, il campo
\emph{Collision Complexity} decide come combinarle; l'opzione \emph{Use Complex
Collision As Simple} usa i triangoli anche per la fisica, ed è comoda per oggetti
statici di forma strana, ma va usata con giudizio, per il costo di cui parla Davide.

% ---------------------------------------------------------------------
\section{Raggruppare e fondere: Merge Actors}
\label{ue3:sec:merge}

``Se ad esempio vogliamo avere tutti e tre gli oggetti assieme'', o semplicemente
spostarli insieme, ci sono due strade.

\paragraph{Raggruppare.} Si selezionano tutti e tre e li si \textbf{raggruppa} con
\textbf{Ctrl+G}; si spostano insieme; con \textbf{Shift+G} si toglie il gruppo.

\paragraph{Fondere.} Oppure li si \textbf{fonde} (\emph{merge}). Si evidenziano i tre
pezzi in scena, si va nel menu \textbf{Actor} in alto e si sceglie \textbf{Merge
Actors}. Unreal chiede dove salvare il nuovo oggetto e con che nome (per esempio
\textbf{SM\_Turret}\footnote{Nella trascrizione ``smTourette''.}), e crea una
\textbf{nuova mesh}, ``che è semplicemente la fusione delle tre mesh precedenti'': una
somma, un \emph{merge}, appunto.

\paragraph{Il merge, di nuovo.} Più tardi uno studente dice di essersi perso come si fa
il merge, e Davide lo rifà con calma. I tre oggetti (il cannone, la base e il pivot)
sono in scena. Con il \textbf{tasto destro} su di essi ci sono delle opzioni, ``però
queste sono opzioni limitate''. Per averle tutte si usa il menu \textbf{Actor} in alto:
mostra gran parte delle stesse opzioni, ma in basso, nella sezione degli strumenti di
Unreal Engine, c'è \textbf{Merge Actors}. Al merge Unreal chiede dove salvare e con che
nome, e ``ti fonde i tre oggetti''. Nello stesso menu c'è anche \textbf{Chaos}, che
``serve per gli oggetti distruttibili'': cose più avanzate, che probabilmente non si
vedranno nel corso, ma da approfondire ``se vi capiterà di voler mettere distruttibilità
all'interno del vostro gioco''.

\paragraph{Approfondimento: gruppo, merge o Blueprint?} Le tre soluzioni non sono
equivalenti. Il \textbf{gruppo} è solo una comodità dell'editor: gli actor restano
separati, e in gioco non cambia nulla. Il \textbf{merge} crea una mesh nuova, una sola:
meno oggetti da disegnare (utile per le prestazioni, per esempio con tanti pezzi di
arredo statici), ma i pezzi \textbf{non si possono più muovere uno rispetto
all'altro}. Per la torretta di Igor, che deve ruotare il supporto e alzare la canna, il
merge è quindi la scelta sbagliata: serve un \textbf{Blueprint} con i tre pezzi come
\textbf{componenti} separati, imparentati tra loro, come si è visto nella seconda
lezione. Il merge della lezione è una dimostrazione dello strumento, e serve anche a
rispondere alla domanda sui materiali (sezione~\ref{ue3:sec:domande}).

% ---------------------------------------------------------------------
\section{I materiali della torretta}
\label{ue3:sec:matturret}

``Ora abbiamo importato le texture. Creiamoci i materiali di queste torrette.'' Tre
materiali, uno per pezzo: \textbf{M\_Turret\_Base}, \textbf{M\_Turret\_Pivot} e
\textbf{M\_Turret\_Cannon}. Si potrebbero fare molto elaborati, ``ma partiamo
semplicemente andando a utilizzare le texture che ci sono state date''.

\paragraph{Quattro Texture Sample.} Nel materiale della base Davide prende \textbf{quattro
Texture Sample} e ci mette le texture della base: \textbf{Base Color},
\textbf{Metallic}, \textbf{Roughness} e \textbf{Normal} (Turret Base Color, Turret Base
Metallic, Turret Base Roughness, Turret Base Normal). Poi le collega ai pin
corrispondenti del nodo risultante: il colore al Base Color, la Metallic al Metallic, la
Roughness alla Roughness, la Normal alla Normal (figura~\ref{ue3:fig:turret}).

\begin{figure}[htbp]
\centering
\begin{tikzpicture}[
  font=\small,
  ue3node/.style={draw, thick, rounded corners=4pt, align=center,
                  minimum height=0.85cm, text width=3.6cm},
  ue3col/.style={-{Stealth[length=2mm]}, thick, blue!60!black},
  ue3flt/.style={-{Stealth[length=2mm]}, thick, green!50!black}
]
  \node[ue3node, fill=orange!20] (c) at (0,2.1) {\textbf{Texture Sample}\\
    \footnotesize T\_Turret\_Base\_Color};
  \node[ue3node, fill=gray!15] (m) at (0,0.7) {\textbf{Texture Sample}\\
    \footnotesize T\_Turret\_Base\_Metallic};
  \node[ue3node, fill=gray!15] (r) at (0,-0.7) {\textbf{Texture Sample}\\
    \footnotesize T\_Turret\_Base\_Roughness};
  \node[ue3node, fill=blue!15] (n) at (0,-2.1) {\textbf{Texture Sample}\\
    \footnotesize T\_Turret\_Base\_Normal};
  \node[draw, very thick, rounded corners=4pt, fill=blue!6, minimum width=3.0cm,
    minimum height=5.2cm] (res) at (6.2,0) {};
  \node[font=\small\bfseries] at (6.2,2.3) {M\_Turret\_Base};
  \foreach \lab/\y in {Base Color/2.1, Metallic/0.7, Specular/0.0, Roughness/-0.7,
                       Emissive Color/-1.4, Normal/-2.1}
    \node[font=\footnotesize, anchor=west] at (4.85,\y) {\lab};
  \draw[ue3col] (c.east) -- (4.8,2.1);
  \draw[ue3flt] (m.east) -- node[above, font=\scriptsize]{R} (4.8,0.7);
  \draw[ue3flt] (r.east) -- node[above, font=\scriptsize]{R} (4.8,-0.7);
  \draw[ue3col] (n.east) -- (4.8,-2.1);
\end{tikzpicture}
\caption{Il materiale della base della torretta, fatto a mano: quattro Texture Sample,
quattro pin. Metallic e Roughness sono mappe in bianco e nero, e basta un canale (per
esempio R). M\_Turret\_Pivot e M\_Turret\_Cannon sono identici, con le loro texture.
Ricostruzione dalla dimostrazione; i nomi esatti dei file sono indicativi.}
\label{ue3:fig:turret}
\end{figure}

\paragraph{Le preview.} ``Mi ero dimenticato di farvi vedere una cosa.'' Nella
\emph{preview} del materiale si può scegliere la forma: il \textbf{cilindro}, la
\textbf{sfera}, il \textbf{piano}, il \textbf{cubo} e il \textbf{custom}. Con il custom,
se si seleziona una mesh nel Content Browser e si clicca il pulsante, la preview mostra
\textbf{quella mesh} con il materiale applicato. La prima volta non la mostra: ``sta
compilando''. Poi compare (Davide aveva collegato una texture ``dalla parte sbagliata'',
e la sistema).

\paragraph{Gli slot.} Il materiale è salvato e compilato. Sulla mesh della torretta ci
sono i \textbf{due slot}: in uno c'è M\_Turret\_Base, ma ``per come è stata strutturata
la torretta'' nell'altro va il materiale del \textbf{pivot}. ``Sta rallentando
veramente tanto il PC.''

\paragraph{Un grafico, un engine.} Davide si ferma su che cosa si è fatto: ``qui non
abbiamo fatto altro che utilizzare le mesh che ci sono state fornite da un grafico''.
Idealmente un grafico, in un software a parte, ha realizzato la mesh e le texture, e ce
le ha passate; noi le abbiamo inserite. Si potrebbe lavorare ancora sull'estetica della
torretta modificando i materiali, ``tuttavia, in questo caso non è necessario''.

\paragraph{Il pivot e il cannone, più furbi.} Per il materiale del pivot Davide è ``un
po' più furbo'': cerca la texture \emph{Turret Pivot Color}, e le altre le copia da lì,
perché nella lista ``stanno direttamente dopo''. Collega Metallic, Roughness e Normal,
salva, ``e abbiamo anche il materiale del nostro pivot''. Poi lo stesso per il cannone:
la texture \emph{Cannon Color}, \textbf{Ctrl+C}, \textbf{Ctrl+V}. Lì Davide ricorda le
scorciatoie classiche: ``tutti i classici shortcut di Control'', Ctrl+C, Ctrl+V, Ctrl+Z,
Ctrl+Y, Ctrl+A, ``sono applicabili anche a Unreal nella maggior parte dei casi''.

\paragraph{Salvate spesso.} Uno studente chiede perché il suo non si è applicato:
``Salvalo.'' E Davide ne fa una regola: ``puntualmente mi capita di vedere gente che sta
lavorando a qualcosa, non salva per qualche motivo, Unreal crasha e perde tipo tutto il
lavoro che stava facendo. E quindi salvate spesso.''

\paragraph{La torretta finita.} Nei due slot vanno il pivot e il cannone. ``Oh, eccoci
qua. Bello lucido.'' Salvato tutto, la torretta ha i suoi materiali.

\paragraph{Ritoccare il colore con un Lerp.} ``Se volete che la torretta sia meno
lucida, o che sia meno nitido l'arancione'', si può intervenire sui singoli materiali.
Davide improvvisa sul cannone: non è sicuro di volere il colore ``così netto'', quindi
prende un \textbf{nero} e un \textbf{Lerp}, collega il nero e la texture del colore alle
due entrate e mette nell'alfa un parametro con cui decidere ``quanto si deve vedere'',
per esempio 0,6. Si può convertire in parametro, creare una \textbf{Material
Instance}, e regolarlo da lì. Al posto del nero si può mettere un altro colore, un
verde: ``viene un po' più verdastra, è troppo verde'', e con l'alfa a 0,8 viene ``questo
verde mezzo, insomma, un po' scuro, un po' cromato''. ``Qui chiaramente sto improvvisando
cose, ma cose anche stupide. Ma nel momento in cui voi avete un'idea un po' più chiara,
un po' più specifica di quello che volete, potete andare a fare un po' più di cose.''

\paragraph{Approfondimento: lo stesso lavoro con un master material.} I tre materiali
della torretta hanno lo stesso identico grafo: cambiano solo le texture. È il caso da
manuale per le Material Instance della prima parte della lezione: un solo master
material (per esempio \texttt{M\_Turret}) con quattro Texture Parameter (colore,
metallic, roughness, normal) e magari un Lerp per la tinta, e tre istanze
(\texttt{MI\_Turret\_Base}, \texttt{MI\_Turret\_Pivot}, \texttt{MI\_Turret\_Cannon}) che
impostano le texture di ciascun pezzo. Un solo shader da compilare, e il ritocco
verdastro del cannone diventa una casella da regolare nell'istanza. Se le texture
fossero state importate insieme alle mesh (Import Materials e Import Textures), Unreal
avrebbe creato da solo i materiali, uno per slot, con un grafo simile a quello della
figura~\ref{ue3:fig:turret}.

% ---------------------------------------------------------------------
\section{Le domande: il Reimport e il merge}
\label{ue3:sec:domande}

``Ora tutto chiaro il discorso dell'import e dei materiali?'' Due domande.

\paragraph{Perché il Reimport con la Z?} Una studentessa non ha capito il reimport
modificando la Z. Davide riapre il cannone. Il reimport è servito a \textbf{modificare la
mesh} alla radice. L'oggetto importato, modellato da qualcun altro, era fatto in un
certo modo, e non è un problema in sé: lo si può abbassare e girare a mano, in scena.
``Però poi, tutte le volte che io prendo e rimetto in scena l'oggetto, l'oggetto sarà
sempre posizionato in quel modo e ruotato in quel modo.'' Il reimport invece ``modifica
proprio alla radice il modo in cui funziona l'oggetto''. E il problema principale non
era la rotazione (``so che l'angolo è 113, allora lo rimodifico''), ma il
\textbf{pivot}, ``in una posizione scomoda'': il reimport è un modo per spostarlo.

\paragraph{La Modeling Mode.} ``Non è l'unico modo.'' Davide apre la \textbf{Modeling
Mode}, cerca lo strumento del \textbf{pivot}, e mostra che lo si può mettere al centro,
a sinistra, a destra, oppure ``nel punto specifico che dico io: lo prendo e me lo
sposto''. Ma questa è ``una soluzione che hanno aggiunto con Unreal Engine 5'': prima la
Modeling Mode non c'era, e in Unreal Engine 4 reimportare era il modo comodo per farlo
``senza avere bisogno di rimodellare l'oggetto a posteriori''. ``Fai la stessa cosa,
semplicemente non hai fatto un'operazione extra.''

\paragraph{Una prova.} Davide stesso è curioso: se si reimporta di nuovo con $-135$ e
113? Il reimport \textbf{sovrascrive} quello che c'era prima, ma il cannone in scena
conserva la rotazione che gli era stata data a mano. Davide lo rimette dove stava (``io
sto facendo a cazzotti da tutta la lezione con la finestra di Zoom, se non si fosse
capito'').

\paragraph{Che cosa succede ai materiali con il merge?} Era la domanda arrivata in chat
durante la costruzione dei materiali: se faccio il merge e i pezzi della torretta hanno
texture diverse? Davide prova: seleziona i tre oggetti, \textbf{Actor}, \textbf{Merge
Actors}, nome \emph{Turret}, salva. Il risultato è un oggetto con \textbf{tre slot} per i
materiali: lo slot 0 è tutta la parte del cannone, lo slot 1 il pivot ``e tutta questa
zona qui'', e così via. Se in uno slot si mette un materiale vuoto, ``diventa tutto
nero'': ``vedete, tutta questa zona adesso viene influenzata da quel materiale lì''.
Risposta: con il merge l'actor viene fuso, e i materiali vengono \textbf{riposizionati
negli slot} della nuova mesh.

% ---------------------------------------------------------------------
\section{Lo stretch, lo z-fighting e i materiali World Aligned}
\label{ue3:sec:world}

``Mi è venuto in mente un'ultima cosa che vi potrei far vedere per quanto riguarda i
materiali, e poi vi vorrei far fare una cosina.'' Manca una mezz'oretta, ``visto che
Igor vi ha fatto fare delle cose, ho visto che siete molto partecipi. Quindi vorrei
fomentare questa cosa.''

\paragraph{M\_Test.} Davide ricrea il materiale della corteccia, cancellato per sbaglio:
\textbf{M\_Test}, con due Texture Sample (Bark e la sua Normal) e un Texture
Coordinate, ``proprio la roba più facile del mondo'', con il tiling a 10.

\paragraph{Ingrandire e stirare.} ``Ora vediamo se abbiamo capito come funziona tutto
quanto.'' I materiali sono collegati alle \textbf{UV} di un oggetto. Quindi se
l'oggetto si \textbf{ingrandisce}, al materiale ``non dovrebbe succedere granché'': si
ingrandisce anche lui (``la risoluzione non ne giova particolarmente''). E se l'oggetto
si \textbf{stira} (\emph{stretch}), tirandolo verso destra? ``Credo si stira anche
lui.'' Esatto: il materiale si allunga. ``In alcuni casi non è neanche particolarmente
brutto'', ma la corteccia ``non sembra più corteccia''; e tirando di più, ``così è
brutto''.

\paragraph{Il pavimento grande.} Se si vuole fare un pavimento per una zona molto
ampia, dover pensare a questo ``può essere una noia'': per non stirare il materiale si
dovrebbero fare tante copie dello stesso pezzo, una attaccata all'altra, ``che
tendenzialmente non è quello che voglio'', perché è scomodo.

\paragraph{Una digressione: lo z-fighting.} ``Visto che è una domanda che arriverà,
anticipo'': se due piani si trovano \textbf{nella stessa posizione}, ``clippano'': la
superficie sfarfalla a chiazze, perché l'engine ``non sa se utilizzare il materiale
dell'oggetto a destra o dell'oggetto a sinistra''. Le soluzioni: separarli davvero,
oppure spostarne uno di pochissimo in \textbf{Z}, per esempio di 0,1: ``è
impercettibile, non si nota, però uno è leggermente più sopra rispetto all'altro'', e il
problema sparisce.

\paragraph{Allineato al mondo.} ``C'è un modo per far sì che il mio materiale non si
stiri quando stiro l'oggetto? Chiaramente sì, perché se no non avrei fatto questa
domanda.'' Bisogna fare un materiale che non è allineato con le coordinate
dell'oggetto, ma con quelle del \textbf{mondo}. Davide lo spiega con il \emph{gizmo}
di spostamento: un oggetto ruotato si può muovere lungo gli assi del mondo o lungo i
suoi assi locali, e il pulsante in alto passa dall'uno all'altro. ``Più o meno il
discorso è lo stesso'': nel caso del materiale non si parla di assi, ma di UV.

\paragraph{I nodi World Aligned.} Davide copia il materiale (\textbf{M\_Test1}) e usa
due nodi, \textbf{World Aligned Texture} e \textbf{World Aligned Normal}. Non hanno
scorciatoie: si cercano per nome. Hanno ``un sacco di parametri'', ma per ora ne servono
due: \textbf{Texture Object} e \textbf{Texture Size}. World Aligned Texture va nel
\textbf{Base Color}, World Aligned Normal nella \textbf{Normal}, usando l'uscita
\textbf{XYZ}, ``perché almeno prendo tutti gli assi'': il materiale è allineato al
mondo su tutti e tre.

\paragraph{Texture Size e Texture Object.} Il \textbf{Texture Size}, un vettore a tre
valori, ``fa la stessa cosa del Texture Coordinate'': regola la grandezza della
ripetizione. Perché allora non tenere il Texture Coordinate sulle UV del Texture
Sample? Perché i nodi World Aligned \textbf{non prendono} un Texture Sample: vogliono un
\textbf{Texture Object}. Si converte: tasto destro sul Texture Sample, \textbf{Convert to
Texture Object}. Il Texture Object ``è un oggetto molto più semplice'': non ha
ingressi, ha una sola uscita, niente RGB; ``è molto più limitato in quello che può
fare'', ma è il ``messaggio'' che serve ai nodi World Aligned.

\paragraph{La prova.} Davide prova diversi valori di Texture Size (1, 2, 10) per avvicinarsi
al risultato dell'altro materiale, con qualche esitazione su quale sia ``troppo'' e quale
``troppo poco'' (``è proprio piccola la texture, allora''), e applica il nuovo
materiale a un secondo oggetto. Poi \textbf{stira} entrambi: il primo si stira; il
secondo ``semplicemente genera nuove texture'', rimane orientato con il mondo,
``rimane sempre della stessa proporzione''.

\paragraph{Dov'è la fregatura.} ``Se questo rimane sempre proporzionato anche quando
stiriamo le cose, perché non dovremmo avere tutti i materiali World Aligned? Dov'è la
fregatura?'' Eccola: essendo allineato al mondo, se l'oggetto si \textbf{muove} o
\textbf{ruota}, il materiale \textbf{resta fermo}. Una \textbf{porta} con un materiale
normale si apre e il materiale resta attaccato alla porta; una porta con un materiale
World Aligned si apre e ``il materiale sembra scorrere sulla porta'', ``come se il
materiale fosse animato''. ``Ci siamo? Ho rotto le palle? Datemi un feedback.''

\begin{figure}[htbp]
\centering
\begin{tikzpicture}[font=\small]
  % materiale UV: stiramento
  \node[font=\footnotesize\bfseries] at (2.6,2.6) {materiale normale (UV)};
  \draw[thick, fill=brown!20] (0,1.0) rectangle (1.2,2.2);
  \foreach \i in {0.2,0.6,1.0} \draw[brown!70!black, thick] (\i,1.15) -- (\i,2.05);
  \draw[-{Stealth[length=2mm]}, thick] (1.5,1.6) -- (2.1,1.6);
  \draw[thick, fill=brown!20] (2.4,1.0) rectangle (5.4,2.2);
  \foreach \i in {0.5,1.5,2.5} \draw[brown!70!black, thick] (2.4+\i,1.15) --
    (2.4+\i,2.05);
  \node[font=\footnotesize, align=center] at (2.7,0.55) {stirato: le venature si
    allargano};
  % materiale world aligned
  \begin{scope}[shift={(6.6,0)}]
    \node[font=\footnotesize\bfseries] at (2.6,2.6) {materiale World Aligned};
    \draw[thick, fill=brown!20] (0,1.0) rectangle (1.2,2.2);
    \foreach \i in {0.2,0.6,1.0} \draw[brown!70!black, thick] (\i,1.15) -- (\i,2.05);
    \draw[-{Stealth[length=2mm]}, thick] (1.5,1.6) -- (2.1,1.6);
    \draw[thick, fill=brown!20] (2.4,1.0) rectangle (5.4,2.2);
    \foreach \i in {0.2,0.6,...,3.0} \draw[brown!70!black, thick] (2.4+\i,1.15) --
      (2.4+\i,2.05);
    \node[font=\footnotesize, align=center] at (2.7,0.55) {stirato: ``genera nuove
      texture'', stessa scala};
  \end{scope}
  % la porta
  \node[font=\footnotesize, align=center, text width=12.5cm] at (5.9,-0.3)
    {la fregatura: se l'oggetto si muove o ruota (una porta che si apre), il materiale
    World Aligned resta fermo nel mondo e ``sembra scorrere'' sull'oggetto};
\end{tikzpicture}
\caption{Materiale legato alle UV e materiale World Aligned a confronto. Il primo
segue l'oggetto e si deforma con lui; il secondo non si deforma, ma non segue
l'oggetto quando si muove. Va bene per pavimenti, muri e rocce statici; non per porte,
veicoli o la torretta.}
\label{ue3:fig:world}
\end{figure}

\paragraph{Le linee nere.} Una studentessa vede delle \textbf{linee nere} al confine tra
una ripetizione della texture e l'altra, sul pavimento. Davide guarda il suo grafo:
stessa texture, Bark, ma ha collegato l'uscita \textbf{XY} invece di \textbf{XYZ}. ``In
quel modo lui aveva due assi che erano allineati col mondo, e l'altro no'', e il
risultato era diverso. ``Ok, grazie.''

\paragraph{Approfondimento: come funziona un World Aligned.} Il nodo World Aligned
Texture, una Material Function del motore, non usa le UV della mesh: per ogni pixel
prende la sua \textbf{posizione nel mondo} e la usa come coordinata della texture. Lo fa
tre volte, proiettando la texture lungo i tre assi (dall'alto, di lato, di fronte), e
mescola le tre proiezioni secondo la direzione della superficie: è la tecnica detta
\textbf{triplanar mapping}. L'uscita XY usa una sola proiezione (dall'alto, buona per i
pavimenti), l'uscita XYZ le mescola tutte; per questo, su un oggetto con facce in tutte
le direzioni, con XY le pareti laterali ricevono una texture ``stirata'' a strisce, che
può dare proprio le linee scure viste dalla studentessa. Il \textbf{Texture Size} è
espresso in \textbf{unità del mondo}, cioè in centimetri: è la grandezza di una
ripetizione della texture. Funziona quindi al contrario del tiling: un tiling più alto
fa ripetere di più la texture, un Texture Size più alto la fa ripetere di meno (una
texture di 100 si ripete ogni metro, una di 10 ogni dieci centimetri). È la ragione
delle esitazioni di Davide su quale valore fosse ``troppo''. Il costo: tre letture della
texture invece di una, quindi un materiale un po' più pesante.

% ---------------------------------------------------------------------
\section{L'esercizio della lava, il nodo Noise e la chiusura}
\label{ue3:sec:lava}

``Allora, ci siamo visti un po' di cose, fatto un miscuglione di cose per quanto riguarda
materiali, texture e import delle mesh.'' Nessuna domanda: ``sembra tutto a posto''.

\paragraph{La consegna.} Per esercitarsi con i materiali, e magari scoprire qualche nodo
non usato a lezione (il nodo \textbf{Noise}, qualche \textbf{maschera}), Davide propone
una prova di una decina di minuti, da continuare poi a casa come le cose di Igor:
\textbf{un materiale di lava}. Di che componenti avrà bisogno? Deve \textbf{muoversi}
un po', quindi bisognerà gestire la velocità del \textbf{Panner}, ``a meno che non vi
cercate proprio la texture della lava''. Ma la richiesta è di improvvisare: ``provate a
vedere cosa riuscite a combinare andando a fare delle operazioni fra colori e fra altre
texture'', con un po' di movimento, un \textbf{Base Color}, ``ma anche un
\textbf{Emissive}, perché la lava comunque deve emettere luce''.

\paragraph{Il ripasso dei nodi.} Una studentessa non ricorda il nome della funzione
``che abbiamo usato per collegare rosso e blu'': il Lerp. Davide ricondivide lo schermo
e fa il riepilogo dei nodi della giornata: la \textbf{costante}, il \textbf{colore}, la
\textbf{texture}, il \textbf{nodo risultante}, il \textbf{Texture Coordinate}, i
\textbf{World Aligned}, le \textbf{operazioni} (con il consiglio sulla moltiplicazione
visto nella sezione~\ref{ue3:sec:costanti}), il \textbf{Panner}, il \textbf{Lerp}.
Divisione e sottrazione servono soprattutto per gli effetti speciali, e Davide vorrebbe
fare una lezione su \textbf{Niagara}, ``il tool che usa Unreal per fare gli effetti
speciali'', se ci sarà tempo.

\paragraph{Le prossime lezioni con Davide.} Davide annuncia come saranno le sue
lezioni finali: ci sarà ``una lezione che sarà un'esercitazione intera'' in cui si
proverà a fare dei \textbf{livelli}, e chiederà agli studenti di condividere lo schermo
e di \emph{streamare} a turno (``se per voi è un problema, magari lo facciamo solo con le
persone per cui non è un problema'', e lo dirà in anticipo). Per esempio, ``un livello
\textbf{platform}'', con le interazioni che si faranno nelle lezioni di Igor e quello che
si vedrà insieme per il \emph{level design}.

\paragraph{Il nodo Noise.} Un nodo ``un po' più complicato'', per chi vuole provare: il
\textbf{Noise}, che ``serve per simulare un \emph{noise}'': il \emph{Voronoi}, il
\emph{White Noise}, un'alternanza di bianco e nero ``che si usa molto spesso per fare
delle distorsioni''. Ha molti parametri: la scala del \emph{noise}, e la possibilità di
filtrarlo con dei vettori o con delle mappe. ``Potete smanettarci un pochino.''

\paragraph{Lo stream della studentessa.} ``C'è qualcuno che, nel mentre sta provando a
fare qualcosa, vuole streamare in questi ultimi cinque minuti?'' Una studentessa ha
provato, ma non riesce a trasformare la texture in World Aligned: è perché usa un
Texture Sample, e serviva un Texture Object. ``Prova a non usare il World Aligned, e
farlo così, diciamo, semplice.'' Senza World Aligned le viene ``questa cosa'', sempre
con il Bark. Davide la guida:
\begin{itemize}
  \item nella lava servono ``delle zone più scure e delle zone più luminose'': quindi
        un \textbf{Lerp} tra un \textbf{colore} e un \textbf{nero}, gestito da un'alfa
        che sia una texture con dei bianchi e dei neri ``posizionati come vuoi'', da
        mandare nell'\textbf{Emissive Color};
  \item l'alfa va cercata tra le texture del progetto, ``qualcosa per dare un effetto
        tipo di distorsione''; le \textbf{Normal}, invece, ``lasciale perdere'': per
        un'alfa si usano le Diffuse o le mappe come Specular e Roughness;
  \item se una texture fa da alfa, nelle entrate del Lerp ci vogliono \textbf{due
        colori}, ``uno che sia molto scuro, uno che sia più chiaro''; ricordando che
        l'alfa ``prende i bianchi e i neri e dice: nei bianchi vai più verso B, nei neri
        vai più verso A'';
  \item il risultato va ``forse più nell'Emissive che non nel Base Color'', perché la
        lava la si vuole ``infuocata, accesa, luccicante''; il Base Color si può anche
        non usare, o usarlo con una texture adatta (``il Bark, per come è fatto, forse
        non è proprio la texture più adatta per fare una lava'');
  \item nell'Emissive si può \textbf{moltiplicare} il risultato, in modo che le zone
        accese siano molto luminose.
\end{itemize}
La studentessa nota che il nero, nell'Emissive, quasi non si vede: ``esatto'', il nero
non emette niente. E usandolo come colore ``non si vede molto'', perché il Bark come alfa
``è molto bianco'', e finisce per scurire un po' il tutto. ``È comunque una cosa
complicata, quindi vedete effettivamente se riuscite a tirar fuori qualcosa.''

\begin{figure}[htbp]
\centering
\begin{tikzpicture}[
  font=\small,
  ue3node/.style={draw, thick, rounded corners=4pt, align=center,
                  minimum height=0.85cm},
  ue3uv/.style={-{Stealth[length=2mm]}, thick, red!60!black},
  ue3flt/.style={-{Stealth[length=2mm]}, thick, green!50!black},
  ue3col/.style={-{Stealth[length=2mm]}, thick, blue!60!black}
]
  \node[ue3node, fill=red!12, text width=1.7cm] (tc) at (0,0) {\textbf{TexCoord}};
  \node[ue3node, fill=gray!10, text width=1.9cm] (pan) at (2.6,0) {\textbf{Panner}\\
    \footnotesize Speed 0,05; 0,02};
  \node[ue3node, fill=gray!15, text width=2.3cm] (msk) at (5.6,0)
    {\textbf{Texture Sample}\\\footnotesize maschera in b/n\\\footnotesize (o un
    Noise)};
  \node[ue3node, fill=black!70, text=white, text width=1.9cm] (dark) at (5.6,2.0)
    {\textbf{scuro}\\\footnotesize quasi nero};
  \node[ue3node, fill=orange!60, text width=1.9cm] (hot) at (5.6,-2.0)
    {\textbf{chiaro}\\\footnotesize arancione};
  \node[ue3node, fill=gray!10, text width=1.3cm, minimum height=2.0cm] (lerp)
    at (8.6,0) {\textbf{Lerp}\\[3pt]\footnotesize A\\\footnotesize B\\
    \footnotesize Alpha};
  \node[ue3node, fill=green!8, text width=1.6cm] (k) at (8.6,-2.4)
    {\footnotesize intensità\\\footnotesize (parametro)};
  \node[ue3node, fill=gray!10, text width=1.4cm] (mul) at (10.9,-0.6)
    {\textbf{Multiply}};
  \node[ue3node, fill=blue!6, text width=2.0cm, minimum height=1.6cm] (res)
    at (13.4,0) {\footnotesize Base Color\\[6pt]\footnotesize\textbf{Emissive}};
  \draw[ue3uv] (tc) -- (pan);
  \draw[ue3uv] (pan) -- node[above, font=\scriptsize]{UVs} (msk);
  \draw[ue3col] (dark.east) -- ++(0.7,0) |- ([yshift=5mm]lerp.west);
  \draw[ue3col] (hot.east) -- ++(0.9,0) |- ([yshift=0mm]lerp.west);
  \draw[ue3flt] (msk.east) -- ++(0.35,0) |- ([yshift=-5mm]lerp.west);
  \draw[ue3col] (lerp.east) -- ++(0.3,0) |- ([yshift=2mm]mul.west);
  \draw[ue3flt] (k.east) -- ++(0.5,0) |- ([yshift=-2mm]mul.west);
  \draw[ue3col] (mul.east) -- ++(0.4,0) |- ([yshift=-3mm]res.west);
  \draw[ue3col, dashed] (lerp.north) |- ([yshift=4mm]res.west);
\end{tikzpicture}
\caption{Una possibile soluzione dell'esercizio della lava, costruita con i soli nodi
della lezione e con i consigli dati alla studentessa: una maschera in bianco e nero che
scorre con il Panner fa da alfa a un Lerp tra uno scuro e un arancione; il risultato,
moltiplicato, va nell'Emissive (e, tratteggiato, eventualmente anche nel Base Color).
Non è un grafo mostrato a lezione: i valori sono indicativi.}
\label{ue3:fig:lava}
\end{figure}

\paragraph{La chiusura.} ``Vabbè, io in realtà adesso ho gettato questa bomba, e penso
che vi lascerò fare. Prendetevi tutto il tempo che volete, nessuna pressione.'' Davide
rivedrà la classe ``a luglio, probabilmente'', per le ultime quattro o cinque lezioni del
corso, sul \textbf{level design} e sulla \textbf{creazione di livelli}. ``Si sono fatte
le 7.'' ``Direi che è stato un piacere, mi sembrate tutti abbastanza interessati,
sicuramente belli, belli, belli attivi.'' E forse ci si rivedrà prima, ``se Igor ha
bisogno di un'altra sostituzione''. ``Ciao a tutti, arrivederci.''

\paragraph{Approfondimento: una lava più convincente.} Partendo dallo schema della
figura~\ref{ue3:fig:lava}, tre accorgimenti, tutti con nodi già visti, rendono la lava
molto più viva. Primo: \textbf{due Panner} con velocità e tiling diversi sulla stessa
maschera, sommati o moltiplicati, così che il movimento non si ripeta mai uguale (lo
stesso trucco dell'acqua). Secondo: un Multiply sull'Emissive con valori ben oltre 1
(10, 20), perché solo così il \emph{bloom} fa ``bruciare'' le zone calde, come nel
materiale con il Glow Intensity. Terzo: usare la \textbf{stessa} maschera per la
\textbf{Roughness}, con un Lerp tra 0,3 (la lava liquida, lucida) e 1 (la crosta
raffreddata, opaca): il materiale cambia così non solo colore ma anche superficie, che è
il senso del PBR. Il nodo Noise di Davide può sostituire la texture della maschera, ma è
calcolato a ogni pixel e costa più di una texture: per un oggetto grande in scena, una
texture di \emph{noise} è di solito la scelta più leggera.

\begin{table}[htbp]
\centering
\small
\begin{tabular}{@{}p{4.4cm}p{8.6cm}@{}}
\toprule
\textbf{Comando} & \textbf{Che cosa fa} \\
\midrule
Tasto destro nel grafo & menu dei nodi, con ricerca \\
T + clic & Texture Sample \\
1 + clic / 2 + clic / 3 + clic & Constant, Constant2Vector, Constant3Vector (il
  colore) \\
A / M / D + clic & Add, Multiply, Divide \\
S + clic & Scalar Parameter (per questo Subtract non ha scorciatoia) \\
U + clic & Texture Coordinate (``lavora sulle UV'') \\
L + clic & Lerp (Linear Interpolate) \\
P + clic & Panner \\
Alt + clic su un pin & stacca i collegamenti \\
Tasto destro su un nodo, Convert to Parameter & trasforma il nodo in un parametro,
  regolabile dalle Material Instance \\
Tasto destro su un Texture Sample, Convert to Texture Object & per i nodi World Aligned \\
Use Selected Asset from Content Browser & mette nel nodo la texture selezionata nel
  Content Browser \\
Tasto destro su un materiale, Create Material Instance & crea un'istanza \\
Reimport Mesh With Dialog & reimporta una mesh riaprendo la finestra di import \\
Ctrl+G / Shift+G & raggruppa e separa gli actor selezionati \\
Actor, Merge Actors & fonde gli actor selezionati in una nuova mesh \\
Ctrl+C, Ctrl+V, Ctrl+Z, Ctrl+Y, Ctrl+A & le scorciatoie classiche, valide quasi
  ovunque \\
Ctrl+S & salvare: ``salvate spesso'' \\
\bottomrule
\end{tabular}
\caption{I comandi della terza lezione nell'editor dei materiali, nel Content Browser e
nel Level Editor di Unreal Engine 5. Le scorciatoie con il clic si usano tenendo premuto
il tasto e cliccando nel grafo; World Aligned Texture e World Aligned Normal non ne
hanno.}
\label{ue3:tab:scorciatoie}
\end{table}

% ---------------------------------------------------------------------
\section*{Punti chiave del capitolo}
\addcontentsline{toc}{section}{Punti chiave del capitolo}

\begin{itemize}
  \item Il grafico lavora soprattutto fuori da Unreal (Blender, Maya); dentro Unreal
        \textbf{importa} gli oggetti e costruisce gli \textbf{shader}.
  \item Una \textbf{texture} è la ``carta regalo'' dell'oggetto: si avvolge sulla mesh
        attraverso le \textbf{UV}, che si creano prima, nel software di modellazione.
        Una texture non pensata per quelle UV si stira.
  \item Le mappe principali: \textbf{Diffuse/Albedo} (il colore, suffisso \_D o \_A),
        \textbf{Normal} (blu: l'illusione del rilievo senza poligoni in più),
        \textbf{Specular}, \textbf{Roughness}, \textbf{Metallic}, la mappa di
        \textbf{altezza} (\emph{displacement} vero, ``una normal sotto steroidi'').
  \item Nelle mappe in bianco e nero il \textbf{bianco è 1}, il \textbf{nero è 0}, il
        grigio sta in mezzo. Ogni canale RGB è una mappa a sé: una texture
        \textbf{AORM}/ORM impacchetta tre mappe in un file.
  \item Il \textbf{materiale} è la ``vernice'', uno \textbf{shader}: si costruisce con
        nodi, come un Blueprint, e combina le texture. Sugli oggetti si applicano
        materiali, non texture.
  \item \textbf{Material Domain} (per che cosa), \textbf{Blend Mode} (Opaque, Masked,
        Translucent) e \textbf{Shading Model} (Default Lit, Unlit) cambiano i pin del
        nodo risultante.
  \item Nodi di base: \textbf{Texture Sample} (T), costanti (1, 2, 3), \textbf{Add},
        \textbf{Multiply}, \textbf{Divide} (A, M, D), Scalar Parameter (S),
        \textbf{Texture Coordinate} (U), \textbf{Lerp} (L), \textbf{Panner} (P). La
        moltiplicazione è l'operazione più usata.
  \item Il rendering è \textbf{PBR}: Roughness 1 è la roccia, 0 lo specchio; Metallic
        e Specular decidono come la superficie riflette; l'\textbf{Emissive}
        moltiplicato fa brillare il materiale.
  \item \textbf{UTiling} e \textbf{VTiling} (o un Multiply sulle coordinate) decidono
        quante volte si ripete la texture: 2 e 2 sono quattro ripetizioni, 0,5 e 0,5 un
        quarto della texture.
  \item Un nodo diventa un \textbf{parametro} con Convert to Parameter. Le
        \textbf{Material Instance} (MI\_) cambiano i parametri senza toccare il grafo
        del \textbf{master material} (M\_), si aggiornano subito e sono più leggere:
        sugli oggetti si applicano sempre le istanze.
  \item Le \textbf{Material Function} sono ``parentesi'' di calcoli riusabili tra
        materiali e tra membri del team.
  \item Il \textbf{Lerp} interpola tra A e B secondo un'\textbf{alfa} da 0 a 1, che può
        essere una texture (``l'arbitro''); il \textbf{Panner} fa scorrere la texture
        (acqua, fluidi, lava).
  \item Importare \textbf{FBX}, controllare le \textbf{unità} (Unreal usa i
        centimetri), texture in PNG o JPEG. La finestra di import genera collisioni,
        corregge il transform, importa materiali e texture.
  \item Una mesh è fatta di \textbf{vertici}, \textbf{spigoli} e \textbf{facce}
        (triangoli). Il \textbf{pivot} deve stare dove il pezzo ruota: con
        \textbf{Reimport Mesh With Dialog} (cannone: Z $-135$, rotazione 113) o con la
        \textbf{Modeling Mode} di Unreal Engine 5.
  \item Le \textbf{simple collision} (box, sfera, capsula: il ``capsulone'' dei
        \emph{character}) sono leggere; le \textbf{complex} usano i triangoli della mesh
        e costano molto di più.
  \item \textbf{Ctrl+G} raggruppa; \textbf{Merge Actors} fonde in una nuova mesh, con
        tanti slot quanti i materiali dei pezzi.
  \item Due piani nella stessa posizione sfarfallano (\textbf{z-fighting}): basta
        alzarne uno di 0,1. I materiali \textbf{World Aligned} (Texture Object, Texture
        Size, uscita XYZ) non si stirano, ma restano fermi se l'oggetto si muove.
  \item \textbf{Salvate spesso.} E per esercizio: un materiale di \textbf{lava}, con
        Panner, Lerp ed Emissive.
\end{itemize}

%
%  Tutte le etichette di questo capitolo hanno il prefisso "ue3:"
%  (ue3:cap:..., ue3:sec:..., ue3:fig:..., ue3:tab:...), e tutti gli
%  stili TikZ hanno il prefisso "ue3" (ue3node, ue3arr, ...) per non
%  collidere con chiavi native di TikZ né con quelli delle prime due
%  lezioni. Il capitolo è autonomo: non contiene \ref verso altri
%  capitoli (i rimandi alle lezioni precedenti sono a parole).
%  Le figure sono ricostruzioni schematiche di quanto mostrato e detto a
%  lezione, non copie delle slide o delle schermate del docente. Dove la
%  trascrizione è troppo rovinata per ricostruire un dettaglio, il
%  capitolo lo dice invece di inventarlo.
% =====================================================================

\chapter[Texture, materiali e mesh]{Texture, materiali e mesh: la grafica dentro
Unreal e la torretta}
\label{ue3:cap:materiali}

Questa è la \textbf{terza lezione} del corso breve di \textbf{introduzione a Unreal
Engine}, e cambia sia docente sia argomento. Al posto di Igor, che ``aveva un
impiccio'', c'è \textbf{Davide}, che si presenta così: ``io sono Davide e in questa
lezione, nello specifico, sostituisco Igor''. Davide tornerà più avanti: ``le ultime
lezioni del corso le fate con me'', e saranno lezioni ``un po' più basate su
\textbf{level design} e la creazione di mappe'', con qualche esercitazione insieme. Dopo
due lezioni sull'editor e sulla logica (piazzare oggetti, i \emph{transform}, i
Blueprint, la pallina sulla piattaforma), questa lezione guarda a ciò che ``fa il
grafico all'interno di Unreal'': le \textbf{mesh}, le \textbf{texture}, i
\textbf{materiali}, e ``qualche nodo, qualche cosa carina che potete fare'' combinando
mesh e materiali.

La lezione ha due metà separate da una pausa. La prima è sui \textbf{materiali}: che
cos'è una texture e come si avvolge su un oggetto attraverso le \textbf{UV}; quali tipi
di mappa esistono (\textbf{Diffuse}, \textbf{Normal}, \textbf{Specular}, la mappa di
altezza per il \emph{displacement}) e perché in queste mappe il bianco vale 1 e il nero
0; che cos'è un \textbf{materiale} e come si costruisce nell'\textbf{editor dei
materiali}, ``nello stesso modo dell'editor dei Blueprint''; le impostazioni principali
(\textbf{Material Domain}, \textbf{Blend Mode}, \textbf{Shading Model}); i nodi di base
(\textbf{Texture Sample}, le costanti, le operazioni, il \textbf{Texture Coordinate},
il \textbf{Lerp}, il \textbf{Panner}); il rendering \textbf{PBR} con Metallic,
Specular e Roughness; e le \textbf{Material Instance}, il modo ``bon ton'' di applicare i
materiali. La seconda metà è sulle \textbf{mesh}: si importa una \textbf{torretta} in
tre pezzi dal materiale didattico, si guarda la finestra di \textbf{import}, si
corregge il \textbf{pivot} del cannone con un \textbf{Reimport}, si esplorano lo Static
Mesh Editor e le \textbf{collisioni}, si fondono i pezzi con \textbf{Merge Actors}, si
creano a mano i materiali della torretta, e si chiude con i materiali \textbf{World
Aligned}, che non si deformano quando si stira l'oggetto. In fondo, un esercizio libero:
un materiale di \textbf{lava}.

La lezione è \textbf{online}, su Zoom, come le precedenti, e finisce poco prima delle
19 (``si sono fatte le 7''). La pausa, una decina di minuti, arriva dopo ``un'oretta e
mezza''. Davide, alla sua prima volta con questa classe, combatte per tutta la lezione
con la condivisione dello schermo (``io sto facendo a cazzotti da tutta la lezione con
la finestra di Zoom'', ``quanto non sono abituato a streamare'') e con un PC di
emergenza: ``il fisso mi è esploso letteralmente tre giorni fa'', e a metà della seconda
parte il computer ``inizia ad accusare la stream'', tanto che Davide gira il
ventilatore verso di lui. Chiede spesso conferma che gli studenti vedano quello che
mostra: ``mi impanicherò ogni tanto a caso e vi chiederò conferma delle robe''.

Tre avvertenze di lettura, le stesse delle prime due lezioni. La prima: la trascrizione
è automatica e storpia moltissimi nomi. Unreal diventa ``ArrayLink'', ``ArrayL4'',
``Arial Engine'' e ``a Real''; la \textbf{Material Instance} è quasi sempre
``materialistance'', ``Material Distance'' e una volta ``mie dire distanze''; il
\textbf{Panner} è ``panel'', ``planner'' e ``banner''; il \textbf{Lerp} è ``l'ERP'';
\emph{World Aligned} è ``word aligned'' e ``word.allocation''; la texture
\textbf{T\_Bark} (una corteccia) è ``Arc'', ``barca'', ``Barc'' e ``park''; ``strecciare''
e ``si straccia'' stanno per lo \emph{stretch}, lo stiramento; ``traccare'' sta per
\emph{tracciare} un bersaglio (\emph{to track}); la \textbf{mappa di altezza}
(\emph{height map}) è quasi sempre ``light map'' o ``lightmap''. Dove la ricostruzione
è sicura il capitolo usa la forma corretta e, la prima volta, lo segnala in nota; dove
non lo è riporta la forma della trascrizione e lo dice. La seconda: durante la pausa e
verso la fine il riconoscimento vocale produce lunghi blocchi di testo senza senso e
ripetuto (``qual è la risposta che ci può essere?'', discorsi sulla ``vita data per
scontata'' e su ``vittime''): sono artefatti, non parole di nessuno, e il capitolo li
ignora. La terza: gran parte della lezione è una dimostrazione dal vivo nell'editor dei
materiali e nello Static Mesh Editor. Il capitolo ricostruisce i grafi di nodi da ciò
che Davide dice mentre li costruisce, e li completa, dove serve, con i nomi esatti dei
nodi e dei comandi di Unreal Engine 5.

\paragraph{Come leggere questo capitolo.} Il capitolo segue l'ordine della lezione. Le
figure ricostruiscono gli schemi spiegati a voce, sulle slide o nell'editor: le mappe
della torretta, i canali di una texture, il rapporto tra texture, materiale e mesh, i
grafi dei materiali costruiti a lezione, le Material Instance, la struttura di una
mesh, il pivot del cannone. Le tabelle riassumono i tipi di mappa, le impostazioni del
materiale, le opzioni di import, le collisioni e le scorciatoie dell'editor dei
materiali. Qua e là, dei paragrafi di \textbf{approfondimento} collegano la lezione al
libro \emph{Blueprints Visual Scripting for Unreal Engine 5} di Marcos Romero e Brenden
Sewell, il testo di riferimento per questa parte del libro, e a qualche nozione tecnica
generale su Unreal Engine 5; sono segnalati come tali e non fanno parte della lezione.

Il percorso del capitolo:
\begin{enumerate}
  \item l'attesa, le presentazioni e il programma; che cosa fa un grafico dentro
        Unreal; il materiale didattico (sezioni \ref{ue3:sec:attesa}
        e~\ref{ue3:sec:didattico});
  \item le \textbf{texture}: la carta regalo, le \textbf{UV}, i tipi di mappa, il
        bianco e il nero, i \textbf{canali RGB} (sezioni da~\ref{ue3:sec:texture}
        a~\ref{ue3:sec:canali});
  \item i \textbf{materiali}: la vernice, l'editor, Material Domain, Blend Mode e
        Shading Model, i nodi (sezioni da~\ref{ue3:sec:materiale}
        a~\ref{ue3:sec:nodi});
  \item in Engine: la mappa vuota, dentro una texture, il \textbf{Texture Sample}, le
        costanti e le scorciatoie, il colore e l'\textbf{emissive} (sezioni
        da~\ref{ue3:sec:mappavuota} a~\ref{ue3:sec:emissive});
  \item il \textbf{PBR} e la Normal (sezione \ref{ue3:sec:pbr});
  \item il \textbf{Texture Coordinate}, il tiling e il pavimento (sezione
        \ref{ue3:sec:tiling});
  \item le \textbf{Material Instance} e le Material Function (sezioni
        \ref{ue3:sec:instance} e~\ref{ue3:sec:function});
  \item il \textbf{Lerp} e il \textbf{Panner}, poi la pausa (sezioni
        da~\ref{ue3:sec:lerp} a~\ref{ue3:sec:pausa});
  \item l'\textbf{import} e la torretta; che cos'è una \textbf{mesh}; dimensioni e
        formati (sezioni da~\ref{ue3:sec:import} a~\ref{ue3:sec:regole});
  \item il \textbf{pivot} del cannone e il \textbf{Reimport} (sezione
        \ref{ue3:sec:pivot});
  \item lo \textbf{Static Mesh Editor} e le \textbf{collisioni}; raggruppare e
        \textbf{Merge Actors} (sezioni \ref{ue3:sec:sme} e~\ref{ue3:sec:merge});
  \item i \textbf{materiali della torretta} e le domande sul Reimport e sul merge
        (sezioni \ref{ue3:sec:matturret} e~\ref{ue3:sec:domande});
  \item lo stretch, lo \textbf{z-fighting} e i materiali \textbf{World Aligned}
        (sezione \ref{ue3:sec:world});
  \item l'esercizio della \textbf{lava}, il nodo \textbf{Noise} e la chiusura
        (sezione \ref{ue3:sec:lava}).
\end{enumerate}

% ---------------------------------------------------------------------
\section{L'attesa, le presentazioni e il programma}
\label{ue3:sec:attesa}

``Ok, buonasera.'' La registrazione parte subito (``sta registrando? Sì, sta
registrando, perfetto''), e Davide aspetta ``qualche minutino'' che entrino tutti. Uno
studente avvisa che dovrà staccarsi venti minuti prima per andare a lavorare: ``nessun
problema'', tanto la lezione resta registrata, e anzi a Davide piacerebbe finire ``con
voi che fate qualcosina'', quindi andare via un po' in anticipo ``non è neanche un
problema''. Davide chiede quante persone c'erano alle lezioni precedenti (``una dozzina,
credo''), ringrazia chi accende la telecamera (``almeno so che non parlo da solo'') e
controlla che non si senta il ventilatore in sottofondo (``inizia a far caldo, inizia a
essere complicato'').

\paragraph{Che cosa avete fatto con Igor?} Prima di cominciare, Davide vuole sapere a
che punto è la classe, per non ripetere cose già viste. Dal Discord ha capito che si è
``ancora all'inizio'': le schede, le finestre, i \emph{transform} sugli oggetti, i
tavolini, ``qualcosa con gli input, con la pallina che si muove''. ``Se vediamo
qualcosa che avete già visto, ditemelo.''

\paragraph{Il programma.} La lezione di oggi guarda a quello che fa ``il grafico
all'interno di Unreal'': ``un po' di nozioni base per quanto riguarda le mesh, un po' di
nozioni base per quanto riguarda i materiali, le texture'', e ``già qualche nodo,
qualche cosa carina che potete fare'' usando mesh e materiali.

\paragraph{Chi c'è.} Davide chiede a che cosa sono iscritti gli studenti, ``per capire
chi ho davanti''. Le risposte: \textbf{game design} (più d'uno), \textbf{3D art}
(tre persone), \textbf{character animation} (due), \textbf{narrative design}. ``Beh, in
realtà siamo abbastanza assortiti.''

\paragraph{Che cosa fa un grafico dentro Unreal.} Ai futuri grafici 3D Davide mette
``le mani avanti'': la maggior parte del lavoro di un grafico \textbf{non} si svolge
dentro Unreal, ma in altri strumenti, che verranno spiegati nei rispettivi corsi.
Quello che si vede oggi è ciò che un grafico fa ``nel momento in cui va a lavorare con
altre persone'' dentro Unreal Engine: \textbf{importare} gli oggetti e creare gli
\textbf{shader} dentro l'engine.

\paragraph{Come si struttura la lezione.} Davide chiede come lavorava Igor: tutta
teoria e poi tutta pratica, o un po' e un po'? ``Andava prima con la presentazione e poi
con Unreal.'' Davide prova a replicare: ``ci facciamo prima una chiacchiera sulla parte
un po' più teorica e poi passiamo al progetto''. Nella pratica, come si vedrà, teoria
ed engine si alterneranno di continuo.

% ---------------------------------------------------------------------
\section{Il materiale didattico}
\label{ue3:sec:didattico}

Sul Discord c'è stato ``un po' di casino'' per il \textbf{materiale didattico}. ``Non è
molta roba'': è un piccolo progetto di Unreal, ``un progetto che Unreal usa per
spiegare alcune cose'', e probabilmente quello che Igor farà usare nelle lezioni
successive ogni volta che si implementa qualcosa. Gli altri file della cartella sono
\textbf{asset} da importare nel progetto poco alla volta: oggi la \textbf{torretta};
quando si farà l'audio, l'audio; ``poi vediamo se aggiungere altre cose''.

Dove si trova, Davide lo spiega più avanti a una studentessa: sul Discord, ``il 19
maggio, se non ricordo male'', c'è un link a una pagina di Drive con un file che si
chiama, più o meno, ``materiale didattico Unreal Engine Intro''\footnote{La trascrizione
dice ``IVE, materiale didattico, a Real Engine Intro'' e, poco prima, ``Hive Unreal
Engine 5 intro'': la parola iniziale non è ricostruibile con certezza.}. Lo si
scarica, lo si \textbf{decomprime}, e dentro c'è ``un progettino con un
\emph{character}'': sono soprattutto cose per le prossime lezioni con Igor, e per oggi
serve poco.

% ---------------------------------------------------------------------
\section{Che cos'è una texture}
\label{ue3:sec:texture}

\paragraph{La carta regalo.} ``Che cos'è una texture? La texture fondamentalmente è il
\emph{wrap}, è la \textbf{carta regalo} in cui è avvolto il nostro oggetto.'' Sulla
slide c'è la torretta del materiale didattico e, accanto, le sue texture: un'immagine
``tutta scomposta e segmentata e divisa in varie parti''. Perché divisa in parti?
Perché quelle parti poi ``si avvolgeranno sopra la mesh''. Le texture sono ``una serie
di mappe, una serie di informazioni che vanno a dare delle qualità specifiche
all'oggetto''. Con un'altra immagine: la texture ``non è altro che una
\textbf{decalcomania}'' dell'oggetto, della \emph{mesh} che si mette in scena.

\paragraph{Le UV.} Perché una texture funzioni bene su un oggetto, deve essere pensata
``con in mente le \textbf{UV} dei modelli su cui si lavora''. Le UV ``sono le
coordinate'', tutto quello che riguarda ``la struttura dell'oggetto'': la \textbf{mappa
delle UV} è la \textbf{prima parte} del processo di creazione di una texture. Dal
modello 3D si estraggono informazioni che poi servono a creare ``tutti i vari tipi di
texture che esistono'': sulla slide le texture della torretta sono diverse, e ognuna
porta informazioni diverse, ``ma sono tutte texture pensate per la stessa torretta''.

\paragraph{L'UV mapping.} L'\textbf{UV mapping} è la tecnica che serve a
\textbf{mappare le texture 2D su modelli 3D}: si assegnano le coordinate dette UV ``ai
vertici e agli \emph{edge} del modello'', e tutte le texture applicate dopo usano quelle
informazioni ``per precisione''. È un lavoro che si fa \textbf{prima} di entrare in
Unreal, in un altro software ``come Blender, come Maya''. Se si usa una texture che non
è stata pensata per le UV di quell'oggetto, ``la texture verrà strecciata''
(\emph{stretched}, stirata) o comunque non darà il risultato voluto.

\begin{figure}[htbp]
\centering
\begin{tikzpicture}[
  font=\small,
  ue3box/.style={draw, thick, align=center, minimum height=0.9cm},
  ue3arr/.style={-{Stealth[length=2.5mm]}, thick}
]
  % il cubo 3D
  \draw[thick, fill=orange!25] (0,0) -- (1.6,0) -- (1.6,1.6) -- (0,1.6) -- cycle;
  \draw[thick, fill=orange!40] (0,1.6) -- (0.7,2.2) -- (2.3,2.2) -- (1.6,1.6) -- cycle;
  \draw[thick, fill=orange!15] (1.6,0) -- (2.3,0.6) -- (2.3,2.2) -- (1.6,1.6) -- cycle;
  \node[font=\footnotesize, align=center] at (1.15,-0.5) {modello 3D\\(la mesh)};
  % l'apertura
  \draw[ue3arr] (2.8,1.1) -- node[above, font=\footnotesize]{si ``apre''}
    node[below, font=\footnotesize]{(UV mapping)} (5.0,1.1);
  % la mappa UV a croce
  \begin{scope}[shift={(5.4,-0.6)}]
    \draw[step=0.75, gray!40, thin] (0,0) grid (3.75,3.0);
    \foreach \x/\y in {0.75/0, 0/0.75, 0.75/0.75, 1.5/0.75, 2.25/0.75, 0.75/1.5}
      \draw[thick, fill=orange!25] (\x,\y) rectangle ++(0.75,0.75);
    \draw[->, thick] (0,0) -- (4.0,0) node[right, font=\footnotesize]{U};
    \draw[->, thick] (0,0) -- (0,3.25) node[above, font=\footnotesize]{V};
    \node[font=\footnotesize, align=center] at (1.9,-0.5) {mappa UV: le facce\\
      stese sul piano $[0,1]^2$};
  \end{scope}
  % la texture
  \draw[ue3arr] (9.4,1.1) -- node[above, font=\footnotesize]{si dipinge} (10.6,1.1);
  \node[ue3box, minimum width=2.3cm, minimum height=2.3cm, fill=orange!10,
    text width=2.1cm] (tx) at (12.0,1.1) {\footnotesize texture 2D\\
    \footnotesize (Diffuse, Normal,\\\footnotesize Roughness\ldots)\\[3pt]
    \footnotesize si ``riavvolge''\\\footnotesize sulla mesh};
\end{tikzpicture}
\caption{La texture come ``carta regalo'': le facce del modello vengono stese su un
piano (la mappa UV, fatta in Blender o in Maya), le texture si dipingono su quel piano,
e poi si ``riavvolgono'' sull'oggetto. È per questo che la texture della torretta, sulla
slide, appare ``tutta scomposta e segmentata''. Schema ricostruito; le lettere U e V
indicano i due assi della texture, per distinguerli da X, Y e Z del mondo.}
\label{ue3:fig:uv}
\end{figure}

\paragraph{Approfondimento: perché U e V.} Le lettere U e V sono scelte solo per non
confondere gli assi della texture con gli assi X, Y e Z dello spazio 3D. Una coordinata
UV è una coppia di numeri, di solito tra 0 e 1: $(0,0)$ è un angolo della texture e
$(1,1)$ l'angolo opposto. Ogni vertice della mesh riceve una coppia $(u,v)$, e per i
punti interni a una faccia il motore interpola le coordinate dei vertici. Se una mesh
non ha UV (o ne ha di sbagliate), il motore non sa ``quale pezzo di carta regalo'' va
su quale faccia: di qui le texture stirate di cui parla Davide.

% ---------------------------------------------------------------------
\section{I tipi di mappa}
\label{ue3:sec:mappe}

Davide passa in rassegna, sulla slide della torretta, le texture una per una, ``piano
piano'', con le informazioni che portano.

\paragraph{La Diffuse, o Albedo.} La prima è la \textbf{Diffuse}, ``in alcuni casi
detta anche \textbf{Albedo}''. Nei file la si riconosce dalla \textbf{nomenclatura}:
``T maiuscolo, \emph{underscore}, nome dell'oggetto, \emph{underscore}, A'', o D: A per
\emph{albedo}, D per \emph{diffuse} (per esempio \texttt{T\_Turret\_D}). La Diffuse è
il \textbf{colore}: ``tutti i dettagli estetici colorati''. Sulla slide è quella in alto
a destra, con l'arancione: è la texture che dà alla torretta l'arancione sulla base, il
nero negli altri dettagli e il bianco in alcune zone. Più avanti Davide ne completa la
descrizione: la Diffuse applica ``i colori e i dettagli di base'' sulle superfici 3D, e
dà l'informazione del colore, del tipo di materiale, dei dettagli (``queste linee che
vedete qua'') e di quanto un oggetto sia ruvido o levigato, se per esempio si applica la
texture di una roccia o di un metallo.

\paragraph{La Normal.} Poi c'è la \textbf{Normal}, la \textbf{Normal Map}, che ``serve
a simulare una tridimensionalità'' dentro un'informazione che è 2D. Si riconosce
subito: è \textbf{blu}, ``sempre mappe blu con delle sfumature di rosso e di verde''.
L'esempio classico è un \textbf{muro di mattoni}: in gioco il muro è \textbf{piatto}, ma
quando gli si applica la Normal il calcolo delle luci e delle ombre ``prende in
considerazione la normal map'' e proietta le ombre in base a lei, dando l'illusione che
i mattoni siano davvero tridimensionali. La Normal dà informazioni sulla
tridimensionalità ``senza appesantire particolarmente il calcolo'', perché non deforma
davvero l'oggetto come fa una mappa di altezza.

Più avanti Davide ci torna con un esempio sulla torretta. Su una delle gambe della base
servono dei \textbf{graffi}. Si potrebbe cambiare il modello aggiungendo i graffi, ma
questo aggiungerebbe \textbf{poligoni}, e quindi ``potenza di calcolo non
necessaria''. Si può invece aggiungere il dettaglio ``direttamente in 2D'', con la
Normal: protuberanze, graffi, rughe. ``Per un singolo oggetto non è un appesantimento
così importante. Ma considerate quanti oggetti ci sono all'interno di una mappa quando
andate a creare un livello'': se ogni oggetto non è creato con la giusta accortezza,
``si fa molto velocemente ad appesantire il progetto più di quanto dovrebbe essere
pesante''.

\paragraph{La Specular.} Lo \textbf{Specular mapping} riguarda ``la brillantezza e la
riflettività dell'oggetto'': come l'oggetto gestisce la luce, come la luce ``rimbalza''
su di lui. Un oggetto può avere parti metalliche, parti in pietra e altre ancora, e ``a
seconda del tipo di materiale di cui l'oggetto è fatto, rifrangerà la luce in una
maniera diversa''. Queste mappe sono quelle ``in bianco e in nero''. La texture
determina come si calcolano i riflessi ``sui singoli pixel proprio dell'oggetto'',
quanta luce viene riflessa; tornerà utile quando si vedranno le luci in Unreal, perché
``cambiando un materiale la luce ha un effetto completamente diverso''.

\paragraph{La mappa di altezza.} Poi c'è la mappa che serve a creare
\textbf{displacement}\footnote{Nella trascrizione è sempre ``light map'' o
``lightmap'', ma la descrizione (rialza e abbassa zone dell'oggetto, crea
\emph{tessellation}, è ``una normal sotto steroidi'') è quella della \textbf{height
map}, la mappa di altezza o di \emph{displacement}. Le \emph{lightmap} vere sono
un'altra cosa: texture in cui il motore salva l'illuminazione precalcolata.}: crea
``una \emph{tessellation} nell'oggetto'', rialza alcune zone, ne abbassa altre,
``modifica leggermente l'oggetto''. Davide la chiama ``una normal sotto steroidi'',
perché fa davvero \emph{displacement}: aggiunge profondità e dettagli tridimensionali
alle superfici piane, e usa le scale di grigio per creare scanalature e protuberanze.
``Nelle parti in cui vedete il bianco è sempre il sì'', la protuberanza; nel nero,
l'effetto opposto.

\paragraph{Il texturing procedurale.} Infine il \textbf{procedural texturing}: il
processo di creare texture, materiali o \textbf{shader} ``attraverso degli
algoritmi''. ``Andando ad utilizzare dei calcoli matematici, noi andiamo a modificare
delle informazioni per avere un risultato.'' È esattamente quello che fa l'editor dei
materiali di Unreal.

\begin{table}[htbp]
\centering
\small
\begin{tabular}{@{}p{2.5cm}p{2.4cm}p{8.1cm}@{}}
\toprule
\textbf{Mappa} & \textbf{Aspetto} & \textbf{Che cosa porta} \\
\midrule
Diffuse (Albedo) & a colori & il colore e i dettagli estetici; suffisso \_D o \_A \\
Normal & blu, con sfumature rosse e verdi & l'illusione del rilievo: la luce e le
  ombre si calcolano come se la superficie avesse graffi, mattoni, rughe, senza
  aggiungere poligoni \\
Specular & bianco e nero & quanto e come la superficie riflette la luce \\
Roughness & bianco e nero & quanto la superficie è ruvida (1) o liscia (0) \\
Metallic & bianco e nero & dove la superficie è metallo (1) e dove no (0) \\
Ambient Occlusion & bianco e nero & le zone in ombra perché ``chiuse'' (pieghe,
  angoli) \\
Height (displacement) & bianco e nero & deforma davvero la superficie: il bianco
  rialza, il nero abbassa \\
\bottomrule
\end{tabular}
\caption{I tipi di mappa nominati a lezione. Diffuse, Normal, Specular e la mappa di
altezza sono quelle della slide; Roughness, Metallic e Ambient Occlusion compaiono
nella parte sulle mappe impacchettate e nei materiali della torretta.}
\label{ue3:tab:mappe}
\end{table}

\paragraph{Il bianco è 1, il nero è 0.} La cosa importante da capire, insiste Davide, è
che nelle mappe in bianco e nero ``dovete sempre considerare il \textbf{bianco come un
1}, il \textbf{nero come uno 0} e il \textbf{grigio come dei valori nel mezzo} fra lo 0 e
l'1''. Che cosa vuol dire? Che in Unreal, per esempio, il valore dello Specular va
soltanto tra 0 e 1: a 1 riflette la luce, a 0 no. Una mappa di Specular, quindi, va
letta come una griglia di numeri: ``ci sono degli 1 in queste parti e degli 0 nel resto
delle parti''. È la chiave per capire quasi tutto il resto della lezione.

\paragraph{Approfondimento: la mappa di altezza, la Normal e Unreal Engine 5.} La
differenza tra Normal e mappa di altezza è la differenza tra \emph{sembrare} e
\emph{essere}. La Normal Map memorizza, per ogni pixel, la direzione in cui ``guarda''
la superficie (la sua \emph{normale}): i tre numeri X, Y e Z della direzione sono
salvati nei canali rosso, verde e blu, e siccome la maggior parte dei pixel guarda
``dritto fuori'' dalla superficie (Z alto) la mappa appare blu. La geometria non
cambia: di profilo, un muro con Normal Map resta piatto. La mappa di altezza invece
sposta davvero i vertici, e per farlo servono vertici da spostare: di qui la
\emph{tessellation}, che suddivide la superficie in molti triangoli. In Unreal Engine 5
il \emph{displacement} vero si fa soprattutto con le mesh \textbf{Nanite}, ed è più
costoso della Normal: per questo, per graffi e piccoli dettagli, la Normal resta la
scelta di default, come dice Davide.

% ---------------------------------------------------------------------
\section{I canali RGB e le mappe impacchettate}
\label{ue3:sec:canali}

``Chiaramente, queste texture utilizzano l'\textbf{RGB}'': non hanno un solo canale, ma
\textbf{tre canali} su cui passano le informazioni. Sulla slide c'è una Diffuse
scomposta in rosso, verde e blu: ognuno porta informazioni separate, che combinate
danno il colore che serve.

Ma c'è un altro modo di vedere la cosa. \textbf{Singolarmente}, ogni canale appare come
un'immagine in \textbf{bianco, nero e grigio}. E bianchi, neri e grigi ``sono
combinazioni di 0 e di 1'': quindi ogni canale si può usare per trasmettere
\textbf{un'informazione diversa}. È il \textbf{multichannel}. Una pratica comune,
nell'esportare e importare texture, è avere una singola mappa ``che spesso e volentieri
viene siglata come \textbf{AORM}'', per \textbf{Ambient Occlusion, Roughness e
Metallic}: tre mappe in bianco e nero, ognuna in un canale. ``Magari noi vogliamo
utilizzare questa mappa qui sopra per l'ambient occlusion, questa mappa qui sotto per
le metallic'', e così via. Una sola mappa con più informazioni invece di tante mappe
diverse: ``va ad alleggerire il progetto, va a rendere più veloce, più modulare''
l'export delle texture.

\begin{figure}[htbp]
\centering
\begin{tikzpicture}[
  font=\small,
  ue3ch/.style={draw, thick, minimum width=2.4cm, minimum height=1.0cm,
                align=center},
  ue3arr/.style={-{Stealth[length=2.5mm]}, thick}
]
  \node[draw, very thick, minimum width=2.6cm, minimum height=3.6cm, align=center,
    fill=gray!10] (tex) at (0,0) {\textbf{una texture}\\\textbf{RGB}\\[4pt]
    \footnotesize es. T\_Turret\_ORM};
  \node[ue3ch, fill=red!15] (r) at (4.2,1.3) {canale \textbf{R}};
  \node[ue3ch, fill=green!15] (g) at (4.2,0) {canale \textbf{G}};
  \node[ue3ch, fill=blue!15] (b) at (4.2,-1.3) {canale \textbf{B}};
  \draw[ue3arr] (tex.east) ++(0,1.0) -- (r.west);
  \draw[ue3arr] (tex.east) -- (g.west);
  \draw[ue3arr] (tex.east) ++(0,-1.0) -- (b.west);
  \node[ue3ch, text width=3.4cm] (ao) at (8.6,1.3) {\textbf{Ambient Occlusion}\\
    \footnotesize grigi tra 0 e 1};
  \node[ue3ch, text width=3.4cm] (ro) at (8.6,0) {\textbf{Roughness}\\
    \footnotesize grigi tra 0 e 1};
  \node[ue3ch, text width=3.4cm] (me) at (8.6,-1.3) {\textbf{Metallic}\\
    \footnotesize grigi tra 0 e 1};
  \draw[ue3arr] (r) -- (ao);
  \draw[ue3arr] (g) -- (ro);
  \draw[ue3arr] (b) -- (me);
  \node[font=\footnotesize, align=center, text width=11cm] at (4.3,-2.5)
    {ogni canale, da solo, è un'immagine in bianco e nero: tre ``mappe'' diverse in
    un solo file};
\end{tikzpicture}
\caption{Una mappa impacchettata (\emph{channel packing}). La lezione la chiama AORM;
nella pratica di Unreal la convenzione più diffusa è ORM, con l'Ambient Occlusion nel
rosso, la Roughness nel verde e il Metallic nel blu, ma l'ordine dipende da chi
esporta le texture.}
\label{ue3:fig:canali}
\end{figure}

\paragraph{Approfondimento: sRGB e compressione.} Quando si importa in Unreal una
texture impacchettata (o qualunque mappa ``di dati'', come Roughness o Metallic), conviene
aprirla e \textbf{togliere la spunta a sRGB}: l'opzione sRGB dice al motore che
l'immagine è un \emph{colore} da correggere per la percezione dell'occhio, mentre una
mappa di dati contiene numeri da leggere così come sono. Nelle stesse impostazioni, il
campo \emph{Compression Settings} si mette su \textbf{Masks} per le mappe impacchettate
e su \textbf{Normalmap} per le Normal (Unreal di solito lo riconosce da solo per le
texture con suffisso \_N). Se ci si dimentica, il materiale funziona lo stesso, ma i
grigi vengono letti un po' sbagliati e la superficie risulta più lucida o più opaca del
previsto.

% ---------------------------------------------------------------------
\section{Il materiale: la vernice}
\label{ue3:sec:materiale}

\paragraph{Le texture non si modificano.} ``Quello che possiamo fare con le texture in
Unreal è relativamente poco, perché non possiamo andare a modificare effettivamente le
texture all'interno dell'engine.'' E allora come si danno qualità visive alle mesh?
``Lo facciamo con i \textbf{materiali}.''

\paragraph{Il materiale è uno shader.} I materiali ``sono degli \textbf{shader}''. E
l'\textbf{editor dei materiali}, quello che si vedrà ``principalmente'' oggi, ``funziona
praticamente nello stesso modo dell'editor dei Blueprint'': una serie di \textbf{nodi}
collegati, che trasportano informazioni fino a un \textbf{nodo risultante}.

\paragraph{Vernice, non carta regalo.} La distinzione è netta: ``il materiale è
effettivamente la \textbf{vernice} che viene applicata sulla mesh, non la texture''. Sugli
oggetti si applicano \textbf{materiali}. Le texture servono a \textbf{costruire} i
materiali: si combinano le texture in un unico materiale, che poi si applica
sull'oggetto finale. Il materiale dice all'oggetto ``come dovrebbe interagire con la
luce, che tipo di informazioni visive dovrebbe portare'', dove ci sono
tridimensionalità e così via: ``tutto quello che abbiamo visto dalle varie texture
precedentemente viene tutto convogliato all'interno dei materiali''. Questo, precisa
Davide, se si parla di \textbf{oggetti fisici}. I materiali hanno anche altre funzioni:
servono per esempio a fare \textbf{effetti speciali}, e un materiale per un effetto
speciale ``dovrà essere costruito in maniera diversa'' da quello che riveste una
torretta.

\begin{figure}[htbp]
\centering
\begin{tikzpicture}[
  font=\small,
  ue3tex/.style={draw, thick, minimum width=2.3cm, minimum height=0.7cm,
                 align=center, fill=gray!10},
  ue3big/.style={draw, very thick, rounded corners=4pt, align=center},
  ue3arr/.style={-{Stealth[length=2.5mm]}, thick}
]
  \node[ue3tex] (t1) at (0,1.5) {T\_Turret\_D};
  \node[ue3tex] (t2) at (0,0.5) {T\_Turret\_N};
  \node[ue3tex] (t3) at (0,-0.5) {T\_Turret\_R};
  \node[ue3tex] (t4) at (0,-1.5) {T\_Turret\_M};
  \node[font=\footnotesize, align=center] at (0,2.4) {\textbf{texture}\\(immagini)};
  \node[ue3big, minimum width=3.4cm, minimum height=3.6cm, fill=blue!6,
    text width=3.2cm] (mat) at (4.4,0) {\textbf{materiale}\\M\_Turret\\[6pt]
    \footnotesize un grafo di nodi\\\footnotesize (uno \emph{shader})\\[6pt]
    \footnotesize la ``vernice''};
  \foreach \tn in {t1,t2,t3,t4} \draw[ue3arr] (\tn.east) -- (mat.west |- \tn);
  \node[ue3big, minimum width=2.8cm, minimum height=1.6cm, fill=orange!12,
    text width=2.6cm] (mi) at (8.4,0) {\textbf{Material}\\\textbf{Instance}\\
    \footnotesize MI\_Turret\\\footnotesize (parametri)};
  \draw[ue3arr] (mat) -- (mi);
  \node[ue3big, minimum width=2.4cm, minimum height=1.6cm, fill=green!8,
    text width=2.2cm] (mesh) at (11.8,0) {\textbf{mesh}\\\footnotesize SM\_Turret\\
    \footnotesize slot del\\\footnotesize materiale};
  \draw[ue3arr] (mi) -- node[above, font=\footnotesize]{si applica} (mesh);
\end{tikzpicture}
\caption{Texture, materiale, istanza e mesh. Le texture non si applicano direttamente
agli oggetti: entrano in un materiale, che è la ``vernice''; dal materiale si crea
un'istanza, ed è l'istanza che si applica allo slot del materiale della mesh
(sezione~\ref{ue3:sec:instance}). I nomi seguono la nomenclatura suggerita a lezione
(T\_, M\_, MI\_, SM\_).}
\label{ue3:fig:pipeline}
\end{figure}

\paragraph{Approfondimento: il materiale nel libro di Romero e Sewell.} Il quinto
capitolo del libro (``Object Interaction with Blueprints'') usa la stessa immagine di
Davide: un materiale è un \emph{asset} che si aggiunge alla mesh di un actor per
crearne l'aspetto, ``come uno strato di vernice''. Il libro crea un materiale rosso,
\texttt{M\_TargetRed}, in una cartella Materials (``non è necessario, ma è buona
pratica tenere ordinata la gerarchia del progetto''), e nota che l'editor dei
materiali ``condivide molte caratteristiche e convenzioni con i Blueprint'': al centro
c'è il grafo; il nodo grande con il nome del materiale è il \textbf{result node}; con il
tasto destro si apre un menu di nodi con una ricerca sensibile al contesto; i fili si
tirano trascinando dai pin, gli ingressi a sinistra e le uscite a destra. Il colore lo
dà un \textbf{VectorParameter} collegato a Base Color; due \textbf{ScalarParameter} a
0,1 danno un po' di Metallic e di Roughness. Poi, in un Blueprint, il nodo \textbf{Set
Material} cambia il materiale di un bersaglio quando viene colpito: un primo esempio di
come la logica (le lezioni di Igor) e la grafica (questa lezione) si incontrino.

% ---------------------------------------------------------------------
\section{Material Domain, Blend Mode e Shading Model}
\label{ue3:sec:dettagli}

Come gli oggetti nella \emph{viewport}, anche il materiale ha i suoi
\textbf{Details}. Dentro l'editor si possono selezionare i singoli nodi, per i dettagli
del nodo, oppure vedere i \textbf{dettagli generali del materiale}. Tra questi ce ne
sono alcuni, evidenziati sulla slide, che ``cambiano proprio effettivamente come
funziona il materiale'' e ne definiscono la struttura e le possibilità.

\paragraph{Material Domain.} Il \textbf{Material Domain} dice \textbf{per che cosa}
verrà usato il materiale: una \textbf{superficie}, l'\textbf{interfaccia} (la UI), il
\textbf{post process}, un effetto.

\paragraph{Blend Mode.} La \textbf{Blend Mode} è il modo in cui il materiale
``\emph{blenda}'' (si mescola) con i pixel calcolati nel resto della scena. La modalità
\textbf{Opaque}, quella di default, \textbf{occlude completamente} ciò che c'è dietro.
Per un effetto speciale, per esempio un \textbf{fuoco}, serve invece la Blend Mode
\textbf{Translucent}, che permette di aggiungere parametri per l'\textbf{opacità} del
materiale.

\paragraph{Shading Model.} Lo \textbf{Shading Model}\footnote{Nella trascrizione una
volta ``scaling model''.} definisce come il materiale \textbf{reagisce alla luce}.
\textbf{Default Lit} vuol dire che segue le regole di base dell'illuminazione: il
materiale viene illuminato. Per un effetto speciale, in cui non si vuole che la luce
influenzi la resa, lo si mette invece su \textbf{Unlit}\footnote{Nella trascrizione
``lo metterò un lit''.}. Ci sono ``diverse tipologie'' di Shading Model, a seconda di
cosa si vuole realizzare.

\paragraph{I pin cambiano.} Queste impostazioni si applicano sul nodo risultante, e
cambiarle cambia i \textbf{pin} del nodo. L'\textbf{Opacity} non si può modificare in
Blend Mode Opaque, perché un materiale opaco ``dà per scontato che l'oggetto sia sempre
opaco'': non c'è bisogno di dire dove e quando sarà trasparente. Viceversa, la Blend
Mode Translucent toglie la possibilità di modificare Metallic, Specular, Roughness e
così via: se il materiale deve essere potenzialmente trasparente, ``non ha senso andare
a calcolare come rifletterà la luce''.

\begin{table}[htbp]
\centering
\small
\begin{tabular}{@{}p{2.6cm}p{3.3cm}p{7.1cm}@{}}
\toprule
\textbf{Impostazione} & \textbf{Valori citati} & \textbf{Che cosa decide} \\
\midrule
Material Domain & Surface, User Interface, Post Process\ldots & per che cosa si usa
  il materiale: una superficie, la UI, il post process, un effetto \\
Blend Mode & \textbf{Opaque} (default), Masked, Translucent & come il materiale si
  mescola con ciò che ha dietro: Opaque nasconde tutto; Masked ritaglia con una
  maschera; Translucent sfuma con l'Opacity (fuoco, fili d'erba) \\
Shading Model & \textbf{Default Lit} (default), Unlit\ldots & come reagisce alla luce:
  Default Lit viene illuminato, Unlit no (effetti speciali) \\
\bottomrule
\end{tabular}
\caption{Le tre impostazioni ``più importanti'' dei dettagli di un materiale. Cambiarle
accende e spegne i pin del nodo risultante.}
\label{ue3:tab:dettagli}
\end{table}

\paragraph{La prova in Engine.} Per far vedere la cosa, Davide apre il progetto del
materiale didattico e invita chi vuole a seguirlo: nel \textbf{Content Browser} si crea
un nuovo materiale (``va bene anche chiamarlo New Material''), lo si apre con un
\textbf{doppio clic} e si entra nell'editor. Le impostazioni di cui si parlava sono nei
dettagli, in basso a sinistra. Se la Blend Mode passa da Opaque a \textbf{Translucent},
il nodo risultante cambia: spariscono Metallic, Specular, Roughness, Anisotropy, e
compare \textbf{Opacity}. È utile, per esempio, per una \textbf{fiamma} che verso l'alto
diventa trasparente (le informazioni grafiche alla base, e in alto ``piano piano va a
dissolversi''), o per un \textbf{filo d'erba} che parte con il suo colore e poi va in
trasparenza. Con la Blend Mode \textbf{Masked}, ``mascherata'', si crea invece una
\textbf{maschera}, un \emph{layer} che funziona in modo diverso dal resto della
superficie. Per ora si resta su \textbf{Opaque}, perché i materiali di oggi andranno su
oggetti 3D. A metà spiegazione arriva uno studente in ritardo: ``Chat, buon pomeriggio,
Lorenzo. Noi qui abbiamo già cominciato, ma in realtà abbiamo parlato di cose
principalmente teoriche.''

\paragraph{Approfondimento: Masked o Translucent?} Le due modalità sembrano simili ma
lavorano in modo opposto. \textbf{Masked} decide pixel per pixel se il materiale
\emph{c'è o non c'è}: l'ingresso Opacity Mask viene confrontato con una soglia (di
default 0,333), e sotto la soglia il pixel non viene disegnato affatto. È la modalità
delle foglie, delle reti e dell'erba ``ritagliata'': economica, compatibile con la
luce e con le ombre, ma senza sfumature. \textbf{Translucent} invece \emph{mescola} il
colore del materiale con ciò che sta dietro, secondo un valore di Opacity tra 0 e 1:
permette le sfumature (fumo, vetro, fiamme), ma costa di più, perché ogni strato
traslucido va sommato a quelli dietro, e gestisce peggio luce e ordinamento. Regola
pratica: Masked quando basta un ``sì o no'', Translucent solo quando serve davvero la
sfumatura.

% ---------------------------------------------------------------------
\section{I nodi: le Material Expression}
\label{ue3:sec:nodi}

``Come per i Blueprint, anche l'editor dei materiali ha dei nodi, delle
\textbf{espressioni}.'' I \textbf{Material Expression Node} sono ``dei piccoli set di
istruzioni preprogrammate che eseguono delle cose nello specifico, esattamente come i
nodi dei Blueprint'', o, con un'altra formula, ``pezzi di codice statico''.

\paragraph{Il primo grafo.} Sulla slide c'è un piccolo grafo
(figura~\ref{ue3:fig:primografo}), e Davide lo legge nodo per nodo:
\begin{itemize}
  \item il \textbf{Texture Sample} comunica l'informazione della texture: ``sta dicendo
        che la texture che c'è in questo nodo è la texture che determina il
        \textbf{base color}''; permette di usare un'immagine e farla diventare la
        texture del materiale;
  \item il \textbf{Texture Coordinate}, il nodo \textbf{rosso}, restituisce le
        coordinate U e V del materiale, e serve a decidere ``quanto è grande il
        \textbf{tiling} della mia texture, quindi quante volte si ripete'';
  \item il \textbf{Multiply} è programmato per moltiplicare due valori: ``moltiplica A
        per B e dà una risultante'';
  \item il quadratino è una \textbf{costante}, ``un singolo \emph{float}'', una variabile
        singola.
\end{itemize}
Il grafo dice: ``moltiplica le coordinate per uno e comunica queste cose alle UV della
texture''. Così com'è, quindi, ``è come se non ci fosse'': moltiplicare per 1 lascia le
coordinate di default. Ma è così che si costruiscono i grafi: combinando nodi si
ottengono risultati molto diversi. Si può moltiplicare, si può fare
un'\textbf{interpolazione}, si possono far \textbf{muovere} le texture per suggerire un
movimento nel materiale, ``eccetera, eccetera, eccetera''.

\begin{figure}[htbp]
\centering
\begin{tikzpicture}[
  font=\small,
  ue3node/.style={draw, thick, rounded corners=4pt, align=center,
                  minimum height=0.9cm},
  ue3uv/.style={-{Stealth[length=2mm]}, thick, red!60!black},
  ue3flt/.style={-{Stealth[length=2mm]}, thick, green!50!black},
  ue3col/.style={-{Stealth[length=2mm]}, thick, blue!60!black}
]
  \node[ue3node, fill=red!12, text width=2.2cm] (tc) at (0,0.6)
    {\textbf{TexCoord}\\\footnotesize UTiling 1, VTiling 1};
  \node[ue3node, fill=green!8, minimum width=1.0cm] (k) at (0,-1.0) {\textbf{1}};
  \node[font=\footnotesize] at (0,-1.75) {costante};
  \node[ue3node, fill=gray!10, text width=1.6cm] (mul) at (3.0,0)
    {\textbf{Multiply}\\\footnotesize A $\times$ B};
  \node[ue3node, fill=gray!10, text width=2.4cm, minimum height=1.8cm] (ts)
    at (6.4,0) {\textbf{Texture}\\\textbf{Sample}\\\footnotesize (T\_Bark)};
  \node[ue3node, fill=blue!6, text width=2.4cm, minimum height=2.6cm] (res)
    at (10.4,0) {\textbf{M\_NewMaterial}\\[3pt]\footnotesize\textbf{Base Color}\\
    \footnotesize Metallic\\\footnotesize Specular\\\footnotesize Roughness\\
    \footnotesize Emissive Color\\\footnotesize Normal\ldots};
  \draw[ue3uv] (tc.east) -- ++(0.4,0) |- node[pos=0.75, above, font=\scriptsize]{A}
    ([yshift=2mm]mul.west);
  \draw[ue3flt] (k.east) -- ++(0.6,0) |- node[pos=0.75, below, font=\scriptsize]{B}
    ([yshift=-2mm]mul.west);
  \draw[ue3uv] (mul.east) -- node[above, font=\scriptsize]{UVs}
    (ts.west |- mul.east);
  \draw[ue3col] (ts.east) -- node[above, font=\scriptsize]{RGB}
    (res.west |- ts.east);
\end{tikzpicture}
\caption{Il primo grafo della slide: le coordinate UV, moltiplicate per una costante,
entrano nel Texture Sample, il cui colore va nel Base Color del nodo risultante. Con la
costante a 1 il Multiply ``è come se non ci fosse''; cambiando la costante cambia il
tiling (sezione~\ref{ue3:sec:tiling}). Schema ricostruito.}
\label{ue3:fig:primografo}
\end{figure}

\paragraph{I parametri.} Non solo le costanti, ma i nodi in generale ``possono essere
utilizzati come \textbf{parametri}''. Per rendere un nodo un parametro si fa
\textbf{tasto destro} sul nodo e \textbf{Convert to Parameter}, gli si dà un nome, e
questo permette di modificare in seguito il valore ``senza dover ricalcolare da capo il
materiale''. I parametri torneranno utili con le \textbf{Material Instance}
(sezione~\ref{ue3:sec:instance}).

\paragraph{Come si chiama un nodo.} ``Tasto destro, scriviamo il nome'', per esempio
\emph{texture sample}, ed ecco il nodo. C'è poi ``tutta una serie di \emph{shortcut}''
per i nodi più usati, che si vedranno tra poco.

% ---------------------------------------------------------------------
\section{Il problema della mappa vuota}
\label{ue3:sec:mappavuota}

Una studentessa interrompe: non riesce più ad aprire la mappa, ``me la apre totalmente
vuota, senza titolo''. Davide le chiede di condividere lo schermo, ma lei segue da due
computer diversi. Prima ipotesi: ha creato un \textbf{nuovo livello}, e un \emph{New
Level} appena caricato è tutto nero, perché non c'è niente in scena (``potremmo andare
a popolarla mettendoci cose''). Ma il problema è un altro: ha aperto il progetto del
materiale didattico, cliccato sulla mappa del personaggio in terza persona
(\emph{Third Person Map}), e non succede niente. Davide sospetta che il progetto sia
ancora in una \textbf{cartella zippata} (``di solito questa cosa succede quando il
progetto è una cartella che è ancora zippata''), e le consiglia di
\textbf{riscaricarlo}. Poco dopo il problema è risolto: ``sono riuscita a risolvere
adesso''. ``Meno male, meno male.''

\paragraph{Approfondimento: aprire un progetto da una cartella.} Un progetto di Unreal
è una cartella che contiene un file \texttt{.uproject} e le sottocartelle
\texttt{Content}, \texttt{Config} e altre. Unreal non può aprire un progetto che sta
ancora dentro un archivio \texttt{.zip}: Windows mostra il contenuto dello zip come se
fosse una cartella, ma l'editor non ci può scrivere, e il risultato tipico è proprio
una mappa vuota o mancante. La soluzione è estrarre tutto (``Estrai tutto'') e aprire il
\texttt{.uproject} estratto con un doppio clic. Se il progetto non compare nell'elenco
dell'Epic Games Launcher o nel Project Browser, basta aprirlo dal \texttt{.uproject}
la prima volta: da lì in poi compare tra i progetti recenti.

% ---------------------------------------------------------------------
\section{Dentro una texture e il Texture Sample}
\label{ue3:sec:sample}

``Facciamo le cose in ordine, se no io mi confondo e mi impiccio.'' Davide torna al
punto saltato: com'è fatta \textbf{dentro} una texture.

\paragraph{La finestra della texture.} Con un doppio clic su una texture nel Content
Browser (Davide ne prende una a caso, una corteccia, \textbf{T\_Bark}) si apre una
finestra in cui ``non possiamo fare grandi modifiche'', ma si possono vedere i
\textbf{canali}: il blu, il verde, il rosso, ciascuno fatto a modo suo. ``Questa texture
non è particolarmente interessante'', ma si rivedrà meglio con la torretta.

\paragraph{Il Texture Sample.} Nel grafo del materiale si mette un \textbf{Texture
Sample}, si seleziona il nodo, e nei dettagli, in basso, si sceglie quale texture
metterci. Cercando T\_Bark e mettendola nel nodo, la \emph{preview} mostra la corteccia:
``il nostro materiale sarà così, perché è effettivamente la texture che ci abbiamo
messo''.

\paragraph{Le uscite separate.} Il Texture Sample non ha una sola uscita: oltre
all'\textbf{RGB} ci sono \textbf{R}, \textbf{G} e \textbf{B} separati (e l'alfa, A). Si
può quindi usare, per esempio, solo il canale \textbf{rosso} (``non quella completa, ma
questa qui'') per trasportare un'informazione: collegandolo allo \textbf{Specular} si
ottiene un materiale che in alcune parti riflette la luce più che in altre. ``Con
questo materiale in realtà è un po' complicato farlo vedere'', ma il principio è quello
dei canali visto sulle slide.

\paragraph{Come si mette una texture nel nodo.} Una studentessa non ha capito come si
aggiunge la texture ``al quadratino nero'' del Texture Sample. Davide rifà il
percorso: si clicca il nodo, e nei suoi dettagli c'è il campo \textbf{Texture}, con il
quadratino e una lista. Due modi:
\begin{itemize}
  \item se una texture è già \textbf{selezionata nel Content Browser}, la freccetta
        \textbf{Use Selected Asset from Content Browser} la inserisce automaticamente
        nel nodo;
  \item più spesso, si clicca sul \textbf{nome} della texture nel campo e si fa un
        \emph{browse}: si cerca il nome (``T underscore Bark'') e lo si sceglie dalla
        lista di tutte le texture disponibili.
\end{itemize}
``Voi lo vedete questo pop-up che esce?'' Sì: perfetto.

% ---------------------------------------------------------------------
\section{Le costanti, le operazioni e le scorciatoie}
\label{ue3:sec:costanti}

``Chiaramente ne vediamo un po', non li vediamo tutti, perché sono tanti, sono
veramente tanti.'' Ci sono ``tantissime possibilità, tante cose che vanno studiate'',
ma per ora servono le basi.

\paragraph{La regola delle scorciatoie.} Il Texture Sample si può chiamare cercandolo
con il tasto destro, oppure \textbf{tenendo premuto il tasto T e facendo clic
sinistro} nel grafo. ``Questa struttura dello shortcut funziona per ogni singolo nodo
che ha uno shortcut'': un tasto, tenuto premuto, più clic sinistro. Il tasto è
``molto spesso l'iniziale del nome del nodo''.

\paragraph{Le costanti.} Cercando \emph{vector} compare, per esempio, il Vector
Noise; cercando \emph{constant} compaiono le costanti: \textbf{Constant},
\textbf{Constant2Vector}, \textbf{Constant3Vector} (e Constant4Vector). Le scorciatoie
sono i numeri: \textbf{1 + clic} per la costante singola, \textbf{2 + clic} per quella
a due valori, \textbf{3 + clic} per quella a tre. La costante a due ``la utilizzerete
relativamente poco''. Le altre due moltissimo:
\begin{itemize}
  \item la \textbf{costante singola} è un \textbf{valore numerico}, che si usa ``per
        andare a fare qualsiasi tipo di calcolo, qualsiasi tipo di costruzione'';
  \item la \textbf{costante a tre} è un \textbf{colore}: l'engine usa i tre valori come
        \textbf{R}, \textbf{G} e \textbf{B}. Si sceglie un colore (Davide prende un
        arancione) e l'engine modifica X, Y e Z per ottenerlo. ``Qualcuno di voi l'aveva
        anche già fatto nel test con la pallina.''
\end{itemize}
Collegando l'arancione al \textbf{Base Color} si ottiene un materiale arancione,
``perché effettivamente è l'unica cosa che gli sto passando''. Ma ``molto raramente''
si usa un colore piatto come Base Color: lo si usa piuttosto per fare
\textbf{operazioni}. Più avanti una studentessa chiederà di nuovo come si crea il colore:
``è come la costante singola, che è uno clic sinistro, ma è a tre, quindi è tre clic
sinistro''.

\paragraph{Staccare i fili.} Per rompere un collegamento: \textbf{Alt + clic sinistro}
sul pin, ``e si rompono''.

\paragraph{Le operazioni.} Le operazioni utili sono ``le classiche'': addizione,
moltiplicazione, divisione, sottrazione. Le scorciatoie: \textbf{A + clic} per
\textbf{Add}, \textbf{M + clic} per \textbf{Multiply}, \textbf{D + clic} per
\textbf{Divide}. La sottrazione (\textbf{Subtract}) non ha scorciatoia: si cerca. Il
motivo lo spiega Davide più avanti: \textbf{S + clic} crea automaticamente un
\textbf{parametro} (uno Scalar Parameter), ``quindi questo potrebbe tornarvi utile''.

\paragraph{Quali si usano davvero.} In chiusura di lezione Davide aggiunge un consiglio
d'esperienza: ``l'operazione che utilizzerete di più in assoluto è la
\textbf{moltiplicazione}''. Moltiplicazione e addizione sono le più usate; divisione e
sottrazione servono soprattutto per cose più complesse, ``tipo sugli effetti
speciali''.

% ---------------------------------------------------------------------
\section{Il colore che brilla: Emissive, Multiply e Glow Intensity}
\label{ue3:sec:emissive}

``A cosa ci servono'' le operazioni? Per esempio ad \textbf{aumentare il valore} di
qualcosa. Il nodo risultante ha diversi pin, ``diverse cose su cui andare a collegare
dei valori'': il \textbf{Base Color}, e più in basso l'\textbf{Emissive Color}, che
controlla ``quali parti del materiale emetteranno una piccola luce'', un
\textbf{\emph{glow}}.

\paragraph{L'arancione nell'Emissive.} Davide stacca l'arancione dal Base Color e lo
collega all'Emissive: il materiale ``sarà leggermente luminoso''. ``Adesso è molto poco
luminoso, ma è comunque un emissive.''

\paragraph{Moltiplicare.} Come si ottiene un'emissione maggiore? ``Moltiplichiamo.''
L'arancione entra nella A di un \textbf{Multiply}, il Multiply va nell'Emissive. Il
secondo valore si può scrivere direttamente nel nodo Multiply (``invece di 1 mettiamo
100''): ``verrà decisamente più luminoso''. Ma così ``è meno controllabile''. Meglio
attaccare alla B una \textbf{costante} separata, darle il valore 100 e convertirla in
\textbf{parametro}, con un nome come \textbf{Glow Intensity}: quando più avanti ci sarà
una Material Instance, il valore si potrà cambiare ``direttamente da lì''. ``Ci andiamo
a estrapolare un parametro.''

\begin{figure}[htbp]
\centering
\begin{tikzpicture}[
  font=\small,
  ue3node/.style={draw, thick, rounded corners=4pt, align=center,
                  minimum height=0.9cm},
  ue3col/.style={-{Stealth[length=2mm]}, thick, blue!60!black},
  ue3flt/.style={-{Stealth[length=2mm]}, thick, green!50!black}
]
  \node[ue3node, fill=orange!30, text width=2.5cm] (col) at (0,0.7)
    {\textbf{Constant3Vector}\\\footnotesize arancione (3 + clic)};
  \node[ue3node, fill=green!8, text width=2.5cm] (gi) at (0,-0.9)
    {\textbf{Glow Intensity}\\\footnotesize parametro = 100};
  \node[ue3node, fill=gray!10, text width=1.6cm] (mul) at (3.6,0)
    {\textbf{Multiply}\\\footnotesize (M + clic)};
  \node[ue3node, fill=blue!6, text width=2.6cm, minimum height=2.2cm] (res)
    at (7.4,0) {\textbf{nodo risultante}\\[3pt]\footnotesize Base Color\\
    \footnotesize Metallic, Specular\\\footnotesize Roughness\\
    \footnotesize\textbf{Emissive Color}};
  \draw[ue3col] (col.east) -- ++(0.4,0) |- node[pos=0.75, above, font=\scriptsize]{A}
    ([yshift=2mm]mul.west);
  \draw[ue3flt] (gi.east) -- ++(0.4,0) |- node[pos=0.75, below, font=\scriptsize]{B}
    ([yshift=-2mm]mul.west);
  \draw[ue3col] (mul.east) -- ++(0.6,0) |- ([yshift=-7mm]res.west);
\end{tikzpicture}
\caption{Un materiale che brilla. Il colore, collegato all'Emissive Color, emette ``una
piccola luce''; moltiplicato per 100 diventa molto più luminoso. Il valore 100 sta in
una costante convertita in parametro (Glow Intensity), così che si possa cambiare da
una Material Instance senza toccare il grafo. Ricostruzione dalla dimostrazione.}
\label{ue3:fig:emissive}
\end{figure}

\paragraph{Approfondimento: perché 100 funziona.} I colori di un materiale non sono
limitati tra 0 e 1. Nel Base Color valori oltre 1 non hanno senso fisico (nessuna
superficie riflette più luce di quella che riceve), ma nell'Emissive sì: un'emissione di
$(1; 0{,}5; 0)$ moltiplicata per 100 diventa $(100; 50; 0)$, una sorgente molto intensa
che il \emph{post process} di Unreal trasforma in un alone luminoso (il \emph{bloom}).
Il Multiply tra un colore e un numero moltiplica ciascun canale per quel numero, quindi
la tinta resta la stessa e cresce solo l'intensità. Con Lumen, il sistema di
illuminazione globale di Unreal Engine 5, un materiale emissivo può anche illuminare un
po' le superfici vicine.

% ---------------------------------------------------------------------
\section{Il PBR: Metallic, Specular, Roughness e Normal}
\label{ue3:sec:pbr}

``Mi rendo conto che sto andando un po' a destra e a sinistra'', ammette Davide, e
torna ai pin del nodo risultante. Base Color ed Emissive si sono visti. Restano
\textbf{Metallic}, \textbf{Specular} e \textbf{Roughness}, ``degli input abbastanza
interessanti'', perché sono quelli che \textbf{determinano la superficie} del
materiale.

\paragraph{Il PBR.} Il rendering dei materiali in Unreal è \textbf{PBR},
\emph{Physically Based Rendering}: per decidere come si comporta una superficie, il
motore mette insieme una serie di informazioni fisiche. Tutti questi parametri vanno
\textbf{da 0 a 1}.

\paragraph{Roughness.} La \textbf{Roughness} dice quanto una superficie è
\textbf{liscia} o \textbf{ruvida}. Una \textbf{roccia} vuole una Roughness molto alta,
cioè 1: il materiale sarà ruvido ``e non rifletterà la luce in nessun modo''. Uno
\textbf{specchio} vuole una Roughness a 0: il più liscio possibile, il meno ruvido, il
più riflettente.

\paragraph{Specular e Metallic.} ``Per aiutare in questa situazione'' ci sono Specular e
Metallic, che ``si alternano a seconda del tipo di materiale''. Lo
\textbf{Specular}, come si è visto con le mappe, determina quanta luce viene riflessa e
l'angolazione a cui viene riflessa. Per creare uno \textbf{specchio} serve un materiale
molto, molto liscio e molto, molto riflettente: Davide abbassa la Roughness e alza la
riflessione, e la \emph{preview} comincia ad assomigliare ``a una parete riflettente'';
con altre cose che si vedranno in seguito ``iniziamo ad avere un riflesso''.

\paragraph{Emissive e Normal.} L'\textbf{Emissive Color} restituisce un colore che è
``effettivamente una luce''. La \textbf{Normal} è l'ultimo parametro utile, e richiede
la texture blu vista prima: si usa sempre un Texture Sample, ma con
\textbf{T\_Bark\_N}, la Normal della corteccia, invece della Diffuse. Davide toglie
tutto il resto e mette a confronto il materiale con e senza Normal: ``vedete la
differenza?'' La Normal aggiunge un dettaglio tridimensionale ``molto semplice da
vedere'', eppure ``la mappa è 2D, non c'è tridimensionalità'': crea solo l'illusione.

\begin{figure}[htbp]
\centering
\begin{tikzpicture}[font=\small]
  \foreach \i/\rr in {0/0, 1/0.25, 2/0.5, 3/0.75, 4/1} {
    \pgfmathsetmacro{\hs}{0.12+0.35*\rr}
    \pgfmathsetmacro{\op}{1-0.85*\rr}
    \shade[ball color=gray!70] (\i*2.4,0) circle (0.8);
    \fill[white, opacity=\op] (\i*2.4-0.28,0.3) circle (\hs);
    \node[font=\footnotesize] at (\i*2.4,-1.2) {Roughness $=\rr$};
  }
  \draw[-{Stealth[length=2.5mm]}, thick] (-0.6,-1.8) -- (10.2,-1.8);
  \node[font=\footnotesize, anchor=north west] at (-0.6,-1.85) {liscio, riflessi netti
    (lo specchio)};
  \node[font=\footnotesize, anchor=north east] at (10.2,-1.85) {ruvido, nessun
    riflesso netto (la roccia)};
\end{tikzpicture}
\caption{La Roughness, da 0 a 1, in uno schema qualitativo: più la superficie è liscia,
più il riflesso della luce è piccolo e netto; più è ruvida, più la luce si sparpaglia
fino a sparire come riflesso. Disegno indicativo, non un rendering di Unreal.}
\label{ue3:fig:roughness}
\end{figure}

\paragraph{Approfondimento: Metallic e Specular nella pratica.} Nel modello PBR di
Unreal i due parametri non hanno lo stesso peso. \textbf{Metallic} è quasi sempre
\emph{binario}: 1 per i metalli, 0 per tutto il resto (legno, pietra, plastica, pelle);
i valori intermedi servono soprattutto per i bordi tra metallo e vernice o per lo
sporco. Lo \textbf{Specular} ha un default di 0,5, che va bene per la grande maggioranza
dei materiali non metallici, e la documentazione di Unreal consiglia di toccarlo poco:
per i metalli viene ignorato, perché il loro riflesso dipende dal Base Color. Lo
specchio della lezione, in termini PBR, è quindi Metallic $=1$, Roughness $=0$, Base
Color chiaro. È il senso dell'``alternarsi'' di cui parla Davide: nei materiali metallici
si lavora con il Metallic, negli altri, se proprio serve, con lo Specular, e quasi
sempre con la Roughness, che è il parametro che cambia di più l'aspetto di una
superficie.

% ---------------------------------------------------------------------
\section{Il Texture Coordinate e il tiling}
\label{ue3:sec:tiling}

``Chiaramente ogni nodo funziona a sé.'' Il \textbf{Texture Coordinate}, il nodo rosso
della prima slide, ha i suoi parametri. La scorciatoia è ``un po' \emph{tricky} da
ricordare'', perché T + clic è già preso dal Texture Sample: è \textbf{U + clic},
``perché il Texture Coordinate lavora sulle UV''.

\paragraph{UTiling e VTiling.} Selezionando il nodo, nei dettagli compaiono
\textbf{UTiling} e \textbf{VTiling}: i parametri ``molto comodi'' con cui si aumenta o
si diminuisce il tiling di un oggetto. Davide collega il Texture Coordinate alle UV
della Normal e della Diffuse, e chiede: se metto UTiling a 2, che cosa succede? ``Vediamo
due volte la texture applicata all'oggetto'', risponde uno studente. ``Sì, però c'è un
\emph{catch}, perché è soltanto sulle U'': la texture si ripete due volte
sull'\textbf{asse orizzontale} e una sola volta su quello verticale. Con anche VTiling a
2, la texture si ripete \textbf{quattro} volte: ``ce n'è una qua, una qua, una qui e una
qua''.

\paragraph{Lo zoom.} Si può fare anche il contrario, per ``zoomare'' sulla texture: con
0,5 la si vede ``la metà delle volte. Anzi, non la metà delle volte, scusate: un
quarto'' (0,5 su entrambi gli assi). E questo non vuol dire che una parte della
texture ``viene mangiata'': viene semplicemente ``applicata in uno spazio differente''.

\begin{figure}[htbp]
\centering
\begin{tikzpicture}[font=\small]
  \newcommand{\uetreTile}[5]{% x0, y0, nu, nv, etichetta
    \begin{scope}[shift={(#1,#2)}]
      \pgfmathsetmacro{\du}{2.2/#3}
      \pgfmathsetmacro{\dv}{2.2/#4}
      \foreach \i in {1,...,#3} \foreach \j in {1,...,#4} {
        \draw[fill=brown!25] ({(\i-1)*\du},{(\j-1)*\dv}) rectangle
          ({\i*\du},{\j*\dv});
        \draw[brown!70!black, thick] ({(\i-1)*\du+0.15*\du},{(\j-1)*\dv+0.2*\dv})
          -- ({(\i-1)*\du+0.85*\du},{(\j-1)*\dv+0.8*\dv});
      }
      \draw[very thick] (0,0) rectangle (2.2,2.2);
      \node[font=\footnotesize, align=center] at (1.1,-0.5) {#5};
    \end{scope}}
  \uetreTile{0}{0}{1}{1}{U = 1, V = 1}
  \uetreTile{3.0}{0}{2}{1}{U = 2, V = 1}
  \uetreTile{6.0}{0}{2}{2}{U = 2, V = 2}
  % zoom: un quarto della texture, ingrandito
  \begin{scope}[shift={(9.0,0)}]
    \draw[fill=brown!25] (0,0) rectangle (2.2,2.2);
    \draw[brown!70!black, thick] (0.66,0.88) -- (2.2,2.2);
    \draw[very thick] (0,0) rectangle (2.2,2.2);
    \node[font=\footnotesize, align=center] at (1.1,-0.5) {U = 0,5, V = 0,5};
  \end{scope}
\end{tikzpicture}
\caption{Il tiling sulla stessa faccia. Con UTiling a 2 la texture si ripete solo in
orizzontale; con entrambi a 2 si ripete quattro volte; con entrambi a 0,5 si vede un
quarto della texture, ingrandito: non è ``mangiata'', è applicata in uno spazio più
grande. Schema ricostruito dalla dimostrazione.}
\label{ue3:fig:tiling}
\end{figure}

\paragraph{Il pavimento.} ``Il tiling, però, va visto.'' Nella mappa c'è un
\textbf{pavimento}, e Davide ne cambia il materiale: lo seleziona, e nei
\textbf{Details}, sotto \textbf{Materials}, al posto della \emph{material instance}
della griglia grigia di prototipazione\footnote{Nella trascrizione ``Materialistance
Prototype Grid Gray'': è l'istanza della griglia grigia dei progetti modello di
Unreal.} trascina il \textbf{master material} appena creato. Il risultato: la corteccia
è ``molto, molto grande'' (``io non so neanche qual è l'utilizzo di questa texture che
sto usando, ci ho preso Bark a caso, ma va bene così''). Davide porta il tiling a 3,
poi a 10: ``iniziamo a vederla più spesso, questa corteccia. Sembrano anche delle dune,
vagamente.'' Il pavimento è probabilmente un cubo allargato, con UV ``molto, molto
piccole'' rispetto alle sue dimensioni: per questo nella \emph{preview} sferica del
materiale la texture si ripete ``diecimila volte'', mentre sul pavimento ``abbiamo un
pavimento che è più dimensionato''.

\paragraph{Una domanda che apre il tema successivo.} Uno studente chiede: si possono
modificare U e V in base all'oggetto su cui è applicata la texture, per esempio 10 su un
pavimento e 0,5 su un altro oggetto? ``Sì, è possibile, non facendolo dallo stesso
materiale. Questo è il momento in cui entrano in gioco le \textbf{Material
Instance}.''

\paragraph{Approfondimento: il tiling in una formula.} Il Texture Coordinate restituisce
per ogni punto la sua coppia $(u, v)$; il Multiply la scala: $(u', v') = (k_u\,u,\;
k_v\,v)$. Una texture è definita sull'intervallo $[0,1]$ e si ripete fuori da lì (è il
comportamento di default, \emph{wrap}): quando $u'$ va da 0 a 2, la texture viene letta
due volte. Con $k < 1$ la faccia ``vede'' solo un pezzo dell'intervallo, e quel pezzo
viene ingrandito. Per questo il tiling dipende dalle UV della mesh: se le UV di un
pavimento enorme occupano comunque tutto il quadrato $[0,1]^2$, una texture con tiling
1 si stende una volta sola su tutta la superficie, ed è ``molto, molto grande''.
Moltiplicare per una costante e usare UTiling e VTiling nel nodo sono due modi di fare
la stessa cosa; il primo ha il vantaggio che la costante può diventare un parametro.

% ---------------------------------------------------------------------
\section{Le Material Instance}
\label{ue3:sec:instance}

Davide torna alle slide (``passiamo un attimo alle Material Instance, visto che ne
stavamo parlando''). Quello costruito finora è un \textbf{master material}: tutto ciò
che sta nel grafo. Se si modifica il grafo, si modifica il materiale. Per avere lo
\textbf{stesso materiale con una piccola differenza} si possono creare due
\textbf{istanze} e cambiare in ciascuna il parametro collegato al Texture Coordinate.

\paragraph{Che cos'è un'istanza.} Una \textbf{Material Instance} è ``un'istanza di un
materiale che noi possiamo utilizzare per modificare quel materiale \textbf{senza
modificare la sorgente}''. Ha diversi vantaggi, ``oltre a essere \emph{bon ton}'', cioè
il modo in cui si dovrebbero applicare i materiali in generale: ``non dovreste prendere
il master material e applicarlo agli oggetti in scena. È sempre consigliabile applicare
un'istanza.''

\paragraph{Perché è più leggera: i tavoli.} L'istanza è ``molto più leggera''. Davide
lo spiega con i \textbf{tavolini} della prima lezione. Tre tavoli in scena con lo
stesso materiale del legno: per ogni tavolo l'engine carica la mesh e carica il
materiale, quindi ``tavolo, tavolo, tavolo'' e ``materiale, materiale, materiale''. Se
invece sui tavoli ci sono tre Material Instance, l'engine carica le tre mesh, ma trova
tre istanze che puntano \textbf{tutte allo stesso master material}, e ``va a
renderizzare un solo materiale''. ``Chiaramente, su tre tavoli non cambia niente. Ora
immaginate che invece di tre tavoli ci siano diecimila oggetti in scena'' e che mille
condividano lo stesso materiale: si capisce perché ``torna comodo avere un'istanza
piuttosto che avere un materiale generale''.

\begin{figure}[htbp]
\centering
\begin{tikzpicture}[
  font=\small,
  ue3tab/.style={draw, thick, minimum width=1.5cm, minimum height=0.7cm,
                 fill=brown!20, align=center},
  ue3mat/.style={draw, thick, rounded corners=3pt, minimum width=1.5cm,
                 minimum height=0.7cm, align=center},
  ue3arr/.style={-{Stealth[length=2mm]}, thick}
]
  % senza istanze
  \node[font=\footnotesize\bfseries] at (2.0,2.2) {come la racconta Davide: master
    material su ogni tavolo};
  \foreach \i in {0,1,2} {
    \node[ue3tab] (t\i) at (\i*2.0,1.2) {tavolo};
    \node[ue3mat, fill=blue!15] (m\i) at (\i*2.0,0) {M\_Legno};
    \draw[ue3arr] (t\i) -- (m\i);
  }
  \node[font=\footnotesize, align=center] at (2.0,-0.8) {tre mesh e ``materiale,
    materiale, materiale''};
  % con istanze
  \begin{scope}[shift={(7.6,0)}]
    \node[font=\footnotesize\bfseries] at (2.0,2.2) {con le Material Instance};
    \foreach \i in {0,1,2} {
      \node[ue3tab] (u\i) at (\i*2.0,1.2) {tavolo};
      \node[ue3mat, fill=orange!15, font=\footnotesize] (i\i) at (\i*2.0,0)
        {MI\_Legno\_\i};
      \draw[ue3arr] (u\i) -- (i\i);
    }
    \node[ue3mat, fill=blue!15] (mm) at (2.0,-1.3) {\textbf{M\_Legno}};
    \foreach \i in {0,1,2} \draw[ue3arr] (i\i) -- (mm);
    \node[font=\footnotesize, align=center] at (2.0,-2.1) {tre mesh, tre istanze, un
      solo materiale da compilare};
  \end{scope}
\end{tikzpicture}
\caption{L'esempio dei tavoli. A sinistra, il master material applicato direttamente;
a destra, tre istanze (che possono avere parametri diversi: un legno più chiaro, un
tiling più fitto) che rimandano tutte allo stesso master material. Schema ricostruito
dalla spiegazione; sul costo reale si veda l'approfondimento.}
\label{ue3:fig:tavoli}
\end{figure}

\paragraph{Come si crea.} ``Come lo facciamo su Unreal?'' Si fa \textbf{tasto destro}
sul materiale nel Content Browser e \textbf{Create Material Instance}, ``qui in
alto''. L'istanza appare \textbf{vuota}: ``non c'è un grafo per modificarlo, perché
chiaramente non posso modificare nessun grafo quando sono qui dentro''. Per avere
qualcosa da cambiare, bisogna estrarre dei \textbf{parametri} dal master material.
Davide torna nel grafo, fa tasto destro sul Texture Sample del colore, \textbf{Convert
to Parameter}, e lo chiama più o meno \emph{Diffuse Color}\footnote{Nella trascrizione
``discolor'': il nome esatto non si ricostruisce, ma è il parametro della texture del
colore.}; lo stesso per l'altro Texture Sample, chiamato \textbf{Normal Map}. Salvato il
materiale, nell'istanza compaiono \textbf{due parametri} selezionabili.

\paragraph{La nomenclatura.} ``Normalmente dovreste chiamarlo'' non \emph{New Material}
ma \textbf{M\_NewMaterial}, e l'istanza \textbf{MI\_NewMaterial}: ``per nomenclatura,
insomma, per ritrovare facilmente tutto quello che ci serve''.

\paragraph{Cambiare la texture dall'istanza.} Davide applica MI\_NewMaterial al
pavimento: è identico a prima. Ma se nell'istanza si cambia il parametro della texture,
da Bark a un'altra, il materiale cambia, senza toccare il grafo. Si potrebbe avere un
master material senza texture, ``delle texture assenti o base'', e tre o quattro
istanze in cui si cambiano le texture: un muro arancione, uno blu, uno rosso e uno verde
sono quattro istanze. ``Quando io renderizzo la scena, non andrò a chiamare cinque
materiali diversi, chiamerò sempre M\_NewMaterial, perché sono tutte istanze dello
stesso materiale.'' ``Non vorrei stare complicando troppo la cosa, però insomma, ditemi
voi se è chiaro il discorso.''

\paragraph{Il Tiling Multiplier.} Ed ecco la risposta alla domanda sul tiling diverso
per ogni oggetto. Nel master material Davide prende un \textbf{Multiply}, ci collega il
Texture Coordinate, collega il Multiply alle UV dei due Texture Sample e attacca alla B
una \textbf{costante}, che converte in parametro con il nome \textbf{Tiling
Multiplier}. Nell'istanza, il Tiling Multiplier portato a 10 dà ``lo stesso risultato di
prima''; ma ora è un parametro. Una copia dell'istanza con il Tiling Multiplier a 5
darebbe un tiling diverso con lo stesso materiale. ``È chiaro il discorso?''
``Perfetto.''

\begin{figure}[htbp]
\centering
\begin{tikzpicture}[
  font=\small,
  ue3node/.style={draw, thick, rounded corners=4pt, align=center,
                  minimum height=0.9cm},
  ue3uv/.style={-{Stealth[length=2mm]}, thick, red!60!black},
  ue3flt/.style={-{Stealth[length=2mm]}, thick, green!50!black},
  ue3col/.style={-{Stealth[length=2mm]}, thick, blue!60!black}
]
  \node[ue3node, fill=red!12, text width=1.9cm] (tc) at (0,0.7)
    {\textbf{TexCoord}\\\footnotesize (U + clic)};
  \node[ue3node, fill=green!8, text width=1.9cm] (tm) at (0,-0.9)
    {\textbf{Tiling}\\\textbf{Multiplier}\\\footnotesize parametro};
  \node[ue3node, fill=gray!10, text width=1.5cm] (mul) at (3.0,0)
    {\textbf{Multiply}};
  \node[ue3node, fill=orange!12, text width=2.4cm] (ts1) at (6.2,1.0)
    {\textbf{Diffuse Color}\\\footnotesize Texture Sample,\\
    \footnotesize parametro};
  \node[ue3node, fill=orange!12, text width=2.4cm] (ts2) at (6.2,-1.0)
    {\textbf{Normal Map}\\\footnotesize Texture Sample,\\\footnotesize parametro};
  \node[ue3node, fill=blue!6, text width=2.2cm, minimum height=2.6cm] (res)
    at (10.0,0) {\textbf{M\_NewMaterial}\\[4pt]\footnotesize Base Color\\[14pt]
    \footnotesize Normal};
  \draw[ue3uv] (tc.east) -- ++(0.4,0) |- ([yshift=2mm]mul.west);
  \draw[ue3flt] (tm.east) -- ++(0.4,0) |- ([yshift=-2mm]mul.west);
  \draw[ue3uv] (mul.east) -- ++(0.5,0) |- node[pos=0.8, above, font=\scriptsize]{UVs}
    (ts1.west);
  \draw[ue3uv] (mul.east) -- ++(0.5,0) |- node[pos=0.8, below, font=\scriptsize]{UVs}
    (ts2.west);
  \draw[ue3col] (ts1.east) -- ++(0.6,0) |- ([yshift=6mm]res.west);
  \draw[ue3col] (ts2.east) -- ++(0.6,0) |- ([yshift=-6mm]res.west);
  % istanze
  \node[ue3node, fill=yellow!15, text width=3.6cm, font=\footnotesize] (mi1)
    at (3.0,-2.9) {\textbf{MI\_NewMaterial}\\Tiling Multiplier = 10\\Diffuse = T\_Bark};
  \node[ue3node, fill=yellow!15, text width=3.6cm, font=\footnotesize] (mi2)
    at (8.0,-2.9) {\textbf{MI\_NewMaterial\_2}\\Tiling Multiplier = 5\\
    Diffuse = un'altra texture};
  \draw[dashed, thick] (mi1.north) -- ++(0,0.4) -| (res.south);
  \draw[dashed, thick] (mi2.north) -- (mi2.north |- res.south);
\end{tikzpicture}
\caption{Il master material con tre parametri (le due texture e il Tiling Multiplier)
e due istanze che li impostano in modo diverso. Le istanze non hanno un grafo: hanno
solo le caselle dei parametri. Ricostruzione dalla dimostrazione.}
\label{ue3:fig:instance}
\end{figure}

\paragraph{Il vantaggio principale.} Ricapitolando dalle slide: le Material Instance
permettono di \textbf{personalizzare velocemente} le superfici ``senza andare a
modificare il grafico che avete fatto'', e le modifiche si vedono
\textbf{immediatamente}, ``senza dover aspettare la ricompilazione del materiale''.

\paragraph{Approfondimento: che cosa si risparmia davvero.} La spiegazione dei tavoli
va presa come un'immagine. Il costo principale di un materiale non è ``caricarlo'' una
volta per oggetto, ma \textbf{compilarlo}: ogni master material diventa uno o più
\emph{shader}, programmi per la scheda video che Unreal deve compilare (e che fanno
comparire, nella prima apertura di un progetto, il famoso contatore ``Compiling
Shaders''). Un'istanza \textbf{non} genera un nuovo shader: riusa quello del genitore e
cambia solo i valori dei parametri. Cento istanze dello stesso master material sono
quindi un solo shader con cento ``schede'' di valori; cento materiali diversi sono cento
shader. Per questo cambiare un parametro in un'istanza è istantaneo, mentre cambiare un
nodo nel master material richiede la ricompilazione, come dice la slide. Le istanze del
Content Browser si chiamano \textbf{Material Instance Constant}; esistono anche le
\textbf{Material Instance Dynamic}, create durante il gioco da un Blueprint (nodo
\emph{Create Dynamic Material Instance}) per cambiare i parametri a \emph{runtime}, per
esempio con \emph{Set Scalar Parameter Value}: un materiale che si accende quando la
torretta di Igor ha un bersaglio nel mirino.

% ---------------------------------------------------------------------
\section{Le Material Function}
\label{ue3:sec:function}

``Per chiudere il discorso materiali'', l'ultima voce delle slide: le \textbf{Material
Function}, che oggi non si useranno. Sono funzioni che raccolgono ``una serie di calcoli
un po' più complessi'': invece di rifare ogni volta dei calcoli che si fanno spesso, li
si mette dentro una funzione e si usa la funzione ``a vostra discrezione''. ``È come se
fossero delle \textbf{parentesi}: voi prendete queste parentesi, le spostate in giro. È
un pacchetto che potete portarvi in giro.'' Servono a velocizzare lo sviluppo quando ci
sono cose da ripetere, e possono essere usate e condivise ``fra membri del team, fra
materiali diversi''.

\paragraph{Approfondimento: Material Function e Material Instance.} Le due cose
risolvono problemi diversi. L'istanza riusa \emph{un intero materiale} cambiandone i
valori; la funzione riusa \emph{un pezzo di grafo} dentro materiali diversi. Unreal
stesso è pieno di Material Function già pronte (si trovano nel menu dei nodi con
l'icona della funzione), e i nodi World Aligned Texture e World Aligned Normal che
compaiono alla fine della lezione (sezione~\ref{ue3:sec:world}) sono, appunto, Material
Function del motore.

% ---------------------------------------------------------------------
\section{Il Lerp: interpolare}
\label{ue3:sec:lerp}

Restano ``un paio di nodi che possono essere carini''. Davide ricapitola quelli visti
(la costante singola, il colore, il Texture Sample, le operazioni, il Texture
Coordinate) e ne aggiunge due: il \textbf{Lerp} e il \textbf{Panner}. ``Se abbiamo
capito la regola degli shortcut, questi saranno quale lettera più clic sinistro?''
``L.'' ``Esatto'': \textbf{L + clic} per il Lerp, \textbf{P + clic} per il Panner. ``Volevo
essere sicuro che era arrivata la cosa.''

\paragraph{Due colori.} Il \textbf{Lerp} si occupa dell'\textbf{interpolazione
lineare} (\emph{linear interpolation}). Per farlo vedere Davide usa due colori, ``una
roba che non farete mai, non con due colori perlomeno'': un \textbf{blu} nella
\textbf{A} e un \textbf{rosso} nella \textbf{B}. Il Lerp, collegato al Base Color, dà il
\textbf{viola}. Perché è inutile farlo con i colori? ``Perché potevo cercarmi il viola
direttamente dal colore.'' Ma mostra come funziona l'interpolazione. Il Lerp si può fare
anche con i numeri, ``per fare operazioni matematiche'', e anche con le texture.

\paragraph{L'alfa.} L'operazione è ``gestita interamente dall'\textbf{alfa}'', il terzo
ingresso. L'alfa va \textbf{da 0 a 1} e decide quale colore prevale, con A che vale 0 e
B che vale 1: con l'alfa a 0 si ha il blu, a 1 il rosso, e in mezzo ``una percentuale di
quel colore'', per esempio 0,7. Modificando l'alfa si va sempre più verso il rosso, o
sempre più verso il blu.

\paragraph{Una texture come alfa.} Ma l'alfa non deve essere un numero fisso. Davide
prende ``il nostro fidato Bark'' e collega \textbf{un suo canale} all'alfa: ``vi ricordate
il discorso dei canali sulle texture?'' Ora quella sequenza di 0 e di 1, di bianco e di
nero, ``è l'\textbf{arbitro} nel decidere dove ci sarà più A e dove ci sarà più B'':
nelle zone più nere il Lerp mette più blu, in quelle più bianche più rosso. Questa
texture è ``estremamente bianca'', quindi il risultato tende al rosso; con altre
texture, o lavorando su questa, si otterrebbero risultati diversi. E collegando il
Multiply delle coordinate anche al Texture Sample dell'alfa, si ottiene un tiling
anche dell'alfa, cioè di dove va il blu e dove il rosso. ``Chiaro questo processo?''
``Sì.''

\begin{figure}[htbp]
\centering
\begin{tikzpicture}[
  font=\small,
  ue3node/.style={draw, thick, rounded corners=4pt, align=center,
                  minimum height=0.9cm},
  ue3col/.style={-{Stealth[length=2mm]}, thick, blue!60!black},
  ue3flt/.style={-{Stealth[length=2mm]}, thick, green!50!black}
]
  \node[ue3node, fill=blue!30, text width=1.9cm] (a) at (0,1.4) {\textbf{blu}\\
    \footnotesize Constant3Vector};
  \node[ue3node, fill=red!30, text width=1.9cm] (b) at (0,0) {\textbf{rosso}\\
    \footnotesize Constant3Vector};
  \node[ue3node, fill=gray!10, text width=1.9cm] (al) at (0,-1.5)
    {\textbf{Alpha}\\\footnotesize 0,7 oppure\\\footnotesize un canale di T\_Bark};
  \node[ue3node, fill=gray!10, text width=1.6cm, minimum height=2.2cm] (lerp)
    at (3.6,0) {\textbf{Lerp}\\\footnotesize (L + clic)\\[4pt]\footnotesize A\\
    \footnotesize B\\\footnotesize Alpha};
  \node[ue3node, fill=violet!30, text width=2.0cm] (res) at (7.0,0)
    {\textbf{Base Color}\\\footnotesize viola, più o\\\footnotesize meno rosso};
  \draw[ue3col] (a.east) -- ++(0.6,0) |- ([yshift=5mm]lerp.west);
  \draw[ue3col] (b.east) -- ([yshift=0mm]lerp.west);
  \draw[ue3flt] (al.east) -- ++(0.6,0) |- ([yshift=-5mm]lerp.west);
  \draw[ue3col] (lerp) -- (res);
  % la scala dell'alfa
  \begin{scope}[shift={(9.4,-0.6)}]
    \shade[left color=blue!60, right color=red!60] (0,0) rectangle (3.4,0.6);
    \draw[thick] (0,0) rectangle (3.4,0.6);
    \node[font=\footnotesize] at (0,-0.3) {0};
    \node[font=\footnotesize] at (3.4,-0.3) {1};
    \node[font=\footnotesize] at (1.7,-0.3) {Alpha};
    \node[font=\footnotesize] at (0,0.95) {A};
    \node[font=\footnotesize] at (3.4,0.95) {B};
    \draw[thick] (2.38,-0.05) -- (2.38,0.65);
    \node[font=\scriptsize] at (2.38,0.95) {0,7};
  \end{scope}
\end{tikzpicture}
\caption{Il Lerp della dimostrazione: con l'alfa a 0 esce A (il blu), a 1 esce B (il
rosso), in mezzo una miscela. Se l'alfa è un canale di una texture, ogni pixel ha la sua
miscela: blu dove la texture è nera, rosso dove è bianca.}
\label{ue3:fig:lerp}
\end{figure}

\paragraph{L'esercizio della pausa.} ``Provateci voi.'' Durante la pausa, chi vuole può
prendere \textbf{due texture}, collegarle a un Lerp e, giocando con l'alfa, vedere quale
parte delle due texture viene presa di più e quale di meno.

\paragraph{Approfondimento: la formula del Lerp.} L'interpolazione lineare tra A e B
con peso $\alpha$ è
\[
  \operatorname{lerp}(A, B, \alpha) = A\,(1-\alpha) + B\,\alpha .
\]
Con colori, il calcolo si fa canale per canale. Blu $=(0,0,1)$ e rosso $=(1,0,0)$ con
$\alpha = 0{,}5$ danno $(0{,}5;\,0;\,0{,}5)$, il viola della lezione; con $\alpha =
0{,}7$ danno $(0{,}7;\,0;\,0{,}3)$, un viola più rosso. Quando l'alfa è una texture, la
formula si applica pixel per pixel, ed è così che si mescolano due materiali su una
stessa superficie: erba e terra su un terreno, vernice e ruggine su un metallo, con una
maschera in bianco e nero a fare da ``arbitro''.

% ---------------------------------------------------------------------
\section{Il Panner: far muovere le texture}
\label{ue3:sec:panner}

Il \textbf{Panner}\footnote{Nella trascrizione ``panel'', ``planner'' e ``banner''.}
``si occupa di far muovere le texture''. Davide riprende il Base Color, stacca il resto,
e collega il Panner alle \textbf{UV} del Texture Sample, ``perché lui lavora sulle UV
della texture''. Selezionando il nodo (o direttamente sul nodo) si vedono
\textbf{Speed X} e \textbf{Speed Y}. Con Speed X a 1 la \emph{preview} comincia a
muoversi sull'asse X, ``anche un po' troppo veloce per i miei gusti''; con 0,3 in X e
0,1 in Y si ha un Panner che si muove su entrambi gli assi.

\paragraph{A che cosa serve.} ``Che materiale potremmo volerci fare con questo nodo?''
``Dell'acqua.'' ``Dell'acqua. Qualsiasi fluido, qualsiasi cosa che si muove.'' Il Panner
ha il parametro di velocità e l'ingresso delle \textbf{coordinate}: per combinarlo con il
tiling si collega il Texture Coordinate (o il Multiply) all'ingresso Coordinate. C'è
poi l'ingresso \textbf{Time}, con cui si può decidere ``quando e come deve muoversi'';
altrimenti, di default, usa il \emph{game time}, ``cioè è sempre in movimento''.

\paragraph{Il master material cambia, l'istanza segue.} Mentre il Panner lavora, il
pavimento nella scena \textbf{si muove}: sul pavimento c'è la Material Instance, e
``nonostante stiamo facendo delle modifiche al materiale originale e non alla Material
Instance, la Material Instance continua a modificarsi''. È il rovescio del rapporto tra
master e istanze: quello che cambia nel genitore cambia in tutte le istanze.

\paragraph{Approfondimento: il Panner in una formula.} Il Panner non fa altro che
sommare alle coordinate un tempo moltiplicato per una velocità:
\[
  (u', v') = (u + s_x\,t,\; v + s_y\,t),
\]
dove $t$ è il tempo in secondi e $s_x, s_y$ le velocità in ``texture al secondo'': con
Speed X a 1 la texture scorre di una larghezza intera ogni secondo, per questo era
troppo veloce. Siccome la texture si ripete, lo scorrimento è infinito e senza salti,
purché la texture sia \emph{tileable}, cioè con i bordi che combaciano. Due Panner con
velocità e tiling diversi sommati tra loro sono il trucco classico per l'acqua e per la
lava: il movimento non si ripete mai uguale.

% ---------------------------------------------------------------------
\section{La pausa: il progetto che non si trova, e Git}
\label{ue3:sec:pausa}

``Visto che comunque ci siamo fatti un'oretta e mezza'', Davide chiude per ora il
discorso materiali e propone una pausa di una decina di minuti. Al ritorno: l'import,
le mesh, e ``una cosina con la torretta che c'è nei file che vi abbiamo passato'', così
che Igor, quando torna, potrà farci ``un piccolo sistema di puntamento''.

\paragraph{Il progetto che Unreal non vede.} Al rientro, in attesa degli altri, una
studentessa racconta il suo problema: ha scaricato e decompresso la cartella del Drive,
ma aprendo Unreal il progetto non compare, e deve aprirlo ``a mano'' da un'altra parte.
Davide spiega che dipende probabilmente dal percorso in cui è salvata la cartella: se la
si mette dove sono salvati gli altri progetti, Unreal dovrebbe trovarla; in ogni caso,
aprire il progetto dal file eseguibile nella cartella funziona. ``Se non è un problema
che ti impedisce di accedere al progetto, fa niente.''

\paragraph{Come si lavora insieme a distanza?} Uno studente, programmatore, fa una
domanda ``fuori dall'argomento'': i programmatori usano \textbf{Git}, ma i file di un
progetto di Unreal sembrano troppo pesanti per lavorarci in contemporanea da lontano.
Davide conferma che \textbf{Git}, ``o diciamo repository di sorta'', è ``la cosa
principale che utilizzerai''. Con \textbf{Unity} ha sempre avuto meno problemi, perché
Unreal ``va gestito in maniera un po' più delicata'', ma entrambi supportano senza
problemi repository e \emph{source control}. Esistono anche altri metodi, che Davide non
ha mai provato: per Unreal, ``una cosa che io sento da una vita'', un modo per lavorare
\textbf{contemporaneamente} sulla stessa cosa, tramite \textbf{VPN} o in \textbf{LAN} e
con dei \emph{plugin} specifici, in modo che due colleghi modifichino attivamente lo
stesso livello. Con Git invece ``ognuno lavora sul suo e poi si \emph{merge} tutto
insieme''. ``Git è quello che ho utilizzato sempre, è il più comodo, quindi dovresti
trovarti abbastanza a casa.''

\paragraph{Approfondimento: Git LFS, Perforce e Multi-User Editing.} I timori dello
studente sono fondati. Git nasce per file di testo, e gli \emph{asset} di Unreal
(\texttt{.uasset}, \texttt{.umap}) sono file \textbf{binari}, spesso grandi, che Git non
sa fondere: se due persone modificano lo stesso livello, una delle due modifiche va
persa. Per questo con Unreal si usa \textbf{Git LFS} (\emph{Large File Storage}), che
tiene i file pesanti fuori dalla cronologia normale e permette di \textbf{bloccare}
(\emph{lock}) un file mentre lo si modifica. Nell'industria, per i progetti Unreal, è
molto diffuso \textbf{Perforce} (Helix Core), il sistema che Epic stessa usa e che
l'editor supporta direttamente, proprio per la gestione dei file binari e dei blocchi.
Il sistema ``sentito da una vita'' di cui parla Davide esiste: è il plugin
\textbf{Multi-User Editing}, con cui più editor connessi in rete (in LAN o via VPN)
condividono una stessa sessione e vedono le modifiche degli altri in tempo reale; nato
per la produzione virtuale nel cinema, si può usare anche nello sviluppo di giochi, ma
non sostituisce il \emph{source control}, a cui le modifiche vanno comunque consegnate
alla fine.

% ---------------------------------------------------------------------
\section{L'import e la torretta}
\label{ue3:sec:import}

Passato ``il nostro quarto d'ora'', si riparte: si importano le prime mesh, la
\textbf{torretta}, ``le tre parti della torretta, perché l'abbiamo divisa in tre
parti''.

\paragraph{La finestra di import, sulla slide.} Quando si importa una mesh compare una
finestra (``che può essere più larga, può essere più stretta'') piena di parametri. Le
cose principali che si possono fare da lì:
\begin{itemize}
  \item \textbf{generare automaticamente le collisioni}, se non sono presenti nel file
        FBX, perché magari gli oggetti non sono stati creati con le collisioni;
  \item \textbf{modificare il transform} dell'oggetto importato ``proprio di base'':
        ruotarlo sul suo \textbf{pivot}, traslarlo. ``Avete già parlato di cos'è il
        pivot con Igor, giusto? Il fulcro che vi permette di traslare e ruotare gli
        oggetti all'interno della scena.'' Con la funzione di \textbf{reimport}, si
        vedrà, questa possibilità diventa molto utile;
  \item adattare le \textbf{unità} e gli \textbf{assi} tra il software di origine e
        Unreal: \textbf{convertire la scena} dalle coordinate FBX a quelle di Unreal,
        \textbf{convertire le unità} (Unreal lavora in centimetri), oppure usare
        \textbf{Force Front X Axis}, che cambia l'asse che Unreal considera come
        ``fronte'' dell'oggetto. Davide l'ha usato poche volte: una volta il modello di
        un \emph{character}, importato normalmente, ``veniva rivoltato, come se il
        dentro fosse fuori'', e cambiare l'asse frontale gli ha restituito la mesh
        com'era stata esportata (``mi pare fosse Blender'');
  \item decidere se importare anche i \textbf{materiali} e le \textbf{texture}
        collegati alla mesh, e \textbf{dove} cercarli: in locale, in un'altra cartella,
        e così via.
\end{itemize}

\paragraph{Il pulsante Import.} Prima di importare, Davide spegne il Panner (``facciamo
che smetto di far muovere il pavimento'') e il pavimento torna normale. Nel Content
Browser c'è ``questo pulsantone con \textbf{Import}'', che apre l'esplora risorse del
computer per cercare i file. Nella cartella del materiale didattico c'è una cartella
\textbf{torretta}, con tre \textbf{static mesh}: \textbf{SM\_Turret\_Base},
\textbf{SM\_Turret\_Cannon} e \textbf{SM\_Turret\_Pivot}\footnote{Nella trascrizione
``Turret Base, Turret Cannon e Turret Pivo''; il prefisso SM\_ è detto a voce.}.

\paragraph{La finestra di import, dal vivo.} Importandole si apre la finestra (``nel
mio caso è così enorme, ma perché io tendo a usarla così''), e si ritrovano le opzioni
della slide: la possibilità di modificare traslazione e rotazione, ``anche
singolarmente sui vari assi'' (``questo è qualcosa che faremo più tardi''); i
\textbf{LOD} (``decidere il gruppo dei LOD''); le collisioni, se e come importarle. Se si
importa una \textbf{Skeletal Mesh}\footnote{Nella trascrizione ``scheda dal Mesh'' e
``Scheletal Mesh'', come nella seconda lezione.}, come un personaggio, la finestra
permette di collegarla a uno \textbf{scheletro} esistente, di generarne uno, oppure di
importare solo delle \textbf{animazioni}: ``faremo una lezione specifica sulle
animazioni, quindi per il momento glissiamo''. In alto ci sono anche le sezioni per
andare più veloci, tra cui \textbf{Import Materials} e \textbf{Import Textures}, con la
\textbf{Search Location} in cui cercarli.

\paragraph{Niente materiali e texture, per ora.} La consegna di Davide: \textbf{togliere
la spunta} a Import Materials e Import Textures, e fare \textbf{Import}. Unreal importa i
tre oggetti (Davide elimina quelli della prova precedente e salva). Perché togliere le
spunte? Se si fossero importati gli oggetti con materiali e texture, Unreal avrebbe
importato automaticamente tutte le texture, creato i materiali adatti e li avrebbe
anche \textbf{applicati}, ``in modo da averli già precisi sulla torretta''. Ed è
proprio quello che si farà, ma \textbf{a mano}, ``così vediamo effettivamente anche il
processo''. ``Ci sono stati problemi? Tutto a posto con la torretta?'' ``Sì, nessun
problema.''

\begin{table}[htbp]
\centering
\small
\begin{tabular}{@{}p{4.2cm}p{8.8cm}@{}}
\toprule
\textbf{Opzione} & \textbf{A che cosa serve} \\
\midrule
Generate Missing Collision & genera le collisioni se il file FBX non le ha \\
Import Translation / Rotation / Uniform Scale & sposta, ruota, scala l'oggetto ``alla
  radice'': cambia dove sta il pivot rispetto alla geometria \\
LOD Group & sceglie un gruppo di livelli di dettaglio (\emph{Level Of Detail}) \\
Convert Scene, Convert Scene Unit & adattano assi e unità di misura del file FBX a
  quelli di Unreal (centimetri) \\
Force Front X Axis & cambia l'asse considerato come ``fronte'' \\
Skeleton / Import Animations & per le Skeletal Mesh: scheletro esistente o nuovo, solo
  animazioni \\
Import Materials, Import Textures & importano e applicano materiali e texture
  collegati \\
Search Location & dove cercare materiali e texture da importare \\
\bottomrule
\end{tabular}
\caption{Le opzioni della finestra di import citate a lezione, con i nomi della
finestra di Unreal Engine 5 (FBX Import Options; il raggruppamento può variare tra le
versioni). Per la torretta: Import Materials e Import Textures senza spunta.}
\label{ue3:tab:import}
\end{table}

% ---------------------------------------------------------------------
\section{Che cos'è una mesh}
\label{ue3:sec:mesh}

Mentre la torretta è importata, Davide torna sulle slide per ``le ultime cose sulle
mesh''. Sulla slide c'è un \textbf{vaso} in \emph{wireframe} e, sotto, una
semplificazione. Una mesh è ``un insieme di \textbf{vertici}, \textbf{spigoli} e
\textbf{facce} che vanno a creare un oggetto in tre dimensioni''.
\begin{itemize}
  \item I \textbf{vertici} sono i \textbf{punti} che costituiscono gli angoli
        dell'oggetto, punti nello spazio.
  \item Gli \textbf{spigoli} (\emph{edge}) sono i \textbf{segmenti} che collegano due
        vertici.
  \item Le \textbf{facce} sono le \textbf{superfici piatte} definite da più segmenti:
        sulla slide, ``tutti i triangoli''.
\end{itemize}
Il \textbf{triangolo} è ``la forma più comoda che esiste a livello di calcolo'', e
tendenzialmente tutti gli oggetti saranno composti da ``un insieme di triangolini''. Più
triangoli insieme, più facce, formano dei \textbf{poligoni}; i poligoni combinati formano
delle \textbf{superfici}; e la superficie determina la mesh. Tutto questo si può vedere
anche nell'engine, con la visualizzazione \textbf{wireframe}: l'oggetto non ``come
appare'', ma ``come griglia, come insieme di tutti i poligoni, vertici e spigoli''.

\begin{figure}[htbp]
\centering
\begin{tikzpicture}[font=\small]
  \coordinate (A) at (0,0);
  \coordinate (B) at (2.6,0.3);
  \coordinate (C) at (3.2,2.2);
  \coordinate (D) at (0.6,2.6);
  \coordinate (E) at (1.7,1.3);
  \fill[orange!15] (A) -- (B) -- (E) -- cycle;
  \fill[orange!30] (B) -- (C) -- (E) -- cycle;
  \fill[orange!20] (C) -- (D) -- (E) -- cycle;
  \fill[orange!10] (D) -- (A) -- (E) -- cycle;
  \draw[very thick] (A) -- (B) -- (C) -- (D) -- cycle;
  \draw[very thick] (A) -- (E) (B) -- (E) (C) -- (E) (D) -- (E);
  \foreach \p in {A,B,C,D,E} \fill (\p) circle (2.5pt);
  \draw[-{Stealth[length=2mm]}] (4.3,2.8) node[right, font=\footnotesize]
    {\textbf{vertice}: un punto} -- (C);
  \draw[-{Stealth[length=2mm]}] (4.3,1.6) node[right, font=\footnotesize]
    {\textbf{spigolo} (\emph{edge}): un segmento tra due vertici} -- ($(B)!0.5!(C)$);
  \draw[-{Stealth[length=2mm]}] (4.3,0.4) node[right, font=\footnotesize]
    {\textbf{faccia}: un triangolo} -- (1.45,0.55);
  \node[font=\footnotesize, align=center, text width=4cm] at (1.6,-0.8)
    {4 facce triangolari\\formano un poligono};
\end{tikzpicture}
\caption{Vertici, spigoli e facce, come nella semplificazione della slide. Il triangolo
è l'unità di calcolo: anche le facce a quattro lati dei programmi di modellazione,
quando arrivano in Unreal, vengono divise in triangoli.}
\label{ue3:fig:mesh}
\end{figure}

% ---------------------------------------------------------------------
\section{Dimensioni, unità di misura e formati}
\label{ue3:sec:regole}

Quando si porta una mesh in Unreal, mentre la si produce o mentre la si cerca, bisogna
rispettare alcune cose.

\paragraph{Le dimensioni.} La cosa più comune: si decide di fare ``una mappa molto
carina'' su Unreal, si cercano mesh su internet o su \textbf{Fab} (``su Fab questo
problema c'è meno, però utilizzando siti esterni può capitare''), e la mesh importata è
\textbf{gigantesca} o \textbf{incredibilmente piccola}, perché chi ha fatto l'export, o
chi ha fatto l'import, non ha rispettato ``le misure di adattamento e transizione da un
software all'altro''. Unreal usa come unità di misura i \textbf{centimetri}; Blender,
per esempio, ``usa altra roba''. Nel passaggio da Blender a Unreal bisogna controllare
che la conversione delle dimensioni sia corretta.

\paragraph{FBX.} ``Altra cosa molto importante'': Unreal ``tendenzialmente legge
soltanto \textbf{FBX}''. Sono supportati anche altri formati, come gli \textbf{OBJ}, ``ma
non assicurano un corretto funzionamento all'interno dell'engine''. Il consiglio:
``importate sempre unicamente FBX''.

\paragraph{I formati delle texture.} Per le texture i formati supportati sono molti
(sulla slide c'è un elenco), ma nella pratica si usano soprattutto \textbf{JPEG} e
\textbf{PNG}, ``i più comuni, a meno che non state facendo qualcosa di molto
specifico''.

\paragraph{Approfondimento: Blender, i metri e il fattore 100.} Blender lavora di default
in metri, con una ``unità'' uguale a un metro; Unreal lavora in centimetri. Un cubo di
lato 2 in Blender è quindi un cubo di 2~m, cioè 200 unità di Unreal: se l'export non
converte, l'oggetto arriva cento volte troppo piccolo (o, con altre impostazioni, cento
volte troppo grande). La soluzione standard è nell'esportatore FBX di Blender
(l'opzione \emph{Apply Scalings} su \emph{FBX Units Scale}) o, lato Unreal, nella
scala uniforme e nella conversione delle unità della finestra di import. Un controllo
rapido: un personaggio umano in Unreal è alto circa 180 unità, come il manichino dei
progetti modello; se la mesh appena importata è molto diversa, c'è un problema di scala.
Oltre a FBX, le versioni recenti di Unreal Engine 5 importano anche \textbf{glTF} e
\textbf{USD}, ma FBX resta il formato più collaudato, come dice Davide.

% ---------------------------------------------------------------------
\section{Il pivot del cannone e il Reimport}
\label{ue3:sec:pivot}

Davide apre le tre mesh (la base, la torretta e il pivot) e le piazza in scena. Sono
state esportate con il \textbf{pivot nello stesso punto}: per sovrapporle senza
spostarle a mano basta \textbf{copiare la X e la Y} di un pezzo negli altri. Ed ecco la
torretta, ``a parte il fatto che non c'ha il materiale''.

\paragraph{I problemi della torretta.} Ma ``sta torretta ha dei problemi: quali sono?''
Il primo si vede: per qualche motivo \textbf{non è allineata}, ed è ``scomodamente''
ruotata: trovare la rotazione giusta sarà scomodo. Il secondo è più sottile. Si
supponga di mettere i tre pezzi insieme in un \textbf{Blueprint}. Che cosa deve fare una
torretta? Deve \textbf{tracciare dei bersagli}: ruotare su sé stessa, ma anche andare su
e giù, se il bersaglio è in alto o in basso. E qui c'è il problema: il \textbf{pivot del
cannone} non è in una posizione comoda per la rotazione verticale, perché sta
\textbf{in basso}, in fondo. Ruotare sul piano orizzontale non dà problemi, ma per
alzare e abbassare la canna il pivot dovrebbe essere ``più o meno all'altezza di questo
\textbf{bullone}'', il fulcro, e da lì il cannone dovrebbe essere manipolato e ruotato
(figura~\ref{ue3:fig:pivot}).

\begin{figure}[htbp]
\centering
\begin{tikzpicture}[font=\small,
  ue3arr/.style={-{Stealth[length=2.5mm]}, thick}]
  \foreach \xs/\titolo/\px/\py/\ang in {0/{pivot in basso (com'era)}/1.25/0.0/0,
                                         6.6/{pivot nel fulcro (dopo il Reimport)}/1.25/1.75/0} {
    \begin{scope}[shift={(\xs,0)}]
      % base
      \draw[thick, fill=orange!35] (0.3,0) -- (2.2,0) -- (1.9,0.5) -- (0.6,0.5) -- cycle;
      % pivot (supporto)
      \draw[thick, fill=gray!30] (0.8,0.5) rectangle (1.7,1.5);
      % cannone
      \draw[thick, fill=orange!20] (0.7,1.5) rectangle (1.8,2.1);
      \draw[thick, fill=orange!20] (1.8,1.6) rectangle (4.0,1.95);
      % bullone
      \fill[black] (1.25,1.75) circle (2pt);
      % il pivot del cannone
      \draw[red, very thick] (\px-0.18,\py) -- (\px+0.18,\py) (\px,\py-0.18) -- (\px,\py+0.18);
      \draw[red, thick] (\px,\py) circle (0.18);
      \node[font=\footnotesize, align=center] at (2.0,-0.6) {\titolo};
    \end{scope}
  }
  % archi di rotazione
  \draw[ue3arr, red] (3.4,0.5) arc[start angle=-20, end angle=25, radius=2.6];
  \node[font=\scriptsize, red, align=center] at (4.6,0.55) {la canna ruota\\attorno al
    piede};
  \draw[ue3arr, red] (6.6+3.6,1.0) arc[start angle=-30, end angle=30, radius=1.6];
  \node[font=\scriptsize, red, align=center] at (6.6+4.75,1.75) {la canna ruota\\
    attorno al bullone};
  % assi
  \draw[ue3arr, blue] (-0.6,0) -- (-0.6,1.2) node[above, font=\scriptsize]{Z};
\end{tikzpicture}
\caption{Il problema del pivot del cannone, visto di lato (schema non in scala). Con il
pivot (il mirino rosso) in fondo, alzare la canna la farebbe girare attorno al piede,
staccandola dal supporto; con il pivot nel fulcro, all'altezza del bullone, la canna
si alza e si abbassa come una vera torretta. La rotazione orizzontale, attorno all'asse
Z, funziona invece in entrambi i casi.}
\label{ue3:fig:pivot}
\end{figure}

\paragraph{Come si sposta un pivot.} ``In realtà possiamo farlo in tutta una miriade di
modi.'' Con la \textbf{Modeling Mode} si potrebbe rimodellare l'oggetto. Ma lo si può
fare anche \textbf{nell'import}: si fa tasto destro sulla mesh e \textbf{Reimport Mesh
With Dialog}, che riapre la finestra di import vista prima. Unreal importa lo stesso
oggetto e lo \textbf{sovrascrive}, ma permette di cambiare come verrà posizionato:
l'\textbf{offset della rotazione} e l'\textbf{offset della traslazione}.

\paragraph{I numeri.} Davide se li era scritti, ma non li ricorda, e li ricava. Per
sapere quale asse toccare, guarda il \emph{gizmo}: il cannone deve essere allineato con
il resto della torretta (e lo stesso il pivot), e il pivot del cannone deve stare non
``qui giù in fondo'', ma allineato con il fulcro, sull'\textbf{asse Z}, ``perché è
questo blu''. Per questa torretta, ``quindi anche per voi'', i parametri sono:
\begin{itemize}
  \item per il \textbf{cannone}: traslazione Z $= -135$, che reimporta il cannone ``più
        attaccato al pivot'', e rotazione $= 113$, che lo riallinea\footnote{La
        trascrizione dice prima ``113 per la traslazione'' e poi si corregge: ``e 113
        sulla rotazione''. L'asse della rotazione non è detto; trattandosi di
        riallineare la torretta nel piano orizzontale, è con ogni probabilità la
        rotazione attorno a Z (il \emph{yaw}).};
  \item per il \textbf{pivot}: ``questa volta non serve modificare la traslazione'', ma
        solo la rotazione, sempre 113.
\end{itemize}
Dopo il reimport il cannone ``sta sotto'' (la sua origine ora è nel fulcro, quindi
nella posizione di prima risulta più basso) e va rimesso a posto; ma ora, ruotandolo, si
alza e si abbassa nel modo giusto. ``E chiaramente poi questo farete più avanti, con dei
Blueprint.'' Si salva tutto.

\paragraph{Approfondimento: perché il pivot conta così tanto.} In Unreal ogni rotazione
e ogni scala di un componente avvengono \textbf{attorno alla sua origine}, cioè al
pivot della mesh. Un Blueprint di torretta, di solito, ha la base come componente
radice, il supporto come figlio che ruota attorno a Z (il \emph{yaw}, per seguire il
bersaglio a destra e a sinistra) e il cannone come figlio del supporto che ruota attorno
a Y (il \emph{pitch}, su e giù). Per questo ogni pezzo deve avere l'origine esattamente
nel punto in cui gira nella realtà: se l'origine è sbagliata, nessun nodo di rotazione
potrà correggerla senza calcoli aggiuntivi. Il ``piccolo sistema di puntamento'' che
Igor farà con questa torretta (con nodi come \emph{Find Look at Rotation}, che calcola
la rotazione necessaria per guardare un punto) si appoggia interamente a questo lavoro
preparatorio.

% ---------------------------------------------------------------------
\section{Mettere ordine: le cartelle}
\label{ue3:sec:cartelle}

Prima di importare le \textbf{texture}, ``per dare un po' di ordine'', Davide crea nel
Content Browser una nuova cartella, \textbf{Torretta}, ci sposta dentro tutto
(``muoviamo, così è un po' più pulito''), e dentro crea un'ulteriore cartella,
\textbf{Texture}, in cui importare le texture della torretta. ``Questo non è
obbligatorio, è solo una questione di \emph{best practice}, cioè è come il discorso
delle Material Instance. Non vi vengo a trovare a casa se non lo fate, però vi risulta
tutto più comodo.''

% ---------------------------------------------------------------------
\section{Lo Static Mesh Editor e le collisioni}
\label{ue3:sec:sme}

Con un doppio clic su una mesh si apre lo \textbf{Static Mesh Editor}: ``andiamoci a
vedere effettivamente un po' qui dentro cosa possiamo vedere delle mesh''.

\paragraph{Gli slot dei materiali.} A destra, nei \textbf{Material Slots}, queste mesh
hanno \textbf{due slot}: ``non è la normalità''. (Una delle tre ne ha uno solo.) Il modo
in cui la torretta è stata suddivisa fa sì che, per rendere bene i materiali, servano
due slot: la base e il cannone condividono in parte il materiale del \textbf{pivot}. Nei
due slot si inseriranno i materiali più avanti. Altre impostazioni (Davide accenna a
quelle sui LOD e sulla \emph{Nanite}\footnote{La trascrizione dice ``per quanto riguarda
la night'': la ricostruzione come Nanite, la tecnologia di geometria virtualizzata di
Unreal Engine 5 che ha una sua sezione nei dettagli della mesh, è probabile ma non
certa.}) non si vedono oggi.

\paragraph{Le modalità di visualizzazione.} Dal menu della \emph{viewport} (la
prospettiva ``non ci interessa'') si cambia la modalità: \textbf{Lit}, quella normale;
\textbf{Unlit}, ``com'è il mio oggetto se non è colpito dalla luce?''; e
\textbf{Wireframe}: ``vedete quello che dicevamo prima?'' Vertici, facce ed \emph{edge}
della mesh, che si potrebbero modificare con la \textbf{Modeling Mode} (``se avanzerà
il tempo con le lezioni, sarebbe carino vederla più avanti''). Ci sono anche
\textbf{Lighting Only}, le \textbf{Reflections} e altre ancora. Dal menu \textbf{Show}
si possono mostrare i \textbf{vertici}: ``dove sono i vertici di questo oggetto?''

\paragraph{Simple e complex.} Un'altra cosa utile ``da un punto di vista di gameplay''
sono le \textbf{collisioni}. Dal menu Show si attiva la \textbf{Simple Collision}: ``queste
sono le collisioni del mio oggetto''. In scena, il \emph{character} interagirà con
queste: ci salterà sopra, o lo dovrà scavalcare. Ma oltre alle \textbf{simple
collision} ci sono le \textbf{complex collision}. La differenza:
\begin{itemize}
  \item le \textbf{simple collision} sono forme che si possono posizionare a mano,
        anche da qui, e rifare da capo;
  \item le \textbf{complex collision} usano la struttura in vertici della mesh stessa
        come collisione. ``È qualcosa che è molto più pesante rispetto a una collisione
        semplice, e torna utile solo in alcuni casi specifici.''
\end{itemize}

\paragraph{Modificare le collisioni.} La collisione semplice si può selezionare,
spostare, anche eliminare; e se ne possono \textbf{aggiungere} altre: una
\textbf{sfera}, di cui si decide la grandezza; una \textbf{capsula}, che ``è
tendenzialmente la collisione che usano i vostri \emph{character}'': il personaggio di
solito non ha una collisione super precisa, ``è un capsulone'', a meno che non si voglia
davvero il contrario; una \textbf{box}. E una box, per questa mesh, ``è più ottimizzata
rispetto alla collisione che questa aveva di base'': all'atto pratico fa la stessa cosa
``con molti meno vertici''. Su una singola mesh cambia poco; su molte mesh aiuta le
prestazioni. Ci sono anche collisioni più complesse, ma ``idealmente, il modo migliore
per fare delle simple collision è letteralmente prendere delle \textbf{forme base}'',
magari più d'una, ``piuttosto che avere tanti vertici sulla stessa collisione''.

\begin{table}[htbp]
\centering
\small
\begin{tabular}{@{}p{3.3cm}p{4.4cm}p{5.3cm}@{}}
\toprule
\textbf{Collisione} & \textbf{Com'è fatta} & \textbf{Quando} \\
\midrule
Box & un parallelepipedo & oggetti squadrati; spesso basta anche per oggetti più
  complessi, ed è la più leggera \\
Sphere & una sfera & oggetti tondi, palle \\
Capsule & un cilindro con due semisfere & i \emph{character}: ``un capsulone'' \\
Convex (più forme) & poliedri convessi, magari più d'uno & quando una forma sola non
  basta; meglio poche forme base che tanti vertici \\
Complex & i triangoli della mesh stessa & casi specifici: precisione massima, costo
  molto più alto \\
\bottomrule
\end{tabular}
\caption{Le collisioni di una static mesh viste nello Static Mesh Editor. Le prime
quattro sono \emph{simple collision}, che si aggiungono e si modificano dal menu
Collision dell'editor; l'ultima è la \emph{complex collision}.}
\label{ue3:tab:collisioni}
\end{table}

\paragraph{Approfondimento: chi usa quale collisione.} Unreal tiene le due collisioni
separate perché le usa per cose diverse. La \textbf{simple} è quella della
\textbf{fisica}: la pallina della seconda lezione che rotola sulla piattaforma, il
\emph{character} che cammina e salta, gli oggetti che cadono; i calcoli fisici su forme
semplici sono molto più veloci. La \textbf{complex} si usa soprattutto per le
\textbf{interrogazioni} precise, come un \emph{line trace} (un ``raggio'' lanciato per
sapere che cosa colpisce un proiettile o che cosa sta guardando il giocatore), quando
la precisione al triangolo serve davvero. Nei dettagli della mesh, il campo
\emph{Collision Complexity} decide come combinarle; l'opzione \emph{Use Complex
Collision As Simple} usa i triangoli anche per la fisica, ed è comoda per oggetti
statici di forma strana, ma va usata con giudizio, per il costo di cui parla Davide.

% ---------------------------------------------------------------------
\section{Raggruppare e fondere: Merge Actors}
\label{ue3:sec:merge}

``Se ad esempio vogliamo avere tutti e tre gli oggetti assieme'', o semplicemente
spostarli insieme, ci sono due strade.

\paragraph{Raggruppare.} Si selezionano tutti e tre e li si \textbf{raggruppa} con
\textbf{Ctrl+G}; si spostano insieme; con \textbf{Shift+G} si toglie il gruppo.

\paragraph{Fondere.} Oppure li si \textbf{fonde} (\emph{merge}). Si evidenziano i tre
pezzi in scena, si va nel menu \textbf{Actor} in alto e si sceglie \textbf{Merge
Actors}. Unreal chiede dove salvare il nuovo oggetto e con che nome (per esempio
\textbf{SM\_Turret}\footnote{Nella trascrizione ``smTourette''.}), e crea una
\textbf{nuova mesh}, ``che è semplicemente la fusione delle tre mesh precedenti'': una
somma, un \emph{merge}, appunto.

\paragraph{Il merge, di nuovo.} Più tardi uno studente dice di essersi perso come si fa
il merge, e Davide lo rifà con calma. I tre oggetti (il cannone, la base e il pivot)
sono in scena. Con il \textbf{tasto destro} su di essi ci sono delle opzioni, ``però
queste sono opzioni limitate''. Per averle tutte si usa il menu \textbf{Actor} in alto:
mostra gran parte delle stesse opzioni, ma in basso, nella sezione degli strumenti di
Unreal Engine, c'è \textbf{Merge Actors}. Al merge Unreal chiede dove salvare e con che
nome, e ``ti fonde i tre oggetti''. Nello stesso menu c'è anche \textbf{Chaos}, che
``serve per gli oggetti distruttibili'': cose più avanzate, che probabilmente non si
vedranno nel corso, ma da approfondire ``se vi capiterà di voler mettere distruttibilità
all'interno del vostro gioco''.

\paragraph{Approfondimento: gruppo, merge o Blueprint?} Le tre soluzioni non sono
equivalenti. Il \textbf{gruppo} è solo una comodità dell'editor: gli actor restano
separati, e in gioco non cambia nulla. Il \textbf{merge} crea una mesh nuova, una sola:
meno oggetti da disegnare (utile per le prestazioni, per esempio con tanti pezzi di
arredo statici), ma i pezzi \textbf{non si possono più muovere uno rispetto
all'altro}. Per la torretta di Igor, che deve ruotare il supporto e alzare la canna, il
merge è quindi la scelta sbagliata: serve un \textbf{Blueprint} con i tre pezzi come
\textbf{componenti} separati, imparentati tra loro, come si è visto nella seconda
lezione. Il merge della lezione è una dimostrazione dello strumento, e serve anche a
rispondere alla domanda sui materiali (sezione~\ref{ue3:sec:domande}).

% ---------------------------------------------------------------------
\section{I materiali della torretta}
\label{ue3:sec:matturret}

``Ora abbiamo importato le texture. Creiamoci i materiali di queste torrette.'' Tre
materiali, uno per pezzo: \textbf{M\_Turret\_Base}, \textbf{M\_Turret\_Pivot} e
\textbf{M\_Turret\_Cannon}. Si potrebbero fare molto elaborati, ``ma partiamo
semplicemente andando a utilizzare le texture che ci sono state date''.

\paragraph{Quattro Texture Sample.} Nel materiale della base Davide prende \textbf{quattro
Texture Sample} e ci mette le texture della base: \textbf{Base Color},
\textbf{Metallic}, \textbf{Roughness} e \textbf{Normal} (Turret Base Color, Turret Base
Metallic, Turret Base Roughness, Turret Base Normal). Poi le collega ai pin
corrispondenti del nodo risultante: il colore al Base Color, la Metallic al Metallic, la
Roughness alla Roughness, la Normal alla Normal (figura~\ref{ue3:fig:turret}).

\begin{figure}[htbp]
\centering
\begin{tikzpicture}[
  font=\small,
  ue3node/.style={draw, thick, rounded corners=4pt, align=center,
                  minimum height=0.85cm, text width=3.6cm},
  ue3col/.style={-{Stealth[length=2mm]}, thick, blue!60!black},
  ue3flt/.style={-{Stealth[length=2mm]}, thick, green!50!black}
]
  \node[ue3node, fill=orange!20] (c) at (0,2.1) {\textbf{Texture Sample}\\
    \footnotesize T\_Turret\_Base\_Color};
  \node[ue3node, fill=gray!15] (m) at (0,0.7) {\textbf{Texture Sample}\\
    \footnotesize T\_Turret\_Base\_Metallic};
  \node[ue3node, fill=gray!15] (r) at (0,-0.7) {\textbf{Texture Sample}\\
    \footnotesize T\_Turret\_Base\_Roughness};
  \node[ue3node, fill=blue!15] (n) at (0,-2.1) {\textbf{Texture Sample}\\
    \footnotesize T\_Turret\_Base\_Normal};
  \node[draw, very thick, rounded corners=4pt, fill=blue!6, minimum width=3.0cm,
    minimum height=5.2cm] (res) at (6.2,0) {};
  \node[font=\small\bfseries] at (6.2,2.3) {M\_Turret\_Base};
  \foreach \lab/\y in {Base Color/2.1, Metallic/0.7, Specular/0.0, Roughness/-0.7,
                       Emissive Color/-1.4, Normal/-2.1}
    \node[font=\footnotesize, anchor=west] at (4.85,\y) {\lab};
  \draw[ue3col] (c.east) -- (4.8,2.1);
  \draw[ue3flt] (m.east) -- node[above, font=\scriptsize]{R} (4.8,0.7);
  \draw[ue3flt] (r.east) -- node[above, font=\scriptsize]{R} (4.8,-0.7);
  \draw[ue3col] (n.east) -- (4.8,-2.1);
\end{tikzpicture}
\caption{Il materiale della base della torretta, fatto a mano: quattro Texture Sample,
quattro pin. Metallic e Roughness sono mappe in bianco e nero, e basta un canale (per
esempio R). M\_Turret\_Pivot e M\_Turret\_Cannon sono identici, con le loro texture.
Ricostruzione dalla dimostrazione; i nomi esatti dei file sono indicativi.}
\label{ue3:fig:turret}
\end{figure}

\paragraph{Le preview.} ``Mi ero dimenticato di farvi vedere una cosa.'' Nella
\emph{preview} del materiale si può scegliere la forma: il \textbf{cilindro}, la
\textbf{sfera}, il \textbf{piano}, il \textbf{cubo} e il \textbf{custom}. Con il custom,
se si seleziona una mesh nel Content Browser e si clicca il pulsante, la preview mostra
\textbf{quella mesh} con il materiale applicato. La prima volta non la mostra: ``sta
compilando''. Poi compare (Davide aveva collegato una texture ``dalla parte sbagliata'',
e la sistema).

\paragraph{Gli slot.} Il materiale è salvato e compilato. Sulla mesh della torretta ci
sono i \textbf{due slot}: in uno c'è M\_Turret\_Base, ma ``per come è stata strutturata
la torretta'' nell'altro va il materiale del \textbf{pivot}. ``Sta rallentando
veramente tanto il PC.''

\paragraph{Un grafico, un engine.} Davide si ferma su che cosa si è fatto: ``qui non
abbiamo fatto altro che utilizzare le mesh che ci sono state fornite da un grafico''.
Idealmente un grafico, in un software a parte, ha realizzato la mesh e le texture, e ce
le ha passate; noi le abbiamo inserite. Si potrebbe lavorare ancora sull'estetica della
torretta modificando i materiali, ``tuttavia, in questo caso non è necessario''.

\paragraph{Il pivot e il cannone, più furbi.} Per il materiale del pivot Davide è ``un
po' più furbo'': cerca la texture \emph{Turret Pivot Color}, e le altre le copia da lì,
perché nella lista ``stanno direttamente dopo''. Collega Metallic, Roughness e Normal,
salva, ``e abbiamo anche il materiale del nostro pivot''. Poi lo stesso per il cannone:
la texture \emph{Cannon Color}, \textbf{Ctrl+C}, \textbf{Ctrl+V}. Lì Davide ricorda le
scorciatoie classiche: ``tutti i classici shortcut di Control'', Ctrl+C, Ctrl+V, Ctrl+Z,
Ctrl+Y, Ctrl+A, ``sono applicabili anche a Unreal nella maggior parte dei casi''.

\paragraph{Salvate spesso.} Uno studente chiede perché il suo non si è applicato:
``Salvalo.'' E Davide ne fa una regola: ``puntualmente mi capita di vedere gente che sta
lavorando a qualcosa, non salva per qualche motivo, Unreal crasha e perde tipo tutto il
lavoro che stava facendo. E quindi salvate spesso.''

\paragraph{La torretta finita.} Nei due slot vanno il pivot e il cannone. ``Oh, eccoci
qua. Bello lucido.'' Salvato tutto, la torretta ha i suoi materiali.

\paragraph{Ritoccare il colore con un Lerp.} ``Se volete che la torretta sia meno
lucida, o che sia meno nitido l'arancione'', si può intervenire sui singoli materiali.
Davide improvvisa sul cannone: non è sicuro di volere il colore ``così netto'', quindi
prende un \textbf{nero} e un \textbf{Lerp}, collega il nero e la texture del colore alle
due entrate e mette nell'alfa un parametro con cui decidere ``quanto si deve vedere'',
per esempio 0,6. Si può convertire in parametro, creare una \textbf{Material
Instance}, e regolarlo da lì. Al posto del nero si può mettere un altro colore, un
verde: ``viene un po' più verdastra, è troppo verde'', e con l'alfa a 0,8 viene ``questo
verde mezzo, insomma, un po' scuro, un po' cromato''. ``Qui chiaramente sto improvvisando
cose, ma cose anche stupide. Ma nel momento in cui voi avete un'idea un po' più chiara,
un po' più specifica di quello che volete, potete andare a fare un po' più di cose.''

\paragraph{Approfondimento: lo stesso lavoro con un master material.} I tre materiali
della torretta hanno lo stesso identico grafo: cambiano solo le texture. È il caso da
manuale per le Material Instance della prima parte della lezione: un solo master
material (per esempio \texttt{M\_Turret}) con quattro Texture Parameter (colore,
metallic, roughness, normal) e magari un Lerp per la tinta, e tre istanze
(\texttt{MI\_Turret\_Base}, \texttt{MI\_Turret\_Pivot}, \texttt{MI\_Turret\_Cannon}) che
impostano le texture di ciascun pezzo. Un solo shader da compilare, e il ritocco
verdastro del cannone diventa una casella da regolare nell'istanza. Se le texture
fossero state importate insieme alle mesh (Import Materials e Import Textures), Unreal
avrebbe creato da solo i materiali, uno per slot, con un grafo simile a quello della
figura~\ref{ue3:fig:turret}.

% ---------------------------------------------------------------------
\section{Le domande: il Reimport e il merge}
\label{ue3:sec:domande}

``Ora tutto chiaro il discorso dell'import e dei materiali?'' Due domande.

\paragraph{Perché il Reimport con la Z?} Una studentessa non ha capito il reimport
modificando la Z. Davide riapre il cannone. Il reimport è servito a \textbf{modificare la
mesh} alla radice. L'oggetto importato, modellato da qualcun altro, era fatto in un
certo modo, e non è un problema in sé: lo si può abbassare e girare a mano, in scena.
``Però poi, tutte le volte che io prendo e rimetto in scena l'oggetto, l'oggetto sarà
sempre posizionato in quel modo e ruotato in quel modo.'' Il reimport invece ``modifica
proprio alla radice il modo in cui funziona l'oggetto''. E il problema principale non
era la rotazione (``so che l'angolo è 113, allora lo rimodifico''), ma il
\textbf{pivot}, ``in una posizione scomoda'': il reimport è un modo per spostarlo.

\paragraph{La Modeling Mode.} ``Non è l'unico modo.'' Davide apre la \textbf{Modeling
Mode}, cerca lo strumento del \textbf{pivot}, e mostra che lo si può mettere al centro,
a sinistra, a destra, oppure ``nel punto specifico che dico io: lo prendo e me lo
sposto''. Ma questa è ``una soluzione che hanno aggiunto con Unreal Engine 5'': prima la
Modeling Mode non c'era, e in Unreal Engine 4 reimportare era il modo comodo per farlo
``senza avere bisogno di rimodellare l'oggetto a posteriori''. ``Fai la stessa cosa,
semplicemente non hai fatto un'operazione extra.''

\paragraph{Una prova.} Davide stesso è curioso: se si reimporta di nuovo con $-135$ e
113? Il reimport \textbf{sovrascrive} quello che c'era prima, ma il cannone in scena
conserva la rotazione che gli era stata data a mano. Davide lo rimette dove stava (``io
sto facendo a cazzotti da tutta la lezione con la finestra di Zoom, se non si fosse
capito'').

\paragraph{Che cosa succede ai materiali con il merge?} Era la domanda arrivata in chat
durante la costruzione dei materiali: se faccio il merge e i pezzi della torretta hanno
texture diverse? Davide prova: seleziona i tre oggetti, \textbf{Actor}, \textbf{Merge
Actors}, nome \emph{Turret}, salva. Il risultato è un oggetto con \textbf{tre slot} per i
materiali: lo slot 0 è tutta la parte del cannone, lo slot 1 il pivot ``e tutta questa
zona qui'', e così via. Se in uno slot si mette un materiale vuoto, ``diventa tutto
nero'': ``vedete, tutta questa zona adesso viene influenzata da quel materiale lì''.
Risposta: con il merge l'actor viene fuso, e i materiali vengono \textbf{riposizionati
negli slot} della nuova mesh.

% ---------------------------------------------------------------------
\section{Lo stretch, lo z-fighting e i materiali World Aligned}
\label{ue3:sec:world}

``Mi è venuto in mente un'ultima cosa che vi potrei far vedere per quanto riguarda i
materiali, e poi vi vorrei far fare una cosina.'' Manca una mezz'oretta, ``visto che
Igor vi ha fatto fare delle cose, ho visto che siete molto partecipi. Quindi vorrei
fomentare questa cosa.''

\paragraph{M\_Test.} Davide ricrea il materiale della corteccia, cancellato per sbaglio:
\textbf{M\_Test}, con due Texture Sample (Bark e la sua Normal) e un Texture
Coordinate, ``proprio la roba più facile del mondo'', con il tiling a 10.

\paragraph{Ingrandire e stirare.} ``Ora vediamo se abbiamo capito come funziona tutto
quanto.'' I materiali sono collegati alle \textbf{UV} di un oggetto. Quindi se
l'oggetto si \textbf{ingrandisce}, al materiale ``non dovrebbe succedere granché'': si
ingrandisce anche lui (``la risoluzione non ne giova particolarmente''). E se l'oggetto
si \textbf{stira} (\emph{stretch}), tirandolo verso destra? ``Credo si stira anche
lui.'' Esatto: il materiale si allunga. ``In alcuni casi non è neanche particolarmente
brutto'', ma la corteccia ``non sembra più corteccia''; e tirando di più, ``così è
brutto''.

\paragraph{Il pavimento grande.} Se si vuole fare un pavimento per una zona molto
ampia, dover pensare a questo ``può essere una noia'': per non stirare il materiale si
dovrebbero fare tante copie dello stesso pezzo, una attaccata all'altra, ``che
tendenzialmente non è quello che voglio'', perché è scomodo.

\paragraph{Una digressione: lo z-fighting.} ``Visto che è una domanda che arriverà,
anticipo'': se due piani si trovano \textbf{nella stessa posizione}, ``clippano'': la
superficie sfarfalla a chiazze, perché l'engine ``non sa se utilizzare il materiale
dell'oggetto a destra o dell'oggetto a sinistra''. Le soluzioni: separarli davvero,
oppure spostarne uno di pochissimo in \textbf{Z}, per esempio di 0,1: ``è
impercettibile, non si nota, però uno è leggermente più sopra rispetto all'altro'', e il
problema sparisce.

\paragraph{Allineato al mondo.} ``C'è un modo per far sì che il mio materiale non si
stiri quando stiro l'oggetto? Chiaramente sì, perché se no non avrei fatto questa
domanda.'' Bisogna fare un materiale che non è allineato con le coordinate
dell'oggetto, ma con quelle del \textbf{mondo}. Davide lo spiega con il \emph{gizmo}
di spostamento: un oggetto ruotato si può muovere lungo gli assi del mondo o lungo i
suoi assi locali, e il pulsante in alto passa dall'uno all'altro. ``Più o meno il
discorso è lo stesso'': nel caso del materiale non si parla di assi, ma di UV.

\paragraph{I nodi World Aligned.} Davide copia il materiale (\textbf{M\_Test1}) e usa
due nodi, \textbf{World Aligned Texture} e \textbf{World Aligned Normal}. Non hanno
scorciatoie: si cercano per nome. Hanno ``un sacco di parametri'', ma per ora ne servono
due: \textbf{Texture Object} e \textbf{Texture Size}. World Aligned Texture va nel
\textbf{Base Color}, World Aligned Normal nella \textbf{Normal}, usando l'uscita
\textbf{XYZ}, ``perché almeno prendo tutti gli assi'': il materiale è allineato al
mondo su tutti e tre.

\paragraph{Texture Size e Texture Object.} Il \textbf{Texture Size}, un vettore a tre
valori, ``fa la stessa cosa del Texture Coordinate'': regola la grandezza della
ripetizione. Perché allora non tenere il Texture Coordinate sulle UV del Texture
Sample? Perché i nodi World Aligned \textbf{non prendono} un Texture Sample: vogliono un
\textbf{Texture Object}. Si converte: tasto destro sul Texture Sample, \textbf{Convert to
Texture Object}. Il Texture Object ``è un oggetto molto più semplice'': non ha
ingressi, ha una sola uscita, niente RGB; ``è molto più limitato in quello che può
fare'', ma è il ``messaggio'' che serve ai nodi World Aligned.

\paragraph{La prova.} Davide prova diversi valori di Texture Size (1, 2, 10) per avvicinarsi
al risultato dell'altro materiale, con qualche esitazione su quale sia ``troppo'' e quale
``troppo poco'' (``è proprio piccola la texture, allora''), e applica il nuovo
materiale a un secondo oggetto. Poi \textbf{stira} entrambi: il primo si stira; il
secondo ``semplicemente genera nuove texture'', rimane orientato con il mondo,
``rimane sempre della stessa proporzione''.

\paragraph{Dov'è la fregatura.} ``Se questo rimane sempre proporzionato anche quando
stiriamo le cose, perché non dovremmo avere tutti i materiali World Aligned? Dov'è la
fregatura?'' Eccola: essendo allineato al mondo, se l'oggetto si \textbf{muove} o
\textbf{ruota}, il materiale \textbf{resta fermo}. Una \textbf{porta} con un materiale
normale si apre e il materiale resta attaccato alla porta; una porta con un materiale
World Aligned si apre e ``il materiale sembra scorrere sulla porta'', ``come se il
materiale fosse animato''. ``Ci siamo? Ho rotto le palle? Datemi un feedback.''

\begin{figure}[htbp]
\centering
\begin{tikzpicture}[font=\small]
  % materiale UV: stiramento
  \node[font=\footnotesize\bfseries] at (2.6,2.6) {materiale normale (UV)};
  \draw[thick, fill=brown!20] (0,1.0) rectangle (1.2,2.2);
  \foreach \i in {0.2,0.6,1.0} \draw[brown!70!black, thick] (\i,1.15) -- (\i,2.05);
  \draw[-{Stealth[length=2mm]}, thick] (1.5,1.6) -- (2.1,1.6);
  \draw[thick, fill=brown!20] (2.4,1.0) rectangle (5.4,2.2);
  \foreach \i in {0.5,1.5,2.5} \draw[brown!70!black, thick] (2.4+\i,1.15) --
    (2.4+\i,2.05);
  \node[font=\footnotesize, align=center] at (2.7,0.55) {stirato: le venature si
    allargano};
  % materiale world aligned
  \begin{scope}[shift={(6.6,0)}]
    \node[font=\footnotesize\bfseries] at (2.6,2.6) {materiale World Aligned};
    \draw[thick, fill=brown!20] (0,1.0) rectangle (1.2,2.2);
    \foreach \i in {0.2,0.6,1.0} \draw[brown!70!black, thick] (\i,1.15) -- (\i,2.05);
    \draw[-{Stealth[length=2mm]}, thick] (1.5,1.6) -- (2.1,1.6);
    \draw[thick, fill=brown!20] (2.4,1.0) rectangle (5.4,2.2);
    \foreach \i in {0.2,0.6,...,3.0} \draw[brown!70!black, thick] (2.4+\i,1.15) --
      (2.4+\i,2.05);
    \node[font=\footnotesize, align=center] at (2.7,0.55) {stirato: ``genera nuove
      texture'', stessa scala};
  \end{scope}
  % la porta
  \node[font=\footnotesize, align=center, text width=12.5cm] at (5.9,-0.3)
    {la fregatura: se l'oggetto si muove o ruota (una porta che si apre), il materiale
    World Aligned resta fermo nel mondo e ``sembra scorrere'' sull'oggetto};
\end{tikzpicture}
\caption{Materiale legato alle UV e materiale World Aligned a confronto. Il primo
segue l'oggetto e si deforma con lui; il secondo non si deforma, ma non segue
l'oggetto quando si muove. Va bene per pavimenti, muri e rocce statici; non per porte,
veicoli o la torretta.}
\label{ue3:fig:world}
\end{figure}

\paragraph{Le linee nere.} Una studentessa vede delle \textbf{linee nere} al confine tra
una ripetizione della texture e l'altra, sul pavimento. Davide guarda il suo grafo:
stessa texture, Bark, ma ha collegato l'uscita \textbf{XY} invece di \textbf{XYZ}. ``In
quel modo lui aveva due assi che erano allineati col mondo, e l'altro no'', e il
risultato era diverso. ``Ok, grazie.''

\paragraph{Approfondimento: come funziona un World Aligned.} Il nodo World Aligned
Texture, una Material Function del motore, non usa le UV della mesh: per ogni pixel
prende la sua \textbf{posizione nel mondo} e la usa come coordinata della texture. Lo fa
tre volte, proiettando la texture lungo i tre assi (dall'alto, di lato, di fronte), e
mescola le tre proiezioni secondo la direzione della superficie: è la tecnica detta
\textbf{triplanar mapping}. L'uscita XY usa una sola proiezione (dall'alto, buona per i
pavimenti), l'uscita XYZ le mescola tutte; per questo, su un oggetto con facce in tutte
le direzioni, con XY le pareti laterali ricevono una texture ``stirata'' a strisce, che
può dare proprio le linee scure viste dalla studentessa. Il \textbf{Texture Size} è
espresso in \textbf{unità del mondo}, cioè in centimetri: è la grandezza di una
ripetizione della texture. Funziona quindi al contrario del tiling: un tiling più alto
fa ripetere di più la texture, un Texture Size più alto la fa ripetere di meno (una
texture di 100 si ripete ogni metro, una di 10 ogni dieci centimetri). È la ragione
delle esitazioni di Davide su quale valore fosse ``troppo''. Il costo: tre letture della
texture invece di una, quindi un materiale un po' più pesante.

% ---------------------------------------------------------------------
\section{L'esercizio della lava, il nodo Noise e la chiusura}
\label{ue3:sec:lava}

``Allora, ci siamo visti un po' di cose, fatto un miscuglione di cose per quanto riguarda
materiali, texture e import delle mesh.'' Nessuna domanda: ``sembra tutto a posto''.

\paragraph{La consegna.} Per esercitarsi con i materiali, e magari scoprire qualche nodo
non usato a lezione (il nodo \textbf{Noise}, qualche \textbf{maschera}), Davide propone
una prova di una decina di minuti, da continuare poi a casa come le cose di Igor:
\textbf{un materiale di lava}. Di che componenti avrà bisogno? Deve \textbf{muoversi}
un po', quindi bisognerà gestire la velocità del \textbf{Panner}, ``a meno che non vi
cercate proprio la texture della lava''. Ma la richiesta è di improvvisare: ``provate a
vedere cosa riuscite a combinare andando a fare delle operazioni fra colori e fra altre
texture'', con un po' di movimento, un \textbf{Base Color}, ``ma anche un
\textbf{Emissive}, perché la lava comunque deve emettere luce''.

\paragraph{Il ripasso dei nodi.} Una studentessa non ricorda il nome della funzione
``che abbiamo usato per collegare rosso e blu'': il Lerp. Davide ricondivide lo schermo
e fa il riepilogo dei nodi della giornata: la \textbf{costante}, il \textbf{colore}, la
\textbf{texture}, il \textbf{nodo risultante}, il \textbf{Texture Coordinate}, i
\textbf{World Aligned}, le \textbf{operazioni} (con il consiglio sulla moltiplicazione
visto nella sezione~\ref{ue3:sec:costanti}), il \textbf{Panner}, il \textbf{Lerp}.
Divisione e sottrazione servono soprattutto per gli effetti speciali, e Davide vorrebbe
fare una lezione su \textbf{Niagara}, ``il tool che usa Unreal per fare gli effetti
speciali'', se ci sarà tempo.

\paragraph{Le prossime lezioni con Davide.} Davide annuncia come saranno le sue
lezioni finali: ci sarà ``una lezione che sarà un'esercitazione intera'' in cui si
proverà a fare dei \textbf{livelli}, e chiederà agli studenti di condividere lo schermo
e di \emph{streamare} a turno (``se per voi è un problema, magari lo facciamo solo con le
persone per cui non è un problema'', e lo dirà in anticipo). Per esempio, ``un livello
\textbf{platform}'', con le interazioni che si faranno nelle lezioni di Igor e quello che
si vedrà insieme per il \emph{level design}.

\paragraph{Il nodo Noise.} Un nodo ``un po' più complicato'', per chi vuole provare: il
\textbf{Noise}, che ``serve per simulare un \emph{noise}'': il \emph{Voronoi}, il
\emph{White Noise}, un'alternanza di bianco e nero ``che si usa molto spesso per fare
delle distorsioni''. Ha molti parametri: la scala del \emph{noise}, e la possibilità di
filtrarlo con dei vettori o con delle mappe. ``Potete smanettarci un pochino.''

\paragraph{Lo stream della studentessa.} ``C'è qualcuno che, nel mentre sta provando a
fare qualcosa, vuole streamare in questi ultimi cinque minuti?'' Una studentessa ha
provato, ma non riesce a trasformare la texture in World Aligned: è perché usa un
Texture Sample, e serviva un Texture Object. ``Prova a non usare il World Aligned, e
farlo così, diciamo, semplice.'' Senza World Aligned le viene ``questa cosa'', sempre
con il Bark. Davide la guida:
\begin{itemize}
  \item nella lava servono ``delle zone più scure e delle zone più luminose'': quindi
        un \textbf{Lerp} tra un \textbf{colore} e un \textbf{nero}, gestito da un'alfa
        che sia una texture con dei bianchi e dei neri ``posizionati come vuoi'', da
        mandare nell'\textbf{Emissive Color};
  \item l'alfa va cercata tra le texture del progetto, ``qualcosa per dare un effetto
        tipo di distorsione''; le \textbf{Normal}, invece, ``lasciale perdere'': per
        un'alfa si usano le Diffuse o le mappe come Specular e Roughness;
  \item se una texture fa da alfa, nelle entrate del Lerp ci vogliono \textbf{due
        colori}, ``uno che sia molto scuro, uno che sia più chiaro''; ricordando che
        l'alfa ``prende i bianchi e i neri e dice: nei bianchi vai più verso B, nei neri
        vai più verso A'';
  \item il risultato va ``forse più nell'Emissive che non nel Base Color'', perché la
        lava la si vuole ``infuocata, accesa, luccicante''; il Base Color si può anche
        non usare, o usarlo con una texture adatta (``il Bark, per come è fatto, forse
        non è proprio la texture più adatta per fare una lava'');
  \item nell'Emissive si può \textbf{moltiplicare} il risultato, in modo che le zone
        accese siano molto luminose.
\end{itemize}
La studentessa nota che il nero, nell'Emissive, quasi non si vede: ``esatto'', il nero
non emette niente. E usandolo come colore ``non si vede molto'', perché il Bark come alfa
``è molto bianco'', e finisce per scurire un po' il tutto. ``È comunque una cosa
complicata, quindi vedete effettivamente se riuscite a tirar fuori qualcosa.''

\begin{figure}[htbp]
\centering
\begin{tikzpicture}[
  font=\small,
  ue3node/.style={draw, thick, rounded corners=4pt, align=center,
                  minimum height=0.85cm},
  ue3uv/.style={-{Stealth[length=2mm]}, thick, red!60!black},
  ue3flt/.style={-{Stealth[length=2mm]}, thick, green!50!black},
  ue3col/.style={-{Stealth[length=2mm]}, thick, blue!60!black}
]
  \node[ue3node, fill=red!12, text width=1.7cm] (tc) at (0,0) {\textbf{TexCoord}};
  \node[ue3node, fill=gray!10, text width=1.9cm] (pan) at (2.6,0) {\textbf{Panner}\\
    \footnotesize Speed 0,05; 0,02};
  \node[ue3node, fill=gray!15, text width=2.3cm] (msk) at (5.6,0)
    {\textbf{Texture Sample}\\\footnotesize maschera in b/n\\\footnotesize (o un
    Noise)};
  \node[ue3node, fill=black!70, text=white, text width=1.9cm] (dark) at (5.6,2.0)
    {\textbf{scuro}\\\footnotesize quasi nero};
  \node[ue3node, fill=orange!60, text width=1.9cm] (hot) at (5.6,-2.0)
    {\textbf{chiaro}\\\footnotesize arancione};
  \node[ue3node, fill=gray!10, text width=1.3cm, minimum height=2.0cm] (lerp)
    at (8.6,0) {\textbf{Lerp}\\[3pt]\footnotesize A\\\footnotesize B\\
    \footnotesize Alpha};
  \node[ue3node, fill=green!8, text width=1.6cm] (k) at (8.6,-2.4)
    {\footnotesize intensità\\\footnotesize (parametro)};
  \node[ue3node, fill=gray!10, text width=1.4cm] (mul) at (10.9,-0.6)
    {\textbf{Multiply}};
  \node[ue3node, fill=blue!6, text width=2.0cm, minimum height=1.6cm] (res)
    at (13.4,0) {\footnotesize Base Color\\[6pt]\footnotesize\textbf{Emissive}};
  \draw[ue3uv] (tc) -- (pan);
  \draw[ue3uv] (pan) -- node[above, font=\scriptsize]{UVs} (msk);
  \draw[ue3col] (dark.east) -- ++(0.7,0) |- ([yshift=5mm]lerp.west);
  \draw[ue3col] (hot.east) -- ++(0.9,0) |- ([yshift=0mm]lerp.west);
  \draw[ue3flt] (msk.east) -- ++(0.35,0) |- ([yshift=-5mm]lerp.west);
  \draw[ue3col] (lerp.east) -- ++(0.3,0) |- ([yshift=2mm]mul.west);
  \draw[ue3flt] (k.east) -- ++(0.5,0) |- ([yshift=-2mm]mul.west);
  \draw[ue3col] (mul.east) -- ++(0.4,0) |- ([yshift=-3mm]res.west);
  \draw[ue3col, dashed] (lerp.north) |- ([yshift=4mm]res.west);
\end{tikzpicture}
\caption{Una possibile soluzione dell'esercizio della lava, costruita con i soli nodi
della lezione e con i consigli dati alla studentessa: una maschera in bianco e nero che
scorre con il Panner fa da alfa a un Lerp tra uno scuro e un arancione; il risultato,
moltiplicato, va nell'Emissive (e, tratteggiato, eventualmente anche nel Base Color).
Non è un grafo mostrato a lezione: i valori sono indicativi.}
\label{ue3:fig:lava}
\end{figure}

\paragraph{La chiusura.} ``Vabbè, io in realtà adesso ho gettato questa bomba, e penso
che vi lascerò fare. Prendetevi tutto il tempo che volete, nessuna pressione.'' Davide
rivedrà la classe ``a luglio, probabilmente'', per le ultime quattro o cinque lezioni del
corso, sul \textbf{level design} e sulla \textbf{creazione di livelli}. ``Si sono fatte
le 7.'' ``Direi che è stato un piacere, mi sembrate tutti abbastanza interessati,
sicuramente belli, belli, belli attivi.'' E forse ci si rivedrà prima, ``se Igor ha
bisogno di un'altra sostituzione''. ``Ciao a tutti, arrivederci.''

\paragraph{Approfondimento: una lava più convincente.} Partendo dallo schema della
figura~\ref{ue3:fig:lava}, tre accorgimenti, tutti con nodi già visti, rendono la lava
molto più viva. Primo: \textbf{due Panner} con velocità e tiling diversi sulla stessa
maschera, sommati o moltiplicati, così che il movimento non si ripeta mai uguale (lo
stesso trucco dell'acqua). Secondo: un Multiply sull'Emissive con valori ben oltre 1
(10, 20), perché solo così il \emph{bloom} fa ``bruciare'' le zone calde, come nel
materiale con il Glow Intensity. Terzo: usare la \textbf{stessa} maschera per la
\textbf{Roughness}, con un Lerp tra 0,3 (la lava liquida, lucida) e 1 (la crosta
raffreddata, opaca): il materiale cambia così non solo colore ma anche superficie, che è
il senso del PBR. Il nodo Noise di Davide può sostituire la texture della maschera, ma è
calcolato a ogni pixel e costa più di una texture: per un oggetto grande in scena, una
texture di \emph{noise} è di solito la scelta più leggera.

\begin{table}[htbp]
\centering
\small
\begin{tabular}{@{}p{4.4cm}p{8.6cm}@{}}
\toprule
\textbf{Comando} & \textbf{Che cosa fa} \\
\midrule
Tasto destro nel grafo & menu dei nodi, con ricerca \\
T + clic & Texture Sample \\
1 + clic / 2 + clic / 3 + clic & Constant, Constant2Vector, Constant3Vector (il
  colore) \\
A / M / D + clic & Add, Multiply, Divide \\
S + clic & Scalar Parameter (per questo Subtract non ha scorciatoia) \\
U + clic & Texture Coordinate (``lavora sulle UV'') \\
L + clic & Lerp (Linear Interpolate) \\
P + clic & Panner \\
Alt + clic su un pin & stacca i collegamenti \\
Tasto destro su un nodo, Convert to Parameter & trasforma il nodo in un parametro,
  regolabile dalle Material Instance \\
Tasto destro su un Texture Sample, Convert to Texture Object & per i nodi World Aligned \\
Use Selected Asset from Content Browser & mette nel nodo la texture selezionata nel
  Content Browser \\
Tasto destro su un materiale, Create Material Instance & crea un'istanza \\
Reimport Mesh With Dialog & reimporta una mesh riaprendo la finestra di import \\
Ctrl+G / Shift+G & raggruppa e separa gli actor selezionati \\
Actor, Merge Actors & fonde gli actor selezionati in una nuova mesh \\
Ctrl+C, Ctrl+V, Ctrl+Z, Ctrl+Y, Ctrl+A & le scorciatoie classiche, valide quasi
  ovunque \\
Ctrl+S & salvare: ``salvate spesso'' \\
\bottomrule
\end{tabular}
\caption{I comandi della terza lezione nell'editor dei materiali, nel Content Browser e
nel Level Editor di Unreal Engine 5. Le scorciatoie con il clic si usano tenendo premuto
il tasto e cliccando nel grafo; World Aligned Texture e World Aligned Normal non ne
hanno.}
\label{ue3:tab:scorciatoie}
\end{table}

% ---------------------------------------------------------------------
\section*{Punti chiave del capitolo}
\addcontentsline{toc}{section}{Punti chiave del capitolo}

\begin{itemize}
  \item Il grafico lavora soprattutto fuori da Unreal (Blender, Maya); dentro Unreal
        \textbf{importa} gli oggetti e costruisce gli \textbf{shader}.
  \item Una \textbf{texture} è la ``carta regalo'' dell'oggetto: si avvolge sulla mesh
        attraverso le \textbf{UV}, che si creano prima, nel software di modellazione.
        Una texture non pensata per quelle UV si stira.
  \item Le mappe principali: \textbf{Diffuse/Albedo} (il colore, suffisso \_D o \_A),
        \textbf{Normal} (blu: l'illusione del rilievo senza poligoni in più),
        \textbf{Specular}, \textbf{Roughness}, \textbf{Metallic}, la mappa di
        \textbf{altezza} (\emph{displacement} vero, ``una normal sotto steroidi'').
  \item Nelle mappe in bianco e nero il \textbf{bianco è 1}, il \textbf{nero è 0}, il
        grigio sta in mezzo. Ogni canale RGB è una mappa a sé: una texture
        \textbf{AORM}/ORM impacchetta tre mappe in un file.
  \item Il \textbf{materiale} è la ``vernice'', uno \textbf{shader}: si costruisce con
        nodi, come un Blueprint, e combina le texture. Sugli oggetti si applicano
        materiali, non texture.
  \item \textbf{Material Domain} (per che cosa), \textbf{Blend Mode} (Opaque, Masked,
        Translucent) e \textbf{Shading Model} (Default Lit, Unlit) cambiano i pin del
        nodo risultante.
  \item Nodi di base: \textbf{Texture Sample} (T), costanti (1, 2, 3), \textbf{Add},
        \textbf{Multiply}, \textbf{Divide} (A, M, D), Scalar Parameter (S),
        \textbf{Texture Coordinate} (U), \textbf{Lerp} (L), \textbf{Panner} (P). La
        moltiplicazione è l'operazione più usata.
  \item Il rendering è \textbf{PBR}: Roughness 1 è la roccia, 0 lo specchio; Metallic
        e Specular decidono come la superficie riflette; l'\textbf{Emissive}
        moltiplicato fa brillare il materiale.
  \item \textbf{UTiling} e \textbf{VTiling} (o un Multiply sulle coordinate) decidono
        quante volte si ripete la texture: 2 e 2 sono quattro ripetizioni, 0,5 e 0,5 un
        quarto della texture.
  \item Un nodo diventa un \textbf{parametro} con Convert to Parameter. Le
        \textbf{Material Instance} (MI\_) cambiano i parametri senza toccare il grafo
        del \textbf{master material} (M\_), si aggiornano subito e sono più leggere:
        sugli oggetti si applicano sempre le istanze.
  \item Le \textbf{Material Function} sono ``parentesi'' di calcoli riusabili tra
        materiali e tra membri del team.
  \item Il \textbf{Lerp} interpola tra A e B secondo un'\textbf{alfa} da 0 a 1, che può
        essere una texture (``l'arbitro''); il \textbf{Panner} fa scorrere la texture
        (acqua, fluidi, lava).
  \item Importare \textbf{FBX}, controllare le \textbf{unità} (Unreal usa i
        centimetri), texture in PNG o JPEG. La finestra di import genera collisioni,
        corregge il transform, importa materiali e texture.
  \item Una mesh è fatta di \textbf{vertici}, \textbf{spigoli} e \textbf{facce}
        (triangoli). Il \textbf{pivot} deve stare dove il pezzo ruota: con
        \textbf{Reimport Mesh With Dialog} (cannone: Z $-135$, rotazione 113) o con la
        \textbf{Modeling Mode} di Unreal Engine 5.
  \item Le \textbf{simple collision} (box, sfera, capsula: il ``capsulone'' dei
        \emph{character}) sono leggere; le \textbf{complex} usano i triangoli della mesh
        e costano molto di più.
  \item \textbf{Ctrl+G} raggruppa; \textbf{Merge Actors} fonde in una nuova mesh, con
        tanti slot quanti i materiali dei pezzi.
  \item Due piani nella stessa posizione sfarfallano (\textbf{z-fighting}): basta
        alzarne uno di 0,1. I materiali \textbf{World Aligned} (Texture Object, Texture
        Size, uscita XYZ) non si stirano, ma restano fermi se l'oggetto si muove.
  \item \textbf{Salvate spesso.} E per esercizio: un materiale di \textbf{lava}, con
        Panner, Lerp ed Emissive.
\end{itemize}

%
%  Tutte le etichette di questo capitolo hanno il prefisso "ue3:"
%  (ue3:cap:..., ue3:sec:..., ue3:fig:..., ue3:tab:...), e tutti gli
%  stili TikZ hanno il prefisso "ue3" (ue3node, ue3arr, ...) per non
%  collidere con chiavi native di TikZ né con quelli delle prime due
%  lezioni. Il capitolo è autonomo: non contiene \ref verso altri
%  capitoli (i rimandi alle lezioni precedenti sono a parole).
%  Le figure sono ricostruzioni schematiche di quanto mostrato e detto a
%  lezione, non copie delle slide o delle schermate del docente. Dove la
%  trascrizione è troppo rovinata per ricostruire un dettaglio, il
%  capitolo lo dice invece di inventarlo.
% =====================================================================

\chapter[Texture, materiali e mesh]{Texture, materiali e mesh: la grafica dentro
Unreal e la torretta}
\label{ue3:cap:materiali}

Questa è la \textbf{terza lezione} del corso breve di \textbf{introduzione a Unreal
Engine}, e cambia sia docente sia argomento. Al posto di Igor, che ``aveva un
impiccio'', c'è \textbf{Davide}, che si presenta così: ``io sono Davide e in questa
lezione, nello specifico, sostituisco Igor''. Davide tornerà più avanti: ``le ultime
lezioni del corso le fate con me'', e saranno lezioni ``un po' più basate su
\textbf{level design} e la creazione di mappe'', con qualche esercitazione insieme. Dopo
due lezioni sull'editor e sulla logica (piazzare oggetti, i \emph{transform}, i
Blueprint, la pallina sulla piattaforma), questa lezione guarda a ciò che ``fa il
grafico all'interno di Unreal'': le \textbf{mesh}, le \textbf{texture}, i
\textbf{materiali}, e ``qualche nodo, qualche cosa carina che potete fare'' combinando
mesh e materiali.

La lezione ha due metà separate da una pausa. La prima è sui \textbf{materiali}: che
cos'è una texture e come si avvolge su un oggetto attraverso le \textbf{UV}; quali tipi
di mappa esistono (\textbf{Diffuse}, \textbf{Normal}, \textbf{Specular}, la mappa di
altezza per il \emph{displacement}) e perché in queste mappe il bianco vale 1 e il nero
0; che cos'è un \textbf{materiale} e come si costruisce nell'\textbf{editor dei
materiali}, ``nello stesso modo dell'editor dei Blueprint''; le impostazioni principali
(\textbf{Material Domain}, \textbf{Blend Mode}, \textbf{Shading Model}); i nodi di base
(\textbf{Texture Sample}, le costanti, le operazioni, il \textbf{Texture Coordinate},
il \textbf{Lerp}, il \textbf{Panner}); il rendering \textbf{PBR} con Metallic,
Specular e Roughness; e le \textbf{Material Instance}, il modo ``bon ton'' di applicare i
materiali. La seconda metà è sulle \textbf{mesh}: si importa una \textbf{torretta} in
tre pezzi dal materiale didattico, si guarda la finestra di \textbf{import}, si
corregge il \textbf{pivot} del cannone con un \textbf{Reimport}, si esplorano lo Static
Mesh Editor e le \textbf{collisioni}, si fondono i pezzi con \textbf{Merge Actors}, si
creano a mano i materiali della torretta, e si chiude con i materiali \textbf{World
Aligned}, che non si deformano quando si stira l'oggetto. In fondo, un esercizio libero:
un materiale di \textbf{lava}.

La lezione è \textbf{online}, su Zoom, come le precedenti, e finisce poco prima delle
19 (``si sono fatte le 7''). La pausa, una decina di minuti, arriva dopo ``un'oretta e
mezza''. Davide, alla sua prima volta con questa classe, combatte per tutta la lezione
con la condivisione dello schermo (``io sto facendo a cazzotti da tutta la lezione con
la finestra di Zoom'', ``quanto non sono abituato a streamare'') e con un PC di
emergenza: ``il fisso mi è esploso letteralmente tre giorni fa'', e a metà della seconda
parte il computer ``inizia ad accusare la stream'', tanto che Davide gira il
ventilatore verso di lui. Chiede spesso conferma che gli studenti vedano quello che
mostra: ``mi impanicherò ogni tanto a caso e vi chiederò conferma delle robe''.

Tre avvertenze di lettura, le stesse delle prime due lezioni. La prima: la trascrizione
è automatica e storpia moltissimi nomi. Unreal diventa ``ArrayLink'', ``ArrayL4'',
``Arial Engine'' e ``a Real''; la \textbf{Material Instance} è quasi sempre
``materialistance'', ``Material Distance'' e una volta ``mie dire distanze''; il
\textbf{Panner} è ``panel'', ``planner'' e ``banner''; il \textbf{Lerp} è ``l'ERP'';
\emph{World Aligned} è ``word aligned'' e ``word.allocation''; la texture
\textbf{T\_Bark} (una corteccia) è ``Arc'', ``barca'', ``Barc'' e ``park''; ``strecciare''
e ``si straccia'' stanno per lo \emph{stretch}, lo stiramento; ``traccare'' sta per
\emph{tracciare} un bersaglio (\emph{to track}); la \textbf{mappa di altezza}
(\emph{height map}) è quasi sempre ``light map'' o ``lightmap''. Dove la ricostruzione
è sicura il capitolo usa la forma corretta e, la prima volta, lo segnala in nota; dove
non lo è riporta la forma della trascrizione e lo dice. La seconda: durante la pausa e
verso la fine il riconoscimento vocale produce lunghi blocchi di testo senza senso e
ripetuto (``qual è la risposta che ci può essere?'', discorsi sulla ``vita data per
scontata'' e su ``vittime''): sono artefatti, non parole di nessuno, e il capitolo li
ignora. La terza: gran parte della lezione è una dimostrazione dal vivo nell'editor dei
materiali e nello Static Mesh Editor. Il capitolo ricostruisce i grafi di nodi da ciò
che Davide dice mentre li costruisce, e li completa, dove serve, con i nomi esatti dei
nodi e dei comandi di Unreal Engine 5.

\paragraph{Come leggere questo capitolo.} Il capitolo segue l'ordine della lezione. Le
figure ricostruiscono gli schemi spiegati a voce, sulle slide o nell'editor: le mappe
della torretta, i canali di una texture, il rapporto tra texture, materiale e mesh, i
grafi dei materiali costruiti a lezione, le Material Instance, la struttura di una
mesh, il pivot del cannone. Le tabelle riassumono i tipi di mappa, le impostazioni del
materiale, le opzioni di import, le collisioni e le scorciatoie dell'editor dei
materiali. Qua e là, dei paragrafi di \textbf{approfondimento} collegano la lezione al
libro \emph{Blueprints Visual Scripting for Unreal Engine 5} di Marcos Romero e Brenden
Sewell, il testo di riferimento per questa parte del libro, e a qualche nozione tecnica
generale su Unreal Engine 5; sono segnalati come tali e non fanno parte della lezione.

Il percorso del capitolo:
\begin{enumerate}
  \item l'attesa, le presentazioni e il programma; che cosa fa un grafico dentro
        Unreal; il materiale didattico (sezioni \ref{ue3:sec:attesa}
        e~\ref{ue3:sec:didattico});
  \item le \textbf{texture}: la carta regalo, le \textbf{UV}, i tipi di mappa, il
        bianco e il nero, i \textbf{canali RGB} (sezioni da~\ref{ue3:sec:texture}
        a~\ref{ue3:sec:canali});
  \item i \textbf{materiali}: la vernice, l'editor, Material Domain, Blend Mode e
        Shading Model, i nodi (sezioni da~\ref{ue3:sec:materiale}
        a~\ref{ue3:sec:nodi});
  \item in Engine: la mappa vuota, dentro una texture, il \textbf{Texture Sample}, le
        costanti e le scorciatoie, il colore e l'\textbf{emissive} (sezioni
        da~\ref{ue3:sec:mappavuota} a~\ref{ue3:sec:emissive});
  \item il \textbf{PBR} e la Normal (sezione \ref{ue3:sec:pbr});
  \item il \textbf{Texture Coordinate}, il tiling e il pavimento (sezione
        \ref{ue3:sec:tiling});
  \item le \textbf{Material Instance} e le Material Function (sezioni
        \ref{ue3:sec:instance} e~\ref{ue3:sec:function});
  \item il \textbf{Lerp} e il \textbf{Panner}, poi la pausa (sezioni
        da~\ref{ue3:sec:lerp} a~\ref{ue3:sec:pausa});
  \item l'\textbf{import} e la torretta; che cos'è una \textbf{mesh}; dimensioni e
        formati (sezioni da~\ref{ue3:sec:import} a~\ref{ue3:sec:regole});
  \item il \textbf{pivot} del cannone e il \textbf{Reimport} (sezione
        \ref{ue3:sec:pivot});
  \item lo \textbf{Static Mesh Editor} e le \textbf{collisioni}; raggruppare e
        \textbf{Merge Actors} (sezioni \ref{ue3:sec:sme} e~\ref{ue3:sec:merge});
  \item i \textbf{materiali della torretta} e le domande sul Reimport e sul merge
        (sezioni \ref{ue3:sec:matturret} e~\ref{ue3:sec:domande});
  \item lo stretch, lo \textbf{z-fighting} e i materiali \textbf{World Aligned}
        (sezione \ref{ue3:sec:world});
  \item l'esercizio della \textbf{lava}, il nodo \textbf{Noise} e la chiusura
        (sezione \ref{ue3:sec:lava}).
\end{enumerate}

% ---------------------------------------------------------------------
\section{L'attesa, le presentazioni e il programma}
\label{ue3:sec:attesa}

``Ok, buonasera.'' La registrazione parte subito (``sta registrando? Sì, sta
registrando, perfetto''), e Davide aspetta ``qualche minutino'' che entrino tutti. Uno
studente avvisa che dovrà staccarsi venti minuti prima per andare a lavorare: ``nessun
problema'', tanto la lezione resta registrata, e anzi a Davide piacerebbe finire ``con
voi che fate qualcosina'', quindi andare via un po' in anticipo ``non è neanche un
problema''. Davide chiede quante persone c'erano alle lezioni precedenti (``una dozzina,
credo''), ringrazia chi accende la telecamera (``almeno so che non parlo da solo'') e
controlla che non si senta il ventilatore in sottofondo (``inizia a far caldo, inizia a
essere complicato'').

\paragraph{Che cosa avete fatto con Igor?} Prima di cominciare, Davide vuole sapere a
che punto è la classe, per non ripetere cose già viste. Dal Discord ha capito che si è
``ancora all'inizio'': le schede, le finestre, i \emph{transform} sugli oggetti, i
tavolini, ``qualcosa con gli input, con la pallina che si muove''. ``Se vediamo
qualcosa che avete già visto, ditemelo.''

\paragraph{Il programma.} La lezione di oggi guarda a quello che fa ``il grafico
all'interno di Unreal'': ``un po' di nozioni base per quanto riguarda le mesh, un po' di
nozioni base per quanto riguarda i materiali, le texture'', e ``già qualche nodo,
qualche cosa carina che potete fare'' usando mesh e materiali.

\paragraph{Chi c'è.} Davide chiede a che cosa sono iscritti gli studenti, ``per capire
chi ho davanti''. Le risposte: \textbf{game design} (più d'uno), \textbf{3D art}
(tre persone), \textbf{character animation} (due), \textbf{narrative design}. ``Beh, in
realtà siamo abbastanza assortiti.''

\paragraph{Che cosa fa un grafico dentro Unreal.} Ai futuri grafici 3D Davide mette
``le mani avanti'': la maggior parte del lavoro di un grafico \textbf{non} si svolge
dentro Unreal, ma in altri strumenti, che verranno spiegati nei rispettivi corsi.
Quello che si vede oggi è ciò che un grafico fa ``nel momento in cui va a lavorare con
altre persone'' dentro Unreal Engine: \textbf{importare} gli oggetti e creare gli
\textbf{shader} dentro l'engine.

\paragraph{Come si struttura la lezione.} Davide chiede come lavorava Igor: tutta
teoria e poi tutta pratica, o un po' e un po'? ``Andava prima con la presentazione e poi
con Unreal.'' Davide prova a replicare: ``ci facciamo prima una chiacchiera sulla parte
un po' più teorica e poi passiamo al progetto''. Nella pratica, come si vedrà, teoria
ed engine si alterneranno di continuo.

% ---------------------------------------------------------------------
\section{Il materiale didattico}
\label{ue3:sec:didattico}

Sul Discord c'è stato ``un po' di casino'' per il \textbf{materiale didattico}. ``Non è
molta roba'': è un piccolo progetto di Unreal, ``un progetto che Unreal usa per
spiegare alcune cose'', e probabilmente quello che Igor farà usare nelle lezioni
successive ogni volta che si implementa qualcosa. Gli altri file della cartella sono
\textbf{asset} da importare nel progetto poco alla volta: oggi la \textbf{torretta};
quando si farà l'audio, l'audio; ``poi vediamo se aggiungere altre cose''.

Dove si trova, Davide lo spiega più avanti a una studentessa: sul Discord, ``il 19
maggio, se non ricordo male'', c'è un link a una pagina di Drive con un file che si
chiama, più o meno, ``materiale didattico Unreal Engine Intro''\footnote{La trascrizione
dice ``IVE, materiale didattico, a Real Engine Intro'' e, poco prima, ``Hive Unreal
Engine 5 intro'': la parola iniziale non è ricostruibile con certezza.}. Lo si
scarica, lo si \textbf{decomprime}, e dentro c'è ``un progettino con un
\emph{character}'': sono soprattutto cose per le prossime lezioni con Igor, e per oggi
serve poco.

% ---------------------------------------------------------------------
\section{Che cos'è una texture}
\label{ue3:sec:texture}

\paragraph{La carta regalo.} ``Che cos'è una texture? La texture fondamentalmente è il
\emph{wrap}, è la \textbf{carta regalo} in cui è avvolto il nostro oggetto.'' Sulla
slide c'è la torretta del materiale didattico e, accanto, le sue texture: un'immagine
``tutta scomposta e segmentata e divisa in varie parti''. Perché divisa in parti?
Perché quelle parti poi ``si avvolgeranno sopra la mesh''. Le texture sono ``una serie
di mappe, una serie di informazioni che vanno a dare delle qualità specifiche
all'oggetto''. Con un'altra immagine: la texture ``non è altro che una
\textbf{decalcomania}'' dell'oggetto, della \emph{mesh} che si mette in scena.

\paragraph{Le UV.} Perché una texture funzioni bene su un oggetto, deve essere pensata
``con in mente le \textbf{UV} dei modelli su cui si lavora''. Le UV ``sono le
coordinate'', tutto quello che riguarda ``la struttura dell'oggetto'': la \textbf{mappa
delle UV} è la \textbf{prima parte} del processo di creazione di una texture. Dal
modello 3D si estraggono informazioni che poi servono a creare ``tutti i vari tipi di
texture che esistono'': sulla slide le texture della torretta sono diverse, e ognuna
porta informazioni diverse, ``ma sono tutte texture pensate per la stessa torretta''.

\paragraph{L'UV mapping.} L'\textbf{UV mapping} è la tecnica che serve a
\textbf{mappare le texture 2D su modelli 3D}: si assegnano le coordinate dette UV ``ai
vertici e agli \emph{edge} del modello'', e tutte le texture applicate dopo usano quelle
informazioni ``per precisione''. È un lavoro che si fa \textbf{prima} di entrare in
Unreal, in un altro software ``come Blender, come Maya''. Se si usa una texture che non
è stata pensata per le UV di quell'oggetto, ``la texture verrà strecciata''
(\emph{stretched}, stirata) o comunque non darà il risultato voluto.

\begin{figure}[htbp]
\centering
\begin{tikzpicture}[
  font=\small,
  ue3box/.style={draw, thick, align=center, minimum height=0.9cm},
  ue3arr/.style={-{Stealth[length=2.5mm]}, thick}
]
  % il cubo 3D
  \draw[thick, fill=orange!25] (0,0) -- (1.6,0) -- (1.6,1.6) -- (0,1.6) -- cycle;
  \draw[thick, fill=orange!40] (0,1.6) -- (0.7,2.2) -- (2.3,2.2) -- (1.6,1.6) -- cycle;
  \draw[thick, fill=orange!15] (1.6,0) -- (2.3,0.6) -- (2.3,2.2) -- (1.6,1.6) -- cycle;
  \node[font=\footnotesize, align=center] at (1.15,-0.5) {modello 3D\\(la mesh)};
  % l'apertura
  \draw[ue3arr] (2.8,1.1) -- node[above, font=\footnotesize]{si ``apre''}
    node[below, font=\footnotesize]{(UV mapping)} (5.0,1.1);
  % la mappa UV a croce
  \begin{scope}[shift={(5.4,-0.6)}]
    \draw[step=0.75, gray!40, thin] (0,0) grid (3.75,3.0);
    \foreach \x/\y in {0.75/0, 0/0.75, 0.75/0.75, 1.5/0.75, 2.25/0.75, 0.75/1.5}
      \draw[thick, fill=orange!25] (\x,\y) rectangle ++(0.75,0.75);
    \draw[->, thick] (0,0) -- (4.0,0) node[right, font=\footnotesize]{U};
    \draw[->, thick] (0,0) -- (0,3.25) node[above, font=\footnotesize]{V};
    \node[font=\footnotesize, align=center] at (1.9,-0.5) {mappa UV: le facce\\
      stese sul piano $[0,1]^2$};
  \end{scope}
  % la texture
  \draw[ue3arr] (9.4,1.1) -- node[above, font=\footnotesize]{si dipinge} (10.6,1.1);
  \node[ue3box, minimum width=2.3cm, minimum height=2.3cm, fill=orange!10,
    text width=2.1cm] (tx) at (12.0,1.1) {\footnotesize texture 2D\\
    \footnotesize (Diffuse, Normal,\\\footnotesize Roughness\ldots)\\[3pt]
    \footnotesize si ``riavvolge''\\\footnotesize sulla mesh};
\end{tikzpicture}
\caption{La texture come ``carta regalo'': le facce del modello vengono stese su un
piano (la mappa UV, fatta in Blender o in Maya), le texture si dipingono su quel piano,
e poi si ``riavvolgono'' sull'oggetto. È per questo che la texture della torretta, sulla
slide, appare ``tutta scomposta e segmentata''. Schema ricostruito; le lettere U e V
indicano i due assi della texture, per distinguerli da X, Y e Z del mondo.}
\label{ue3:fig:uv}
\end{figure}

\paragraph{Approfondimento: perché U e V.} Le lettere U e V sono scelte solo per non
confondere gli assi della texture con gli assi X, Y e Z dello spazio 3D. Una coordinata
UV è una coppia di numeri, di solito tra 0 e 1: $(0,0)$ è un angolo della texture e
$(1,1)$ l'angolo opposto. Ogni vertice della mesh riceve una coppia $(u,v)$, e per i
punti interni a una faccia il motore interpola le coordinate dei vertici. Se una mesh
non ha UV (o ne ha di sbagliate), il motore non sa ``quale pezzo di carta regalo'' va
su quale faccia: di qui le texture stirate di cui parla Davide.

% ---------------------------------------------------------------------
\section{I tipi di mappa}
\label{ue3:sec:mappe}

Davide passa in rassegna, sulla slide della torretta, le texture una per una, ``piano
piano'', con le informazioni che portano.

\paragraph{La Diffuse, o Albedo.} La prima è la \textbf{Diffuse}, ``in alcuni casi
detta anche \textbf{Albedo}''. Nei file la si riconosce dalla \textbf{nomenclatura}:
``T maiuscolo, \emph{underscore}, nome dell'oggetto, \emph{underscore}, A'', o D: A per
\emph{albedo}, D per \emph{diffuse} (per esempio \texttt{T\_Turret\_D}). La Diffuse è
il \textbf{colore}: ``tutti i dettagli estetici colorati''. Sulla slide è quella in alto
a destra, con l'arancione: è la texture che dà alla torretta l'arancione sulla base, il
nero negli altri dettagli e il bianco in alcune zone. Più avanti Davide ne completa la
descrizione: la Diffuse applica ``i colori e i dettagli di base'' sulle superfici 3D, e
dà l'informazione del colore, del tipo di materiale, dei dettagli (``queste linee che
vedete qua'') e di quanto un oggetto sia ruvido o levigato, se per esempio si applica la
texture di una roccia o di un metallo.

\paragraph{La Normal.} Poi c'è la \textbf{Normal}, la \textbf{Normal Map}, che ``serve
a simulare una tridimensionalità'' dentro un'informazione che è 2D. Si riconosce
subito: è \textbf{blu}, ``sempre mappe blu con delle sfumature di rosso e di verde''.
L'esempio classico è un \textbf{muro di mattoni}: in gioco il muro è \textbf{piatto}, ma
quando gli si applica la Normal il calcolo delle luci e delle ombre ``prende in
considerazione la normal map'' e proietta le ombre in base a lei, dando l'illusione che
i mattoni siano davvero tridimensionali. La Normal dà informazioni sulla
tridimensionalità ``senza appesantire particolarmente il calcolo'', perché non deforma
davvero l'oggetto come fa una mappa di altezza.

Più avanti Davide ci torna con un esempio sulla torretta. Su una delle gambe della base
servono dei \textbf{graffi}. Si potrebbe cambiare il modello aggiungendo i graffi, ma
questo aggiungerebbe \textbf{poligoni}, e quindi ``potenza di calcolo non
necessaria''. Si può invece aggiungere il dettaglio ``direttamente in 2D'', con la
Normal: protuberanze, graffi, rughe. ``Per un singolo oggetto non è un appesantimento
così importante. Ma considerate quanti oggetti ci sono all'interno di una mappa quando
andate a creare un livello'': se ogni oggetto non è creato con la giusta accortezza,
``si fa molto velocemente ad appesantire il progetto più di quanto dovrebbe essere
pesante''.

\paragraph{La Specular.} Lo \textbf{Specular mapping} riguarda ``la brillantezza e la
riflettività dell'oggetto'': come l'oggetto gestisce la luce, come la luce ``rimbalza''
su di lui. Un oggetto può avere parti metalliche, parti in pietra e altre ancora, e ``a
seconda del tipo di materiale di cui l'oggetto è fatto, rifrangerà la luce in una
maniera diversa''. Queste mappe sono quelle ``in bianco e in nero''. La texture
determina come si calcolano i riflessi ``sui singoli pixel proprio dell'oggetto'',
quanta luce viene riflessa; tornerà utile quando si vedranno le luci in Unreal, perché
``cambiando un materiale la luce ha un effetto completamente diverso''.

\paragraph{La mappa di altezza.} Poi c'è la mappa che serve a creare
\textbf{displacement}\footnote{Nella trascrizione è sempre ``light map'' o
``lightmap'', ma la descrizione (rialza e abbassa zone dell'oggetto, crea
\emph{tessellation}, è ``una normal sotto steroidi'') è quella della \textbf{height
map}, la mappa di altezza o di \emph{displacement}. Le \emph{lightmap} vere sono
un'altra cosa: texture in cui il motore salva l'illuminazione precalcolata.}: crea
``una \emph{tessellation} nell'oggetto'', rialza alcune zone, ne abbassa altre,
``modifica leggermente l'oggetto''. Davide la chiama ``una normal sotto steroidi'',
perché fa davvero \emph{displacement}: aggiunge profondità e dettagli tridimensionali
alle superfici piane, e usa le scale di grigio per creare scanalature e protuberanze.
``Nelle parti in cui vedete il bianco è sempre il sì'', la protuberanza; nel nero,
l'effetto opposto.

\paragraph{Il texturing procedurale.} Infine il \textbf{procedural texturing}: il
processo di creare texture, materiali o \textbf{shader} ``attraverso degli
algoritmi''. ``Andando ad utilizzare dei calcoli matematici, noi andiamo a modificare
delle informazioni per avere un risultato.'' È esattamente quello che fa l'editor dei
materiali di Unreal.

\begin{table}[htbp]
\centering
\small
\begin{tabular}{@{}p{2.5cm}p{2.4cm}p{8.1cm}@{}}
\toprule
\textbf{Mappa} & \textbf{Aspetto} & \textbf{Che cosa porta} \\
\midrule
Diffuse (Albedo) & a colori & il colore e i dettagli estetici; suffisso \_D o \_A \\
Normal & blu, con sfumature rosse e verdi & l'illusione del rilievo: la luce e le
  ombre si calcolano come se la superficie avesse graffi, mattoni, rughe, senza
  aggiungere poligoni \\
Specular & bianco e nero & quanto e come la superficie riflette la luce \\
Roughness & bianco e nero & quanto la superficie è ruvida (1) o liscia (0) \\
Metallic & bianco e nero & dove la superficie è metallo (1) e dove no (0) \\
Ambient Occlusion & bianco e nero & le zone in ombra perché ``chiuse'' (pieghe,
  angoli) \\
Height (displacement) & bianco e nero & deforma davvero la superficie: il bianco
  rialza, il nero abbassa \\
\bottomrule
\end{tabular}
\caption{I tipi di mappa nominati a lezione. Diffuse, Normal, Specular e la mappa di
altezza sono quelle della slide; Roughness, Metallic e Ambient Occlusion compaiono
nella parte sulle mappe impacchettate e nei materiali della torretta.}
\label{ue3:tab:mappe}
\end{table}

\paragraph{Il bianco è 1, il nero è 0.} La cosa importante da capire, insiste Davide, è
che nelle mappe in bianco e nero ``dovete sempre considerare il \textbf{bianco come un
1}, il \textbf{nero come uno 0} e il \textbf{grigio come dei valori nel mezzo} fra lo 0 e
l'1''. Che cosa vuol dire? Che in Unreal, per esempio, il valore dello Specular va
soltanto tra 0 e 1: a 1 riflette la luce, a 0 no. Una mappa di Specular, quindi, va
letta come una griglia di numeri: ``ci sono degli 1 in queste parti e degli 0 nel resto
delle parti''. È la chiave per capire quasi tutto il resto della lezione.

\paragraph{Approfondimento: la mappa di altezza, la Normal e Unreal Engine 5.} La
differenza tra Normal e mappa di altezza è la differenza tra \emph{sembrare} e
\emph{essere}. La Normal Map memorizza, per ogni pixel, la direzione in cui ``guarda''
la superficie (la sua \emph{normale}): i tre numeri X, Y e Z della direzione sono
salvati nei canali rosso, verde e blu, e siccome la maggior parte dei pixel guarda
``dritto fuori'' dalla superficie (Z alto) la mappa appare blu. La geometria non
cambia: di profilo, un muro con Normal Map resta piatto. La mappa di altezza invece
sposta davvero i vertici, e per farlo servono vertici da spostare: di qui la
\emph{tessellation}, che suddivide la superficie in molti triangoli. In Unreal Engine 5
il \emph{displacement} vero si fa soprattutto con le mesh \textbf{Nanite}, ed è più
costoso della Normal: per questo, per graffi e piccoli dettagli, la Normal resta la
scelta di default, come dice Davide.

% ---------------------------------------------------------------------
\section{I canali RGB e le mappe impacchettate}
\label{ue3:sec:canali}

``Chiaramente, queste texture utilizzano l'\textbf{RGB}'': non hanno un solo canale, ma
\textbf{tre canali} su cui passano le informazioni. Sulla slide c'è una Diffuse
scomposta in rosso, verde e blu: ognuno porta informazioni separate, che combinate
danno il colore che serve.

Ma c'è un altro modo di vedere la cosa. \textbf{Singolarmente}, ogni canale appare come
un'immagine in \textbf{bianco, nero e grigio}. E bianchi, neri e grigi ``sono
combinazioni di 0 e di 1'': quindi ogni canale si può usare per trasmettere
\textbf{un'informazione diversa}. È il \textbf{multichannel}. Una pratica comune,
nell'esportare e importare texture, è avere una singola mappa ``che spesso e volentieri
viene siglata come \textbf{AORM}'', per \textbf{Ambient Occlusion, Roughness e
Metallic}: tre mappe in bianco e nero, ognuna in un canale. ``Magari noi vogliamo
utilizzare questa mappa qui sopra per l'ambient occlusion, questa mappa qui sotto per
le metallic'', e così via. Una sola mappa con più informazioni invece di tante mappe
diverse: ``va ad alleggerire il progetto, va a rendere più veloce, più modulare''
l'export delle texture.

\begin{figure}[htbp]
\centering
\begin{tikzpicture}[
  font=\small,
  ue3ch/.style={draw, thick, minimum width=2.4cm, minimum height=1.0cm,
                align=center},
  ue3arr/.style={-{Stealth[length=2.5mm]}, thick}
]
  \node[draw, very thick, minimum width=2.6cm, minimum height=3.6cm, align=center,
    fill=gray!10] (tex) at (0,0) {\textbf{una texture}\\\textbf{RGB}\\[4pt]
    \footnotesize es. T\_Turret\_ORM};
  \node[ue3ch, fill=red!15] (r) at (4.2,1.3) {canale \textbf{R}};
  \node[ue3ch, fill=green!15] (g) at (4.2,0) {canale \textbf{G}};
  \node[ue3ch, fill=blue!15] (b) at (4.2,-1.3) {canale \textbf{B}};
  \draw[ue3arr] (tex.east) ++(0,1.0) -- (r.west);
  \draw[ue3arr] (tex.east) -- (g.west);
  \draw[ue3arr] (tex.east) ++(0,-1.0) -- (b.west);
  \node[ue3ch, text width=3.4cm] (ao) at (8.6,1.3) {\textbf{Ambient Occlusion}\\
    \footnotesize grigi tra 0 e 1};
  \node[ue3ch, text width=3.4cm] (ro) at (8.6,0) {\textbf{Roughness}\\
    \footnotesize grigi tra 0 e 1};
  \node[ue3ch, text width=3.4cm] (me) at (8.6,-1.3) {\textbf{Metallic}\\
    \footnotesize grigi tra 0 e 1};
  \draw[ue3arr] (r) -- (ao);
  \draw[ue3arr] (g) -- (ro);
  \draw[ue3arr] (b) -- (me);
  \node[font=\footnotesize, align=center, text width=11cm] at (4.3,-2.5)
    {ogni canale, da solo, è un'immagine in bianco e nero: tre ``mappe'' diverse in
    un solo file};
\end{tikzpicture}
\caption{Una mappa impacchettata (\emph{channel packing}). La lezione la chiama AORM;
nella pratica di Unreal la convenzione più diffusa è ORM, con l'Ambient Occlusion nel
rosso, la Roughness nel verde e il Metallic nel blu, ma l'ordine dipende da chi
esporta le texture.}
\label{ue3:fig:canali}
\end{figure}

\paragraph{Approfondimento: sRGB e compressione.} Quando si importa in Unreal una
texture impacchettata (o qualunque mappa ``di dati'', come Roughness o Metallic), conviene
aprirla e \textbf{togliere la spunta a sRGB}: l'opzione sRGB dice al motore che
l'immagine è un \emph{colore} da correggere per la percezione dell'occhio, mentre una
mappa di dati contiene numeri da leggere così come sono. Nelle stesse impostazioni, il
campo \emph{Compression Settings} si mette su \textbf{Masks} per le mappe impacchettate
e su \textbf{Normalmap} per le Normal (Unreal di solito lo riconosce da solo per le
texture con suffisso \_N). Se ci si dimentica, il materiale funziona lo stesso, ma i
grigi vengono letti un po' sbagliati e la superficie risulta più lucida o più opaca del
previsto.

% ---------------------------------------------------------------------
\section{Il materiale: la vernice}
\label{ue3:sec:materiale}

\paragraph{Le texture non si modificano.} ``Quello che possiamo fare con le texture in
Unreal è relativamente poco, perché non possiamo andare a modificare effettivamente le
texture all'interno dell'engine.'' E allora come si danno qualità visive alle mesh?
``Lo facciamo con i \textbf{materiali}.''

\paragraph{Il materiale è uno shader.} I materiali ``sono degli \textbf{shader}''. E
l'\textbf{editor dei materiali}, quello che si vedrà ``principalmente'' oggi, ``funziona
praticamente nello stesso modo dell'editor dei Blueprint'': una serie di \textbf{nodi}
collegati, che trasportano informazioni fino a un \textbf{nodo risultante}.

\paragraph{Vernice, non carta regalo.} La distinzione è netta: ``il materiale è
effettivamente la \textbf{vernice} che viene applicata sulla mesh, non la texture''. Sugli
oggetti si applicano \textbf{materiali}. Le texture servono a \textbf{costruire} i
materiali: si combinano le texture in un unico materiale, che poi si applica
sull'oggetto finale. Il materiale dice all'oggetto ``come dovrebbe interagire con la
luce, che tipo di informazioni visive dovrebbe portare'', dove ci sono
tridimensionalità e così via: ``tutto quello che abbiamo visto dalle varie texture
precedentemente viene tutto convogliato all'interno dei materiali''. Questo, precisa
Davide, se si parla di \textbf{oggetti fisici}. I materiali hanno anche altre funzioni:
servono per esempio a fare \textbf{effetti speciali}, e un materiale per un effetto
speciale ``dovrà essere costruito in maniera diversa'' da quello che riveste una
torretta.

\begin{figure}[htbp]
\centering
\begin{tikzpicture}[
  font=\small,
  ue3tex/.style={draw, thick, minimum width=2.3cm, minimum height=0.7cm,
                 align=center, fill=gray!10},
  ue3big/.style={draw, very thick, rounded corners=4pt, align=center},
  ue3arr/.style={-{Stealth[length=2.5mm]}, thick}
]
  \node[ue3tex] (t1) at (0,1.5) {T\_Turret\_D};
  \node[ue3tex] (t2) at (0,0.5) {T\_Turret\_N};
  \node[ue3tex] (t3) at (0,-0.5) {T\_Turret\_R};
  \node[ue3tex] (t4) at (0,-1.5) {T\_Turret\_M};
  \node[font=\footnotesize, align=center] at (0,2.4) {\textbf{texture}\\(immagini)};
  \node[ue3big, minimum width=3.4cm, minimum height=3.6cm, fill=blue!6,
    text width=3.2cm] (mat) at (4.4,0) {\textbf{materiale}\\M\_Turret\\[6pt]
    \footnotesize un grafo di nodi\\\footnotesize (uno \emph{shader})\\[6pt]
    \footnotesize la ``vernice''};
  \foreach \tn in {t1,t2,t3,t4} \draw[ue3arr] (\tn.east) -- (mat.west |- \tn);
  \node[ue3big, minimum width=2.8cm, minimum height=1.6cm, fill=orange!12,
    text width=2.6cm] (mi) at (8.4,0) {\textbf{Material}\\\textbf{Instance}\\
    \footnotesize MI\_Turret\\\footnotesize (parametri)};
  \draw[ue3arr] (mat) -- (mi);
  \node[ue3big, minimum width=2.4cm, minimum height=1.6cm, fill=green!8,
    text width=2.2cm] (mesh) at (11.8,0) {\textbf{mesh}\\\footnotesize SM\_Turret\\
    \footnotesize slot del\\\footnotesize materiale};
  \draw[ue3arr] (mi) -- node[above, font=\footnotesize]{si applica} (mesh);
\end{tikzpicture}
\caption{Texture, materiale, istanza e mesh. Le texture non si applicano direttamente
agli oggetti: entrano in un materiale, che è la ``vernice''; dal materiale si crea
un'istanza, ed è l'istanza che si applica allo slot del materiale della mesh
(sezione~\ref{ue3:sec:instance}). I nomi seguono la nomenclatura suggerita a lezione
(T\_, M\_, MI\_, SM\_).}
\label{ue3:fig:pipeline}
\end{figure}

\paragraph{Approfondimento: il materiale nel libro di Romero e Sewell.} Il quinto
capitolo del libro (``Object Interaction with Blueprints'') usa la stessa immagine di
Davide: un materiale è un \emph{asset} che si aggiunge alla mesh di un actor per
crearne l'aspetto, ``come uno strato di vernice''. Il libro crea un materiale rosso,
\texttt{M\_TargetRed}, in una cartella Materials (``non è necessario, ma è buona
pratica tenere ordinata la gerarchia del progetto''), e nota che l'editor dei
materiali ``condivide molte caratteristiche e convenzioni con i Blueprint'': al centro
c'è il grafo; il nodo grande con il nome del materiale è il \textbf{result node}; con il
tasto destro si apre un menu di nodi con una ricerca sensibile al contesto; i fili si
tirano trascinando dai pin, gli ingressi a sinistra e le uscite a destra. Il colore lo
dà un \textbf{VectorParameter} collegato a Base Color; due \textbf{ScalarParameter} a
0,1 danno un po' di Metallic e di Roughness. Poi, in un Blueprint, il nodo \textbf{Set
Material} cambia il materiale di un bersaglio quando viene colpito: un primo esempio di
come la logica (le lezioni di Igor) e la grafica (questa lezione) si incontrino.

% ---------------------------------------------------------------------
\section{Material Domain, Blend Mode e Shading Model}
\label{ue3:sec:dettagli}

Come gli oggetti nella \emph{viewport}, anche il materiale ha i suoi
\textbf{Details}. Dentro l'editor si possono selezionare i singoli nodi, per i dettagli
del nodo, oppure vedere i \textbf{dettagli generali del materiale}. Tra questi ce ne
sono alcuni, evidenziati sulla slide, che ``cambiano proprio effettivamente come
funziona il materiale'' e ne definiscono la struttura e le possibilità.

\paragraph{Material Domain.} Il \textbf{Material Domain} dice \textbf{per che cosa}
verrà usato il materiale: una \textbf{superficie}, l'\textbf{interfaccia} (la UI), il
\textbf{post process}, un effetto.

\paragraph{Blend Mode.} La \textbf{Blend Mode} è il modo in cui il materiale
``\emph{blenda}'' (si mescola) con i pixel calcolati nel resto della scena. La modalità
\textbf{Opaque}, quella di default, \textbf{occlude completamente} ciò che c'è dietro.
Per un effetto speciale, per esempio un \textbf{fuoco}, serve invece la Blend Mode
\textbf{Translucent}, che permette di aggiungere parametri per l'\textbf{opacità} del
materiale.

\paragraph{Shading Model.} Lo \textbf{Shading Model}\footnote{Nella trascrizione una
volta ``scaling model''.} definisce come il materiale \textbf{reagisce alla luce}.
\textbf{Default Lit} vuol dire che segue le regole di base dell'illuminazione: il
materiale viene illuminato. Per un effetto speciale, in cui non si vuole che la luce
influenzi la resa, lo si mette invece su \textbf{Unlit}\footnote{Nella trascrizione
``lo metterò un lit''.}. Ci sono ``diverse tipologie'' di Shading Model, a seconda di
cosa si vuole realizzare.

\paragraph{I pin cambiano.} Queste impostazioni si applicano sul nodo risultante, e
cambiarle cambia i \textbf{pin} del nodo. L'\textbf{Opacity} non si può modificare in
Blend Mode Opaque, perché un materiale opaco ``dà per scontato che l'oggetto sia sempre
opaco'': non c'è bisogno di dire dove e quando sarà trasparente. Viceversa, la Blend
Mode Translucent toglie la possibilità di modificare Metallic, Specular, Roughness e
così via: se il materiale deve essere potenzialmente trasparente, ``non ha senso andare
a calcolare come rifletterà la luce''.

\begin{table}[htbp]
\centering
\small
\begin{tabular}{@{}p{2.6cm}p{3.3cm}p{7.1cm}@{}}
\toprule
\textbf{Impostazione} & \textbf{Valori citati} & \textbf{Che cosa decide} \\
\midrule
Material Domain & Surface, User Interface, Post Process\ldots & per che cosa si usa
  il materiale: una superficie, la UI, il post process, un effetto \\
Blend Mode & \textbf{Opaque} (default), Masked, Translucent & come il materiale si
  mescola con ciò che ha dietro: Opaque nasconde tutto; Masked ritaglia con una
  maschera; Translucent sfuma con l'Opacity (fuoco, fili d'erba) \\
Shading Model & \textbf{Default Lit} (default), Unlit\ldots & come reagisce alla luce:
  Default Lit viene illuminato, Unlit no (effetti speciali) \\
\bottomrule
\end{tabular}
\caption{Le tre impostazioni ``più importanti'' dei dettagli di un materiale. Cambiarle
accende e spegne i pin del nodo risultante.}
\label{ue3:tab:dettagli}
\end{table}

\paragraph{La prova in Engine.} Per far vedere la cosa, Davide apre il progetto del
materiale didattico e invita chi vuole a seguirlo: nel \textbf{Content Browser} si crea
un nuovo materiale (``va bene anche chiamarlo New Material''), lo si apre con un
\textbf{doppio clic} e si entra nell'editor. Le impostazioni di cui si parlava sono nei
dettagli, in basso a sinistra. Se la Blend Mode passa da Opaque a \textbf{Translucent},
il nodo risultante cambia: spariscono Metallic, Specular, Roughness, Anisotropy, e
compare \textbf{Opacity}. È utile, per esempio, per una \textbf{fiamma} che verso l'alto
diventa trasparente (le informazioni grafiche alla base, e in alto ``piano piano va a
dissolversi''), o per un \textbf{filo d'erba} che parte con il suo colore e poi va in
trasparenza. Con la Blend Mode \textbf{Masked}, ``mascherata'', si crea invece una
\textbf{maschera}, un \emph{layer} che funziona in modo diverso dal resto della
superficie. Per ora si resta su \textbf{Opaque}, perché i materiali di oggi andranno su
oggetti 3D. A metà spiegazione arriva uno studente in ritardo: ``Chat, buon pomeriggio,
Lorenzo. Noi qui abbiamo già cominciato, ma in realtà abbiamo parlato di cose
principalmente teoriche.''

\paragraph{Approfondimento: Masked o Translucent?} Le due modalità sembrano simili ma
lavorano in modo opposto. \textbf{Masked} decide pixel per pixel se il materiale
\emph{c'è o non c'è}: l'ingresso Opacity Mask viene confrontato con una soglia (di
default 0,333), e sotto la soglia il pixel non viene disegnato affatto. È la modalità
delle foglie, delle reti e dell'erba ``ritagliata'': economica, compatibile con la
luce e con le ombre, ma senza sfumature. \textbf{Translucent} invece \emph{mescola} il
colore del materiale con ciò che sta dietro, secondo un valore di Opacity tra 0 e 1:
permette le sfumature (fumo, vetro, fiamme), ma costa di più, perché ogni strato
traslucido va sommato a quelli dietro, e gestisce peggio luce e ordinamento. Regola
pratica: Masked quando basta un ``sì o no'', Translucent solo quando serve davvero la
sfumatura.

% ---------------------------------------------------------------------
\section{I nodi: le Material Expression}
\label{ue3:sec:nodi}

``Come per i Blueprint, anche l'editor dei materiali ha dei nodi, delle
\textbf{espressioni}.'' I \textbf{Material Expression Node} sono ``dei piccoli set di
istruzioni preprogrammate che eseguono delle cose nello specifico, esattamente come i
nodi dei Blueprint'', o, con un'altra formula, ``pezzi di codice statico''.

\paragraph{Il primo grafo.} Sulla slide c'è un piccolo grafo
(figura~\ref{ue3:fig:primografo}), e Davide lo legge nodo per nodo:
\begin{itemize}
  \item il \textbf{Texture Sample} comunica l'informazione della texture: ``sta dicendo
        che la texture che c'è in questo nodo è la texture che determina il
        \textbf{base color}''; permette di usare un'immagine e farla diventare la
        texture del materiale;
  \item il \textbf{Texture Coordinate}, il nodo \textbf{rosso}, restituisce le
        coordinate U e V del materiale, e serve a decidere ``quanto è grande il
        \textbf{tiling} della mia texture, quindi quante volte si ripete'';
  \item il \textbf{Multiply} è programmato per moltiplicare due valori: ``moltiplica A
        per B e dà una risultante'';
  \item il quadratino è una \textbf{costante}, ``un singolo \emph{float}'', una variabile
        singola.
\end{itemize}
Il grafo dice: ``moltiplica le coordinate per uno e comunica queste cose alle UV della
texture''. Così com'è, quindi, ``è come se non ci fosse'': moltiplicare per 1 lascia le
coordinate di default. Ma è così che si costruiscono i grafi: combinando nodi si
ottengono risultati molto diversi. Si può moltiplicare, si può fare
un'\textbf{interpolazione}, si possono far \textbf{muovere} le texture per suggerire un
movimento nel materiale, ``eccetera, eccetera, eccetera''.

\begin{figure}[htbp]
\centering
\begin{tikzpicture}[
  font=\small,
  ue3node/.style={draw, thick, rounded corners=4pt, align=center,
                  minimum height=0.9cm},
  ue3uv/.style={-{Stealth[length=2mm]}, thick, red!60!black},
  ue3flt/.style={-{Stealth[length=2mm]}, thick, green!50!black},
  ue3col/.style={-{Stealth[length=2mm]}, thick, blue!60!black}
]
  \node[ue3node, fill=red!12, text width=2.2cm] (tc) at (0,0.6)
    {\textbf{TexCoord}\\\footnotesize UTiling 1, VTiling 1};
  \node[ue3node, fill=green!8, minimum width=1.0cm] (k) at (0,-1.0) {\textbf{1}};
  \node[font=\footnotesize] at (0,-1.75) {costante};
  \node[ue3node, fill=gray!10, text width=1.6cm] (mul) at (3.0,0)
    {\textbf{Multiply}\\\footnotesize A $\times$ B};
  \node[ue3node, fill=gray!10, text width=2.4cm, minimum height=1.8cm] (ts)
    at (6.4,0) {\textbf{Texture}\\\textbf{Sample}\\\footnotesize (T\_Bark)};
  \node[ue3node, fill=blue!6, text width=2.4cm, minimum height=2.6cm] (res)
    at (10.4,0) {\textbf{M\_NewMaterial}\\[3pt]\footnotesize\textbf{Base Color}\\
    \footnotesize Metallic\\\footnotesize Specular\\\footnotesize Roughness\\
    \footnotesize Emissive Color\\\footnotesize Normal\ldots};
  \draw[ue3uv] (tc.east) -- ++(0.4,0) |- node[pos=0.75, above, font=\scriptsize]{A}
    ([yshift=2mm]mul.west);
  \draw[ue3flt] (k.east) -- ++(0.6,0) |- node[pos=0.75, below, font=\scriptsize]{B}
    ([yshift=-2mm]mul.west);
  \draw[ue3uv] (mul.east) -- node[above, font=\scriptsize]{UVs}
    (ts.west |- mul.east);
  \draw[ue3col] (ts.east) -- node[above, font=\scriptsize]{RGB}
    (res.west |- ts.east);
\end{tikzpicture}
\caption{Il primo grafo della slide: le coordinate UV, moltiplicate per una costante,
entrano nel Texture Sample, il cui colore va nel Base Color del nodo risultante. Con la
costante a 1 il Multiply ``è come se non ci fosse''; cambiando la costante cambia il
tiling (sezione~\ref{ue3:sec:tiling}). Schema ricostruito.}
\label{ue3:fig:primografo}
\end{figure}

\paragraph{I parametri.} Non solo le costanti, ma i nodi in generale ``possono essere
utilizzati come \textbf{parametri}''. Per rendere un nodo un parametro si fa
\textbf{tasto destro} sul nodo e \textbf{Convert to Parameter}, gli si dà un nome, e
questo permette di modificare in seguito il valore ``senza dover ricalcolare da capo il
materiale''. I parametri torneranno utili con le \textbf{Material Instance}
(sezione~\ref{ue3:sec:instance}).

\paragraph{Come si chiama un nodo.} ``Tasto destro, scriviamo il nome'', per esempio
\emph{texture sample}, ed ecco il nodo. C'è poi ``tutta una serie di \emph{shortcut}''
per i nodi più usati, che si vedranno tra poco.

% ---------------------------------------------------------------------
\section{Il problema della mappa vuota}
\label{ue3:sec:mappavuota}

Una studentessa interrompe: non riesce più ad aprire la mappa, ``me la apre totalmente
vuota, senza titolo''. Davide le chiede di condividere lo schermo, ma lei segue da due
computer diversi. Prima ipotesi: ha creato un \textbf{nuovo livello}, e un \emph{New
Level} appena caricato è tutto nero, perché non c'è niente in scena (``potremmo andare
a popolarla mettendoci cose''). Ma il problema è un altro: ha aperto il progetto del
materiale didattico, cliccato sulla mappa del personaggio in terza persona
(\emph{Third Person Map}), e non succede niente. Davide sospetta che il progetto sia
ancora in una \textbf{cartella zippata} (``di solito questa cosa succede quando il
progetto è una cartella che è ancora zippata''), e le consiglia di
\textbf{riscaricarlo}. Poco dopo il problema è risolto: ``sono riuscita a risolvere
adesso''. ``Meno male, meno male.''

\paragraph{Approfondimento: aprire un progetto da una cartella.} Un progetto di Unreal
è una cartella che contiene un file \texttt{.uproject} e le sottocartelle
\texttt{Content}, \texttt{Config} e altre. Unreal non può aprire un progetto che sta
ancora dentro un archivio \texttt{.zip}: Windows mostra il contenuto dello zip come se
fosse una cartella, ma l'editor non ci può scrivere, e il risultato tipico è proprio
una mappa vuota o mancante. La soluzione è estrarre tutto (``Estrai tutto'') e aprire il
\texttt{.uproject} estratto con un doppio clic. Se il progetto non compare nell'elenco
dell'Epic Games Launcher o nel Project Browser, basta aprirlo dal \texttt{.uproject}
la prima volta: da lì in poi compare tra i progetti recenti.

% ---------------------------------------------------------------------
\section{Dentro una texture e il Texture Sample}
\label{ue3:sec:sample}

``Facciamo le cose in ordine, se no io mi confondo e mi impiccio.'' Davide torna al
punto saltato: com'è fatta \textbf{dentro} una texture.

\paragraph{La finestra della texture.} Con un doppio clic su una texture nel Content
Browser (Davide ne prende una a caso, una corteccia, \textbf{T\_Bark}) si apre una
finestra in cui ``non possiamo fare grandi modifiche'', ma si possono vedere i
\textbf{canali}: il blu, il verde, il rosso, ciascuno fatto a modo suo. ``Questa texture
non è particolarmente interessante'', ma si rivedrà meglio con la torretta.

\paragraph{Il Texture Sample.} Nel grafo del materiale si mette un \textbf{Texture
Sample}, si seleziona il nodo, e nei dettagli, in basso, si sceglie quale texture
metterci. Cercando T\_Bark e mettendola nel nodo, la \emph{preview} mostra la corteccia:
``il nostro materiale sarà così, perché è effettivamente la texture che ci abbiamo
messo''.

\paragraph{Le uscite separate.} Il Texture Sample non ha una sola uscita: oltre
all'\textbf{RGB} ci sono \textbf{R}, \textbf{G} e \textbf{B} separati (e l'alfa, A). Si
può quindi usare, per esempio, solo il canale \textbf{rosso} (``non quella completa, ma
questa qui'') per trasportare un'informazione: collegandolo allo \textbf{Specular} si
ottiene un materiale che in alcune parti riflette la luce più che in altre. ``Con
questo materiale in realtà è un po' complicato farlo vedere'', ma il principio è quello
dei canali visto sulle slide.

\paragraph{Come si mette una texture nel nodo.} Una studentessa non ha capito come si
aggiunge la texture ``al quadratino nero'' del Texture Sample. Davide rifà il
percorso: si clicca il nodo, e nei suoi dettagli c'è il campo \textbf{Texture}, con il
quadratino e una lista. Due modi:
\begin{itemize}
  \item se una texture è già \textbf{selezionata nel Content Browser}, la freccetta
        \textbf{Use Selected Asset from Content Browser} la inserisce automaticamente
        nel nodo;
  \item più spesso, si clicca sul \textbf{nome} della texture nel campo e si fa un
        \emph{browse}: si cerca il nome (``T underscore Bark'') e lo si sceglie dalla
        lista di tutte le texture disponibili.
\end{itemize}
``Voi lo vedete questo pop-up che esce?'' Sì: perfetto.

% ---------------------------------------------------------------------
\section{Le costanti, le operazioni e le scorciatoie}
\label{ue3:sec:costanti}

``Chiaramente ne vediamo un po', non li vediamo tutti, perché sono tanti, sono
veramente tanti.'' Ci sono ``tantissime possibilità, tante cose che vanno studiate'',
ma per ora servono le basi.

\paragraph{La regola delle scorciatoie.} Il Texture Sample si può chiamare cercandolo
con il tasto destro, oppure \textbf{tenendo premuto il tasto T e facendo clic
sinistro} nel grafo. ``Questa struttura dello shortcut funziona per ogni singolo nodo
che ha uno shortcut'': un tasto, tenuto premuto, più clic sinistro. Il tasto è
``molto spesso l'iniziale del nome del nodo''.

\paragraph{Le costanti.} Cercando \emph{vector} compare, per esempio, il Vector
Noise; cercando \emph{constant} compaiono le costanti: \textbf{Constant},
\textbf{Constant2Vector}, \textbf{Constant3Vector} (e Constant4Vector). Le scorciatoie
sono i numeri: \textbf{1 + clic} per la costante singola, \textbf{2 + clic} per quella
a due valori, \textbf{3 + clic} per quella a tre. La costante a due ``la utilizzerete
relativamente poco''. Le altre due moltissimo:
\begin{itemize}
  \item la \textbf{costante singola} è un \textbf{valore numerico}, che si usa ``per
        andare a fare qualsiasi tipo di calcolo, qualsiasi tipo di costruzione'';
  \item la \textbf{costante a tre} è un \textbf{colore}: l'engine usa i tre valori come
        \textbf{R}, \textbf{G} e \textbf{B}. Si sceglie un colore (Davide prende un
        arancione) e l'engine modifica X, Y e Z per ottenerlo. ``Qualcuno di voi l'aveva
        anche già fatto nel test con la pallina.''
\end{itemize}
Collegando l'arancione al \textbf{Base Color} si ottiene un materiale arancione,
``perché effettivamente è l'unica cosa che gli sto passando''. Ma ``molto raramente''
si usa un colore piatto come Base Color: lo si usa piuttosto per fare
\textbf{operazioni}. Più avanti una studentessa chiederà di nuovo come si crea il colore:
``è come la costante singola, che è uno clic sinistro, ma è a tre, quindi è tre clic
sinistro''.

\paragraph{Staccare i fili.} Per rompere un collegamento: \textbf{Alt + clic sinistro}
sul pin, ``e si rompono''.

\paragraph{Le operazioni.} Le operazioni utili sono ``le classiche'': addizione,
moltiplicazione, divisione, sottrazione. Le scorciatoie: \textbf{A + clic} per
\textbf{Add}, \textbf{M + clic} per \textbf{Multiply}, \textbf{D + clic} per
\textbf{Divide}. La sottrazione (\textbf{Subtract}) non ha scorciatoia: si cerca. Il
motivo lo spiega Davide più avanti: \textbf{S + clic} crea automaticamente un
\textbf{parametro} (uno Scalar Parameter), ``quindi questo potrebbe tornarvi utile''.

\paragraph{Quali si usano davvero.} In chiusura di lezione Davide aggiunge un consiglio
d'esperienza: ``l'operazione che utilizzerete di più in assoluto è la
\textbf{moltiplicazione}''. Moltiplicazione e addizione sono le più usate; divisione e
sottrazione servono soprattutto per cose più complesse, ``tipo sugli effetti
speciali''.

% ---------------------------------------------------------------------
\section{Il colore che brilla: Emissive, Multiply e Glow Intensity}
\label{ue3:sec:emissive}

``A cosa ci servono'' le operazioni? Per esempio ad \textbf{aumentare il valore} di
qualcosa. Il nodo risultante ha diversi pin, ``diverse cose su cui andare a collegare
dei valori'': il \textbf{Base Color}, e più in basso l'\textbf{Emissive Color}, che
controlla ``quali parti del materiale emetteranno una piccola luce'', un
\textbf{\emph{glow}}.

\paragraph{L'arancione nell'Emissive.} Davide stacca l'arancione dal Base Color e lo
collega all'Emissive: il materiale ``sarà leggermente luminoso''. ``Adesso è molto poco
luminoso, ma è comunque un emissive.''

\paragraph{Moltiplicare.} Come si ottiene un'emissione maggiore? ``Moltiplichiamo.''
L'arancione entra nella A di un \textbf{Multiply}, il Multiply va nell'Emissive. Il
secondo valore si può scrivere direttamente nel nodo Multiply (``invece di 1 mettiamo
100''): ``verrà decisamente più luminoso''. Ma così ``è meno controllabile''. Meglio
attaccare alla B una \textbf{costante} separata, darle il valore 100 e convertirla in
\textbf{parametro}, con un nome come \textbf{Glow Intensity}: quando più avanti ci sarà
una Material Instance, il valore si potrà cambiare ``direttamente da lì''. ``Ci andiamo
a estrapolare un parametro.''

\begin{figure}[htbp]
\centering
\begin{tikzpicture}[
  font=\small,
  ue3node/.style={draw, thick, rounded corners=4pt, align=center,
                  minimum height=0.9cm},
  ue3col/.style={-{Stealth[length=2mm]}, thick, blue!60!black},
  ue3flt/.style={-{Stealth[length=2mm]}, thick, green!50!black}
]
  \node[ue3node, fill=orange!30, text width=2.5cm] (col) at (0,0.7)
    {\textbf{Constant3Vector}\\\footnotesize arancione (3 + clic)};
  \node[ue3node, fill=green!8, text width=2.5cm] (gi) at (0,-0.9)
    {\textbf{Glow Intensity}\\\footnotesize parametro = 100};
  \node[ue3node, fill=gray!10, text width=1.6cm] (mul) at (3.6,0)
    {\textbf{Multiply}\\\footnotesize (M + clic)};
  \node[ue3node, fill=blue!6, text width=2.6cm, minimum height=2.2cm] (res)
    at (7.4,0) {\textbf{nodo risultante}\\[3pt]\footnotesize Base Color\\
    \footnotesize Metallic, Specular\\\footnotesize Roughness\\
    \footnotesize\textbf{Emissive Color}};
  \draw[ue3col] (col.east) -- ++(0.4,0) |- node[pos=0.75, above, font=\scriptsize]{A}
    ([yshift=2mm]mul.west);
  \draw[ue3flt] (gi.east) -- ++(0.4,0) |- node[pos=0.75, below, font=\scriptsize]{B}
    ([yshift=-2mm]mul.west);
  \draw[ue3col] (mul.east) -- ++(0.6,0) |- ([yshift=-7mm]res.west);
\end{tikzpicture}
\caption{Un materiale che brilla. Il colore, collegato all'Emissive Color, emette ``una
piccola luce''; moltiplicato per 100 diventa molto più luminoso. Il valore 100 sta in
una costante convertita in parametro (Glow Intensity), così che si possa cambiare da
una Material Instance senza toccare il grafo. Ricostruzione dalla dimostrazione.}
\label{ue3:fig:emissive}
\end{figure}

\paragraph{Approfondimento: perché 100 funziona.} I colori di un materiale non sono
limitati tra 0 e 1. Nel Base Color valori oltre 1 non hanno senso fisico (nessuna
superficie riflette più luce di quella che riceve), ma nell'Emissive sì: un'emissione di
$(1; 0{,}5; 0)$ moltiplicata per 100 diventa $(100; 50; 0)$, una sorgente molto intensa
che il \emph{post process} di Unreal trasforma in un alone luminoso (il \emph{bloom}).
Il Multiply tra un colore e un numero moltiplica ciascun canale per quel numero, quindi
la tinta resta la stessa e cresce solo l'intensità. Con Lumen, il sistema di
illuminazione globale di Unreal Engine 5, un materiale emissivo può anche illuminare un
po' le superfici vicine.

% ---------------------------------------------------------------------
\section{Il PBR: Metallic, Specular, Roughness e Normal}
\label{ue3:sec:pbr}

``Mi rendo conto che sto andando un po' a destra e a sinistra'', ammette Davide, e
torna ai pin del nodo risultante. Base Color ed Emissive si sono visti. Restano
\textbf{Metallic}, \textbf{Specular} e \textbf{Roughness}, ``degli input abbastanza
interessanti'', perché sono quelli che \textbf{determinano la superficie} del
materiale.

\paragraph{Il PBR.} Il rendering dei materiali in Unreal è \textbf{PBR},
\emph{Physically Based Rendering}: per decidere come si comporta una superficie, il
motore mette insieme una serie di informazioni fisiche. Tutti questi parametri vanno
\textbf{da 0 a 1}.

\paragraph{Roughness.} La \textbf{Roughness} dice quanto una superficie è
\textbf{liscia} o \textbf{ruvida}. Una \textbf{roccia} vuole una Roughness molto alta,
cioè 1: il materiale sarà ruvido ``e non rifletterà la luce in nessun modo''. Uno
\textbf{specchio} vuole una Roughness a 0: il più liscio possibile, il meno ruvido, il
più riflettente.

\paragraph{Specular e Metallic.} ``Per aiutare in questa situazione'' ci sono Specular e
Metallic, che ``si alternano a seconda del tipo di materiale''. Lo
\textbf{Specular}, come si è visto con le mappe, determina quanta luce viene riflessa e
l'angolazione a cui viene riflessa. Per creare uno \textbf{specchio} serve un materiale
molto, molto liscio e molto, molto riflettente: Davide abbassa la Roughness e alza la
riflessione, e la \emph{preview} comincia ad assomigliare ``a una parete riflettente'';
con altre cose che si vedranno in seguito ``iniziamo ad avere un riflesso''.

\paragraph{Emissive e Normal.} L'\textbf{Emissive Color} restituisce un colore che è
``effettivamente una luce''. La \textbf{Normal} è l'ultimo parametro utile, e richiede
la texture blu vista prima: si usa sempre un Texture Sample, ma con
\textbf{T\_Bark\_N}, la Normal della corteccia, invece della Diffuse. Davide toglie
tutto il resto e mette a confronto il materiale con e senza Normal: ``vedete la
differenza?'' La Normal aggiunge un dettaglio tridimensionale ``molto semplice da
vedere'', eppure ``la mappa è 2D, non c'è tridimensionalità'': crea solo l'illusione.

\begin{figure}[htbp]
\centering
\begin{tikzpicture}[font=\small]
  \foreach \i/\rr in {0/0, 1/0.25, 2/0.5, 3/0.75, 4/1} {
    \pgfmathsetmacro{\hs}{0.12+0.35*\rr}
    \pgfmathsetmacro{\op}{1-0.85*\rr}
    \shade[ball color=gray!70] (\i*2.4,0) circle (0.8);
    \fill[white, opacity=\op] (\i*2.4-0.28,0.3) circle (\hs);
    \node[font=\footnotesize] at (\i*2.4,-1.2) {Roughness $=\rr$};
  }
  \draw[-{Stealth[length=2.5mm]}, thick] (-0.6,-1.8) -- (10.2,-1.8);
  \node[font=\footnotesize, anchor=north west] at (-0.6,-1.85) {liscio, riflessi netti
    (lo specchio)};
  \node[font=\footnotesize, anchor=north east] at (10.2,-1.85) {ruvido, nessun
    riflesso netto (la roccia)};
\end{tikzpicture}
\caption{La Roughness, da 0 a 1, in uno schema qualitativo: più la superficie è liscia,
più il riflesso della luce è piccolo e netto; più è ruvida, più la luce si sparpaglia
fino a sparire come riflesso. Disegno indicativo, non un rendering di Unreal.}
\label{ue3:fig:roughness}
\end{figure}

\paragraph{Approfondimento: Metallic e Specular nella pratica.} Nel modello PBR di
Unreal i due parametri non hanno lo stesso peso. \textbf{Metallic} è quasi sempre
\emph{binario}: 1 per i metalli, 0 per tutto il resto (legno, pietra, plastica, pelle);
i valori intermedi servono soprattutto per i bordi tra metallo e vernice o per lo
sporco. Lo \textbf{Specular} ha un default di 0,5, che va bene per la grande maggioranza
dei materiali non metallici, e la documentazione di Unreal consiglia di toccarlo poco:
per i metalli viene ignorato, perché il loro riflesso dipende dal Base Color. Lo
specchio della lezione, in termini PBR, è quindi Metallic $=1$, Roughness $=0$, Base
Color chiaro. È il senso dell'``alternarsi'' di cui parla Davide: nei materiali metallici
si lavora con il Metallic, negli altri, se proprio serve, con lo Specular, e quasi
sempre con la Roughness, che è il parametro che cambia di più l'aspetto di una
superficie.

% ---------------------------------------------------------------------
\section{Il Texture Coordinate e il tiling}
\label{ue3:sec:tiling}

``Chiaramente ogni nodo funziona a sé.'' Il \textbf{Texture Coordinate}, il nodo rosso
della prima slide, ha i suoi parametri. La scorciatoia è ``un po' \emph{tricky} da
ricordare'', perché T + clic è già preso dal Texture Sample: è \textbf{U + clic},
``perché il Texture Coordinate lavora sulle UV''.

\paragraph{UTiling e VTiling.} Selezionando il nodo, nei dettagli compaiono
\textbf{UTiling} e \textbf{VTiling}: i parametri ``molto comodi'' con cui si aumenta o
si diminuisce il tiling di un oggetto. Davide collega il Texture Coordinate alle UV
della Normal e della Diffuse, e chiede: se metto UTiling a 2, che cosa succede? ``Vediamo
due volte la texture applicata all'oggetto'', risponde uno studente. ``Sì, però c'è un
\emph{catch}, perché è soltanto sulle U'': la texture si ripete due volte
sull'\textbf{asse orizzontale} e una sola volta su quello verticale. Con anche VTiling a
2, la texture si ripete \textbf{quattro} volte: ``ce n'è una qua, una qua, una qui e una
qua''.

\paragraph{Lo zoom.} Si può fare anche il contrario, per ``zoomare'' sulla texture: con
0,5 la si vede ``la metà delle volte. Anzi, non la metà delle volte, scusate: un
quarto'' (0,5 su entrambi gli assi). E questo non vuol dire che una parte della
texture ``viene mangiata'': viene semplicemente ``applicata in uno spazio differente''.

\begin{figure}[htbp]
\centering
\begin{tikzpicture}[font=\small]
  \newcommand{\uetreTile}[5]{% x0, y0, nu, nv, etichetta
    \begin{scope}[shift={(#1,#2)}]
      \pgfmathsetmacro{\du}{2.2/#3}
      \pgfmathsetmacro{\dv}{2.2/#4}
      \foreach \i in {1,...,#3} \foreach \j in {1,...,#4} {
        \draw[fill=brown!25] ({(\i-1)*\du},{(\j-1)*\dv}) rectangle
          ({\i*\du},{\j*\dv});
        \draw[brown!70!black, thick] ({(\i-1)*\du+0.15*\du},{(\j-1)*\dv+0.2*\dv})
          -- ({(\i-1)*\du+0.85*\du},{(\j-1)*\dv+0.8*\dv});
      }
      \draw[very thick] (0,0) rectangle (2.2,2.2);
      \node[font=\footnotesize, align=center] at (1.1,-0.5) {#5};
    \end{scope}}
  \uetreTile{0}{0}{1}{1}{U = 1, V = 1}
  \uetreTile{3.0}{0}{2}{1}{U = 2, V = 1}
  \uetreTile{6.0}{0}{2}{2}{U = 2, V = 2}
  % zoom: un quarto della texture, ingrandito
  \begin{scope}[shift={(9.0,0)}]
    \draw[fill=brown!25] (0,0) rectangle (2.2,2.2);
    \draw[brown!70!black, thick] (0.66,0.88) -- (2.2,2.2);
    \draw[very thick] (0,0) rectangle (2.2,2.2);
    \node[font=\footnotesize, align=center] at (1.1,-0.5) {U = 0,5, V = 0,5};
  \end{scope}
\end{tikzpicture}
\caption{Il tiling sulla stessa faccia. Con UTiling a 2 la texture si ripete solo in
orizzontale; con entrambi a 2 si ripete quattro volte; con entrambi a 0,5 si vede un
quarto della texture, ingrandito: non è ``mangiata'', è applicata in uno spazio più
grande. Schema ricostruito dalla dimostrazione.}
\label{ue3:fig:tiling}
\end{figure}

\paragraph{Il pavimento.} ``Il tiling, però, va visto.'' Nella mappa c'è un
\textbf{pavimento}, e Davide ne cambia il materiale: lo seleziona, e nei
\textbf{Details}, sotto \textbf{Materials}, al posto della \emph{material instance}
della griglia grigia di prototipazione\footnote{Nella trascrizione ``Materialistance
Prototype Grid Gray'': è l'istanza della griglia grigia dei progetti modello di
Unreal.} trascina il \textbf{master material} appena creato. Il risultato: la corteccia
è ``molto, molto grande'' (``io non so neanche qual è l'utilizzo di questa texture che
sto usando, ci ho preso Bark a caso, ma va bene così''). Davide porta il tiling a 3,
poi a 10: ``iniziamo a vederla più spesso, questa corteccia. Sembrano anche delle dune,
vagamente.'' Il pavimento è probabilmente un cubo allargato, con UV ``molto, molto
piccole'' rispetto alle sue dimensioni: per questo nella \emph{preview} sferica del
materiale la texture si ripete ``diecimila volte'', mentre sul pavimento ``abbiamo un
pavimento che è più dimensionato''.

\paragraph{Una domanda che apre il tema successivo.} Uno studente chiede: si possono
modificare U e V in base all'oggetto su cui è applicata la texture, per esempio 10 su un
pavimento e 0,5 su un altro oggetto? ``Sì, è possibile, non facendolo dallo stesso
materiale. Questo è il momento in cui entrano in gioco le \textbf{Material
Instance}.''

\paragraph{Approfondimento: il tiling in una formula.} Il Texture Coordinate restituisce
per ogni punto la sua coppia $(u, v)$; il Multiply la scala: $(u', v') = (k_u\,u,\;
k_v\,v)$. Una texture è definita sull'intervallo $[0,1]$ e si ripete fuori da lì (è il
comportamento di default, \emph{wrap}): quando $u'$ va da 0 a 2, la texture viene letta
due volte. Con $k < 1$ la faccia ``vede'' solo un pezzo dell'intervallo, e quel pezzo
viene ingrandito. Per questo il tiling dipende dalle UV della mesh: se le UV di un
pavimento enorme occupano comunque tutto il quadrato $[0,1]^2$, una texture con tiling
1 si stende una volta sola su tutta la superficie, ed è ``molto, molto grande''.
Moltiplicare per una costante e usare UTiling e VTiling nel nodo sono due modi di fare
la stessa cosa; il primo ha il vantaggio che la costante può diventare un parametro.

% ---------------------------------------------------------------------
\section{Le Material Instance}
\label{ue3:sec:instance}

Davide torna alle slide (``passiamo un attimo alle Material Instance, visto che ne
stavamo parlando''). Quello costruito finora è un \textbf{master material}: tutto ciò
che sta nel grafo. Se si modifica il grafo, si modifica il materiale. Per avere lo
\textbf{stesso materiale con una piccola differenza} si possono creare due
\textbf{istanze} e cambiare in ciascuna il parametro collegato al Texture Coordinate.

\paragraph{Che cos'è un'istanza.} Una \textbf{Material Instance} è ``un'istanza di un
materiale che noi possiamo utilizzare per modificare quel materiale \textbf{senza
modificare la sorgente}''. Ha diversi vantaggi, ``oltre a essere \emph{bon ton}'', cioè
il modo in cui si dovrebbero applicare i materiali in generale: ``non dovreste prendere
il master material e applicarlo agli oggetti in scena. È sempre consigliabile applicare
un'istanza.''

\paragraph{Perché è più leggera: i tavoli.} L'istanza è ``molto più leggera''. Davide
lo spiega con i \textbf{tavolini} della prima lezione. Tre tavoli in scena con lo
stesso materiale del legno: per ogni tavolo l'engine carica la mesh e carica il
materiale, quindi ``tavolo, tavolo, tavolo'' e ``materiale, materiale, materiale''. Se
invece sui tavoli ci sono tre Material Instance, l'engine carica le tre mesh, ma trova
tre istanze che puntano \textbf{tutte allo stesso master material}, e ``va a
renderizzare un solo materiale''. ``Chiaramente, su tre tavoli non cambia niente. Ora
immaginate che invece di tre tavoli ci siano diecimila oggetti in scena'' e che mille
condividano lo stesso materiale: si capisce perché ``torna comodo avere un'istanza
piuttosto che avere un materiale generale''.

\begin{figure}[htbp]
\centering
\begin{tikzpicture}[
  font=\small,
  ue3tab/.style={draw, thick, minimum width=1.5cm, minimum height=0.7cm,
                 fill=brown!20, align=center},
  ue3mat/.style={draw, thick, rounded corners=3pt, minimum width=1.5cm,
                 minimum height=0.7cm, align=center},
  ue3arr/.style={-{Stealth[length=2mm]}, thick}
]
  % senza istanze
  \node[font=\footnotesize\bfseries] at (2.0,2.2) {come la racconta Davide: master
    material su ogni tavolo};
  \foreach \i in {0,1,2} {
    \node[ue3tab] (t\i) at (\i*2.0,1.2) {tavolo};
    \node[ue3mat, fill=blue!15] (m\i) at (\i*2.0,0) {M\_Legno};
    \draw[ue3arr] (t\i) -- (m\i);
  }
  \node[font=\footnotesize, align=center] at (2.0,-0.8) {tre mesh e ``materiale,
    materiale, materiale''};
  % con istanze
  \begin{scope}[shift={(7.6,0)}]
    \node[font=\footnotesize\bfseries] at (2.0,2.2) {con le Material Instance};
    \foreach \i in {0,1,2} {
      \node[ue3tab] (u\i) at (\i*2.0,1.2) {tavolo};
      \node[ue3mat, fill=orange!15, font=\footnotesize] (i\i) at (\i*2.0,0)
        {MI\_Legno\_\i};
      \draw[ue3arr] (u\i) -- (i\i);
    }
    \node[ue3mat, fill=blue!15] (mm) at (2.0,-1.3) {\textbf{M\_Legno}};
    \foreach \i in {0,1,2} \draw[ue3arr] (i\i) -- (mm);
    \node[font=\footnotesize, align=center] at (2.0,-2.1) {tre mesh, tre istanze, un
      solo materiale da compilare};
  \end{scope}
\end{tikzpicture}
\caption{L'esempio dei tavoli. A sinistra, il master material applicato direttamente;
a destra, tre istanze (che possono avere parametri diversi: un legno più chiaro, un
tiling più fitto) che rimandano tutte allo stesso master material. Schema ricostruito
dalla spiegazione; sul costo reale si veda l'approfondimento.}
\label{ue3:fig:tavoli}
\end{figure}

\paragraph{Come si crea.} ``Come lo facciamo su Unreal?'' Si fa \textbf{tasto destro}
sul materiale nel Content Browser e \textbf{Create Material Instance}, ``qui in
alto''. L'istanza appare \textbf{vuota}: ``non c'è un grafo per modificarlo, perché
chiaramente non posso modificare nessun grafo quando sono qui dentro''. Per avere
qualcosa da cambiare, bisogna estrarre dei \textbf{parametri} dal master material.
Davide torna nel grafo, fa tasto destro sul Texture Sample del colore, \textbf{Convert
to Parameter}, e lo chiama più o meno \emph{Diffuse Color}\footnote{Nella trascrizione
``discolor'': il nome esatto non si ricostruisce, ma è il parametro della texture del
colore.}; lo stesso per l'altro Texture Sample, chiamato \textbf{Normal Map}. Salvato il
materiale, nell'istanza compaiono \textbf{due parametri} selezionabili.

\paragraph{La nomenclatura.} ``Normalmente dovreste chiamarlo'' non \emph{New Material}
ma \textbf{M\_NewMaterial}, e l'istanza \textbf{MI\_NewMaterial}: ``per nomenclatura,
insomma, per ritrovare facilmente tutto quello che ci serve''.

\paragraph{Cambiare la texture dall'istanza.} Davide applica MI\_NewMaterial al
pavimento: è identico a prima. Ma se nell'istanza si cambia il parametro della texture,
da Bark a un'altra, il materiale cambia, senza toccare il grafo. Si potrebbe avere un
master material senza texture, ``delle texture assenti o base'', e tre o quattro
istanze in cui si cambiano le texture: un muro arancione, uno blu, uno rosso e uno verde
sono quattro istanze. ``Quando io renderizzo la scena, non andrò a chiamare cinque
materiali diversi, chiamerò sempre M\_NewMaterial, perché sono tutte istanze dello
stesso materiale.'' ``Non vorrei stare complicando troppo la cosa, però insomma, ditemi
voi se è chiaro il discorso.''

\paragraph{Il Tiling Multiplier.} Ed ecco la risposta alla domanda sul tiling diverso
per ogni oggetto. Nel master material Davide prende un \textbf{Multiply}, ci collega il
Texture Coordinate, collega il Multiply alle UV dei due Texture Sample e attacca alla B
una \textbf{costante}, che converte in parametro con il nome \textbf{Tiling
Multiplier}. Nell'istanza, il Tiling Multiplier portato a 10 dà ``lo stesso risultato di
prima''; ma ora è un parametro. Una copia dell'istanza con il Tiling Multiplier a 5
darebbe un tiling diverso con lo stesso materiale. ``È chiaro il discorso?''
``Perfetto.''

\begin{figure}[htbp]
\centering
\begin{tikzpicture}[
  font=\small,
  ue3node/.style={draw, thick, rounded corners=4pt, align=center,
                  minimum height=0.9cm},
  ue3uv/.style={-{Stealth[length=2mm]}, thick, red!60!black},
  ue3flt/.style={-{Stealth[length=2mm]}, thick, green!50!black},
  ue3col/.style={-{Stealth[length=2mm]}, thick, blue!60!black}
]
  \node[ue3node, fill=red!12, text width=1.9cm] (tc) at (0,0.7)
    {\textbf{TexCoord}\\\footnotesize (U + clic)};
  \node[ue3node, fill=green!8, text width=1.9cm] (tm) at (0,-0.9)
    {\textbf{Tiling}\\\textbf{Multiplier}\\\footnotesize parametro};
  \node[ue3node, fill=gray!10, text width=1.5cm] (mul) at (3.0,0)
    {\textbf{Multiply}};
  \node[ue3node, fill=orange!12, text width=2.4cm] (ts1) at (6.2,1.0)
    {\textbf{Diffuse Color}\\\footnotesize Texture Sample,\\
    \footnotesize parametro};
  \node[ue3node, fill=orange!12, text width=2.4cm] (ts2) at (6.2,-1.0)
    {\textbf{Normal Map}\\\footnotesize Texture Sample,\\\footnotesize parametro};
  \node[ue3node, fill=blue!6, text width=2.2cm, minimum height=2.6cm] (res)
    at (10.0,0) {\textbf{M\_NewMaterial}\\[4pt]\footnotesize Base Color\\[14pt]
    \footnotesize Normal};
  \draw[ue3uv] (tc.east) -- ++(0.4,0) |- ([yshift=2mm]mul.west);
  \draw[ue3flt] (tm.east) -- ++(0.4,0) |- ([yshift=-2mm]mul.west);
  \draw[ue3uv] (mul.east) -- ++(0.5,0) |- node[pos=0.8, above, font=\scriptsize]{UVs}
    (ts1.west);
  \draw[ue3uv] (mul.east) -- ++(0.5,0) |- node[pos=0.8, below, font=\scriptsize]{UVs}
    (ts2.west);
  \draw[ue3col] (ts1.east) -- ++(0.6,0) |- ([yshift=6mm]res.west);
  \draw[ue3col] (ts2.east) -- ++(0.6,0) |- ([yshift=-6mm]res.west);
  % istanze
  \node[ue3node, fill=yellow!15, text width=3.6cm, font=\footnotesize] (mi1)
    at (3.0,-2.9) {\textbf{MI\_NewMaterial}\\Tiling Multiplier = 10\\Diffuse = T\_Bark};
  \node[ue3node, fill=yellow!15, text width=3.6cm, font=\footnotesize] (mi2)
    at (8.0,-2.9) {\textbf{MI\_NewMaterial\_2}\\Tiling Multiplier = 5\\
    Diffuse = un'altra texture};
  \draw[dashed, thick] (mi1.north) -- ++(0,0.4) -| (res.south);
  \draw[dashed, thick] (mi2.north) -- (mi2.north |- res.south);
\end{tikzpicture}
\caption{Il master material con tre parametri (le due texture e il Tiling Multiplier)
e due istanze che li impostano in modo diverso. Le istanze non hanno un grafo: hanno
solo le caselle dei parametri. Ricostruzione dalla dimostrazione.}
\label{ue3:fig:instance}
\end{figure}

\paragraph{Il vantaggio principale.} Ricapitolando dalle slide: le Material Instance
permettono di \textbf{personalizzare velocemente} le superfici ``senza andare a
modificare il grafico che avete fatto'', e le modifiche si vedono
\textbf{immediatamente}, ``senza dover aspettare la ricompilazione del materiale''.

\paragraph{Approfondimento: che cosa si risparmia davvero.} La spiegazione dei tavoli
va presa come un'immagine. Il costo principale di un materiale non è ``caricarlo'' una
volta per oggetto, ma \textbf{compilarlo}: ogni master material diventa uno o più
\emph{shader}, programmi per la scheda video che Unreal deve compilare (e che fanno
comparire, nella prima apertura di un progetto, il famoso contatore ``Compiling
Shaders''). Un'istanza \textbf{non} genera un nuovo shader: riusa quello del genitore e
cambia solo i valori dei parametri. Cento istanze dello stesso master material sono
quindi un solo shader con cento ``schede'' di valori; cento materiali diversi sono cento
shader. Per questo cambiare un parametro in un'istanza è istantaneo, mentre cambiare un
nodo nel master material richiede la ricompilazione, come dice la slide. Le istanze del
Content Browser si chiamano \textbf{Material Instance Constant}; esistono anche le
\textbf{Material Instance Dynamic}, create durante il gioco da un Blueprint (nodo
\emph{Create Dynamic Material Instance}) per cambiare i parametri a \emph{runtime}, per
esempio con \emph{Set Scalar Parameter Value}: un materiale che si accende quando la
torretta di Igor ha un bersaglio nel mirino.

% ---------------------------------------------------------------------
\section{Le Material Function}
\label{ue3:sec:function}

``Per chiudere il discorso materiali'', l'ultima voce delle slide: le \textbf{Material
Function}, che oggi non si useranno. Sono funzioni che raccolgono ``una serie di calcoli
un po' più complessi'': invece di rifare ogni volta dei calcoli che si fanno spesso, li
si mette dentro una funzione e si usa la funzione ``a vostra discrezione''. ``È come se
fossero delle \textbf{parentesi}: voi prendete queste parentesi, le spostate in giro. È
un pacchetto che potete portarvi in giro.'' Servono a velocizzare lo sviluppo quando ci
sono cose da ripetere, e possono essere usate e condivise ``fra membri del team, fra
materiali diversi''.

\paragraph{Approfondimento: Material Function e Material Instance.} Le due cose
risolvono problemi diversi. L'istanza riusa \emph{un intero materiale} cambiandone i
valori; la funzione riusa \emph{un pezzo di grafo} dentro materiali diversi. Unreal
stesso è pieno di Material Function già pronte (si trovano nel menu dei nodi con
l'icona della funzione), e i nodi World Aligned Texture e World Aligned Normal che
compaiono alla fine della lezione (sezione~\ref{ue3:sec:world}) sono, appunto, Material
Function del motore.

% ---------------------------------------------------------------------
\section{Il Lerp: interpolare}
\label{ue3:sec:lerp}

Restano ``un paio di nodi che possono essere carini''. Davide ricapitola quelli visti
(la costante singola, il colore, il Texture Sample, le operazioni, il Texture
Coordinate) e ne aggiunge due: il \textbf{Lerp} e il \textbf{Panner}. ``Se abbiamo
capito la regola degli shortcut, questi saranno quale lettera più clic sinistro?''
``L.'' ``Esatto'': \textbf{L + clic} per il Lerp, \textbf{P + clic} per il Panner. ``Volevo
essere sicuro che era arrivata la cosa.''

\paragraph{Due colori.} Il \textbf{Lerp} si occupa dell'\textbf{interpolazione
lineare} (\emph{linear interpolation}). Per farlo vedere Davide usa due colori, ``una
roba che non farete mai, non con due colori perlomeno'': un \textbf{blu} nella
\textbf{A} e un \textbf{rosso} nella \textbf{B}. Il Lerp, collegato al Base Color, dà il
\textbf{viola}. Perché è inutile farlo con i colori? ``Perché potevo cercarmi il viola
direttamente dal colore.'' Ma mostra come funziona l'interpolazione. Il Lerp si può fare
anche con i numeri, ``per fare operazioni matematiche'', e anche con le texture.

\paragraph{L'alfa.} L'operazione è ``gestita interamente dall'\textbf{alfa}'', il terzo
ingresso. L'alfa va \textbf{da 0 a 1} e decide quale colore prevale, con A che vale 0 e
B che vale 1: con l'alfa a 0 si ha il blu, a 1 il rosso, e in mezzo ``una percentuale di
quel colore'', per esempio 0,7. Modificando l'alfa si va sempre più verso il rosso, o
sempre più verso il blu.

\paragraph{Una texture come alfa.} Ma l'alfa non deve essere un numero fisso. Davide
prende ``il nostro fidato Bark'' e collega \textbf{un suo canale} all'alfa: ``vi ricordate
il discorso dei canali sulle texture?'' Ora quella sequenza di 0 e di 1, di bianco e di
nero, ``è l'\textbf{arbitro} nel decidere dove ci sarà più A e dove ci sarà più B'':
nelle zone più nere il Lerp mette più blu, in quelle più bianche più rosso. Questa
texture è ``estremamente bianca'', quindi il risultato tende al rosso; con altre
texture, o lavorando su questa, si otterrebbero risultati diversi. E collegando il
Multiply delle coordinate anche al Texture Sample dell'alfa, si ottiene un tiling
anche dell'alfa, cioè di dove va il blu e dove il rosso. ``Chiaro questo processo?''
``Sì.''

\begin{figure}[htbp]
\centering
\begin{tikzpicture}[
  font=\small,
  ue3node/.style={draw, thick, rounded corners=4pt, align=center,
                  minimum height=0.9cm},
  ue3col/.style={-{Stealth[length=2mm]}, thick, blue!60!black},
  ue3flt/.style={-{Stealth[length=2mm]}, thick, green!50!black}
]
  \node[ue3node, fill=blue!30, text width=1.9cm] (a) at (0,1.4) {\textbf{blu}\\
    \footnotesize Constant3Vector};
  \node[ue3node, fill=red!30, text width=1.9cm] (b) at (0,0) {\textbf{rosso}\\
    \footnotesize Constant3Vector};
  \node[ue3node, fill=gray!10, text width=1.9cm] (al) at (0,-1.5)
    {\textbf{Alpha}\\\footnotesize 0,7 oppure\\\footnotesize un canale di T\_Bark};
  \node[ue3node, fill=gray!10, text width=1.6cm, minimum height=2.2cm] (lerp)
    at (3.6,0) {\textbf{Lerp}\\\footnotesize (L + clic)\\[4pt]\footnotesize A\\
    \footnotesize B\\\footnotesize Alpha};
  \node[ue3node, fill=violet!30, text width=2.0cm] (res) at (7.0,0)
    {\textbf{Base Color}\\\footnotesize viola, più o\\\footnotesize meno rosso};
  \draw[ue3col] (a.east) -- ++(0.6,0) |- ([yshift=5mm]lerp.west);
  \draw[ue3col] (b.east) -- ([yshift=0mm]lerp.west);
  \draw[ue3flt] (al.east) -- ++(0.6,0) |- ([yshift=-5mm]lerp.west);
  \draw[ue3col] (lerp) -- (res);
  % la scala dell'alfa
  \begin{scope}[shift={(9.4,-0.6)}]
    \shade[left color=blue!60, right color=red!60] (0,0) rectangle (3.4,0.6);
    \draw[thick] (0,0) rectangle (3.4,0.6);
    \node[font=\footnotesize] at (0,-0.3) {0};
    \node[font=\footnotesize] at (3.4,-0.3) {1};
    \node[font=\footnotesize] at (1.7,-0.3) {Alpha};
    \node[font=\footnotesize] at (0,0.95) {A};
    \node[font=\footnotesize] at (3.4,0.95) {B};
    \draw[thick] (2.38,-0.05) -- (2.38,0.65);
    \node[font=\scriptsize] at (2.38,0.95) {0,7};
  \end{scope}
\end{tikzpicture}
\caption{Il Lerp della dimostrazione: con l'alfa a 0 esce A (il blu), a 1 esce B (il
rosso), in mezzo una miscela. Se l'alfa è un canale di una texture, ogni pixel ha la sua
miscela: blu dove la texture è nera, rosso dove è bianca.}
\label{ue3:fig:lerp}
\end{figure}

\paragraph{L'esercizio della pausa.} ``Provateci voi.'' Durante la pausa, chi vuole può
prendere \textbf{due texture}, collegarle a un Lerp e, giocando con l'alfa, vedere quale
parte delle due texture viene presa di più e quale di meno.

\paragraph{Approfondimento: la formula del Lerp.} L'interpolazione lineare tra A e B
con peso $\alpha$ è
\[
  \operatorname{lerp}(A, B, \alpha) = A\,(1-\alpha) + B\,\alpha .
\]
Con colori, il calcolo si fa canale per canale. Blu $=(0,0,1)$ e rosso $=(1,0,0)$ con
$\alpha = 0{,}5$ danno $(0{,}5;\,0;\,0{,}5)$, il viola della lezione; con $\alpha =
0{,}7$ danno $(0{,}7;\,0;\,0{,}3)$, un viola più rosso. Quando l'alfa è una texture, la
formula si applica pixel per pixel, ed è così che si mescolano due materiali su una
stessa superficie: erba e terra su un terreno, vernice e ruggine su un metallo, con una
maschera in bianco e nero a fare da ``arbitro''.

% ---------------------------------------------------------------------
\section{Il Panner: far muovere le texture}
\label{ue3:sec:panner}

Il \textbf{Panner}\footnote{Nella trascrizione ``panel'', ``planner'' e ``banner''.}
``si occupa di far muovere le texture''. Davide riprende il Base Color, stacca il resto,
e collega il Panner alle \textbf{UV} del Texture Sample, ``perché lui lavora sulle UV
della texture''. Selezionando il nodo (o direttamente sul nodo) si vedono
\textbf{Speed X} e \textbf{Speed Y}. Con Speed X a 1 la \emph{preview} comincia a
muoversi sull'asse X, ``anche un po' troppo veloce per i miei gusti''; con 0,3 in X e
0,1 in Y si ha un Panner che si muove su entrambi gli assi.

\paragraph{A che cosa serve.} ``Che materiale potremmo volerci fare con questo nodo?''
``Dell'acqua.'' ``Dell'acqua. Qualsiasi fluido, qualsiasi cosa che si muove.'' Il Panner
ha il parametro di velocità e l'ingresso delle \textbf{coordinate}: per combinarlo con il
tiling si collega il Texture Coordinate (o il Multiply) all'ingresso Coordinate. C'è
poi l'ingresso \textbf{Time}, con cui si può decidere ``quando e come deve muoversi'';
altrimenti, di default, usa il \emph{game time}, ``cioè è sempre in movimento''.

\paragraph{Il master material cambia, l'istanza segue.} Mentre il Panner lavora, il
pavimento nella scena \textbf{si muove}: sul pavimento c'è la Material Instance, e
``nonostante stiamo facendo delle modifiche al materiale originale e non alla Material
Instance, la Material Instance continua a modificarsi''. È il rovescio del rapporto tra
master e istanze: quello che cambia nel genitore cambia in tutte le istanze.

\paragraph{Approfondimento: il Panner in una formula.} Il Panner non fa altro che
sommare alle coordinate un tempo moltiplicato per una velocità:
\[
  (u', v') = (u + s_x\,t,\; v + s_y\,t),
\]
dove $t$ è il tempo in secondi e $s_x, s_y$ le velocità in ``texture al secondo'': con
Speed X a 1 la texture scorre di una larghezza intera ogni secondo, per questo era
troppo veloce. Siccome la texture si ripete, lo scorrimento è infinito e senza salti,
purché la texture sia \emph{tileable}, cioè con i bordi che combaciano. Due Panner con
velocità e tiling diversi sommati tra loro sono il trucco classico per l'acqua e per la
lava: il movimento non si ripete mai uguale.

% ---------------------------------------------------------------------
\section{La pausa: il progetto che non si trova, e Git}
\label{ue3:sec:pausa}

``Visto che comunque ci siamo fatti un'oretta e mezza'', Davide chiude per ora il
discorso materiali e propone una pausa di una decina di minuti. Al ritorno: l'import,
le mesh, e ``una cosina con la torretta che c'è nei file che vi abbiamo passato'', così
che Igor, quando torna, potrà farci ``un piccolo sistema di puntamento''.

\paragraph{Il progetto che Unreal non vede.} Al rientro, in attesa degli altri, una
studentessa racconta il suo problema: ha scaricato e decompresso la cartella del Drive,
ma aprendo Unreal il progetto non compare, e deve aprirlo ``a mano'' da un'altra parte.
Davide spiega che dipende probabilmente dal percorso in cui è salvata la cartella: se la
si mette dove sono salvati gli altri progetti, Unreal dovrebbe trovarla; in ogni caso,
aprire il progetto dal file eseguibile nella cartella funziona. ``Se non è un problema
che ti impedisce di accedere al progetto, fa niente.''

\paragraph{Come si lavora insieme a distanza?} Uno studente, programmatore, fa una
domanda ``fuori dall'argomento'': i programmatori usano \textbf{Git}, ma i file di un
progetto di Unreal sembrano troppo pesanti per lavorarci in contemporanea da lontano.
Davide conferma che \textbf{Git}, ``o diciamo repository di sorta'', è ``la cosa
principale che utilizzerai''. Con \textbf{Unity} ha sempre avuto meno problemi, perché
Unreal ``va gestito in maniera un po' più delicata'', ma entrambi supportano senza
problemi repository e \emph{source control}. Esistono anche altri metodi, che Davide non
ha mai provato: per Unreal, ``una cosa che io sento da una vita'', un modo per lavorare
\textbf{contemporaneamente} sulla stessa cosa, tramite \textbf{VPN} o in \textbf{LAN} e
con dei \emph{plugin} specifici, in modo che due colleghi modifichino attivamente lo
stesso livello. Con Git invece ``ognuno lavora sul suo e poi si \emph{merge} tutto
insieme''. ``Git è quello che ho utilizzato sempre, è il più comodo, quindi dovresti
trovarti abbastanza a casa.''

\paragraph{Approfondimento: Git LFS, Perforce e Multi-User Editing.} I timori dello
studente sono fondati. Git nasce per file di testo, e gli \emph{asset} di Unreal
(\texttt{.uasset}, \texttt{.umap}) sono file \textbf{binari}, spesso grandi, che Git non
sa fondere: se due persone modificano lo stesso livello, una delle due modifiche va
persa. Per questo con Unreal si usa \textbf{Git LFS} (\emph{Large File Storage}), che
tiene i file pesanti fuori dalla cronologia normale e permette di \textbf{bloccare}
(\emph{lock}) un file mentre lo si modifica. Nell'industria, per i progetti Unreal, è
molto diffuso \textbf{Perforce} (Helix Core), il sistema che Epic stessa usa e che
l'editor supporta direttamente, proprio per la gestione dei file binari e dei blocchi.
Il sistema ``sentito da una vita'' di cui parla Davide esiste: è il plugin
\textbf{Multi-User Editing}, con cui più editor connessi in rete (in LAN o via VPN)
condividono una stessa sessione e vedono le modifiche degli altri in tempo reale; nato
per la produzione virtuale nel cinema, si può usare anche nello sviluppo di giochi, ma
non sostituisce il \emph{source control}, a cui le modifiche vanno comunque consegnate
alla fine.

% ---------------------------------------------------------------------
\section{L'import e la torretta}
\label{ue3:sec:import}

Passato ``il nostro quarto d'ora'', si riparte: si importano le prime mesh, la
\textbf{torretta}, ``le tre parti della torretta, perché l'abbiamo divisa in tre
parti''.

\paragraph{La finestra di import, sulla slide.} Quando si importa una mesh compare una
finestra (``che può essere più larga, può essere più stretta'') piena di parametri. Le
cose principali che si possono fare da lì:
\begin{itemize}
  \item \textbf{generare automaticamente le collisioni}, se non sono presenti nel file
        FBX, perché magari gli oggetti non sono stati creati con le collisioni;
  \item \textbf{modificare il transform} dell'oggetto importato ``proprio di base'':
        ruotarlo sul suo \textbf{pivot}, traslarlo. ``Avete già parlato di cos'è il
        pivot con Igor, giusto? Il fulcro che vi permette di traslare e ruotare gli
        oggetti all'interno della scena.'' Con la funzione di \textbf{reimport}, si
        vedrà, questa possibilità diventa molto utile;
  \item adattare le \textbf{unità} e gli \textbf{assi} tra il software di origine e
        Unreal: \textbf{convertire la scena} dalle coordinate FBX a quelle di Unreal,
        \textbf{convertire le unità} (Unreal lavora in centimetri), oppure usare
        \textbf{Force Front X Axis}, che cambia l'asse che Unreal considera come
        ``fronte'' dell'oggetto. Davide l'ha usato poche volte: una volta il modello di
        un \emph{character}, importato normalmente, ``veniva rivoltato, come se il
        dentro fosse fuori'', e cambiare l'asse frontale gli ha restituito la mesh
        com'era stata esportata (``mi pare fosse Blender'');
  \item decidere se importare anche i \textbf{materiali} e le \textbf{texture}
        collegati alla mesh, e \textbf{dove} cercarli: in locale, in un'altra cartella,
        e così via.
\end{itemize}

\paragraph{Il pulsante Import.} Prima di importare, Davide spegne il Panner (``facciamo
che smetto di far muovere il pavimento'') e il pavimento torna normale. Nel Content
Browser c'è ``questo pulsantone con \textbf{Import}'', che apre l'esplora risorse del
computer per cercare i file. Nella cartella del materiale didattico c'è una cartella
\textbf{torretta}, con tre \textbf{static mesh}: \textbf{SM\_Turret\_Base},
\textbf{SM\_Turret\_Cannon} e \textbf{SM\_Turret\_Pivot}\footnote{Nella trascrizione
``Turret Base, Turret Cannon e Turret Pivo''; il prefisso SM\_ è detto a voce.}.

\paragraph{La finestra di import, dal vivo.} Importandole si apre la finestra (``nel
mio caso è così enorme, ma perché io tendo a usarla così''), e si ritrovano le opzioni
della slide: la possibilità di modificare traslazione e rotazione, ``anche
singolarmente sui vari assi'' (``questo è qualcosa che faremo più tardi''); i
\textbf{LOD} (``decidere il gruppo dei LOD''); le collisioni, se e come importarle. Se si
importa una \textbf{Skeletal Mesh}\footnote{Nella trascrizione ``scheda dal Mesh'' e
``Scheletal Mesh'', come nella seconda lezione.}, come un personaggio, la finestra
permette di collegarla a uno \textbf{scheletro} esistente, di generarne uno, oppure di
importare solo delle \textbf{animazioni}: ``faremo una lezione specifica sulle
animazioni, quindi per il momento glissiamo''. In alto ci sono anche le sezioni per
andare più veloci, tra cui \textbf{Import Materials} e \textbf{Import Textures}, con la
\textbf{Search Location} in cui cercarli.

\paragraph{Niente materiali e texture, per ora.} La consegna di Davide: \textbf{togliere
la spunta} a Import Materials e Import Textures, e fare \textbf{Import}. Unreal importa i
tre oggetti (Davide elimina quelli della prova precedente e salva). Perché togliere le
spunte? Se si fossero importati gli oggetti con materiali e texture, Unreal avrebbe
importato automaticamente tutte le texture, creato i materiali adatti e li avrebbe
anche \textbf{applicati}, ``in modo da averli già precisi sulla torretta''. Ed è
proprio quello che si farà, ma \textbf{a mano}, ``così vediamo effettivamente anche il
processo''. ``Ci sono stati problemi? Tutto a posto con la torretta?'' ``Sì, nessun
problema.''

\begin{table}[htbp]
\centering
\small
\begin{tabular}{@{}p{4.2cm}p{8.8cm}@{}}
\toprule
\textbf{Opzione} & \textbf{A che cosa serve} \\
\midrule
Generate Missing Collision & genera le collisioni se il file FBX non le ha \\
Import Translation / Rotation / Uniform Scale & sposta, ruota, scala l'oggetto ``alla
  radice'': cambia dove sta il pivot rispetto alla geometria \\
LOD Group & sceglie un gruppo di livelli di dettaglio (\emph{Level Of Detail}) \\
Convert Scene, Convert Scene Unit & adattano assi e unità di misura del file FBX a
  quelli di Unreal (centimetri) \\
Force Front X Axis & cambia l'asse considerato come ``fronte'' \\
Skeleton / Import Animations & per le Skeletal Mesh: scheletro esistente o nuovo, solo
  animazioni \\
Import Materials, Import Textures & importano e applicano materiali e texture
  collegati \\
Search Location & dove cercare materiali e texture da importare \\
\bottomrule
\end{tabular}
\caption{Le opzioni della finestra di import citate a lezione, con i nomi della
finestra di Unreal Engine 5 (FBX Import Options; il raggruppamento può variare tra le
versioni). Per la torretta: Import Materials e Import Textures senza spunta.}
\label{ue3:tab:import}
\end{table}

% ---------------------------------------------------------------------
\section{Che cos'è una mesh}
\label{ue3:sec:mesh}

Mentre la torretta è importata, Davide torna sulle slide per ``le ultime cose sulle
mesh''. Sulla slide c'è un \textbf{vaso} in \emph{wireframe} e, sotto, una
semplificazione. Una mesh è ``un insieme di \textbf{vertici}, \textbf{spigoli} e
\textbf{facce} che vanno a creare un oggetto in tre dimensioni''.
\begin{itemize}
  \item I \textbf{vertici} sono i \textbf{punti} che costituiscono gli angoli
        dell'oggetto, punti nello spazio.
  \item Gli \textbf{spigoli} (\emph{edge}) sono i \textbf{segmenti} che collegano due
        vertici.
  \item Le \textbf{facce} sono le \textbf{superfici piatte} definite da più segmenti:
        sulla slide, ``tutti i triangoli''.
\end{itemize}
Il \textbf{triangolo} è ``la forma più comoda che esiste a livello di calcolo'', e
tendenzialmente tutti gli oggetti saranno composti da ``un insieme di triangolini''. Più
triangoli insieme, più facce, formano dei \textbf{poligoni}; i poligoni combinati formano
delle \textbf{superfici}; e la superficie determina la mesh. Tutto questo si può vedere
anche nell'engine, con la visualizzazione \textbf{wireframe}: l'oggetto non ``come
appare'', ma ``come griglia, come insieme di tutti i poligoni, vertici e spigoli''.

\begin{figure}[htbp]
\centering
\begin{tikzpicture}[font=\small]
  \coordinate (A) at (0,0);
  \coordinate (B) at (2.6,0.3);
  \coordinate (C) at (3.2,2.2);
  \coordinate (D) at (0.6,2.6);
  \coordinate (E) at (1.7,1.3);
  \fill[orange!15] (A) -- (B) -- (E) -- cycle;
  \fill[orange!30] (B) -- (C) -- (E) -- cycle;
  \fill[orange!20] (C) -- (D) -- (E) -- cycle;
  \fill[orange!10] (D) -- (A) -- (E) -- cycle;
  \draw[very thick] (A) -- (B) -- (C) -- (D) -- cycle;
  \draw[very thick] (A) -- (E) (B) -- (E) (C) -- (E) (D) -- (E);
  \foreach \p in {A,B,C,D,E} \fill (\p) circle (2.5pt);
  \draw[-{Stealth[length=2mm]}] (4.3,2.8) node[right, font=\footnotesize]
    {\textbf{vertice}: un punto} -- (C);
  \draw[-{Stealth[length=2mm]}] (4.3,1.6) node[right, font=\footnotesize]
    {\textbf{spigolo} (\emph{edge}): un segmento tra due vertici} -- ($(B)!0.5!(C)$);
  \draw[-{Stealth[length=2mm]}] (4.3,0.4) node[right, font=\footnotesize]
    {\textbf{faccia}: un triangolo} -- (1.45,0.55);
  \node[font=\footnotesize, align=center, text width=4cm] at (1.6,-0.8)
    {4 facce triangolari\\formano un poligono};
\end{tikzpicture}
\caption{Vertici, spigoli e facce, come nella semplificazione della slide. Il triangolo
è l'unità di calcolo: anche le facce a quattro lati dei programmi di modellazione,
quando arrivano in Unreal, vengono divise in triangoli.}
\label{ue3:fig:mesh}
\end{figure}

% ---------------------------------------------------------------------
\section{Dimensioni, unità di misura e formati}
\label{ue3:sec:regole}

Quando si porta una mesh in Unreal, mentre la si produce o mentre la si cerca, bisogna
rispettare alcune cose.

\paragraph{Le dimensioni.} La cosa più comune: si decide di fare ``una mappa molto
carina'' su Unreal, si cercano mesh su internet o su \textbf{Fab} (``su Fab questo
problema c'è meno, però utilizzando siti esterni può capitare''), e la mesh importata è
\textbf{gigantesca} o \textbf{incredibilmente piccola}, perché chi ha fatto l'export, o
chi ha fatto l'import, non ha rispettato ``le misure di adattamento e transizione da un
software all'altro''. Unreal usa come unità di misura i \textbf{centimetri}; Blender,
per esempio, ``usa altra roba''. Nel passaggio da Blender a Unreal bisogna controllare
che la conversione delle dimensioni sia corretta.

\paragraph{FBX.} ``Altra cosa molto importante'': Unreal ``tendenzialmente legge
soltanto \textbf{FBX}''. Sono supportati anche altri formati, come gli \textbf{OBJ}, ``ma
non assicurano un corretto funzionamento all'interno dell'engine''. Il consiglio:
``importate sempre unicamente FBX''.

\paragraph{I formati delle texture.} Per le texture i formati supportati sono molti
(sulla slide c'è un elenco), ma nella pratica si usano soprattutto \textbf{JPEG} e
\textbf{PNG}, ``i più comuni, a meno che non state facendo qualcosa di molto
specifico''.

\paragraph{Approfondimento: Blender, i metri e il fattore 100.} Blender lavora di default
in metri, con una ``unità'' uguale a un metro; Unreal lavora in centimetri. Un cubo di
lato 2 in Blender è quindi un cubo di 2~m, cioè 200 unità di Unreal: se l'export non
converte, l'oggetto arriva cento volte troppo piccolo (o, con altre impostazioni, cento
volte troppo grande). La soluzione standard è nell'esportatore FBX di Blender
(l'opzione \emph{Apply Scalings} su \emph{FBX Units Scale}) o, lato Unreal, nella
scala uniforme e nella conversione delle unità della finestra di import. Un controllo
rapido: un personaggio umano in Unreal è alto circa 180 unità, come il manichino dei
progetti modello; se la mesh appena importata è molto diversa, c'è un problema di scala.
Oltre a FBX, le versioni recenti di Unreal Engine 5 importano anche \textbf{glTF} e
\textbf{USD}, ma FBX resta il formato più collaudato, come dice Davide.

% ---------------------------------------------------------------------
\section{Il pivot del cannone e il Reimport}
\label{ue3:sec:pivot}

Davide apre le tre mesh (la base, la torretta e il pivot) e le piazza in scena. Sono
state esportate con il \textbf{pivot nello stesso punto}: per sovrapporle senza
spostarle a mano basta \textbf{copiare la X e la Y} di un pezzo negli altri. Ed ecco la
torretta, ``a parte il fatto che non c'ha il materiale''.

\paragraph{I problemi della torretta.} Ma ``sta torretta ha dei problemi: quali sono?''
Il primo si vede: per qualche motivo \textbf{non è allineata}, ed è ``scomodamente''
ruotata: trovare la rotazione giusta sarà scomodo. Il secondo è più sottile. Si
supponga di mettere i tre pezzi insieme in un \textbf{Blueprint}. Che cosa deve fare una
torretta? Deve \textbf{tracciare dei bersagli}: ruotare su sé stessa, ma anche andare su
e giù, se il bersaglio è in alto o in basso. E qui c'è il problema: il \textbf{pivot del
cannone} non è in una posizione comoda per la rotazione verticale, perché sta
\textbf{in basso}, in fondo. Ruotare sul piano orizzontale non dà problemi, ma per
alzare e abbassare la canna il pivot dovrebbe essere ``più o meno all'altezza di questo
\textbf{bullone}'', il fulcro, e da lì il cannone dovrebbe essere manipolato e ruotato
(figura~\ref{ue3:fig:pivot}).

\begin{figure}[htbp]
\centering
\begin{tikzpicture}[font=\small,
  ue3arr/.style={-{Stealth[length=2.5mm]}, thick}]
  \foreach \xs/\titolo/\px/\py/\ang in {0/{pivot in basso (com'era)}/1.25/0.0/0,
                                         6.6/{pivot nel fulcro (dopo il Reimport)}/1.25/1.75/0} {
    \begin{scope}[shift={(\xs,0)}]
      % base
      \draw[thick, fill=orange!35] (0.3,0) -- (2.2,0) -- (1.9,0.5) -- (0.6,0.5) -- cycle;
      % pivot (supporto)
      \draw[thick, fill=gray!30] (0.8,0.5) rectangle (1.7,1.5);
      % cannone
      \draw[thick, fill=orange!20] (0.7,1.5) rectangle (1.8,2.1);
      \draw[thick, fill=orange!20] (1.8,1.6) rectangle (4.0,1.95);
      % bullone
      \fill[black] (1.25,1.75) circle (2pt);
      % il pivot del cannone
      \draw[red, very thick] (\px-0.18,\py) -- (\px+0.18,\py) (\px,\py-0.18) -- (\px,\py+0.18);
      \draw[red, thick] (\px,\py) circle (0.18);
      \node[font=\footnotesize, align=center] at (2.0,-0.6) {\titolo};
    \end{scope}
  }
  % archi di rotazione
  \draw[ue3arr, red] (3.4,0.5) arc[start angle=-20, end angle=25, radius=2.6];
  \node[font=\scriptsize, red, align=center] at (4.6,0.55) {la canna ruota\\attorno al
    piede};
  \draw[ue3arr, red] (6.6+3.6,1.0) arc[start angle=-30, end angle=30, radius=1.6];
  \node[font=\scriptsize, red, align=center] at (6.6+4.75,1.75) {la canna ruota\\
    attorno al bullone};
  % assi
  \draw[ue3arr, blue] (-0.6,0) -- (-0.6,1.2) node[above, font=\scriptsize]{Z};
\end{tikzpicture}
\caption{Il problema del pivot del cannone, visto di lato (schema non in scala). Con il
pivot (il mirino rosso) in fondo, alzare la canna la farebbe girare attorno al piede,
staccandola dal supporto; con il pivot nel fulcro, all'altezza del bullone, la canna
si alza e si abbassa come una vera torretta. La rotazione orizzontale, attorno all'asse
Z, funziona invece in entrambi i casi.}
\label{ue3:fig:pivot}
\end{figure}

\paragraph{Come si sposta un pivot.} ``In realtà possiamo farlo in tutta una miriade di
modi.'' Con la \textbf{Modeling Mode} si potrebbe rimodellare l'oggetto. Ma lo si può
fare anche \textbf{nell'import}: si fa tasto destro sulla mesh e \textbf{Reimport Mesh
With Dialog}, che riapre la finestra di import vista prima. Unreal importa lo stesso
oggetto e lo \textbf{sovrascrive}, ma permette di cambiare come verrà posizionato:
l'\textbf{offset della rotazione} e l'\textbf{offset della traslazione}.

\paragraph{I numeri.} Davide se li era scritti, ma non li ricorda, e li ricava. Per
sapere quale asse toccare, guarda il \emph{gizmo}: il cannone deve essere allineato con
il resto della torretta (e lo stesso il pivot), e il pivot del cannone deve stare non
``qui giù in fondo'', ma allineato con il fulcro, sull'\textbf{asse Z}, ``perché è
questo blu''. Per questa torretta, ``quindi anche per voi'', i parametri sono:
\begin{itemize}
  \item per il \textbf{cannone}: traslazione Z $= -135$, che reimporta il cannone ``più
        attaccato al pivot'', e rotazione $= 113$, che lo riallinea\footnote{La
        trascrizione dice prima ``113 per la traslazione'' e poi si corregge: ``e 113
        sulla rotazione''. L'asse della rotazione non è detto; trattandosi di
        riallineare la torretta nel piano orizzontale, è con ogni probabilità la
        rotazione attorno a Z (il \emph{yaw}).};
  \item per il \textbf{pivot}: ``questa volta non serve modificare la traslazione'', ma
        solo la rotazione, sempre 113.
\end{itemize}
Dopo il reimport il cannone ``sta sotto'' (la sua origine ora è nel fulcro, quindi
nella posizione di prima risulta più basso) e va rimesso a posto; ma ora, ruotandolo, si
alza e si abbassa nel modo giusto. ``E chiaramente poi questo farete più avanti, con dei
Blueprint.'' Si salva tutto.

\paragraph{Approfondimento: perché il pivot conta così tanto.} In Unreal ogni rotazione
e ogni scala di un componente avvengono \textbf{attorno alla sua origine}, cioè al
pivot della mesh. Un Blueprint di torretta, di solito, ha la base come componente
radice, il supporto come figlio che ruota attorno a Z (il \emph{yaw}, per seguire il
bersaglio a destra e a sinistra) e il cannone come figlio del supporto che ruota attorno
a Y (il \emph{pitch}, su e giù). Per questo ogni pezzo deve avere l'origine esattamente
nel punto in cui gira nella realtà: se l'origine è sbagliata, nessun nodo di rotazione
potrà correggerla senza calcoli aggiuntivi. Il ``piccolo sistema di puntamento'' che
Igor farà con questa torretta (con nodi come \emph{Find Look at Rotation}, che calcola
la rotazione necessaria per guardare un punto) si appoggia interamente a questo lavoro
preparatorio.

% ---------------------------------------------------------------------
\section{Mettere ordine: le cartelle}
\label{ue3:sec:cartelle}

Prima di importare le \textbf{texture}, ``per dare un po' di ordine'', Davide crea nel
Content Browser una nuova cartella, \textbf{Torretta}, ci sposta dentro tutto
(``muoviamo, così è un po' più pulito''), e dentro crea un'ulteriore cartella,
\textbf{Texture}, in cui importare le texture della torretta. ``Questo non è
obbligatorio, è solo una questione di \emph{best practice}, cioè è come il discorso
delle Material Instance. Non vi vengo a trovare a casa se non lo fate, però vi risulta
tutto più comodo.''

% ---------------------------------------------------------------------
\section{Lo Static Mesh Editor e le collisioni}
\label{ue3:sec:sme}

Con un doppio clic su una mesh si apre lo \textbf{Static Mesh Editor}: ``andiamoci a
vedere effettivamente un po' qui dentro cosa possiamo vedere delle mesh''.

\paragraph{Gli slot dei materiali.} A destra, nei \textbf{Material Slots}, queste mesh
hanno \textbf{due slot}: ``non è la normalità''. (Una delle tre ne ha uno solo.) Il modo
in cui la torretta è stata suddivisa fa sì che, per rendere bene i materiali, servano
due slot: la base e il cannone condividono in parte il materiale del \textbf{pivot}. Nei
due slot si inseriranno i materiali più avanti. Altre impostazioni (Davide accenna a
quelle sui LOD e sulla \emph{Nanite}\footnote{La trascrizione dice ``per quanto riguarda
la night'': la ricostruzione come Nanite, la tecnologia di geometria virtualizzata di
Unreal Engine 5 che ha una sua sezione nei dettagli della mesh, è probabile ma non
certa.}) non si vedono oggi.

\paragraph{Le modalità di visualizzazione.} Dal menu della \emph{viewport} (la
prospettiva ``non ci interessa'') si cambia la modalità: \textbf{Lit}, quella normale;
\textbf{Unlit}, ``com'è il mio oggetto se non è colpito dalla luce?''; e
\textbf{Wireframe}: ``vedete quello che dicevamo prima?'' Vertici, facce ed \emph{edge}
della mesh, che si potrebbero modificare con la \textbf{Modeling Mode} (``se avanzerà
il tempo con le lezioni, sarebbe carino vederla più avanti''). Ci sono anche
\textbf{Lighting Only}, le \textbf{Reflections} e altre ancora. Dal menu \textbf{Show}
si possono mostrare i \textbf{vertici}: ``dove sono i vertici di questo oggetto?''

\paragraph{Simple e complex.} Un'altra cosa utile ``da un punto di vista di gameplay''
sono le \textbf{collisioni}. Dal menu Show si attiva la \textbf{Simple Collision}: ``queste
sono le collisioni del mio oggetto''. In scena, il \emph{character} interagirà con
queste: ci salterà sopra, o lo dovrà scavalcare. Ma oltre alle \textbf{simple
collision} ci sono le \textbf{complex collision}. La differenza:
\begin{itemize}
  \item le \textbf{simple collision} sono forme che si possono posizionare a mano,
        anche da qui, e rifare da capo;
  \item le \textbf{complex collision} usano la struttura in vertici della mesh stessa
        come collisione. ``È qualcosa che è molto più pesante rispetto a una collisione
        semplice, e torna utile solo in alcuni casi specifici.''
\end{itemize}

\paragraph{Modificare le collisioni.} La collisione semplice si può selezionare,
spostare, anche eliminare; e se ne possono \textbf{aggiungere} altre: una
\textbf{sfera}, di cui si decide la grandezza; una \textbf{capsula}, che ``è
tendenzialmente la collisione che usano i vostri \emph{character}'': il personaggio di
solito non ha una collisione super precisa, ``è un capsulone'', a meno che non si voglia
davvero il contrario; una \textbf{box}. E una box, per questa mesh, ``è più ottimizzata
rispetto alla collisione che questa aveva di base'': all'atto pratico fa la stessa cosa
``con molti meno vertici''. Su una singola mesh cambia poco; su molte mesh aiuta le
prestazioni. Ci sono anche collisioni più complesse, ma ``idealmente, il modo migliore
per fare delle simple collision è letteralmente prendere delle \textbf{forme base}'',
magari più d'una, ``piuttosto che avere tanti vertici sulla stessa collisione''.

\begin{table}[htbp]
\centering
\small
\begin{tabular}{@{}p{3.3cm}p{4.4cm}p{5.3cm}@{}}
\toprule
\textbf{Collisione} & \textbf{Com'è fatta} & \textbf{Quando} \\
\midrule
Box & un parallelepipedo & oggetti squadrati; spesso basta anche per oggetti più
  complessi, ed è la più leggera \\
Sphere & una sfera & oggetti tondi, palle \\
Capsule & un cilindro con due semisfere & i \emph{character}: ``un capsulone'' \\
Convex (più forme) & poliedri convessi, magari più d'uno & quando una forma sola non
  basta; meglio poche forme base che tanti vertici \\
Complex & i triangoli della mesh stessa & casi specifici: precisione massima, costo
  molto più alto \\
\bottomrule
\end{tabular}
\caption{Le collisioni di una static mesh viste nello Static Mesh Editor. Le prime
quattro sono \emph{simple collision}, che si aggiungono e si modificano dal menu
Collision dell'editor; l'ultima è la \emph{complex collision}.}
\label{ue3:tab:collisioni}
\end{table}

\paragraph{Approfondimento: chi usa quale collisione.} Unreal tiene le due collisioni
separate perché le usa per cose diverse. La \textbf{simple} è quella della
\textbf{fisica}: la pallina della seconda lezione che rotola sulla piattaforma, il
\emph{character} che cammina e salta, gli oggetti che cadono; i calcoli fisici su forme
semplici sono molto più veloci. La \textbf{complex} si usa soprattutto per le
\textbf{interrogazioni} precise, come un \emph{line trace} (un ``raggio'' lanciato per
sapere che cosa colpisce un proiettile o che cosa sta guardando il giocatore), quando
la precisione al triangolo serve davvero. Nei dettagli della mesh, il campo
\emph{Collision Complexity} decide come combinarle; l'opzione \emph{Use Complex
Collision As Simple} usa i triangoli anche per la fisica, ed è comoda per oggetti
statici di forma strana, ma va usata con giudizio, per il costo di cui parla Davide.

% ---------------------------------------------------------------------
\section{Raggruppare e fondere: Merge Actors}
\label{ue3:sec:merge}

``Se ad esempio vogliamo avere tutti e tre gli oggetti assieme'', o semplicemente
spostarli insieme, ci sono due strade.

\paragraph{Raggruppare.} Si selezionano tutti e tre e li si \textbf{raggruppa} con
\textbf{Ctrl+G}; si spostano insieme; con \textbf{Shift+G} si toglie il gruppo.

\paragraph{Fondere.} Oppure li si \textbf{fonde} (\emph{merge}). Si evidenziano i tre
pezzi in scena, si va nel menu \textbf{Actor} in alto e si sceglie \textbf{Merge
Actors}. Unreal chiede dove salvare il nuovo oggetto e con che nome (per esempio
\textbf{SM\_Turret}\footnote{Nella trascrizione ``smTourette''.}), e crea una
\textbf{nuova mesh}, ``che è semplicemente la fusione delle tre mesh precedenti'': una
somma, un \emph{merge}, appunto.

\paragraph{Il merge, di nuovo.} Più tardi uno studente dice di essersi perso come si fa
il merge, e Davide lo rifà con calma. I tre oggetti (il cannone, la base e il pivot)
sono in scena. Con il \textbf{tasto destro} su di essi ci sono delle opzioni, ``però
queste sono opzioni limitate''. Per averle tutte si usa il menu \textbf{Actor} in alto:
mostra gran parte delle stesse opzioni, ma in basso, nella sezione degli strumenti di
Unreal Engine, c'è \textbf{Merge Actors}. Al merge Unreal chiede dove salvare e con che
nome, e ``ti fonde i tre oggetti''. Nello stesso menu c'è anche \textbf{Chaos}, che
``serve per gli oggetti distruttibili'': cose più avanzate, che probabilmente non si
vedranno nel corso, ma da approfondire ``se vi capiterà di voler mettere distruttibilità
all'interno del vostro gioco''.

\paragraph{Approfondimento: gruppo, merge o Blueprint?} Le tre soluzioni non sono
equivalenti. Il \textbf{gruppo} è solo una comodità dell'editor: gli actor restano
separati, e in gioco non cambia nulla. Il \textbf{merge} crea una mesh nuova, una sola:
meno oggetti da disegnare (utile per le prestazioni, per esempio con tanti pezzi di
arredo statici), ma i pezzi \textbf{non si possono più muovere uno rispetto
all'altro}. Per la torretta di Igor, che deve ruotare il supporto e alzare la canna, il
merge è quindi la scelta sbagliata: serve un \textbf{Blueprint} con i tre pezzi come
\textbf{componenti} separati, imparentati tra loro, come si è visto nella seconda
lezione. Il merge della lezione è una dimostrazione dello strumento, e serve anche a
rispondere alla domanda sui materiali (sezione~\ref{ue3:sec:domande}).

% ---------------------------------------------------------------------
\section{I materiali della torretta}
\label{ue3:sec:matturret}

``Ora abbiamo importato le texture. Creiamoci i materiali di queste torrette.'' Tre
materiali, uno per pezzo: \textbf{M\_Turret\_Base}, \textbf{M\_Turret\_Pivot} e
\textbf{M\_Turret\_Cannon}. Si potrebbero fare molto elaborati, ``ma partiamo
semplicemente andando a utilizzare le texture che ci sono state date''.

\paragraph{Quattro Texture Sample.} Nel materiale della base Davide prende \textbf{quattro
Texture Sample} e ci mette le texture della base: \textbf{Base Color},
\textbf{Metallic}, \textbf{Roughness} e \textbf{Normal} (Turret Base Color, Turret Base
Metallic, Turret Base Roughness, Turret Base Normal). Poi le collega ai pin
corrispondenti del nodo risultante: il colore al Base Color, la Metallic al Metallic, la
Roughness alla Roughness, la Normal alla Normal (figura~\ref{ue3:fig:turret}).

\begin{figure}[htbp]
\centering
\begin{tikzpicture}[
  font=\small,
  ue3node/.style={draw, thick, rounded corners=4pt, align=center,
                  minimum height=0.85cm, text width=3.6cm},
  ue3col/.style={-{Stealth[length=2mm]}, thick, blue!60!black},
  ue3flt/.style={-{Stealth[length=2mm]}, thick, green!50!black}
]
  \node[ue3node, fill=orange!20] (c) at (0,2.1) {\textbf{Texture Sample}\\
    \footnotesize T\_Turret\_Base\_Color};
  \node[ue3node, fill=gray!15] (m) at (0,0.7) {\textbf{Texture Sample}\\
    \footnotesize T\_Turret\_Base\_Metallic};
  \node[ue3node, fill=gray!15] (r) at (0,-0.7) {\textbf{Texture Sample}\\
    \footnotesize T\_Turret\_Base\_Roughness};
  \node[ue3node, fill=blue!15] (n) at (0,-2.1) {\textbf{Texture Sample}\\
    \footnotesize T\_Turret\_Base\_Normal};
  \node[draw, very thick, rounded corners=4pt, fill=blue!6, minimum width=3.0cm,
    minimum height=5.2cm] (res) at (6.2,0) {};
  \node[font=\small\bfseries] at (6.2,2.3) {M\_Turret\_Base};
  \foreach \lab/\y in {Base Color/2.1, Metallic/0.7, Specular/0.0, Roughness/-0.7,
                       Emissive Color/-1.4, Normal/-2.1}
    \node[font=\footnotesize, anchor=west] at (4.85,\y) {\lab};
  \draw[ue3col] (c.east) -- (4.8,2.1);
  \draw[ue3flt] (m.east) -- node[above, font=\scriptsize]{R} (4.8,0.7);
  \draw[ue3flt] (r.east) -- node[above, font=\scriptsize]{R} (4.8,-0.7);
  \draw[ue3col] (n.east) -- (4.8,-2.1);
\end{tikzpicture}
\caption{Il materiale della base della torretta, fatto a mano: quattro Texture Sample,
quattro pin. Metallic e Roughness sono mappe in bianco e nero, e basta un canale (per
esempio R). M\_Turret\_Pivot e M\_Turret\_Cannon sono identici, con le loro texture.
Ricostruzione dalla dimostrazione; i nomi esatti dei file sono indicativi.}
\label{ue3:fig:turret}
\end{figure}

\paragraph{Le preview.} ``Mi ero dimenticato di farvi vedere una cosa.'' Nella
\emph{preview} del materiale si può scegliere la forma: il \textbf{cilindro}, la
\textbf{sfera}, il \textbf{piano}, il \textbf{cubo} e il \textbf{custom}. Con il custom,
se si seleziona una mesh nel Content Browser e si clicca il pulsante, la preview mostra
\textbf{quella mesh} con il materiale applicato. La prima volta non la mostra: ``sta
compilando''. Poi compare (Davide aveva collegato una texture ``dalla parte sbagliata'',
e la sistema).

\paragraph{Gli slot.} Il materiale è salvato e compilato. Sulla mesh della torretta ci
sono i \textbf{due slot}: in uno c'è M\_Turret\_Base, ma ``per come è stata strutturata
la torretta'' nell'altro va il materiale del \textbf{pivot}. ``Sta rallentando
veramente tanto il PC.''

\paragraph{Un grafico, un engine.} Davide si ferma su che cosa si è fatto: ``qui non
abbiamo fatto altro che utilizzare le mesh che ci sono state fornite da un grafico''.
Idealmente un grafico, in un software a parte, ha realizzato la mesh e le texture, e ce
le ha passate; noi le abbiamo inserite. Si potrebbe lavorare ancora sull'estetica della
torretta modificando i materiali, ``tuttavia, in questo caso non è necessario''.

\paragraph{Il pivot e il cannone, più furbi.} Per il materiale del pivot Davide è ``un
po' più furbo'': cerca la texture \emph{Turret Pivot Color}, e le altre le copia da lì,
perché nella lista ``stanno direttamente dopo''. Collega Metallic, Roughness e Normal,
salva, ``e abbiamo anche il materiale del nostro pivot''. Poi lo stesso per il cannone:
la texture \emph{Cannon Color}, \textbf{Ctrl+C}, \textbf{Ctrl+V}. Lì Davide ricorda le
scorciatoie classiche: ``tutti i classici shortcut di Control'', Ctrl+C, Ctrl+V, Ctrl+Z,
Ctrl+Y, Ctrl+A, ``sono applicabili anche a Unreal nella maggior parte dei casi''.

\paragraph{Salvate spesso.} Uno studente chiede perché il suo non si è applicato:
``Salvalo.'' E Davide ne fa una regola: ``puntualmente mi capita di vedere gente che sta
lavorando a qualcosa, non salva per qualche motivo, Unreal crasha e perde tipo tutto il
lavoro che stava facendo. E quindi salvate spesso.''

\paragraph{La torretta finita.} Nei due slot vanno il pivot e il cannone. ``Oh, eccoci
qua. Bello lucido.'' Salvato tutto, la torretta ha i suoi materiali.

\paragraph{Ritoccare il colore con un Lerp.} ``Se volete che la torretta sia meno
lucida, o che sia meno nitido l'arancione'', si può intervenire sui singoli materiali.
Davide improvvisa sul cannone: non è sicuro di volere il colore ``così netto'', quindi
prende un \textbf{nero} e un \textbf{Lerp}, collega il nero e la texture del colore alle
due entrate e mette nell'alfa un parametro con cui decidere ``quanto si deve vedere'',
per esempio 0,6. Si può convertire in parametro, creare una \textbf{Material
Instance}, e regolarlo da lì. Al posto del nero si può mettere un altro colore, un
verde: ``viene un po' più verdastra, è troppo verde'', e con l'alfa a 0,8 viene ``questo
verde mezzo, insomma, un po' scuro, un po' cromato''. ``Qui chiaramente sto improvvisando
cose, ma cose anche stupide. Ma nel momento in cui voi avete un'idea un po' più chiara,
un po' più specifica di quello che volete, potete andare a fare un po' più di cose.''

\paragraph{Approfondimento: lo stesso lavoro con un master material.} I tre materiali
della torretta hanno lo stesso identico grafo: cambiano solo le texture. È il caso da
manuale per le Material Instance della prima parte della lezione: un solo master
material (per esempio \texttt{M\_Turret}) con quattro Texture Parameter (colore,
metallic, roughness, normal) e magari un Lerp per la tinta, e tre istanze
(\texttt{MI\_Turret\_Base}, \texttt{MI\_Turret\_Pivot}, \texttt{MI\_Turret\_Cannon}) che
impostano le texture di ciascun pezzo. Un solo shader da compilare, e il ritocco
verdastro del cannone diventa una casella da regolare nell'istanza. Se le texture
fossero state importate insieme alle mesh (Import Materials e Import Textures), Unreal
avrebbe creato da solo i materiali, uno per slot, con un grafo simile a quello della
figura~\ref{ue3:fig:turret}.

% ---------------------------------------------------------------------
\section{Le domande: il Reimport e il merge}
\label{ue3:sec:domande}

``Ora tutto chiaro il discorso dell'import e dei materiali?'' Due domande.

\paragraph{Perché il Reimport con la Z?} Una studentessa non ha capito il reimport
modificando la Z. Davide riapre il cannone. Il reimport è servito a \textbf{modificare la
mesh} alla radice. L'oggetto importato, modellato da qualcun altro, era fatto in un
certo modo, e non è un problema in sé: lo si può abbassare e girare a mano, in scena.
``Però poi, tutte le volte che io prendo e rimetto in scena l'oggetto, l'oggetto sarà
sempre posizionato in quel modo e ruotato in quel modo.'' Il reimport invece ``modifica
proprio alla radice il modo in cui funziona l'oggetto''. E il problema principale non
era la rotazione (``so che l'angolo è 113, allora lo rimodifico''), ma il
\textbf{pivot}, ``in una posizione scomoda'': il reimport è un modo per spostarlo.

\paragraph{La Modeling Mode.} ``Non è l'unico modo.'' Davide apre la \textbf{Modeling
Mode}, cerca lo strumento del \textbf{pivot}, e mostra che lo si può mettere al centro,
a sinistra, a destra, oppure ``nel punto specifico che dico io: lo prendo e me lo
sposto''. Ma questa è ``una soluzione che hanno aggiunto con Unreal Engine 5'': prima la
Modeling Mode non c'era, e in Unreal Engine 4 reimportare era il modo comodo per farlo
``senza avere bisogno di rimodellare l'oggetto a posteriori''. ``Fai la stessa cosa,
semplicemente non hai fatto un'operazione extra.''

\paragraph{Una prova.} Davide stesso è curioso: se si reimporta di nuovo con $-135$ e
113? Il reimport \textbf{sovrascrive} quello che c'era prima, ma il cannone in scena
conserva la rotazione che gli era stata data a mano. Davide lo rimette dove stava (``io
sto facendo a cazzotti da tutta la lezione con la finestra di Zoom, se non si fosse
capito'').

\paragraph{Che cosa succede ai materiali con il merge?} Era la domanda arrivata in chat
durante la costruzione dei materiali: se faccio il merge e i pezzi della torretta hanno
texture diverse? Davide prova: seleziona i tre oggetti, \textbf{Actor}, \textbf{Merge
Actors}, nome \emph{Turret}, salva. Il risultato è un oggetto con \textbf{tre slot} per i
materiali: lo slot 0 è tutta la parte del cannone, lo slot 1 il pivot ``e tutta questa
zona qui'', e così via. Se in uno slot si mette un materiale vuoto, ``diventa tutto
nero'': ``vedete, tutta questa zona adesso viene influenzata da quel materiale lì''.
Risposta: con il merge l'actor viene fuso, e i materiali vengono \textbf{riposizionati
negli slot} della nuova mesh.

% ---------------------------------------------------------------------
\section{Lo stretch, lo z-fighting e i materiali World Aligned}
\label{ue3:sec:world}

``Mi è venuto in mente un'ultima cosa che vi potrei far vedere per quanto riguarda i
materiali, e poi vi vorrei far fare una cosina.'' Manca una mezz'oretta, ``visto che
Igor vi ha fatto fare delle cose, ho visto che siete molto partecipi. Quindi vorrei
fomentare questa cosa.''

\paragraph{M\_Test.} Davide ricrea il materiale della corteccia, cancellato per sbaglio:
\textbf{M\_Test}, con due Texture Sample (Bark e la sua Normal) e un Texture
Coordinate, ``proprio la roba più facile del mondo'', con il tiling a 10.

\paragraph{Ingrandire e stirare.} ``Ora vediamo se abbiamo capito come funziona tutto
quanto.'' I materiali sono collegati alle \textbf{UV} di un oggetto. Quindi se
l'oggetto si \textbf{ingrandisce}, al materiale ``non dovrebbe succedere granché'': si
ingrandisce anche lui (``la risoluzione non ne giova particolarmente''). E se l'oggetto
si \textbf{stira} (\emph{stretch}), tirandolo verso destra? ``Credo si stira anche
lui.'' Esatto: il materiale si allunga. ``In alcuni casi non è neanche particolarmente
brutto'', ma la corteccia ``non sembra più corteccia''; e tirando di più, ``così è
brutto''.

\paragraph{Il pavimento grande.} Se si vuole fare un pavimento per una zona molto
ampia, dover pensare a questo ``può essere una noia'': per non stirare il materiale si
dovrebbero fare tante copie dello stesso pezzo, una attaccata all'altra, ``che
tendenzialmente non è quello che voglio'', perché è scomodo.

\paragraph{Una digressione: lo z-fighting.} ``Visto che è una domanda che arriverà,
anticipo'': se due piani si trovano \textbf{nella stessa posizione}, ``clippano'': la
superficie sfarfalla a chiazze, perché l'engine ``non sa se utilizzare il materiale
dell'oggetto a destra o dell'oggetto a sinistra''. Le soluzioni: separarli davvero,
oppure spostarne uno di pochissimo in \textbf{Z}, per esempio di 0,1: ``è
impercettibile, non si nota, però uno è leggermente più sopra rispetto all'altro'', e il
problema sparisce.

\paragraph{Allineato al mondo.} ``C'è un modo per far sì che il mio materiale non si
stiri quando stiro l'oggetto? Chiaramente sì, perché se no non avrei fatto questa
domanda.'' Bisogna fare un materiale che non è allineato con le coordinate
dell'oggetto, ma con quelle del \textbf{mondo}. Davide lo spiega con il \emph{gizmo}
di spostamento: un oggetto ruotato si può muovere lungo gli assi del mondo o lungo i
suoi assi locali, e il pulsante in alto passa dall'uno all'altro. ``Più o meno il
discorso è lo stesso'': nel caso del materiale non si parla di assi, ma di UV.

\paragraph{I nodi World Aligned.} Davide copia il materiale (\textbf{M\_Test1}) e usa
due nodi, \textbf{World Aligned Texture} e \textbf{World Aligned Normal}. Non hanno
scorciatoie: si cercano per nome. Hanno ``un sacco di parametri'', ma per ora ne servono
due: \textbf{Texture Object} e \textbf{Texture Size}. World Aligned Texture va nel
\textbf{Base Color}, World Aligned Normal nella \textbf{Normal}, usando l'uscita
\textbf{XYZ}, ``perché almeno prendo tutti gli assi'': il materiale è allineato al
mondo su tutti e tre.

\paragraph{Texture Size e Texture Object.} Il \textbf{Texture Size}, un vettore a tre
valori, ``fa la stessa cosa del Texture Coordinate'': regola la grandezza della
ripetizione. Perché allora non tenere il Texture Coordinate sulle UV del Texture
Sample? Perché i nodi World Aligned \textbf{non prendono} un Texture Sample: vogliono un
\textbf{Texture Object}. Si converte: tasto destro sul Texture Sample, \textbf{Convert to
Texture Object}. Il Texture Object ``è un oggetto molto più semplice'': non ha
ingressi, ha una sola uscita, niente RGB; ``è molto più limitato in quello che può
fare'', ma è il ``messaggio'' che serve ai nodi World Aligned.

\paragraph{La prova.} Davide prova diversi valori di Texture Size (1, 2, 10) per avvicinarsi
al risultato dell'altro materiale, con qualche esitazione su quale sia ``troppo'' e quale
``troppo poco'' (``è proprio piccola la texture, allora''), e applica il nuovo
materiale a un secondo oggetto. Poi \textbf{stira} entrambi: il primo si stira; il
secondo ``semplicemente genera nuove texture'', rimane orientato con il mondo,
``rimane sempre della stessa proporzione''.

\paragraph{Dov'è la fregatura.} ``Se questo rimane sempre proporzionato anche quando
stiriamo le cose, perché non dovremmo avere tutti i materiali World Aligned? Dov'è la
fregatura?'' Eccola: essendo allineato al mondo, se l'oggetto si \textbf{muove} o
\textbf{ruota}, il materiale \textbf{resta fermo}. Una \textbf{porta} con un materiale
normale si apre e il materiale resta attaccato alla porta; una porta con un materiale
World Aligned si apre e ``il materiale sembra scorrere sulla porta'', ``come se il
materiale fosse animato''. ``Ci siamo? Ho rotto le palle? Datemi un feedback.''

\begin{figure}[htbp]
\centering
\begin{tikzpicture}[font=\small]
  % materiale UV: stiramento
  \node[font=\footnotesize\bfseries] at (2.6,2.6) {materiale normale (UV)};
  \draw[thick, fill=brown!20] (0,1.0) rectangle (1.2,2.2);
  \foreach \i in {0.2,0.6,1.0} \draw[brown!70!black, thick] (\i,1.15) -- (\i,2.05);
  \draw[-{Stealth[length=2mm]}, thick] (1.5,1.6) -- (2.1,1.6);
  \draw[thick, fill=brown!20] (2.4,1.0) rectangle (5.4,2.2);
  \foreach \i in {0.5,1.5,2.5} \draw[brown!70!black, thick] (2.4+\i,1.15) --
    (2.4+\i,2.05);
  \node[font=\footnotesize, align=center] at (2.7,0.55) {stirato: le venature si
    allargano};
  % materiale world aligned
  \begin{scope}[shift={(6.6,0)}]
    \node[font=\footnotesize\bfseries] at (2.6,2.6) {materiale World Aligned};
    \draw[thick, fill=brown!20] (0,1.0) rectangle (1.2,2.2);
    \foreach \i in {0.2,0.6,1.0} \draw[brown!70!black, thick] (\i,1.15) -- (\i,2.05);
    \draw[-{Stealth[length=2mm]}, thick] (1.5,1.6) -- (2.1,1.6);
    \draw[thick, fill=brown!20] (2.4,1.0) rectangle (5.4,2.2);
    \foreach \i in {0.2,0.6,...,3.0} \draw[brown!70!black, thick] (2.4+\i,1.15) --
      (2.4+\i,2.05);
    \node[font=\footnotesize, align=center] at (2.7,0.55) {stirato: ``genera nuove
      texture'', stessa scala};
  \end{scope}
  % la porta
  \node[font=\footnotesize, align=center, text width=12.5cm] at (5.9,-0.3)
    {la fregatura: se l'oggetto si muove o ruota (una porta che si apre), il materiale
    World Aligned resta fermo nel mondo e ``sembra scorrere'' sull'oggetto};
\end{tikzpicture}
\caption{Materiale legato alle UV e materiale World Aligned a confronto. Il primo
segue l'oggetto e si deforma con lui; il secondo non si deforma, ma non segue
l'oggetto quando si muove. Va bene per pavimenti, muri e rocce statici; non per porte,
veicoli o la torretta.}
\label{ue3:fig:world}
\end{figure}

\paragraph{Le linee nere.} Una studentessa vede delle \textbf{linee nere} al confine tra
una ripetizione della texture e l'altra, sul pavimento. Davide guarda il suo grafo:
stessa texture, Bark, ma ha collegato l'uscita \textbf{XY} invece di \textbf{XYZ}. ``In
quel modo lui aveva due assi che erano allineati col mondo, e l'altro no'', e il
risultato era diverso. ``Ok, grazie.''

\paragraph{Approfondimento: come funziona un World Aligned.} Il nodo World Aligned
Texture, una Material Function del motore, non usa le UV della mesh: per ogni pixel
prende la sua \textbf{posizione nel mondo} e la usa come coordinata della texture. Lo fa
tre volte, proiettando la texture lungo i tre assi (dall'alto, di lato, di fronte), e
mescola le tre proiezioni secondo la direzione della superficie: è la tecnica detta
\textbf{triplanar mapping}. L'uscita XY usa una sola proiezione (dall'alto, buona per i
pavimenti), l'uscita XYZ le mescola tutte; per questo, su un oggetto con facce in tutte
le direzioni, con XY le pareti laterali ricevono una texture ``stirata'' a strisce, che
può dare proprio le linee scure viste dalla studentessa. Il \textbf{Texture Size} è
espresso in \textbf{unità del mondo}, cioè in centimetri: è la grandezza di una
ripetizione della texture. Funziona quindi al contrario del tiling: un tiling più alto
fa ripetere di più la texture, un Texture Size più alto la fa ripetere di meno (una
texture di 100 si ripete ogni metro, una di 10 ogni dieci centimetri). È la ragione
delle esitazioni di Davide su quale valore fosse ``troppo''. Il costo: tre letture della
texture invece di una, quindi un materiale un po' più pesante.

% ---------------------------------------------------------------------
\section{L'esercizio della lava, il nodo Noise e la chiusura}
\label{ue3:sec:lava}

``Allora, ci siamo visti un po' di cose, fatto un miscuglione di cose per quanto riguarda
materiali, texture e import delle mesh.'' Nessuna domanda: ``sembra tutto a posto''.

\paragraph{La consegna.} Per esercitarsi con i materiali, e magari scoprire qualche nodo
non usato a lezione (il nodo \textbf{Noise}, qualche \textbf{maschera}), Davide propone
una prova di una decina di minuti, da continuare poi a casa come le cose di Igor:
\textbf{un materiale di lava}. Di che componenti avrà bisogno? Deve \textbf{muoversi}
un po', quindi bisognerà gestire la velocità del \textbf{Panner}, ``a meno che non vi
cercate proprio la texture della lava''. Ma la richiesta è di improvvisare: ``provate a
vedere cosa riuscite a combinare andando a fare delle operazioni fra colori e fra altre
texture'', con un po' di movimento, un \textbf{Base Color}, ``ma anche un
\textbf{Emissive}, perché la lava comunque deve emettere luce''.

\paragraph{Il ripasso dei nodi.} Una studentessa non ricorda il nome della funzione
``che abbiamo usato per collegare rosso e blu'': il Lerp. Davide ricondivide lo schermo
e fa il riepilogo dei nodi della giornata: la \textbf{costante}, il \textbf{colore}, la
\textbf{texture}, il \textbf{nodo risultante}, il \textbf{Texture Coordinate}, i
\textbf{World Aligned}, le \textbf{operazioni} (con il consiglio sulla moltiplicazione
visto nella sezione~\ref{ue3:sec:costanti}), il \textbf{Panner}, il \textbf{Lerp}.
Divisione e sottrazione servono soprattutto per gli effetti speciali, e Davide vorrebbe
fare una lezione su \textbf{Niagara}, ``il tool che usa Unreal per fare gli effetti
speciali'', se ci sarà tempo.

\paragraph{Le prossime lezioni con Davide.} Davide annuncia come saranno le sue
lezioni finali: ci sarà ``una lezione che sarà un'esercitazione intera'' in cui si
proverà a fare dei \textbf{livelli}, e chiederà agli studenti di condividere lo schermo
e di \emph{streamare} a turno (``se per voi è un problema, magari lo facciamo solo con le
persone per cui non è un problema'', e lo dirà in anticipo). Per esempio, ``un livello
\textbf{platform}'', con le interazioni che si faranno nelle lezioni di Igor e quello che
si vedrà insieme per il \emph{level design}.

\paragraph{Il nodo Noise.} Un nodo ``un po' più complicato'', per chi vuole provare: il
\textbf{Noise}, che ``serve per simulare un \emph{noise}'': il \emph{Voronoi}, il
\emph{White Noise}, un'alternanza di bianco e nero ``che si usa molto spesso per fare
delle distorsioni''. Ha molti parametri: la scala del \emph{noise}, e la possibilità di
filtrarlo con dei vettori o con delle mappe. ``Potete smanettarci un pochino.''

\paragraph{Lo stream della studentessa.} ``C'è qualcuno che, nel mentre sta provando a
fare qualcosa, vuole streamare in questi ultimi cinque minuti?'' Una studentessa ha
provato, ma non riesce a trasformare la texture in World Aligned: è perché usa un
Texture Sample, e serviva un Texture Object. ``Prova a non usare il World Aligned, e
farlo così, diciamo, semplice.'' Senza World Aligned le viene ``questa cosa'', sempre
con il Bark. Davide la guida:
\begin{itemize}
  \item nella lava servono ``delle zone più scure e delle zone più luminose'': quindi
        un \textbf{Lerp} tra un \textbf{colore} e un \textbf{nero}, gestito da un'alfa
        che sia una texture con dei bianchi e dei neri ``posizionati come vuoi'', da
        mandare nell'\textbf{Emissive Color};
  \item l'alfa va cercata tra le texture del progetto, ``qualcosa per dare un effetto
        tipo di distorsione''; le \textbf{Normal}, invece, ``lasciale perdere'': per
        un'alfa si usano le Diffuse o le mappe come Specular e Roughness;
  \item se una texture fa da alfa, nelle entrate del Lerp ci vogliono \textbf{due
        colori}, ``uno che sia molto scuro, uno che sia più chiaro''; ricordando che
        l'alfa ``prende i bianchi e i neri e dice: nei bianchi vai più verso B, nei neri
        vai più verso A'';
  \item il risultato va ``forse più nell'Emissive che non nel Base Color'', perché la
        lava la si vuole ``infuocata, accesa, luccicante''; il Base Color si può anche
        non usare, o usarlo con una texture adatta (``il Bark, per come è fatto, forse
        non è proprio la texture più adatta per fare una lava'');
  \item nell'Emissive si può \textbf{moltiplicare} il risultato, in modo che le zone
        accese siano molto luminose.
\end{itemize}
La studentessa nota che il nero, nell'Emissive, quasi non si vede: ``esatto'', il nero
non emette niente. E usandolo come colore ``non si vede molto'', perché il Bark come alfa
``è molto bianco'', e finisce per scurire un po' il tutto. ``È comunque una cosa
complicata, quindi vedete effettivamente se riuscite a tirar fuori qualcosa.''

\begin{figure}[htbp]
\centering
\begin{tikzpicture}[
  font=\small,
  ue3node/.style={draw, thick, rounded corners=4pt, align=center,
                  minimum height=0.85cm},
  ue3uv/.style={-{Stealth[length=2mm]}, thick, red!60!black},
  ue3flt/.style={-{Stealth[length=2mm]}, thick, green!50!black},
  ue3col/.style={-{Stealth[length=2mm]}, thick, blue!60!black}
]
  \node[ue3node, fill=red!12, text width=1.7cm] (tc) at (0,0) {\textbf{TexCoord}};
  \node[ue3node, fill=gray!10, text width=1.9cm] (pan) at (2.6,0) {\textbf{Panner}\\
    \footnotesize Speed 0,05; 0,02};
  \node[ue3node, fill=gray!15, text width=2.3cm] (msk) at (5.6,0)
    {\textbf{Texture Sample}\\\footnotesize maschera in b/n\\\footnotesize (o un
    Noise)};
  \node[ue3node, fill=black!70, text=white, text width=1.9cm] (dark) at (5.6,2.0)
    {\textbf{scuro}\\\footnotesize quasi nero};
  \node[ue3node, fill=orange!60, text width=1.9cm] (hot) at (5.6,-2.0)
    {\textbf{chiaro}\\\footnotesize arancione};
  \node[ue3node, fill=gray!10, text width=1.3cm, minimum height=2.0cm] (lerp)
    at (8.6,0) {\textbf{Lerp}\\[3pt]\footnotesize A\\\footnotesize B\\
    \footnotesize Alpha};
  \node[ue3node, fill=green!8, text width=1.6cm] (k) at (8.6,-2.4)
    {\footnotesize intensità\\\footnotesize (parametro)};
  \node[ue3node, fill=gray!10, text width=1.4cm] (mul) at (10.9,-0.6)
    {\textbf{Multiply}};
  \node[ue3node, fill=blue!6, text width=2.0cm, minimum height=1.6cm] (res)
    at (13.4,0) {\footnotesize Base Color\\[6pt]\footnotesize\textbf{Emissive}};
  \draw[ue3uv] (tc) -- (pan);
  \draw[ue3uv] (pan) -- node[above, font=\scriptsize]{UVs} (msk);
  \draw[ue3col] (dark.east) -- ++(0.7,0) |- ([yshift=5mm]lerp.west);
  \draw[ue3col] (hot.east) -- ++(0.9,0) |- ([yshift=0mm]lerp.west);
  \draw[ue3flt] (msk.east) -- ++(0.35,0) |- ([yshift=-5mm]lerp.west);
  \draw[ue3col] (lerp.east) -- ++(0.3,0) |- ([yshift=2mm]mul.west);
  \draw[ue3flt] (k.east) -- ++(0.5,0) |- ([yshift=-2mm]mul.west);
  \draw[ue3col] (mul.east) -- ++(0.4,0) |- ([yshift=-3mm]res.west);
  \draw[ue3col, dashed] (lerp.north) |- ([yshift=4mm]res.west);
\end{tikzpicture}
\caption{Una possibile soluzione dell'esercizio della lava, costruita con i soli nodi
della lezione e con i consigli dati alla studentessa: una maschera in bianco e nero che
scorre con il Panner fa da alfa a un Lerp tra uno scuro e un arancione; il risultato,
moltiplicato, va nell'Emissive (e, tratteggiato, eventualmente anche nel Base Color).
Non è un grafo mostrato a lezione: i valori sono indicativi.}
\label{ue3:fig:lava}
\end{figure}

\paragraph{La chiusura.} ``Vabbè, io in realtà adesso ho gettato questa bomba, e penso
che vi lascerò fare. Prendetevi tutto il tempo che volete, nessuna pressione.'' Davide
rivedrà la classe ``a luglio, probabilmente'', per le ultime quattro o cinque lezioni del
corso, sul \textbf{level design} e sulla \textbf{creazione di livelli}. ``Si sono fatte
le 7.'' ``Direi che è stato un piacere, mi sembrate tutti abbastanza interessati,
sicuramente belli, belli, belli attivi.'' E forse ci si rivedrà prima, ``se Igor ha
bisogno di un'altra sostituzione''. ``Ciao a tutti, arrivederci.''

\paragraph{Approfondimento: una lava più convincente.} Partendo dallo schema della
figura~\ref{ue3:fig:lava}, tre accorgimenti, tutti con nodi già visti, rendono la lava
molto più viva. Primo: \textbf{due Panner} con velocità e tiling diversi sulla stessa
maschera, sommati o moltiplicati, così che il movimento non si ripeta mai uguale (lo
stesso trucco dell'acqua). Secondo: un Multiply sull'Emissive con valori ben oltre 1
(10, 20), perché solo così il \emph{bloom} fa ``bruciare'' le zone calde, come nel
materiale con il Glow Intensity. Terzo: usare la \textbf{stessa} maschera per la
\textbf{Roughness}, con un Lerp tra 0,3 (la lava liquida, lucida) e 1 (la crosta
raffreddata, opaca): il materiale cambia così non solo colore ma anche superficie, che è
il senso del PBR. Il nodo Noise di Davide può sostituire la texture della maschera, ma è
calcolato a ogni pixel e costa più di una texture: per un oggetto grande in scena, una
texture di \emph{noise} è di solito la scelta più leggera.

\begin{table}[htbp]
\centering
\small
\begin{tabular}{@{}p{4.4cm}p{8.6cm}@{}}
\toprule
\textbf{Comando} & \textbf{Che cosa fa} \\
\midrule
Tasto destro nel grafo & menu dei nodi, con ricerca \\
T + clic & Texture Sample \\
1 + clic / 2 + clic / 3 + clic & Constant, Constant2Vector, Constant3Vector (il
  colore) \\
A / M / D + clic & Add, Multiply, Divide \\
S + clic & Scalar Parameter (per questo Subtract non ha scorciatoia) \\
U + clic & Texture Coordinate (``lavora sulle UV'') \\
L + clic & Lerp (Linear Interpolate) \\
P + clic & Panner \\
Alt + clic su un pin & stacca i collegamenti \\
Tasto destro su un nodo, Convert to Parameter & trasforma il nodo in un parametro,
  regolabile dalle Material Instance \\
Tasto destro su un Texture Sample, Convert to Texture Object & per i nodi World Aligned \\
Use Selected Asset from Content Browser & mette nel nodo la texture selezionata nel
  Content Browser \\
Tasto destro su un materiale, Create Material Instance & crea un'istanza \\
Reimport Mesh With Dialog & reimporta una mesh riaprendo la finestra di import \\
Ctrl+G / Shift+G & raggruppa e separa gli actor selezionati \\
Actor, Merge Actors & fonde gli actor selezionati in una nuova mesh \\
Ctrl+C, Ctrl+V, Ctrl+Z, Ctrl+Y, Ctrl+A & le scorciatoie classiche, valide quasi
  ovunque \\
Ctrl+S & salvare: ``salvate spesso'' \\
\bottomrule
\end{tabular}
\caption{I comandi della terza lezione nell'editor dei materiali, nel Content Browser e
nel Level Editor di Unreal Engine 5. Le scorciatoie con il clic si usano tenendo premuto
il tasto e cliccando nel grafo; World Aligned Texture e World Aligned Normal non ne
hanno.}
\label{ue3:tab:scorciatoie}
\end{table}

% ---------------------------------------------------------------------
\section*{Punti chiave del capitolo}
\addcontentsline{toc}{section}{Punti chiave del capitolo}

\begin{itemize}
  \item Il grafico lavora soprattutto fuori da Unreal (Blender, Maya); dentro Unreal
        \textbf{importa} gli oggetti e costruisce gli \textbf{shader}.
  \item Una \textbf{texture} è la ``carta regalo'' dell'oggetto: si avvolge sulla mesh
        attraverso le \textbf{UV}, che si creano prima, nel software di modellazione.
        Una texture non pensata per quelle UV si stira.
  \item Le mappe principali: \textbf{Diffuse/Albedo} (il colore, suffisso \_D o \_A),
        \textbf{Normal} (blu: l'illusione del rilievo senza poligoni in più),
        \textbf{Specular}, \textbf{Roughness}, \textbf{Metallic}, la mappa di
        \textbf{altezza} (\emph{displacement} vero, ``una normal sotto steroidi'').
  \item Nelle mappe in bianco e nero il \textbf{bianco è 1}, il \textbf{nero è 0}, il
        grigio sta in mezzo. Ogni canale RGB è una mappa a sé: una texture
        \textbf{AORM}/ORM impacchetta tre mappe in un file.
  \item Il \textbf{materiale} è la ``vernice'', uno \textbf{shader}: si costruisce con
        nodi, come un Blueprint, e combina le texture. Sugli oggetti si applicano
        materiali, non texture.
  \item \textbf{Material Domain} (per che cosa), \textbf{Blend Mode} (Opaque, Masked,
        Translucent) e \textbf{Shading Model} (Default Lit, Unlit) cambiano i pin del
        nodo risultante.
  \item Nodi di base: \textbf{Texture Sample} (T), costanti (1, 2, 3), \textbf{Add},
        \textbf{Multiply}, \textbf{Divide} (A, M, D), Scalar Parameter (S),
        \textbf{Texture Coordinate} (U), \textbf{Lerp} (L), \textbf{Panner} (P). La
        moltiplicazione è l'operazione più usata.
  \item Il rendering è \textbf{PBR}: Roughness 1 è la roccia, 0 lo specchio; Metallic
        e Specular decidono come la superficie riflette; l'\textbf{Emissive}
        moltiplicato fa brillare il materiale.
  \item \textbf{UTiling} e \textbf{VTiling} (o un Multiply sulle coordinate) decidono
        quante volte si ripete la texture: 2 e 2 sono quattro ripetizioni, 0,5 e 0,5 un
        quarto della texture.
  \item Un nodo diventa un \textbf{parametro} con Convert to Parameter. Le
        \textbf{Material Instance} (MI\_) cambiano i parametri senza toccare il grafo
        del \textbf{master material} (M\_), si aggiornano subito e sono più leggere:
        sugli oggetti si applicano sempre le istanze.
  \item Le \textbf{Material Function} sono ``parentesi'' di calcoli riusabili tra
        materiali e tra membri del team.
  \item Il \textbf{Lerp} interpola tra A e B secondo un'\textbf{alfa} da 0 a 1, che può
        essere una texture (``l'arbitro''); il \textbf{Panner} fa scorrere la texture
        (acqua, fluidi, lava).
  \item Importare \textbf{FBX}, controllare le \textbf{unità} (Unreal usa i
        centimetri), texture in PNG o JPEG. La finestra di import genera collisioni,
        corregge il transform, importa materiali e texture.
  \item Una mesh è fatta di \textbf{vertici}, \textbf{spigoli} e \textbf{facce}
        (triangoli). Il \textbf{pivot} deve stare dove il pezzo ruota: con
        \textbf{Reimport Mesh With Dialog} (cannone: Z $-135$, rotazione 113) o con la
        \textbf{Modeling Mode} di Unreal Engine 5.
  \item Le \textbf{simple collision} (box, sfera, capsula: il ``capsulone'' dei
        \emph{character}) sono leggere; le \textbf{complex} usano i triangoli della mesh
        e costano molto di più.
  \item \textbf{Ctrl+G} raggruppa; \textbf{Merge Actors} fonde in una nuova mesh, con
        tanti slot quanti i materiali dei pezzi.
  \item Due piani nella stessa posizione sfarfallano (\textbf{z-fighting}): basta
        alzarne uno di 0,1. I materiali \textbf{World Aligned} (Texture Object, Texture
        Size, uscita XYZ) non si stirano, ma restano fermi se l'oggetto si muove.
  \item \textbf{Salvate spesso.} E per esercizio: un materiale di \textbf{lava}, con
        Panner, Lerp ed Emissive.
\end{itemize}
